\documentclass[11pt,twoside,openany]{book}
\usepackage{optitab-livre}
\usepackage[version=4]{mhchem}
\usepackage{siunitx}
\usepackage{adjustbox}
% Adaptation collège de la charte éditoriale utilisée en Terminale.
\renewcommand{\otNiveau}{Troisième · Brevet · Collège}
\renewcommand{\otNbChapitres}{18}
\renewcommand{\otNbExosChapitres}{270}
\renewcommand{\otNbDefis}{18}
\renewcommand{\otNbExosTotal}{288}
\renewcommand{\otSousTitre}{Cours et exercices corrigés}
\renewcommand{\otCleardoublepage}{\clearpage}
\setlength{\parskip}{2.2mm}
\setlength{\headheight}{15pt}
\usepackage{xurl}
\urlstyle{same}
\linespread{1.05}
\tcbset{otEncadre/.append style={fontupper=\fontsize{11.5}{15}\selectfont,
 before skip=3mm,after skip=3mm,top=2.8mm,bottom=2.8mm}}
\newenvironment{pcEncadre}[2]{%
 \def\pcCol{otNavy}\def\pcFond{otNavyTint}%
 \ifstrequal{#1}{method}{\def\pcCol{otGreenDark}\def\pcFond{otGreenTint}}{}%
 \ifstrequal{#1}{attention}{\def\pcCol{otOrange!80!black}\def\pcFond{otOrange!10}}{}%
 \begin{tcolorbox}[otEncadre,colback=\pcFond,borderline west={1.1mm}{0pt}{\pcCol}]
 \otEtiq{\pcCol}{#2}\par\vspace{1mm}}{\end{tcolorbox}}

\newcommand{\pcPage}[2]{%
 \clearpage\markright{#1}\otPageTitre{#1}{#2}\vspace{-2mm}}
\newcommand{\pcSousTitre}[1]{%
 \par\addvspace{3.5mm}\noindent{\otSemi\fontsize{12}{15}\selectfont\color{otNavy}#1}\par\nopagebreak\vspace{1mm}}
\newcommand{\pcRetour}[2]{%
 \par{\otMedium\fontsize{8.5}{11}\selectfont\color{otGreenDark}\hyperref[#1]{#2\enspace·\enspace p.~\pageref*{#1}}}\par}
\newcommand{\pcFigure}[2]{%
 \par\vspace{2mm}\noindent\begin{minipage}{\linewidth}\centering
 \adjincludegraphics[max width=\linewidth,max height=57mm]{#1}\par
 \ifstrempty{#2}{}{\vspace{1.3mm}{\otLight\fontsize{8.7}{11.5}\selectfont\color{otMuted}\raggedright #2\par}}
 \end{minipage}\par\vspace{1.5mm}}
% La largeur maximale est prise en charge par adjustbox, sans déformer l'image.

\newsavebox{\pcFicheBox}
\newenvironment{pcFiche}[1]{%
 \def\pcFicheId{#1}\par\noindent
 \begin{lrbox}{\pcFicheBox}\begin{minipage}{\linewidth}
}{%
 \end{minipage}\end{lrbox}%
 \dimen0=\dimexpr\pagegoal-\pagetotal-4mm\relax
 \ifdim\dimen0>\textheight\dimen0=\textheight\fi
 \typeout{PCFICHE|\pcFicheId|\the\dimexpr\ht\pcFicheBox+\dp\pcFicheBox\relax|\the\dimen0}%
 \ifdim\dimexpr\ht\pcFicheBox+\dp\pcFicheBox\relax>\dimen0
   \PackageError{cinquieme}{PCFICHE-TROP|\pcFicheId}{Repartir le contenu sans reduire le texte.}%
 \else\noindent\usebox{\pcFicheBox}\fi\par}

\newcommand{\pcSomPartie}[2]{%
 \par\addvspace{3.2mm}{\otSemi\fontsize{8}{10}\selectfont\color{otGreenDark}\otCaps[12]{Partie #1}}\enspace
 {\otSemi\fontsize{10}{12}\selectfont\color{otNavy}#2}\par\vspace{1.8mm}}
\newcommand{\pcSomLigne}[3]{%
 \noindent\makebox[9mm][l]{\otSemi\fontsize{9}{11}\selectfont\color{otNavy}#1}%
 \hyperref[#3]{{\otRegular\fontsize{9.4}{12}\selectfont #2}%
 \otPoints{\otSemi\fontsize{9}{11}\selectfont\color{otNavy}\pageref*{#3}}}\par\vspace{0.5mm}}

\newcommand{\pcMission}[3]{%
 \phantomsection\label{mission-#1-#2}%
 \otEntete{otNavy}{#1.#2}{#3}{0}{\otRenvoi{Corrigé}{cor-mission-#1-#2}}}
\newcommand{\pcSolutionMission}[3]{%
 \phantomsection\label{cor-mission-#1-#2}%
 \otEntete{otGreen}{#1.#2}{#3}{0}{\otRenvoi{Énoncé}{mission-#1-#2}}}
\newenvironment{pcLexique}[2]{%
 \par\addvspace{3.8mm}\Needspace{5\baselineskip}\noindent
 {\otSemi\fontsize{11}{14}\selectfont\color{otNavy}#1}\hfill
 {\otMedium\fontsize{8}{11}\selectfont\color{otGreenDark}\hyperref[chap-#2]{Chapitre #2\enspace·\enspace p.~\pageref*{chap-#2}}}\par\nopagebreak\vspace{1mm}
 \fontsize{11.5}{15}\selectfont}{\par}

% Les ouvertures conservent le bandeau bleu, le grand numéro, les deux
% colonnes d'objectifs et de repères, ainsi que les onglets de Terminale.
\renewcommand{\otQRCours}{%
 \begin{tikzpicture}\node[fill=otNavyTint,rounded corners=2mm,inner sep=3.5mm,text width=48mm]
 {{\otSemi\fontsize{8.5}{11}\selectfont\color{otNavy}Mon carnet de sciences}\par\vspace{1mm}
  {\otLight\fontsize{8}{11}\selectfont Une question, un schéma légendé et une conclusion à garder.}};\end{tikzpicture}}
\renewcommand{\otLegende}{%
 {\otLight\fontsize{7.8}{10}\selectfont\color{otMuted}Difficulté\hspace{3mm}%
 \otDots{1}\enspace débuter\hspace{4mm}\otDots{2}\enspace s’entraîner\hspace{4mm}\otDots{3}\enspace raisonner}}
\renewcommand{\otDansCe}{%
 \pcOuvertureLigne{Le cours}{ess-\otChapNum}%
 \pcOuvertureLigne{Je démarre}{rub-\otChapNum-1}%
 \pcOuvertureLigne{Je m’entraîne}{rub-\otChapNum-2}%
 \pcOuvertureLigne{Je raisonne}{rub-\otChapNum-3}%
 \pcOuvertureLigne{Défi}{defi-ex-\otChapNum}%
 {\color{otRule}\rule{\linewidth}{0.4pt}}\par\vspace{1mm}%
 \pcOuvertureLigne{Corrigés}{corchap-\otChapNum}}
\newcommand{\pcOuvertureLigne}[2]{%
 \noindent{\otSemi\fontsize{9}{12}\selectfont\color{otNavy}#1}\hfill
 {\otMedium\fontsize{8}{11}\selectfont\hyperref[#2]{p.~\pageref*{#2}}}\par\vspace{2mm}}

% Décors propres aux quatre thèmes du collège.
\expandafter\newcommand\csname otDecocollege1\endcsname{\otDecoGrille
 \foreach \x/\y in {151/70,160/72,169/69,154/80,164/82,174/79,157/90,167/93,178/90}{%
 \fill[white,opacity=0.26] (\x mm,\y mm) circle (1.3mm);}
 \draw[white,opacity=0.4,line width=0.8pt] (145mm,99mm) -- (145mm,61mm) -- (185mm,61mm) -- (185mm,99mm);}
\expandafter\newcommand\csname otDecocollege2\endcsname{\otDecoGrille
 \draw[white,opacity=0.3,line width=1pt] (138mm,70mm) .. controls (159mm,103mm) and (172mm,54mm) .. (196mm,84mm);
 \foreach \x/\y in {138/70,151/80,164/81,177/78,190/79}{\fill[otGreen] (\x mm,\y mm) circle (1.8mm);}}
\expandafter\newcommand\csname otDecocollege3\endcsname{\otDecoGrille
 \draw[white,opacity=0.4,line width=1pt] (141mm,67mm) -- (141mm,91mm) -- (189mm,91mm) -- (189mm,67mm) -- cycle;
 \draw[white,opacity=0.6,line width=1pt] (141mm,77mm) -- (141mm,81mm);
 \fill[otNavy] (162mm,86mm) rectangle (172mm,96mm);
 \draw[otGreen,line width=1pt] (167mm,91mm) circle (4.5mm);
 \draw[otGreen,line width=0.7pt] (164mm,88mm) -- (170mm,94mm) (164mm,94mm) -- (170mm,88mm);}
\expandafter\newcommand\csname otDecocollege4\endcsname{\otDecoGrille
 \draw[white,opacity=0.45,line width=1pt] plot[domain=0:6,samples=110] ({135mm+\x*11mm},{79mm+sin(\x*3 r)*9mm});}

\renewcommand{\otTitreLivre}{Physique-chimie}
\renewcommand{\otSousTitre}{Cours et exercices corrigés}
\clubpenalty=10000\widowpenalty=10000\displaywidowpenalty=10000
\renewcommand{\pcPage}[2]{\clearpage\markright{#1}\otPageTitre{#1}{#2}\vspace{-4mm}}
\renewcommand{\otFootText}[1]{{\otLight\fontsize{6.5}{8}\selectfont\color{otMuted}#1}}
\fancypagestyle{otMain}{\fancyhf{}\renewcommand{\headrulewidth}{0.4pt}
\fancyhead[LE]{\otHeadNum\hspace{.8em}\otHeadText{\leftmark}}
\fancyhead[RO]{\otHeadText{\rightmark}\hspace{.8em}\otHeadNum}
\fancyfoot[LE]{\otThumb\otFootNum\hspace{1.2em}\otFootText{OptiTAB · Troisième · Brevet 2027}}
\fancyfoot[RO]{\otThumb\otFootText{Physique-chimie · Cours et exercices corrigés}\hspace{1.2em}\otFootNum}}
\patchcmd{\otEntete}{\addvspace{6.5mm}}{\addvspace{4mm}}{}{}
\patchcmd{\otEntete}{\addvspace{6.5mm}}{\addvspace{4mm}}{}{}
\patchcmd{\otEntete}{\otBesoinPlace{9\baselineskip}}{\otBesoinPlace{5\baselineskip}}{}{}
\newcommand{\unite}[1]{\ensuremath{\,\symup{#1}}}
\setlength{\emergencystretch}{2em}
\Urlmuskip=0mu plus 1mu\relax
\newcommand{\pcAnnexe}[2]{\clearpage\gdef\otChapNum{0}\gdef\otHeadTag{}\markboth{#1}{#1}\pcPage{#1}{#2}}
\newcommand{\pcDefiTitre}[2]{\otEntete{otNavy}{D#1}{#2}{3}{\otRenvoi{Corrigé}{cor-defi-#1}}}
\newcommand{\pcSolutionDefi}[2]{\otEntete{otGreen}{D#1}{#2}{0}{\otRenvoi{Énoncé}{defi-ex-#1}}}
\newcommand{\pcGrille}{\par\vspace{4mm}\begin{tikzpicture}\draw[step=5mm,otRule!70,line width=.3pt] (.2pt,0) grid (\dimexpr\linewidth-.2pt\relax,60mm);\end{tikzpicture}}
\newcommand{\pcRepere}[2]{\par\noindent\hyperref[#1]{{\otMedium\fontsize{9}{12}\selectfont\color{otGreenDark}#2 · p.~\pageref*{#1}}}\par\vspace{2mm}}
\newcommand{\pcMot}[1]{{\otSemi\color{otNavy}#1}}
\newcommand{\pcCase}{\tikz[baseline=-.5mm]\draw[otMuted,line width=.4pt] (0,0) rectangle (2mm,2mm);}
\newcommand{\pcFigureEx}[2]{\par\vspace{1.5mm}\noindent\begin{minipage}{\linewidth}\centering\adjincludegraphics[max width=\linewidth,max height=46mm]{#1}\par\ifstrempty{#2}{}{\vspace{1mm}{\otLight\fontsize{8.7}{11.5}\selectfont\color{otMuted}\raggedright #2\par}}\end{minipage}\par\vspace{1.5mm}}
\newcommand{\pcFigureExact}[2]{\par\vspace{1.5mm}\noindent\begin{minipage}{\linewidth}\centering\includegraphics{#1}\par\ifstrempty{#2}{}{\vspace{1mm}{\otLight\fontsize{8.7}{11.5}\selectfont\color{otMuted}\raggedright #2\par}}\end{minipage}\par\vspace{1.5mm}}
\renewcommand{\corrigesChapitre}[2]{\clearpage\gdef\otChapNum{#1}\gdef\otChapTitle{#2}\gdef\otHeadTag{\otTwoDigits{#1}}
\global\otCorrigestrue\global\otSoupletrue\gdef\otAccent{otGreenDark}\markboth{Solutions · #2}{Solutions · #2}
\phantomsection\label{corchap-#1}\pdfbookmark[1]{#1 · #2}{bm-corchap-#1}\otPageTitre{Corrigés · #1}{#2}
\fontsize{11}{14}\selectfont\setlength{\parskip}{1mm}\setlist[questions,1]{label=\otLabel{\arabic*.},leftmargin=1.7em,labelsep=.45em,itemsep=.25em,topsep=.2em,parsep=.1em}}
\sisetup{locale=FR,per-mode=symbol}

\hypersetup{pdftitle={OptiTAB - Physique-chimie - Troisième - Brevet 2027 - Manuel complet},pdfauthor={OptiTAB},pdfsubject={France - Cycle 4 - Édition 2026-2027 - Série générale},pdfkeywords={cours, schémas, exercices, corrigés, programme français, DNB, brevets blancs}}
\begin{document}\pagenumbering{alph}\thispagestyle{empty}
\begin{tikzpicture}[remember picture,overlay]\begin{scope}[shift={(current page.south west)}]
\fill[otNavy] (0,0) rectangle (210mm,297mm);
\shade[bottom color=otNavyDark,top color=otNavy] (0,0) rectangle (210mm,165mm);
\draw[white,opacity=.045,step=8mm] (0,0) grid (210mm,297mm);
\node[anchor=north west,inner sep=0pt] at (22mm,279mm){\includegraphics[width=60mm]{logo-blanc}};
\node[anchor=west,text=otGreen,font=\otSemi\fontsize{11}{14}\selectfont] at (24mm,242mm){\otCaps[15]{Troisième · Série générale}};
\node[anchor=west,text=white,font=\otBold\fontsize{36}{42}\selectfont] at (22mm,217mm){Physique-chimie};
\node[anchor=west,text=white,font=\otLight\fontsize{32}{38}\selectfont] at (23mm,194mm){Objectif Brevet 2027};
\fill[otGreen] (24mm,180mm) rectangle (64mm,181.4mm);
\node[anchor=north west,text=white,text width=151mm,font=\otRegular\fontsize{11}{16}\selectfont] at (24mm,166mm)
{Tout le parcours pour comprendre, s’entraîner\\et rédiger une réponse scientifique solide.};
\node[anchor=north west,text=white,text width=151mm,font=\otRegular\fontsize{10.5}{15}\selectfont] at (24mm,132mm)
{\otCheck\hspace{3mm}18 chapitres · 270 exercices corrigés\\[3mm]
\otCheck\hspace{3mm}Énoncés et corrigés illustrés\\[3mm]
\otCheck\hspace{3mm}15 sujets de physique-chimie · format Brevet\\[3mm]
\otCheck\hspace{3mm}Barèmes, méthodes et révisions du cycle 4};
\draw[otGreen,line width=2pt] (24mm,49mm)--(60mm,49mm);
\node[anchor=west,text=white,font=\otSemi\fontsize{9}{12}\selectfont] at (24mm,39mm){France · Édition 2026-2027 · Manuel complet};
\node[anchor=west,text=white,opacity=.7,font=\otLight\fontsize{8}{11}\selectfont] at (24mm,29mm){Fondé sur le programme français et les modalités du DNB 2027};
\node[anchor=west,text=white,font=\otMedium\fontsize{10}{12}\selectfont] at (24mm,15mm){optitab.net};
\end{scope}\end{tikzpicture}\null\clearpage
\pagenumbering{arabic}\pagestyle{otFront}
\otPageTitre{Comprendre, puis réussir}{Un cours de Troisième et une préparation au Brevet}
Ce manuel réunit les apprentissages de fin de cycle et les révisions nécessaires de Cinquième et de Quatrième. Il s’adresse à l’élève qui veut construire ses raisonnements, à la famille qui l’accompagne et à l’enseignant qui cherche un parcours structuré.
\begin{pcEncadre}{method}{Quatre étapes pour apprendre}
\textbf{Comprendre.} Lire le cours et décrire les schémas à voix haute ; refaire l’exemple sans le regarder.\par
\textbf{Essayer.} Chercher un exercice avec le cours fermé et conserver les essais, même incomplets.\par
\textbf{Corriger.} Comparer méthode, conversions et conclusion, puis écrire la cause de l’erreur.\par
\textbf{Reprendre.} Refaire une question quelques jours plus tard, puis composer un sujet blanc en temps limité.
\end{pcEncadre}
\pcSousTitre{Un parcours progressif}
\otDots{1}\quad\textbf{Je démarre} : reconnaître, nommer et appliquer directement.\par
\otDots{2}\quad\textbf{Je m’entraîne} : calculer, représenter et relier les notions.\par
\otDots{3}\quad\textbf{Je raisonne} : choisir, croiser les données et argumenter.
\pcSousTitre{Ce que contient le manuel}
18 cours illustrés, 18 investigations et 18 défis ; 270 exercices avec réponses détaillées ; quinze sujets originaux de physique-chimie avec documents, points et remédiation, dont cinq dossiers supplémentaires au format des annales. Les liens internes permettent de passer de l’énoncé à la solution. Les fiches finales réunissent méthodes, unités, tests et concordance du programme.
\vfill
{\small OptiTAB · Octobre 2026. Manuel indépendant du ministère. Cours et exercices originaux ; certains parcours de révision reprennent des activités de la collection OptiTAB. Aucune homologation ni aucun numéro ISBN fictif n’est revendiqué.}
\clearpage\otPageTitre{Sommaire}{Le programme, puis les entraînements au Brevet}
\begingroup\setlength{\parskip}{0pt}
\pcSomPartie{1}{Organisation et transformations de la matière}
\pcSomLigne{01}{Matière : mesurer, séparer, identifier}{chap-1}
\pcSomLigne{02}{Atomes, noyaux et ions}{chap-2}
\pcSomLigne{03}{Transformations chimiques et équations}{chap-3}
\pcSomLigne{04}{Solutions, ions et pH}{chap-4}
\pcSomLigne{05}{Métaux : réactions et corrosion}{chap-5}
\pcSomLigne{06}{Combustions et conservation de la masse}{chap-6}
\pcSomLigne{07}{La matière dans l’Univers}{chap-7}
\pcSomPartie{2}{Mouvements et interactions}
\pcSomLigne{08}{Décrire et mesurer un mouvement}{chap-8}
\pcSomLigne{09}{Interactions, forces et poids}{chap-9}
\pcSomLigne{10}{Exploiter la gravitation universelle}{chap-10}
\pcSomPartie{3}{L’énergie, ses transferts et ses conversions}
\pcSomLigne{11}{Énergie cinétique et bilan mécanique}{chap-11}
\pcSomLigne{12}{L’énergie : transferts et conversions}{chap-12}
\pcSomLigne{13}{Circuits : intensités et tensions}{chap-13}
\pcSomLigne{14}{Loi d’Ohm et puissance électrique}{chap-14}
\pcSomLigne{15}{Puissance et énergie électrique}{chap-15}
\pcSomPartie{4}{Des signaux pour observer et communiquer}
\pcSomLigne{16}{Lumière : propagation et distances}{chap-16}
\pcSomLigne{17}{Le son : fréquence, propagation et écho}{chap-17}
\pcSomLigne{18}{Signaux et information}{chap-18}
\vspace{2mm}
\pcSomLigne{C}{Corrigés des exercices, investigations et défis}{corriges-debut}
\pcSomLigne{B}{Dix brevets blancs · Énoncés et documents}{brevets}
\pcSomLigne{S}{Corrigés, barèmes et reprise des sujets}{cor-brevets}
\pcSomLigne{A}{Cinq sujets supplémentaires au format des annales}{annales-originales}
\pcSomLigne{C}{Corrigés pédagogiques des sujets 11 à 15}{cor-annales-originales}
\pcSomLigne{M}{Méthodes, formules, unités et lexique}{methodes}
\pcSomLigne{P}{Concordance du programme et références}{programme}
\endgroup
\clearpage\otPageTitre{Le cadre de cette édition}{France · Cycle 4 · Année scolaire 2026-2027}
Le programme en application en physique-chimie est celui du cycle 4 publié au BO du 30 juillet 2020. Il fixe un ensemble pour les trois années du collège ; l’organisation annuelle de ce manuel est un choix éditorial. Les notions déjà rencontrées sont reprises pour préparer l’ensemble du cycle.
\begin{pcEncadre}{definition}{Le Brevet 2027 : le volet sciences}
L’épreuve de sciences dure une heure et porte sur deux disciplines retenues parmi physique-chimie, SVT et technologie. La durée indicative est de 30 minutes par discipline. La note totale est sur 20, soit 10 par discipline, dont 1 point pour la maîtrise de la langue. Les disciplines retenues sont annoncées deux mois avant l’épreuve ; le présent livre ne présume pas ce choix.
\end{pcEncadre}
\pcSousTitre{Des brevets blancs au bon périmètre}
Les quinze sujets sont des \textbf{volets de physique-chimie}, chacun prévu pour 30 minutes et noté sur 10. Ils constituent une préparation à cette discipline ; l’autre volet de l’épreuve de sciences n’est pas traité ici. Les exercices et documents sont créés pour le manuel et ne portent pas le statut de sujets ministériels.
\pcSousTitre{Rédiger fait partie des apprentissages}
Nommer les grandeurs, écrire les relations, convertir avant de calculer, puis conclure dans le contexte. Une question ouverte peut admettre plusieurs démarches cohérentes. La précision du vocabulaire et l’organisation d’une réponse sont travaillées dans les cours et dans chaque sujet blanc.
\pcRepere{programme}{Concordance et textes de référence, vérifiés le 24 octobre 2026}
\clearpage\otPageTitre{Expérimenter et raisonner}{Un cadre pour les investigations en classe}
\begin{pcEncadre}{attention}{Un protocole validé par le professeur}
Les activités de chimie, chauffage, électricité et tests de gaz nécessitent du matériel adapté et un encadrement. Les données et documents peuvent être étudiés sans reproduire une manipulation à domicile.
\end{pcEncadre}
\begin{questions}
\item Lire le protocole, identifier le système étudié et les risques, puis faire vérifier le matériel et le montage.
\item Ne pas goûter ni sentir directement un produit. Employer les protections prescrites, prélever de petites quantités et gérer les déchets selon la consigne.
\item Ne pas mélanger des produits ménagers. Ne pas fermer un dégagement gazeux dans un récipient rigide improvisé. Les tests à la flamme restent encadrés.
\item Couper l’alimentation avant de modifier un circuit ; utiliser seulement l’alimentation scolaire autorisée, sans manipuler le secteur domestique.
\item Ne pas regarder le Soleil ou une source intense ; ne pas augmenter le niveau sonore pour tester son audition.
\end{questions}
\begin{pcEncadre}{method}{La trace d’une démarche scientifique}
Question $\rightarrow$ hypothèse $\rightarrow$ protocole avec témoin $\rightarrow$ résultats et unités $\rightarrow$ interprétation $\rightarrow$ conclusion limitée aux conditions. Distinguer ce qui est observé, calculé et supposé.
\end{pcEncadre}
\pcSousTitre{Pourquoi répéter une mesure ?}
Des répétitions permettent de repérer la dispersion ; elles ne corrigent pas automatiquement un appareil décalé ou une méthode inadaptée. Annoncer la précision utile, comparer les résultats et signaler les limites avant de généraliser.
\clearpage\otPageTitre{Diagnostic de départ}{Dix questions pour choisir ses premières révisions}
Répondre sans le cours. Le diagnostic sert à repérer les besoins ; il ne remplace pas une évaluation de toutes les compétences.\par
\begin{questions}
\item Convertir $0{,}25\unite L$ en mL et $350\unite g$ en kg.
\item Un liquide transparent est-il forcément un corps pur ? Justifier.
\item Compter les atomes H et O dans $2\,\ce{H2O}$.
\item Un atome neutre a Z = 11. Combien possède-t-il d’électrons ?
\item Classer une solution de pH 4 vers 25 °C.
\item Un vélo parcourt 120 m en 30 s. Calculer sa vitesse moyenne.
\item Nommer l’unité du poids et la distinguer de celle de la masse.
\item Indiquer le branchement d’un ampèremètre et d’un voltmètre.
\item Convertir 20 mA en A et 15 min en secondes.
\item Un écho a fait un aller-retour : pourquoi diviser la distance totale par deux ?
\end{questions}
\pcRepere{cor-diagnostic}{Vérifier les réponses et choisir un chapitre}
\begin{pcEncadre}{method}{Un objectif réaliste}
Choisir deux priorités, refaire quelques exercices de niveau 1 et revenir au diagnostic après la correction. Une erreur de conversion appelle un entraînement différent d’une confusion entre force et énergie.
\end{pcEncadre}
\clearpage\otPageTitre{Organiser sa préparation}{Un rythme régulier avant les sujets blancs}
\pcSousTitre{Pendant l’année}
Après chaque cours : refaire l’exemple, chercher cinq exercices puis reprendre les erreurs. Chaque semaine : un court calcul, une lecture de graphique et une réponse rédigée. Conserver une page « mes erreurs fréquentes » avec un exemple corrigé.
\pcSousTitre{Un parcours de six semaines adaptable}
\begin{questions}
\item Matière et transformations : chapitres 1 à 6 ; sujet 2 puis reprise des questions manquées.
\item Atomes, ions et métaux : chapitres 2, 4 et 5 ; sujets 5 et 7.
\item Mouvements, forces et énergie : chapitres 8 à 11 ; sujet 1 et questions mécaniques du sujet 4.
\item Électricité et bilans : chapitres 12 à 15 ; sujets 3 et 8.
\item Univers et signaux : chapitres 7 et 16 à 18 ; sujets 6 et 9.
\item Synthèse : sujets 4 et 10 sans aide ; refaire deux questions mal réussies et une réponse ouverte.
\end{questions}
\begin{pcEncadre}{definition}{Avant de rendre la copie}
Ai-je traité toutes les questions accessibles ? Mes unités sont-elles cohérentes ? Les axes et les flèches ont-ils une légende ? La conclusion répond-elle à la question ? Ai-je écrit une démarche même si le calcul n’est pas terminé ?
\end{pcEncadre}
\pcRepere{brevets}{Accéder aux dix sujets blancs}
\clearpage\pagestyle{otMain}
\begin{ouverture}[college1]{1}{Matière : mesurer, séparer, identifier}{Partie 1 · Organisation et transformations de la matière}{Comment distinguer deux liquides transparents ?}
\objectif{Décrire états, mélanges et changements d’état.}
\objectif{Exploiter une masse volumique et des mesures.}
\objectif{Choisir une méthode de séparation et contrôler sa conclusion.}
\prerequis{Masse en g ou kg ; volume en mL, L ou m³.}
\prerequis{Proportionnalité et division.}
\end{ouverture}
\pcPage{Le cours}{Matière : mesurer, séparer, identifier}
\phantomsection\label{ess-1}
\essSection{1}{États et modèle microscopique}
Un solide a une forme et un volume propres ; un liquide a un volume propre et prend la forme du récipient ; un gaz occupe tout le volume disponible. Dans le modèle, les entités sont proches dans un solide ou un liquide, très espacées dans un gaz. Une fusion, une solidification, une vaporisation ou une liquéfaction change leur organisation, sans former une nouvelle espèce. La sublimation et la condensation solide relient directement solide et gaz.\par
\pcFigure{figures/b01-etats.pdf}{Même espèce ; organisation différente. Les particules ne sont pas à l’échelle.}
\essSection{2}{Mesurer un changement d’état}
À pression fixée, un corps pur présente une température de changement d’état déterminée. Sur une courbe, un palier peut correspondre à la coexistence de deux états, alors que de l’énergie est transférée. Dans un système fermé sans perte de matière, la masse se conserve ; le volume peut changer. Un palier seul ne prouve pas toute la composition d’un échantillon : il faut connaître le protocole et sa précision.\par
\clearpage
\essSection{3}{Masse volumique}
La masse volumique est $\rho=m/V$, soit $m=\rho V$ et $V=m/\rho$. En unités SI : m en kg, V en m³, $\rho$ en kg/m³. On peut aussi employer g et mL, ou g et cm³, si les unités sont cohérentes. $1\unite{mL}=1\unite{cm^3}$ et $1\unite{g/mL}=1\,000\unite{kg/m^3}$. La température, les bulles et une tare incorrecte influencent la mesure.\par
\pcFigure{figures/b01-rho.pdf}{Pour un matériau homogène à température fixée, masse et volume sont proportionnels.}
\essSection{4}{Corps purs, mélanges et séparation}
Un corps pur contient une seule espèce chimique ; un mélange en contient plusieurs. Homogène signifie qu’on ne distingue pas plusieurs phases à l’œil nu. Deux liquides miscibles forment une phase. Une filtration retient des particules solides, mais pas le sel dissous ; une distillation permet de recueillir un liquide vaporisé puis liquéfié. La solubilité est une limite de dissolution à température donnée. Une eau limpide n’est pas nécessairement potable.\par
\begin{pcEncadre}{method}{Un exemple expliqué}Un solide de $156\unite g$ déplace l’eau de 40 à $60\unite{mL}$. Son volume est $20\unite{cm^3}$ ; $\rho=156/20=7{,}8\unite{g/cm^3}$. C’est compatible avec le fer, sans identification unique. Il faut une immersion complète, sans dissolution ni bulle.\end{pcEncadre}
\pcPage{Expérimenter et retenir}{Matière : mesurer, séparer, identifier}
\phantomsection\label{activite-1}
\begin{pcEncadre}{method}{Investigation en classe}En classe, mesurer la masse de trois volumes d’un même liquide avec tare. Tracer m en fonction de V au tableur, comparer les quotients et discuter les écarts. Concevoir ensuite un protocole distinct pour un solide insoluble. Indiquer ce qui est mesuré directement, calculé et supposé.\end{pcEncadre}
\pcSousTitre{Ce que je dois savoir expliquer}\begin{questions}
\item Distinguer espèce, état et mélange.
\item Expliquer conservation de masse et variation possible du volume.
\item Choisir des unités cohérentes pour $
ho=m/V$.
\end{questions}
\begin{pcEncadre}{attention}{Une erreur à éviter}La masse volumique n’est pas la masse : doubler le volume d’un matériau homogène double sa masse, sans doubler sa masse volumique.\end{pcEncadre}
\pcRepere{cor-activite-1}{Pistes pour l’investigation}\pcGrille
\clearpage\markright{Je démarre}
\rubrique{Je démarre}
\par\noindent\begin{minipage}{\linewidth}
\begin{exercice}[1]{Trois états}
\begin{questions}
\item Nommer l’état qui conserve sa forme.
\item Quel état remplit tout un récipient fermé ?
\end{questions}
\end{exercice}
\end{minipage}
\par\noindent\begin{minipage}{\linewidth}
\begin{exercice}[1]{Une eau salée}
\begin{questions}
\item Corps pur ou mélange ?
\item Homogène signifie-t-il pur ?
\end{questions}
\end{exercice}
\end{minipage}
\par\noindent\begin{minipage}{\linewidth}
\begin{exercice}[1]{Les bons mots}
\begin{questions}
\item Nommer le passage liquide-gaz.
\item Nommer le passage gaz-liquide.
\end{questions}
\end{exercice}
\end{minipage}
\par\noindent\begin{minipage}{\linewidth}
\begin{exercice}[1]{Volumes}
\begin{questions}
\item Convertir $0{,}35\unite L$ en mL.
\item Convertir $18\unite{cm^3}$ en mL.
\end{questions}
\end{exercice}
\end{minipage}
\par\noindent\begin{minipage}{\linewidth}
\begin{exercice}[1]{Le quotient}
\begin{questions}
\item Calculer $\rho$ pour 40 g et 50 mL.
\item Donner le résultat en kg/m³.
\end{questions}
\end{exercice}
\end{minipage}
\clearpage\markright{Je m’entraîne}
\rubrique{Je m’entraîne}
\par\noindent\begin{minipage}{\linewidth}
\begin{exercice}[2]{Solide immergé}
\pcFigureEx{figures/rev-ex-immersion.pdf}{}
\begin{questions}
\item Une éprouvette passe de 28 à 43 mL. Trouver le volume du solide.
\item Sa masse est 40,5 g. Calculer sa masse volumique.
\end{questions}
\end{exercice}
\end{minipage}
\par\noindent\begin{minipage}{\linewidth}
\begin{exercice}[2]{Une masse prévue}
\begin{questions}
\item Une huile a $\rho=0{,}90\unite{g/mL}$. Quelle masse pour 250 mL ?
\item La même huile chauffée occupe plus de volume à masse conservée. Comment évolue $\rho$ ?
\end{questions}
\end{exercice}
\end{minipage}
\par\noindent\begin{minipage}{\linewidth}
\begin{exercice}[2]{Retrouver un volume}
\begin{questions}
\item Un morceau de cuivre de 178 g a $\rho=8{,}9\unite{g/cm^3}$. Trouver V.
\item Pourquoi une pièce creuse ne suit-elle pas le même calcul avec son volume extérieur ?
\end{questions}
\end{exercice}
\end{minipage}
\par\noindent\begin{minipage}{\linewidth}
\begin{exercice}[2]{Une limite de dissolution}
\begin{questions}
\item À température fixée, 100 mL d’eau dissolvent au plus 36 g de sel. Que devient un ajout de 40 g ?
\item Estimer la masse maximale pour 250 mL d’eau dans les mêmes conditions.
\end{questions}
\end{exercice}
\end{minipage}
\par\noindent\begin{minipage}{\linewidth}
\begin{exercice}[2]{Interpréter un palier}
\pcFigureEx{figures/rev-ex-ebullition.pdf}{Courbe illustrant le palier de trois minutes décrit dans la question.}
\begin{questions}
\item Pendant trois minutes, une eau pure bout vers 100 °C à pression usuelle. Pourquoi la température ne monte-t-elle pas malgré le chauffage ?
\item Sa masse baisse dans le bécher ouvert. La conservation est-elle violée ?
\end{questions}
\end{exercice}
\end{minipage}
\clearpage\markright{Je raisonne}
\rubrique{Je raisonne}
\par\noindent\begin{minipage}{\linewidth}
\begin{exercice}[3]{Choisir une séparation}
\begin{questions}
\item Un mélange contient eau, sel dissous et sable. Proposer deux opérations pour recueillir de l’eau séparée.
\item La première opération suffit-elle à prouver la potabilité ?
\end{questions}
\end{exercice}
\end{minipage}
\par\noindent\begin{minipage}{\linewidth}
\begin{exercice}[3]{Une mesure discutable}
\begin{questions}
\item Un bloc de 54 g paraît déplacer 25 mL, alors que son volume réel est 20 mL. Quel $\rho$ calcule-t-on ?
\item Une bulle accrochée peut-elle expliquer le sens de l’erreur ?
\end{questions}
\end{exercice}
\end{minipage}
\par\noindent\begin{minipage}{\linewidth}
\begin{exercice}[3]{Identifier avec une plage}
\begin{questions}
\item Une mesure donne $2{,}68\unite{g/cm^3}$ avec une incertitude de lecture estimée à $0{,}05\unite{g/cm^3}$. Comparer à 2,70 pour l’aluminium.
\item Peut-on conclure « aluminium pur certain » ?
\end{questions}
\end{exercice}
\end{minipage}
\par\noindent\begin{minipage}{\linewidth}
\begin{exercice}[3]{Glace et océans}
\begin{questions}
\item Expliquer pourquoi la fonte d’une glace continentale peut augmenter le niveau marin.
\item Citer un autre mécanisme physique lié au réchauffement.
\end{questions}
\end{exercice}
\end{minipage}
\par\noindent\begin{minipage}{\linewidth}
\begin{exercice}[3]{Construire une preuve}
\begin{questions}
\item Deux liquides ont le même aspect. Proposer une comparaison quantitative sans les goûter.
\item Pourquoi conserver température et protocole comparables ?
\end{questions}
\end{exercice}
\end{minipage}
\pcPage{Défi du chapitre}{Matière : mesurer, séparer, identifier}
\phantomsection\label{defi-ex-1}
\pcDefiTitre{1}{Le flacon inconnu}
\begin{questions}
\item Un flacon vide pèse 45 g ; avec 80 mL de liquide, 109 g. Déterminer $\rho$.
\item Références : eau 1,00 ; huile 0,90 ; éthanol 0,79 g/mL. Quelle hypothèse est la plus compatible ?
\item Proposer une limite et une amélioration du protocole.
\end{questions}
\pcGrille

\begin{ouverture}[college1]{2}{Atomes, noyaux et ions}{Partie 1 · Organisation et transformations de la matière}{Comment une entité peut-elle devenir chargée ?}
\objectif{Décrire les constituants de l’atome.}
\objectif{Lire Z et un inventaire de nucléons.}
\objectif{Interpréter une formule moléculaire et une charge ionique.}
\prerequis{Symboles chimiques et nombres entiers.}
\prerequis{Charges positives et négatives.}
\end{ouverture}
\pcPage{Le cours}{Atomes, noyaux et ions}
\phantomsection\label{ess-2}
\essSection{1}{Un noyau et des électrons}
Un atome comporte un noyau central et un nuage d’électrons. Le noyau contient des protons, de charge positive, et généralement des neutrons, sans charge électrique. Les électrons portent une charge négative opposée à celle du proton. Un atome neutre possède autant de protons que d’électrons. Sa masse est presque entièrement concentrée dans le noyau. Le dessin de collège ne représente pas des orbites réelles.\par
\pcFigure{figures/b02-atome.pdf}{Modèle schématique : noyau très agrandi ; les électrons sont figurés par des points.}
\essSection{2}{Le numéro atomique}
Le numéro atomique Z est le nombre de protons et identifie l’élément. Si A est le nombre total de nucléons, le nombre de neutrons vaut $A-Z$. Cette notation est exploitée lorsqu’elle est donnée. Pour $Z=6$ et $A=12$, le noyau contient six protons et six neutrons ; l’atome neutre possède six électrons. Un tableau périodique permet de retrouver nom, symbole et Z.\par
\clearpage
\essSection{3}{Former un ion}
Un ion est une entité chargée ayant perdu ou gagné un ou plusieurs électrons. Une perte d’électrons donne une charge positive ; un gain donne une charge négative. Lors d’une ionisation chimique usuelle, le noyau conserve son nombre de protons : l’élément ne change pas. $\ce{Na+}$ a perdu un électron ; $\ce{Cl-}$ en a gagné un ; $\ce{Cu^{2+}}$ en a perdu deux.\par
\pcFigure{figures/b02-ions.pdf}{Les nombres d’électrons indiqués concernent l’entité entière ; aucun proton n’est déplacé.}
\essSection{4}{Molécules, indices et coefficients}
Une molécule est un assemblage neutre d’atomes liés. $\ce{H2O}$ contient deux H et un O ; $\ce{CO2}$ un C et deux O ; $\ce{CH4}$ un C et quatre H ; $\ce{N2O}$ deux N et un O. Dans $3\,\ce{CO2}$, le coefficient multiplie toute la formule : trois C et six O. Un indice change l’espèce ; un coefficient change sa quantité d’entités. Un ion peut aussi être polyatomique.\par
\begin{pcEncadre}{method}{Un exemple expliqué}Le noyau d’un ion aluminium contient 13 protons et 14 neutrons ; l’entité possède 10 électrons. $Z=13$, $A=27$. La charge totale correspond à $13-10=3$ charges positives : $\ce{Al^{3+}}$. L’atome neutre d’aluminium possède 13 électrons.\end{pcEncadre}
\pcPage{Expérimenter et retenir}{Atomes, noyaux et ions}
\phantomsection\label{activite-2}
\begin{pcEncadre}{method}{Investigation en classe}À partir de cartes « noyau » et de jetons électrons, construire un atome neutre, puis un ion du même élément. Justifier chaque inventaire. Utiliser ensuite un tableau périodique pour vérifier les symboles. Séparer explicitement les limites du modèle graphique et les conclusions sur les charges.\end{pcEncadre}
\pcSousTitre{Ce que je dois savoir expliquer}\begin{questions}
\item Définir Z et distinguer proton, neutron et électron.
\item Retrouver les électrons d’un ion simple avec Z et sa charge.
\item Compter les atomes sans confondre indice et coefficient.
\end{questions}
\begin{pcEncadre}{attention}{Une erreur à éviter}Un ion positif n’a pas gagné des protons : dans les transformations chimiques étudiées, ce sont les électrons qui peuvent être échangés.\end{pcEncadre}
\pcRepere{cor-activite-2}{Pistes pour l’investigation}\pcGrille
\clearpage\markright{Je démarre}
\rubrique{Je démarre}
\par\noindent\begin{minipage}{\linewidth}
\begin{exercice}[1]{Constituants}
\begin{questions}
\item Où se trouvent protons et neutrons ?
\item Quelle particule porte une charge négative ?
\end{questions}
\end{exercice}
\end{minipage}
\par\noindent\begin{minipage}{\linewidth}
\begin{exercice}[1]{Neutralité}
\begin{questions}
\item Un atome possède 8 protons. Combien d’électrons ?
\item Quel est son numéro atomique ?
\end{questions}
\end{exercice}
\end{minipage}
\par\noindent\begin{minipage}{\linewidth}
\begin{exercice}[1]{Un noyau}
\begin{questions}
\item Un noyau contient 11 protons et 12 neutrons. Donner Z.
\item Donner A.
\end{questions}
\end{exercice}
\end{minipage}
\par\noindent\begin{minipage}{\linewidth}
\begin{exercice}[1]{L’eau}
\begin{questions}
\item Compter H et O dans une molécule d’eau.
\item Compter ces atomes dans trois molécules.
\end{questions}
\end{exercice}
\end{minipage}
\par\noindent\begin{minipage}{\linewidth}
\begin{exercice}[1]{Le signe}
\begin{questions}
\item L’ion $\ce{Cl-}$ a-t-il gagné ou perdu un électron ?
\item Même question pour $\ce{Na+}$.
\end{questions}
\end{exercice}
\end{minipage}
\clearpage\markright{Je m’entraîne}
\rubrique{Je m’entraîne}
\par\noindent\begin{minipage}{\linewidth}
\begin{exercice}[2]{Un ion cuivre}
\begin{questions}
\item Cu a $Z=29$. Combien d’électrons dans $\ce{Cu^{2+}}$ ?
\item Combien de protons dans cet ion ?
\end{questions}
\end{exercice}
\end{minipage}
\par\noindent\begin{minipage}{\linewidth}
\begin{exercice}[2]{Un ion chlorure}
\begin{questions}
\item Cl a $Z=17$. Compter les électrons de $\ce{Cl-}$.
\item Cet ion est-il une molécule de dichlore $\ce{Cl2}$ ?
\end{questions}
\end{exercice}
\end{minipage}
\par\noindent\begin{minipage}{\linewidth}
\begin{exercice}[2]{Inventaire complet}
\begin{questions}
\item Un atome a Z = 13 et A = 27. Compter neutrons et électrons.
\item Il perd trois électrons. Donner la charge et le symbole, sachant qu’il s’agit d’aluminium.
\end{questions}
\end{exercice}
\end{minipage}
\par\noindent\begin{minipage}{\linewidth}
\begin{exercice}[2]{Indices et coefficient}
\begin{questions}
\item Compter les atomes dans $2\,\ce{CH4}$.
\item Compter les atomes dans $3\,\ce{N2O}$.
\end{questions}
\end{exercice}
\end{minipage}
\par\noindent\begin{minipage}{\linewidth}
\begin{exercice}[2]{Trois écritures}
\begin{questions}
\item Distinguer O, $\ce{O2}$ et $\ce{O^{2-}}$.
\item Les trois appartiennent-ils au même élément ?
\end{questions}
\end{exercice}
\end{minipage}
\clearpage\markright{Je raisonne}
\rubrique{Je raisonne}
\par\noindent\begin{minipage}{\linewidth}
\begin{exercice}[3]{Retrouver un élément}
\begin{questions}
\item Une entité possède 12 protons et 10 électrons. Trouver sa charge.
\item Le tableau donne Mg pour Z = 12. Écrire l’ion.
\end{questions}
\end{exercice}
\end{minipage}
\par\noindent\begin{minipage}{\linewidth}
\begin{exercice}[3]{Une fausse justification}
\begin{questions}
\item « Le sodium devient chlorure parce qu’il gagne un électron. » Corriger cette phrase.
\item Quel nombre devrait changer pour changer d’élément ?
\end{questions}
\end{exercice}
\end{minipage}
\par\noindent\begin{minipage}{\linewidth}
\begin{exercice}[3]{Un changement d’échelle}
\begin{questions}
\item On prend un diamètre atomique de $10^{-10}\unite m$ et un diamètre nucléaire de $10^{-15}\unite m$. Calculer le rapport.
\item Si l’atome est représenté par 10 m, quel diamètre aurait le noyau ?
\end{questions}
\end{exercice}
\end{minipage}
\par\noindent\begin{minipage}{\linewidth}
\begin{exercice}[3]{Une solution électriquement neutre}
\begin{questions}
\item On représente deux ions $\ce{Na+}$ et deux $\ce{Cl-}$. Quelle charge totale ?
\item Avec deux $\ce{Ca^{2+}}$, combien de $\ce{Cl-}$ faut-il pour compenser ?
\end{questions}
\end{exercice}
\end{minipage}
\par\noindent\begin{minipage}{\linewidth}
\begin{exercice}[3]{Deux noyaux du même élément}
\begin{questions}
\item Deux noyaux ont Z = 6, mais A = 12 et A = 14. Comparer les nombres de neutrons.
\item Pourquoi sont-ils du même élément malgré cette différence ?
\end{questions}
\end{exercice}
\end{minipage}
\pcPage{Défi du chapitre}{Atomes, noyaux et ions}
\phantomsection\label{defi-ex-2}
\pcDefiTitre{2}{L’entité inconnue}
\begin{questions}
\item Un noyau contient 26 protons et 30 neutrons ; l’entité possède 23 électrons. Donner Z et A.
\item Le tableau donne Fe pour Z = 26. Écrire le symbole de l’ion.
\item Expliquer en trois phrases pourquoi il ne s’agit ni d’un nouvel élément ni d’une molécule.
\end{questions}
\pcGrille

\begin{ouverture}[college1]{3}{Transformations chimiques et équations}{Partie 1 · Organisation et transformations de la matière}{Comment rendre compte de ce qui disparaît et apparaît ?}
\objectif{Identifier réactifs et produits par des observations et des tests.}
\objectif{Exploiter une équation fournie et compter les atomes.}
\objectif{Utiliser la conservation de la masse avec un système défini.}
\prerequis{Formules moléculaires et coefficients.}
\prerequis{Différence entre physique et chimique.}
\end{ouverture}
\pcPage{Le cours}{Transformations chimiques et équations}
\phantomsection\label{ess-3}
\essSection{1}{Reconnaître une transformation}
Des réactifs sont consommés et des produits se forment. Une couleur, des bulles, un dépôt ou un changement de température sont des indices à interpréter : des bulles peuvent aussi venir d’une ébullition ou d’un gaz dissous. Un test caractéristique, avec réactif et résultat de référence, apporte une information sur l’espèce. Il faut comparer l’état initial et l’état final.\par
\pcFigure{figures/b03-tests.pdf}{Banque de tests : ils sont effectués sous la responsabilité du professeur.}
\essSection{2}{Lire l’équation fournie}
Dans $\ce{CH4 + 2 O2 -> CO2 + 2 H2O}$, le méthane et le dioxygène sont les réactifs ; l’eau et le dioxyde de carbone les produits. Le modèle traduit une redistribution des atomes : un C, quatre H et quatre O sont présents de chaque côté. Les coefficients donnent des proportions d’entités ; ils ne donnent pas directement des masses en grammes.\par
\clearpage
\essSection{3}{Ce qui se conserve}
Au cours d’une réaction chimique, les atomes de chaque élément et la masse totale se conservent dans un bilan matériel complet. Les espèces, le nombre de molécules et les volumes ne sont généralement pas conservés. Dans un système ouvert, une pesée partielle peut varier parce que de la matière entre ou sort. Toujours nommer ce que l’on pèse avant de calculer.\par
\pcFigure{figures/b03-bilan.pdf}{Même masse totale pour un système fermé ; les molécules ont changé.}
\essSection{4}{Construire une conclusion}
Une conclusion relie une donnée au modèle : « l’eau de chaux se trouble ; le test fourni indique du dioxyde de carbone ». Une absence d’effet ne prouve pas toujours une absence totale d’espèce : quantité, sensibilité du test et protocole interviennent. Au Brevet, une démarche partielle pertinente doit être explicitée, avec relation, conversions et unités.\par
\begin{pcEncadre}{method}{Un exemple expliqué}Dans un système fermé, 4 g de dihydrogène et 32 g de dioxygène réagissent entièrement. L’équation fournie est $\ce{2 H2 + O2 -> 2 H2O}$. La masse d’eau formée est $4+32=36\unite g$. On ne multiplie pas 4 g par le coefficient 2 : les masses données concernent déjà les réactifs utilisés.\end{pcEncadre}
\pcPage{Expérimenter et retenir}{Transformations chimiques et équations}
\phantomsection\label{activite-3}
\begin{pcEncadre}{method}{Investigation en classe}Analyser des observations de dissolution, fusion, dégazage et réaction carbonate-acide, avec témoins. Choisir une banque de tests et justifier les conclusions. Le professeur réalise les manipulations autorisées ; prévoir un dispositif adapté au gaz dégagé et ne jamais fermer un récipient rigide improvisé.\end{pcEncadre}
\pcSousTitre{Ce que je dois savoir expliquer}\begin{questions}
\item Définir réactif et produit.
\item Contrôler chaque élément d’une équation fournie.
\item Expliquer une variation de masse apparente en système ouvert.
\end{questions}
\begin{pcEncadre}{attention}{Une erreur à éviter}Une équation se lit avec ses formules : modifier un indice pour faire le compte revient à changer une espèce chimique.\end{pcEncadre}
\pcRepere{cor-activite-3}{Pistes pour l’investigation}\pcGrille
\clearpage\markright{Je démarre}
\rubrique{Je démarre}
\par\noindent\begin{minipage}{\linewidth}
\begin{exercice}[1]{Fondre ou réagir}
\begin{questions}
\item La glace fond. De nouvelles espèces apparaissent-elles ?
\item Le fer rouille. Pourquoi parle-t-on de transformation chimique ?
\end{questions}
\end{exercice}
\end{minipage}
\par\noindent\begin{minipage}{\linewidth}
\begin{exercice}[1]{Avant et après}
\begin{questions}
\item Dans $\ce{C + O2 -> CO2}$, nommer les réactifs.
\item Nommer le produit.
\end{questions}
\end{exercice}
\end{minipage}
\par\noindent\begin{minipage}{\linewidth}
\begin{exercice}[1]{Un test fourni}
\begin{questions}
\item L’eau de chaux se trouble au contact d’un gaz. Quelle espèce est mise en évidence selon le test usuel ?
\item Des bulles seules suffisent-elles ?
\end{questions}
\end{exercice}
\end{minipage}
\par\noindent\begin{minipage}{\linewidth}
\begin{exercice}[1]{Compter}
\begin{questions}
\item Dans $\ce{2 H2 + O2 -> 2 H2O}$, compter H avant et après.
\item Compter O avant et après.
\end{questions}
\end{exercice}
\end{minipage}
\par\noindent\begin{minipage}{\linewidth}
\begin{exercice}[1]{Masse totale}
\begin{questions}
\item Deux réactifs de 12 g et 32 g sont intégralement transformés sans échange de matière. Quelle masse de produits ?
\item Le nombre de molécules doit-il être conservé ?
\end{questions}
\end{exercice}
\end{minipage}
\clearpage\markright{Je m’entraîne}
\rubrique{Je m’entraîne}
\par\noindent\begin{minipage}{\linewidth}
\begin{exercice}[2]{Un gaz s’échappe}
\begin{questions}
\item Un ensemble ouvert pèse 150 g avant et 147,8 g après une réaction. Quelle masse nette a quitté le système si rien d’autre n’entre ni ne sort ?
\item La matière a-t-elle été détruite ?
\end{questions}
\end{exercice}
\end{minipage}
\par\noindent\begin{minipage}{\linewidth}
\begin{exercice}[2]{Une masse inconnue}
\begin{questions}
\item 6 g de carbone réagissent entièrement avec du dioxygène pour former 22 g de $\ce{CO2}$. Trouver la masse de dioxygène consommée.
\item Pourquoi préciser « entièrement » ?
\end{questions}
\end{exercice}
\end{minipage}
\par\noindent\begin{minipage}{\linewidth}
\begin{exercice}[2]{Le bon inventaire}
\begin{questions}
\item Vérifier C, H et O dans $\ce{C2H6O + 3 O2 -> 2 CO2 + 3 H2O}$.
\item Les coefficients 1, 3, 2, 3 sont-ils des rapports de masses ?
\end{questions}
\end{exercice}
\end{minipage}
\par\noindent\begin{minipage}{\linewidth}
\begin{exercice}[2]{Une équation impossible}
\begin{questions}
\item Pourquoi $\ce{H2 + O2 -> H2O}$ ne convient-elle pas comme équation conservant les atomes ?
\item Comparer à $\ce{2 H2 + O2 -> 2 H2O}$.
\end{questions}
\end{exercice}
\end{minipage}
\par\noindent\begin{minipage}{\linewidth}
\begin{exercice}[2]{Des réactifs en excès}
\begin{questions}
\item Un système fermé contient initialement 10 g et 25 g de réactifs. Après, il reste 5 g d’un réactif. Quelle masse totale de produits si aucun autre constituant n’intervient ?
\item Pourquoi ne pas annoncer 35 g de produits ?
\end{questions}
\end{exercice}
\end{minipage}
\clearpage\markright{Je raisonne}
\rubrique{Je raisonne}
\par\noindent\begin{minipage}{\linewidth}
\begin{exercice}[3]{Comparer des contrôles}
\begin{questions}
\item Un réactif de test change de couleur même sans échantillon. Que montre le témoin ?
\item Peut-on alors attribuer la couleur à l’espèce recherchée ?
\end{questions}
\end{exercice}
\end{minipage}
\par\noindent\begin{minipage}{\linewidth}
\begin{exercice}[3]{Mélange et réaction}
\begin{questions}
\item Dissoudre du sucre sans le transformer est-il une réaction ?
\item Pourquoi « mélanger » n’exclut-il pas qu’une réaction se produise dans une autre situation ?
\end{questions}
\end{exercice}
\end{minipage}
\par\noindent\begin{minipage}{\linewidth}
\begin{exercice}[3]{Un solide gagne de la masse}
\begin{questions}
\item Une laine de fer chauffée dans le dioxygène gagne 0,8 g. Quelle matière supplémentaire doit figurer dans le bilan ?
\item Si le fer initial pesait 2,1 g et réagit entièrement, quelle masse finale obtient-on ?
\end{questions}
\end{exercice}
\end{minipage}
\par\noindent\begin{minipage}{\linewidth}
\begin{exercice}[3]{Un protocole de gaz}
\begin{questions}
\item Proposer comment distinguer un gaz dégagé par ébullition et un produit de réaction, sans se fier aux bulles seules.
\item Citer une limite d’une conclusion négative de test.
\end{questions}
\end{exercice}
\end{minipage}
\par\noindent\begin{minipage}{\linewidth}
\begin{exercice}[3]{Prise d’initiative}
\begin{questions}
\item Une réaction donne 18 g d’eau et 44 g d’un gaz. Les seuls réactifs consommés pèsent au total 62 g. Peut-on identifier ce gaz avec les seules masses ?
\item Quel renseignement complémentaire demander ?
\end{questions}
\end{exercice}
\end{minipage}
\pcPage{Défi du chapitre}{Transformations chimiques et équations}
\phantomsection\label{defi-ex-3}
\pcDefiTitre{3}{Le bilan complet}
\begin{questions}
\item Une enceinte contient 8 g de méthane et 40 g de dioxygène. Après réaction, il reste 8 g de dioxygène et il s’est formé 22 g de $\ce{CO2}$. Trouver la masse d’eau formée.
\item Vérifier le bilan matériel.
\item Expliquer pourquoi on ne doit pas compter l’enceinte comme réactif.
\end{questions}
\pcGrille

\begin{ouverture}[college1]{4}{Solutions, ions et pH}{Partie 1 · Organisation et transformations de la matière}{Que signifie vraiment une solution acide ?}
\objectif{Relier pH, acidité et ions H+ ou OH−.}
\objectif{Distinguer dissolution, dilution et réaction.}
\objectif{Exploiter une banque de tests ioniques.}
\prerequis{Ion et charge.}
\prerequis{Soluté, solvant et solution.}
\end{ouverture}
\pcPage{Le cours}{Solutions, ions et pH}
\phantomsection\label{ess-4}
\essSection{1}{Des ions dans l’eau}
Une solution aqueuse contient de l’eau et des espèces dissoutes, molécules ou ions. Dans le modèle d’une solution de chlorure de sodium, les ions $\ce{Na+}$ et $\ce{Cl-}$ sont dispersés ; le sel n’a pas disparu. Une solution ionique peut conduire le courant grâce au déplacement des ions. Le courant dans un métal est lié à des électrons : les porteurs dépendent du milieu.\par
\pcFigure{figures/b04-solution.pdf}{Dissolution d’un solide ionique ; nombres d’ions dessinés choisis pour la lisibilité.}
\essSection{2}{Interpréter le pH}
Pour les solutions aqueuses usuelles vers 25 °C, pH inférieur à 7 : acide ; égal à 7 : neutre ; supérieur à 7 : basique. Les ions $\ce{H+}$ prédominent sur $\ce{OH-}$ dans une solution acide ; c’est l’inverse dans une solution basique. Les deux types d’ions peuvent être présents. Le pH est sans unité. Le papier pH donne une estimation ; le pH-mètre doit être réglé et utilisé selon son protocole.\par
\clearpage
\essSection{3}{Diluer ou faire réagir}
Diluer consiste à ajouter du solvant sans ajouter de soluté : une solution acide usuelle devient moins acide, son pH se rapproche de 7 ; une solution basique devient moins basique. Le pH n’est pas proportionnel à la quantité de solvant. Une réaction acide-base peut consommer les ions $\ce{H+}$ et $\ce{OH-}$ : $\ce{H+ + OH- -> H2O}$. La solution finale dépend des quantités introduites ; elle n’est pas automatiquement neutre.\par
\pcFigure{figures/b04-ph.pdf}{Échelle indicative pour des solutions aqueuses usuelles ; aucun goût ni contact direct.}
\essSection{4}{Utiliser des tests fournis}
Une banque de tests peut indiquer : avec nitrate d’argent, précipité blanc pour les ions chlorure ; avec soude, bleu pour $\ce{Cu^{2+}}$, vert pour $\ce{Fe^{2+}}$ ou brun pour $\ce{Fe^{3+}}$. On utilise des prélèvements séparés pour éviter de modifier tout le flacon. Le résultat s’interprète dans les conditions du test. Produits, protections et déchets sont gérés selon les consignes du professeur.\par
\begin{pcEncadre}{method}{Un exemple expliqué}Un liquide a pH 3. Il est acide ; dans le modèle, les ions H+ sont plus abondants que les ions OH−. Après dilution à l’eau, on mesure pH 4 : l’acidité diminue, mais le liquide reste acide. On ne déduit pas le nouveau pH par une moyenne des valeurs.\end{pcEncadre}
\pcPage{Expérimenter et retenir}{Solutions, ions et pH}
\phantomsection\label{activite-4}
\begin{pcEncadre}{method}{Investigation en classe}Le professeur fournit des solutions diluées autorisées et une banque de tests. Sur des prélèvements distincts, comparer papier pH et mesure instrumentale, puis identifier un ion par précipitation. Prévoir un témoin, noter la couleur observée et séparer observation, identification et limite du test.\end{pcEncadre}
\pcSousTitre{Ce que je dois savoir expliquer}\begin{questions}
\item Classer une solution avec son pH et citer les ions concernés.
\item Expliquer l’effet qualitatif d’une dilution.
\item Décrire un test avec réactif, observation et conclusion.
\end{questions}
\begin{pcEncadre}{attention}{Une erreur à éviter}Un pH de 4 n’est pas deux fois moins acide qu’un pH de 2. Au collège, on compare et on exploite les informations fournies sans traiter le pH comme une grandeur proportionnelle.\end{pcEncadre}
\pcRepere{cor-activite-4}{Pistes pour l’investigation}\pcGrille
\clearpage\markright{Je démarre}
\rubrique{Je démarre}
\par\noindent\begin{minipage}{\linewidth}
\begin{exercice}[1]{Classer}
\begin{questions}
\item Classer pH 2 et pH 11.
\item Que signifie pH 7 vers 25 °C ?
\end{questions}
\end{exercice}
\end{minipage}
\par\noindent\begin{minipage}{\linewidth}
\begin{exercice}[1]{L’unité}
\begin{questions}
\item Quelle est l’unité du pH ?
\item Quelle grandeur mesure-t-on avec une éprouvette ?
\end{questions}
\end{exercice}
\end{minipage}
\par\noindent\begin{minipage}{\linewidth}
\begin{exercice}[1]{Ions et acidité}
\begin{questions}
\item Quels ions prédominent dans une solution acide : $\ce{H+}$ ou $\ce{OH-}$ ?
\item Même question pour une solution basique.
\end{questions}
\end{exercice}
\end{minipage}
\par\noindent\begin{minipage}{\linewidth}
\begin{exercice}[1]{Du sel dissous}
\begin{questions}
\item Dans une solution de chlorure de sodium, citer les deux ions.
\item Pourquoi la filtration ordinaire ne les sépare-t-elle pas de l’eau ?
\end{questions}
\end{exercice}
\end{minipage}
\par\noindent\begin{minipage}{\linewidth}
\begin{exercice}[1]{Un test bleu}
\begin{questions}
\item La banque fournie associe soude + ions cuivre(II) à un précipité bleu. Que conclure si ce précipité apparaît ?
\item Faut-il verser le réactif dans le flacon d’origine ?
\end{questions}
\end{exercice}
\end{minipage}
\clearpage\markright{Je m’entraîne}
\rubrique{Je m’entraîne}
\par\noindent\begin{minipage}{\linewidth}
\begin{exercice}[2]{Une dilution}
\begin{questions}
\item Un acide usuel passe de pH 2 à pH 4 par ajout d’eau. Est-il encore acide ?
\item Son acidité a-t-elle augmenté ?
\end{questions}
\end{exercice}
\end{minipage}
\par\noindent\begin{minipage}{\linewidth}
\begin{exercice}[2]{Une base diluée}
\begin{questions}
\item Une base passe de pH 12 à pH 10. Comment évolue son caractère basique ?
\item Peut-on obtenir pH 5 par simple dilution idéale de cette base avec de l’eau neutre ?
\end{questions}
\end{exercice}
\end{minipage}
\par\noindent\begin{minipage}{\linewidth}
\begin{exercice}[2]{Trois inconnues}
\begin{questions}
\item A : pH 4 ; B : pH 7 ; C : pH 9. Classer.
\item Une solution contient des ions : peut-on déduire son pH sans les connaître ?
\end{questions}
\end{exercice}
\end{minipage}
\par\noindent\begin{minipage}{\linewidth}
\begin{exercice}[2]{Un résultat vert}
\begin{questions}
\item Le test à la soude donne un précipité vert, associé dans la banque à Fe2+. Nommer l’ion.
\item Combien d’électrons a-t-il perdus par rapport à l’atome ?
\end{questions}
\end{exercice}
\end{minipage}
\par\noindent\begin{minipage}{\linewidth}
\begin{exercice}[2]{Neutraliser}
\begin{questions}
\item Lire $\ce{H+ + OH- -> H2O}$ : quels ions sont consommés ?
\item Pourquoi un mélange acide-base ne finit-il pas toujours à pH 7 ?
\end{questions}
\end{exercice}
\end{minipage}
\clearpage\markright{Je raisonne}
\rubrique{Je raisonne}
\par\noindent\begin{minipage}{\linewidth}
\begin{exercice}[3]{Échantillons comparables}
\begin{questions}
\item Deux bandelettes indiquent environ 5 et environ 6. Peut-on annoncer une différence précise de 1,00 unité de pH ?
\item Proposer une amélioration.
\end{questions}
\end{exercice}
\end{minipage}
\par\noindent\begin{minipage}{\linewidth}
\begin{exercice}[3]{Conduction}
\begin{questions}
\item Comparer les porteurs du courant dans un fil métallique et dans une solution ionique.
\item La solution peut-elle rester globalement neutre tout en conduisant ?
\end{questions}
\end{exercice}
\end{minipage}
\par\noindent\begin{minipage}{\linewidth}
\begin{exercice}[3]{Sans moyenne}
\begin{questions}
\item Un élève annonce pH 7 en mélangeant des volumes égaux de solutions pH 2 et pH 12. Pourquoi ce calcul est-il insuffisant ?
\item Quelle donnée ou action manque ?
\end{questions}
\end{exercice}
\end{minipage}
\par\noindent\begin{minipage}{\linewidth}
\begin{exercice}[3]{L’eau et le climat}
\begin{questions}
\item Une eau passe de pH 8,2 à pH 8,0. Est-elle devenue acide ?
\item Peut-on parler d’acidification ?
\end{questions}
\end{exercice}
\end{minipage}
\par\noindent\begin{minipage}{\linewidth}
\begin{exercice}[3]{Concevoir une analyse}
\begin{questions}
\item Une solution peut contenir des ions chlorure et cuivre(II). Décrire les deux prélèvements et les tests nécessaires d’après la banque du cours.
\item Pourquoi un seul test ne suffit-il pas à conclure à la présence des deux ions ?
\end{questions}
\end{exercice}
\end{minipage}
\pcPage{Défi du chapitre}{Solutions, ions et pH}
\phantomsection\label{defi-ex-4}
\pcDefiTitre{4}{Le laboratoire enquête}
\begin{questions}
\item A : pH 3, test chlorure positif ; B : pH 10, test cuivre négatif. Classer A et B.
\item Un camarade veut mélanger A et B pour identifier tous les ions d’un coup. Pourquoi refuser ce protocole ?
\item Rédiger une conclusion de trois phrases limitée aux observations.
\end{questions}
\pcGrille

\begin{ouverture}[college1]{5}{Métaux : réactions et corrosion}{Partie 1 · Organisation et transformations de la matière}{Pourquoi certains métaux s’altèrent-ils ?}
\objectif{Exploiter une réaction acide-métal avec un modèle fourni.}
\objectif{Identifier gaz et ions à l’aide d’une banque.}
\objectif{Interpréter un protocole de corrosion et une protection.}
\prerequis{Équations fournies, ions et pH.}
\prerequis{Conservation de la masse.}
\end{ouverture}
\pcPage{Le cours}{Métaux : réactions et corrosion}
\phantomsection\label{ess-5}
\essSection{1}{Un acide peut réagir avec un métal}
Certains métaux réagissent avec une solution acide : du métal et des ions H+ sont consommés, des ions métalliques et du dihydrogène peuvent apparaître. Exemple fourni : $\ce{Zn + 2 H+ -> Zn^{2+} + H2}$. Les atomes et la charge totale sont conservés. Tous les métaux ne réagissent pas de la même manière avec un même acide ; on ne généralise pas un seul essai.\par
\pcFigure{figures/b05-reaction.pdf}{Modèle zinc-acide : les ions spectateurs ne figurent pas dans l’équation simplifiée.}
\essSection{2}{Interpréter les tests}
Un test du dihydrogène peut produire une petite détonation au contact d’une flamme, dans un dispositif pédagogique sécurisé par le professeur. Les ions formés peuvent être identifiés sur un autre prélèvement avec une banque de précipitation. La diminution des ions H+ peut s’accompagner d’une hausse du pH ; elle ne signifie pas forcément que la solution finale est neutre.\par
\clearpage
\essSection{3}{Comprendre la corrosion}
La corrosion est une transformation chimique d’un métal avec son environnement. La rouille du fer est un ensemble de composés ; sa formation habituelle nécessite notamment eau et dioxygène. Comparer des expériences en ne faisant varier qu’un facteur permet de rechercher les conditions. Une eau salée peut accélérer le phénomène ; il faut des mesures comparables pour quantifier cet effet.\par
\pcFigure{figures/b05-corrosion.pdf}{Expériences comparées : même métal, même durée ; résultats qualitatifs annoncés.}
\essSection{4}{Protéger et recycler}
Une peinture ou un revêtement isolant limite le contact du fer avec l’eau et le dioxygène tant que la barrière reste efficace. Certains métaux forment une couche superficielle protectrice. La protection dépend du matériau et du milieu. Recycler réduit certains prélèvements de ressources, mais demande tri et énergie. Ne pas confondre altération chimique, fusion du métal et simple fragmentation.\par
\begin{pcEncadre}{method}{Un exemple expliqué}Le fer réagit selon $\ce{Fe + 2 H+ -> Fe^{2+} + H2}$. On observe des bulles, puis le test du gaz indique H2 ; un test à la soude donne un précipité vert associé à Fe2+. Les deux résultats sont cohérents avec la réaction fournie. Les bulles seules n’auraient pas permis cette conclusion.\end{pcEncadre}
\pcPage{Expérimenter et retenir}{Métaux : réactions et corrosion}
\phantomsection\label{activite-5}
\begin{pcEncadre}{method}{Investigation en classe}Exploiter des photographies datées de clous identiques : air sec, eau au contact de l’air, eau salée au contact de l’air, eau désaérée protégée de l’air. Construire un tableau d’observations, choisir les comparaisons pertinentes et proposer une expérience pour tester un revêtement sans changer tous les paramètres.\end{pcEncadre}
\pcSousTitre{Ce que je dois savoir expliquer}\begin{questions}
\item Lire une équation acide-métal et repérer les espèces spectatrices.
\item Distinguer observation et identification du gaz.
\item Relier une protection aux conditions de corrosion.
\end{questions}
\begin{pcEncadre}{attention}{Une erreur à éviter}La rouille n’est pas « du fer sale ». C’est le résultat d’une transformation produisant d’autres espèces ; sa masse peut inclure des éléments venus de l’environnement.\end{pcEncadre}
\pcRepere{cor-activite-5}{Pistes pour l’investigation}\pcGrille
\clearpage\markright{Je démarre}
\rubrique{Je démarre}
\par\noindent\begin{minipage}{\linewidth}
\begin{exercice}[1]{Corrosion}
\begin{questions}
\item Le fer qui rouille subit-il une transformation physique ou chimique ?
\item Citer deux facteurs usuels nécessaires dans le modèle étudié.
\end{questions}
\end{exercice}
\end{minipage}
\par\noindent\begin{minipage}{\linewidth}
\begin{exercice}[1]{Un gaz}
\begin{questions}
\item Quel gaz est formé dans $\ce{Zn + 2 H+ -> Zn^{2+} + H2}$ ?
\item Quel ion métallique apparaît ?
\end{questions}
\end{exercice}
\end{minipage}
\par\noindent\begin{minipage}{\linewidth}
\begin{exercice}[1]{Les réactifs}
\begin{questions}
\item Dans $\ce{Fe + 2 H+ -> Fe^{2+} + H2}$, quelles espèces sont consommées ?
\item L’ion chlorure non écrit peut-il être spectateur ?
\end{questions}
\end{exercice}
\end{minipage}
\par\noindent\begin{minipage}{\linewidth}
\begin{exercice}[1]{Protection}
\begin{questions}
\item Quel rôle joue une peinture sur le fer ?
\item Pourquoi une rayure peut-elle réduire son efficacité ?
\end{questions}
\end{exercice}
\end{minipage}
\par\noindent\begin{minipage}{\linewidth}
\begin{exercice}[1]{Le bon test}
\begin{questions}
\item Une banque associe le précipité vert à Fe2+. Quel ion identifie-t-on ?
\item Pourquoi éviter de réaliser soi-même un test à la flamme à domicile ?
\end{questions}
\end{exercice}
\end{minipage}
\clearpage\markright{Je m’entraîne}
\rubrique{Je m’entraîne}
\par\noindent\begin{minipage}{\linewidth}
\begin{exercice}[2]{Conserver la charge}
\begin{questions}
\item Comparer la charge totale des deux côtés de $\ce{Zn + 2 H+ -> Zn^{2+} + H2}$.
\item Compter les atomes H de chaque côté.
\end{questions}
\end{exercice}
\end{minipage}
\par\noindent\begin{minipage}{\linewidth}
\begin{exercice}[2]{Le pH évolue}
\begin{questions}
\item Une réaction acide-fer fait passer le pH de 2 à 4. Que dire de l’acidité ?
\item Est-ce cohérent avec une consommation d’ions H+ ?
\end{questions}
\end{exercice}
\end{minipage}
\par\noindent\begin{minipage}{\linewidth}
\begin{exercice}[2]{Une pesée ouverte}
\begin{questions}
\item Un dispositif perd 0,06 g pendant une réaction produisant uniquement un gaz qui s’échappe. Quelle masse de gaz s’est échappée si aucune autre perte ne se produit ?
\item Pourquoi faut-il écarter les projections ?
\end{questions}
\end{exercice}
\end{minipage}
\par\noindent\begin{minipage}{\linewidth}
\begin{exercice}[2]{Choisir un couple}
\begin{questions}
\item A : air sec sans rouille ; B : eau + air avec rouille. Que suggère la comparaison ?
\item Quel essai supplémentaire testerait le rôle du dioxygène ?
\end{questions}
\end{exercice}
\end{minipage}
\par\noindent\begin{minipage}{\linewidth}
\begin{exercice}[2]{Ne pas tout changer}
\begin{questions}
\item Deux clous diffèrent par taille, métal, salinité et durée d’exposition. Peut-on isoler l’effet du sel ?
\item Décrire une meilleure comparaison.
\end{questions}
\end{exercice}
\end{minipage}
\clearpage\markright{Je raisonne}
\rubrique{Je raisonne}
\par\noindent\begin{minipage}{\linewidth}
\begin{exercice}[3]{Un bilan de masse}
\begin{questions}
\item Un métal de 5,6 g incorpore 2,4 g de matière venant du milieu. Quelle masse de produit si tout est conservé ?
\item Pourquoi cette augmentation ne viole-t-elle pas la conservation ?
\end{questions}
\end{exercice}
\end{minipage}
\par\noindent\begin{minipage}{\linewidth}
\begin{exercice}[3]{Une conclusion trop large}
\begin{questions}
\item Le cuivre ne produit pas de bulles dans une solution acide donnée pendant un essai. Peut-on conclure que « les acides ne réagissent jamais avec les métaux » ?
\item Formuler une conclusion limitée.
\end{questions}
\end{exercice}
\end{minipage}
\par\noindent\begin{minipage}{\linewidth}
\begin{exercice}[3]{Comparer des résultats}
\begin{questions}
\item Deux séries donnent des masses de fer corrodé de 0,20 g et 0,24 g après sept jours. L’écart seul prouve-t-il un effet du sel ?
\item Que faudrait-il pour rendre l’argument plus solide ?
\end{questions}
\end{exercice}
\end{minipage}
\par\noindent\begin{minipage}{\linewidth}
\begin{exercice}[3]{Choisir un modèle}
\begin{questions}
\item L’équation $\ce{Fe + 2 H+ -> Fe^{2+} + H2}$ décrit-elle à elle seule la rouille ordinaire dans l’air humide ?
\item Pourquoi distinguer les deux situations ?
\end{questions}
\end{exercice}
\end{minipage}
\par\noindent\begin{minipage}{\linewidth}
\begin{exercice}[3]{Réparer plutôt que jeter}
\begin{questions}
\item Proposer deux actions réduisant la consommation de métal pour un objet corrodé.
\item Pourquoi le recyclage ne dispense-t-il pas de toute consommation d’énergie ?
\end{questions}
\end{exercice}
\end{minipage}
\pcPage{Défi du chapitre}{Métaux : réactions et corrosion}
\phantomsection\label{defi-ex-5}
\pcDefiTitre{5}{Le portail près de la mer}
\begin{questions}
\item Un portail en fer rouille surtout aux rayures. Expliquer le rôle du revêtement.
\item Proposer une expérience isolant l’effet de la salinité.
\item Rédiger une recommandation justifiée et une limite de votre étude.
\end{questions}
\pcGrille

\begin{ouverture}[college1]{6}{Combustions et conservation de la masse}{Partie 1 · Organisation et transformations de la matière}{Brûler, est-ce faire disparaître la matière ?}
\objectif{Identifier combustible, comburant et énergie d’activation.}
\objectif{Lire un bilan de combustion et vérifier les atomes.}
\objectif{Raisonner sur un système ouvert ou fermé et sur les émissions.}
\prerequis{Réactifs, produits et formules.}
\end{ouverture}
\pcPage{Le cours}{Combustions et conservation de la masse}
\phantomsection\label{ess-6}
\essSection{1}{Les conditions d’une combustion}
Une combustion est une transformation chimique qui transfère de l’énergie. Elle nécessite un combustible, un comburant et une énergie d’activation. Dans les exemples usuels, le comburant est le dioxygène. Retirer une condition peut empêcher la combustion, mais les gestes d’extinction dépendent du contexte : on suit les consignes et on ne réalise pas d’essai de feu à la maison.\par
\pcFigure{figures/c07-triangle.pdf}{Le triangle du feu est un schéma de conditions, pas un protocole d’extinction à appliquer sans formation.}
\essSection{2}{Lire un bilan simple}
Pour le carbone : $\ce{C + O2 -> CO2}$. Pour la combustion complète du méthane : $\ce{CH4 + 2O2 -> CO2 + 2H2O}$. Ces bilans supposent les conditions de combustion complète annoncées. Les coefficients permettent de vérifier les atomes C, H et O. Ils ne donnent pas directement des masses égales de chaque réactif.\par
\clearpage
\essSection{3}{Une masse totale conservée}
Dans un système fermé qui ne perd pas de matière, la masse totale avant réaction est égale à celle après réaction, à la précision de la mesure. Dans un système ouvert, un gaz peut entrer ou sortir : la seule partie posée sur la balance ne constitue pas toujours le bilan complet. Une combustion peut consommer le dioxygène de l’air et libérer des produits gazeux invisibles.\par
\pcFigure{figures/c07-bilan.pdf}{Le carbone, le dioxygène consommé et le dioxyde de carbone formé appartiennent au bilan complet.}
\essSection{4}{Risques et climat}
Une combustion incomplète d’un combustible carboné peut produire du monoxyde de carbone $\ce{CO}$ et de la suie. Le monoxyde de carbone est très toxique, incolore et inodore : l’odorat ne le détecte pas. Le dioxyde de carbone $\ce{CO2}$ est un gaz à effet de serre ; ses émissions d’origine fossile contribuent au réchauffement. Réduire un gaspillage énergétique diminue des besoins, mais comparer des solutions exige un périmètre cohérent et des données.\par
\begin{pcEncadre}{method}{Un exemple expliqué}Dans un bilan idéal, $3\unite g$ de carbone réagissent avec $8\unite g$ de dioxygène et forment $11\unite g$ de dioxyde de carbone. La somme $3+8=11$ respecte la conservation. Si un dispositif ouvert semble perdre de la masse, on doit inclure les gaz qui entrent et sortent avant d’interpréter.\end{pcEncadre}
\pcPage{Expérimenter et retenir}{Combustions et conservation de la masse}
\phantomsection\label{activite-6}
\begin{pcEncadre}{method}{Investigation en classe}\textbf{Investigation documentaire.} Le professeur fournit les masses d’un montage fermé adapté avant et après une transformation, puis celles d’un montage ouvert. Définir précisément ce qui est pesé, repérer les échanges de matière et construire un bilan. La combustion éventuelle est une démonstration du professeur avec matériel adapté, sans improviser un récipient chauffé fermé. \textbf{À rendre :} schéma du système, tableau des masses, comparaison à la précision de la balance, et phrase distinguant conservation des atomes et changement des molécules.\end{pcEncadre}
\pcSousTitre{Ce que je dois savoir expliquer}\begin{questions}
\item Combustible et comburant sont deux rôles différents.
\item La combustion complète du méthane forme dioxyde de carbone et eau.
\item Une combustion réorganise les atomes.
\item La masse totale se conserve lorsque tout le système matériel est pris en compte.
\item Monoxyde de carbone et dioxyde de carbone ont des propriétés différentes.
\end{questions}
\begin{pcEncadre}{attention}{Une erreur à éviter}« Le combustible a perdu $2\unite g$, donc seulement $2\unite g$ de produits se sont formés. » Une partie de la masse des produits provient du dioxygène consommé.\end{pcEncadre}
\pcRepere{cor-activite-6}{Pistes pour l’investigation}\pcGrille
\clearpage\markright{Je démarre}
\rubrique{Je démarre}
\par\noindent\begin{minipage}{\linewidth}
\begin{exercice}[1]{Deux rôles}
\begin{questions}
\item Dans une combustion de charbon dans l’air, nommer combustible et comburant.
\item L’air entier est-il consommé comme une seule espèce ?
\end{questions}
\end{exercice}
\end{minipage}
\par\noindent\begin{minipage}{\linewidth}
\begin{exercice}[1]{Les produits du méthane}
\begin{questions}
\item Quels produits figurent dans $\ce{CH4 + 2O2 -> CO2 + 2H2O}$ ?
\item La réaction transfère-t-elle de l’énergie ?
\end{questions}
\end{exercice}
\end{minipage}
\par\noindent\begin{minipage}{\linewidth}
\begin{exercice}[1]{Une somme de masses}
\begin{questions}
\item $2\unite g$ et $5\unite g$ de réactifs sont entièrement transformés sans perte. Quelle masse totale de produits ?
\item Pourquoi préciser « sans perte » ?
\end{questions}
\end{exercice}
\end{minipage}
\par\noindent\begin{minipage}{\linewidth}
\begin{exercice}[1]{Deux gaz à distinguer}
\begin{questions}
\item Comparer les formules de CO et de $\ce{CO2}$.
\item Lequel peut être détecté par l’odorat ?
\end{questions}
\end{exercice}
\end{minipage}
\par\noindent\begin{minipage}{\linewidth}
\begin{exercice}[1]{Une condition manquante}
\begin{questions}
\item Nommer les trois composantes du triangle du feu.
\item Peut-on déclencher volontairement une combustion pour essayer ?
\end{questions}
\end{exercice}
\end{minipage}
\clearpage\markright{Je m’entraîne}
\rubrique{Je m’entraîne}
\par\noindent\begin{minipage}{\linewidth}
\begin{exercice}[2]{Compter les atomes}
\begin{questions}
\item Dans le bilan du méthane fourni au cours, compter H et O avant et après.
\item Pourquoi les deux molécules d’eau sont-elles nécessaires ?
\end{questions}
\end{exercice}
\end{minipage}
\par\noindent\begin{minipage}{\linewidth}
\begin{exercice}[2]{Une masse manquante}
\begin{questions}
\item On forme $22\unite g$ de produit à partir de $6\unite g$ d’un réactif et d’un second entièrement consommé. Trouver sa masse.
\item Quelle hypothèse permet ce calcul ?
\end{questions}
\end{exercice}
\end{minipage}
\par\noindent\begin{minipage}{\linewidth}
\begin{exercice}[2]{Le produit est plus lourd}
\begin{questions}
\item $12\unite g$ de carbone forment $44\unite g$ de dioxyde de carbone dans un bilan complet. Quelle masse de dioxygène a réagi ?
\item La masse a-t-elle été créée ?
\end{questions}
\end{exercice}
\end{minipage}
\par\noindent\begin{minipage}{\linewidth}
\begin{exercice}[2]{Un montage ouvert}
\begin{questions}
\item Une bougie allumée devient plus légère. Cela réfute-t-il la conservation ?
\item Quels échanges faudrait-il considérer ?
\end{questions}
\end{exercice}
\end{minipage}
\par\noindent\begin{minipage}{\linewidth}
\begin{exercice}[2]{Un bilan proportionnel}
\begin{questions}
\item Le bilan fourni est $3\unite g$ de C et $8\unite g$ de dioxygène pour $11\unite g$ de produit. Pour $9\unite g$ de C, quelles masses correspondantes ?
\item Ce calcul donne-t-il la composition de l’air nécessaire ?
\end{questions}
\end{exercice}
\end{minipage}
\clearpage\markright{Je raisonne}
\rubrique{Je raisonne}
\par\noindent\begin{minipage}{\linewidth}
\begin{exercice}[3]{Une affirmation à corriger}
\begin{questions}
\item « La combustion incomplète ne produit que du dioxyde de carbone. » Corriger.
\item Pourquoi ne peut-on pas se fier à l’absence d’odeur ?
\end{questions}
\end{exercice}
\end{minipage}
\par\noindent\begin{minipage}{\linewidth}
\begin{exercice}[3]{Une mesure presque égale}
\begin{questions}
\item Le montage fermé pèse $245{,}3\unite g$ avant et $245{,}2\unite g$ après. La balance affiche au dixième de gramme. Peut-on annoncer une destruction de matière ?
\item Que vérifier ?
\end{questions}
\end{exercice}
\end{minipage}
\par\noindent\begin{minipage}{\linewidth}
\begin{exercice}[3]{Un réactif en excès}
\begin{questions}
\item On dispose de $6\unite g$ de C et $24\unite g$ de dioxygène. Le bilan fourni exige $8\unite g$ de dioxygène pour $3\unite g$ de C. Que reste-t-il ?
\item Quelle masse de dioxyde de carbone se forme ?
\end{questions}
\end{exercice}
\end{minipage}
\par\noindent\begin{minipage}{\linewidth}
\begin{exercice}[3]{Émissions comparées}
\begin{questions}
\item Une donnée fictive indique $200\unite g$ de CO2 par unité d’énergie pour A et $80\unite g$ pour B. Pour 5 unités identiques, calculer les émissions.
\item Peut-on conclure sur tous les impacts écologiques ?
\end{questions}
\end{exercice}
\end{minipage}
\par\noindent\begin{minipage}{\linewidth}
\begin{exercice}[3]{Conserver quoi ?}
\begin{questions}
\item Lors d’une réaction équilibrée, le nombre de molécules est-il toujours conservé ?
\item Quelle conservation faut-il vérifier ?
\end{questions}
\end{exercice}
\end{minipage}
\pcPage{Défi du chapitre}{Combustions et conservation de la masse}
\phantomsection\label{defi-ex-6}
\pcDefiTitre{6}{Le bilan sous contrôle}
\begin{questions}
\item On utilise $9\unite g$ de carbone et $20\unite g$ de dioxygène. La donnée est $3\unite g$ de C pour $8\unite g$ de dioxygène. Quel réactif limite la transformation complète annoncée ?
\item Calculer les masses de carbone consommé et restant, puis de produit.
\item Vérifier la masse totale et expliquer la limite du modèle.
\end{questions}
\pcGrille

\begin{ouverture}[college1]{7}{La matière dans l’Univers}{Partie 1 · Organisation et transformations de la matière}{De quoi les étoiles et la Terre sont-elles faites ?}
\objectif{Situer Terre, Système solaire et galaxie.}
\objectif{Comparer des échelles de distance et de temps.}
\objectif{Relier les éléments de la Terre à l’histoire de la matière.}
\prerequis{Proportionnalité et lecture des symboles chimiques.}
\end{ouverture}
\pcPage{Le cours}{La matière dans l’Univers}
\phantomsection\label{ess-7}
\essSection{1}{Des objets emboîtés}
La Terre est une planète du Système solaire ; le Soleil est son étoile. Le Système solaire appartient à la Voie lactée, une galaxie contenant un très grand nombre d’étoiles. L’Univers observable contient de nombreuses galaxies. Une planète ne produit pas sa lumière visible comme une étoile : elle peut diffuser la lumière reçue. Un satellite naturel tourne autour d’une planète dans le modèle étudié.\par
\pcFigure{figures/c08-systeme.pdf}{Ordre des huit planètes : les tailles et les espacements ne sont pas à l’échelle.}
\essSection{2}{Comparer des distances}
Pour le voisinage terrestre, on utilise souvent le kilomètre. L’unité astronomique, notée ua, vaut ici environ $150$ millions de kilomètres et correspond à l’échelle de la distance Terre-Soleil. L’année-lumière est une distance parcourue par la lumière en un an, environ $9{,}46\times10^{12}\unite{km}$. Ce n’est pas une durée. La vitesse de propagation sera étudiée au chapitre 16.\par
\clearpage
\essSection{3}{Une histoire de la matière}
Le Système solaire s’est formé il y a environ $4{,}6$ milliards d’années à partir de gaz et de poussières. L’Univers a évolué avant sa formation : les premiers éléments légers comprennent principalement hydrogène et hélium. Des étoiles et des événements stellaires ont ensuite produit et dispersé de nombreux éléments plus lourds. Carbone, oxygène, fer ou silicium ne sont donc pas propres à la Terre. Le modèle cosmologique situe l’âge de l’Univers vers 13,8 milliards d’années ; cet âge ne se confond pas avec celui de la Terre. Une frise doit distinguer ces échelles de temps.\par
\pcFigure{figures/c08-echelles.pdf}{De la planète à la galaxie : comparer des ordres de grandeur exige de préciser l’unité et le type de distance.}
\essSection{4}{Des ressources inégalement réparties}
La matière observable étudiée dans les étoiles et sur Terre obéit aux mêmes lois physiques. Les proportions d’éléments varient cependant selon les milieux. La présence d’un élément dans la croûte terrestre ne garantit pas un gisement facile à exploiter. Extraire, recycler et utiliser des métaux soulève des questions de ressources, d’énergie et d’environnement. Pour vérifier un document, identifier l’auteur, la date, l’unité et le sens du pourcentage.\par
\begin{pcEncadre}{method}{Un exemple expliqué}À l’échelle $1\unite{cm}$ pour $10$ millions de kilomètres, une distance de $150$ millions de kilomètres devient $15\unite{cm}$. Une distance de $750$ millions devient $75\unite{cm}$. Cette échelle des distances ne permet pas de dessiner simultanément des planètes lisibles à leur vraie taille relative : il faut annoncer une autre échelle pour leurs diamètres.\end{pcEncadre}
\pcPage{Expérimenter et retenir}{La matière dans l’Univers}
\phantomsection\label{activite-7}
\begin{pcEncadre}{method}{Investigation en classe}\textbf{Mission.} Construire une frise des distances moyennes au Soleil à partir d’un tableau fourni par le professeur. Choisir une échelle compatible avec la salle ; calculer les positions avant de tracer. Représenter les planètes par des symboles de taille fixe, avec la mention « diamètres non à l’échelle ». Comparer un document de vulgarisation à une source d’agence spatiale : distinguer données mesurées, valeurs arrondies et illustration artistique. Préparer une explication orale de la différence entre galaxie et système planétaire.\end{pcEncadre}
\pcSousTitre{Ce que je dois savoir expliquer}\begin{questions}
\item Le Soleil est une étoile, la Terre une planète et la Lune un satellite naturel.
\item Le Système solaire appartient à la Voie lactée.
\item Une année-lumière mesure une distance.
\item Les éléments terrestres s’inscrivent dans une histoire cosmique.
\item Un schéma astronomique doit préciser ses échelles.
\end{questions}
\begin{pcEncadre}{attention}{Une erreur à éviter}Une image montrant les planètes alignées ne signifie pas qu’elles restent alignées : c’est une convention de présentation de leur ordre.\end{pcEncadre}
\pcRepere{cor-activite-7}{Pistes pour l’investigation}\pcGrille
\clearpage\markright{Je démarre}
\rubrique{Je démarre}
\par\noindent\begin{minipage}{\linewidth}
\begin{exercice}[1]{Classer des objets}
\begin{questions}
\item Classer Soleil, Terre et Lune en étoile, planète ou satellite naturel.
\item Le Soleil appartient-il à la Voie lactée ?
\end{questions}
\end{exercice}
\end{minipage}
\par\noindent\begin{minipage}{\linewidth}
\begin{exercice}[1]{Le bon emboîtement}
\begin{questions}
\item Ranger : Système solaire, Terre, Univers observable, Voie lactée.
\item Une galaxie est-elle une planète géante ?
\end{questions}
\end{exercice}
\end{minipage}
\par\noindent\begin{minipage}{\linewidth}
\begin{exercice}[1]{Une unité trompeuse}
\begin{questions}
\item L’année-lumière est-elle une durée ou une distance ?
\item Pourquoi son nom contient-il « année » ?
\end{questions}
\end{exercice}
\end{minipage}
\par\noindent\begin{minipage}{\linewidth}
\begin{exercice}[1]{L’ordre des planètes}
\begin{questions}
\item Quelles planètes encadrent la Terre dans l’ordre des distances moyennes au Soleil ?
\item Nommer la plus éloignée des huit planètes.
\end{questions}
\end{exercice}
\end{minipage}
\par\noindent\begin{minipage}{\linewidth}
\begin{exercice}[1]{Lire une échelle}
\begin{questions}
\item À $1\unite{cm}$ pour $10$ millions de kilomètres, représenter $50$ millions de kilomètres.
\item Représenter $150$ millions de kilomètres.
\end{questions}
\end{exercice}
\end{minipage}
\clearpage\markright{Je m’entraîne}
\rubrique{Je m’entraîne}
\par\noindent\begin{minipage}{\linewidth}
\begin{exercice}[2]{Comparer des distances}
\begin{questions}
\item Deux distances valent $150$ et $600$ millions de kilomètres. Combien de fois la seconde est-elle plus grande ?
\item Quel écart les sépare ?
\end{questions}
\end{exercice}
\end{minipage}
\par\noindent\begin{minipage}{\linewidth}
\begin{exercice}[2]{Des milliards d’années}
\begin{questions}
\item Écrire $4{,}6$ milliards d’années en millions d’années.
\item Est-ce l’âge d’une personne ou une échelle géologique ?
\end{questions}
\end{exercice}
\end{minipage}
\par\noindent\begin{minipage}{\linewidth}
\begin{exercice}[2]{Les mêmes éléments}
\begin{questions}
\item On détecte de l’hydrogène dans une étoile. Cet hydrogène correspond-il à un élément propre à cette étoile ?
\item Pourquoi les proportions peuvent-elles différer de celles de la Terre ?
\end{questions}
\end{exercice}
\end{minipage}
\par\noindent\begin{minipage}{\linewidth}
\begin{exercice}[2]{Les limites d’une affiche}
\begin{questions}
\item Une affiche donne la même taille à toutes les planètes. Est-elle forcément inutile ?
\item Quelle légende doit-elle porter ?
\end{questions}
\end{exercice}
\end{minipage}
\par\noindent\begin{minipage}{\linewidth}
\begin{exercice}[2]{Passer à une autre échelle}
\begin{questions}
\item Une frise utilise $1\unite{cm}$ pour $25$ millions de kilomètres. Placer une planète à $750$ millions de kilomètres.
\item Quelle distance réelle représente $8\unite{cm}$ ?
\end{questions}
\end{exercice}
\end{minipage}
\clearpage\markright{Je raisonne}
\rubrique{Je raisonne}
\par\noindent\begin{minipage}{\linewidth}
\begin{exercice}[3]{Une preuve d’autorité ?}
\begin{questions}
\item Deux sites donnent des valeurs légèrement différentes pour une distance planétaire. Nommer deux raisons possibles autres qu’une erreur.
\item Quelles informations chercher ?
\end{questions}
\end{exercice}
\end{minipage}
\par\noindent\begin{minipage}{\linewidth}
\begin{exercice}[3]{Un métal « rare »}
\begin{questions}
\item Une fiche annonce qu’un élément est présent dans beaucoup de roches. Cela garantit-il une extraction facile ?
\item Que faut-il encore examiner ?
\end{questions}
\end{exercice}
\end{minipage}
\par\noindent\begin{minipage}{\linewidth}
\begin{exercice}[3]{Une autre unité}
\begin{questions}
\item On prend $1\unite{ua}=150$ millions de kilomètres. Convertir $450$ millions de kilomètres en ua.
\item Convertir $5\unite{ua}$ en millions de kilomètres.
\end{questions}
\end{exercice}
\end{minipage}
\par\noindent\begin{minipage}{\linewidth}
\begin{exercice}[3]{Un alignement dessiné}
\begin{questions}
\item Un schéma aligne les huit planètes. Peut-on l’utiliser pour prédire leur position demain ?
\item Que représente-t-il utilement ?
\end{questions}
\end{exercice}
\end{minipage}
\par\noindent\begin{minipage}{\linewidth}
\begin{exercice}[3]{Rédiger une histoire}
\begin{questions}
\item Pourquoi trouve-t-on du carbone à la fois sur Terre et dans certaines étoiles ?
\item Pourquoi ne faut-il pas dire que tout le carbone a été créé sur Terre ?
\end{questions}
\end{exercice}
\end{minipage}
\pcPage{Défi du chapitre}{La matière dans l’Univers}
\phantomsection\label{defi-ex-7}
\pcDefiTitre{7}{La frise impossible ?}
\begin{questions}
\item La distance moyenne maximale retenue est $4\,500$ millions de kilomètres. La salle offre $9\unite m$. Choisir une échelle simple utilisant toute la longueur.
\item Placer la Terre à $150$ millions de kilomètres.
\item Pourquoi ne peut-on pas exiger aussi des diamètres lisibles à la même échelle ?
\end{questions}
\pcGrille

\begin{ouverture}[college2]{8}{Décrire et mesurer un mouvement}{Partie 2 · Mouvements et interactions}{Un passager assis est-il immobile ?}
\objectif{Choisir un objet et un référentiel.}
\objectif{Décrire trajectoire, direction, sens et évolution de la vitesse.}
\objectif{Calculer distance, durée et vitesse avec des unités cohérentes.}
\prerequis{Durées, distances et proportionnalité.}
\end{ouverture}
\pcPage{Le cours}{Décrire et mesurer un mouvement}
\phantomsection\label{ess-8}
\essSection{1}{Le mouvement dépend du référentiel}
Pour décrire un mouvement, on choisit l’objet étudié et un référentiel, c’est-à-dire un objet de référence auquel on associe des repères d’espace et de temps. Un passager assis est immobile par rapport à son siège, mais en mouvement par rapport au quai lorsque le train roule. Ces descriptions sont compatibles parce qu’elles utilisent des références différentes.\par
\pcFigure{figures/c09-chrono.pdf}{Des positions relevées à intervalles de temps égaux : l’écartement renseigne sur l’évolution de la valeur de la vitesse.}
\essSection{2}{Décrire la trajectoire et la vitesse}
La trajectoire est l’ensemble des positions successives du point étudié dans le référentiel choisi. Elle est rectiligne, circulaire ou plus généralement curviligne. La vitesse possède une direction, un sens et une valeur. Un mouvement circulaire à valeur de vitesse constante change néanmoins de direction. « Uniforme » signifie que la valeur de la vitesse reste constante ; cela ne signifie pas forcément rectiligne.\par
\clearpage
\essSection{3}{Relier distance, durée et vitesse}
La vitesse moyenne sur un trajet est $v_{\mathrm{moy}}=d/\Delta t$. Pour un mouvement uniforme, on utilise $d=v\Delta t$ et $\Delta t=d/v$. Si $d$ est en m et $\Delta t$ en s, $v$ est en m/s. Pour passer de m/s à km/h, multiplier par $3{,}6$ ; pour le passage inverse, diviser par $3{,}6$. Une vitesse moyenne ne décrit pas tous les instants du trajet.\par
\pcFigure{figures/c09-graphique.pdf}{Sur ce graphique de position rectiligne, une pente plus forte correspond à une valeur de vitesse plus élevée.}
\essSection{4}{Mesurer sans inventer une précision}
Une chronophotographie ou une vidéo fournit des positions et des instants. Pour exploiter l’image, il faut une échelle de distance, une durée entre images et un point suivi cohérent. La perspective peut déformer les longueurs. Une mesure avec un chronomètre manuel comporte une variabilité due au déclenchement. Répéter et comparer aide à estimer la qualité du résultat.\par
\begin{pcEncadre}{method}{Un exemple expliqué}Une cycliste parcourt $180\unite m$ en $30\unite s$. Sa vitesse moyenne vaut $180/30=6{,}0\unite{m/s}$, soit $6{,}0\times3{,}6=21{,}6\unite{km/h}$. Si elle conserve cette vitesse pendant $50\unite s$, elle parcourt $d=6{,}0\times50=300\unite m$. La dernière conclusion suppose bien un mouvement uniforme sur la durée annoncée.\end{pcEncadre}
\pcPage{Expérimenter et retenir}{Décrire et mesurer un mouvement}
\phantomsection\label{activite-8}
\begin{pcEncadre}{method}{Investigation en classe}\textbf{Question.} Un petit chariot avance-t-il uniformément ? Filmer un déplacement rectiligne sur une table, avec un repère de longueur dans le plan du mouvement. Choisir le centre d’une même roue ou un point repéré du chariot. Relever les positions toutes les mêmes durées, puis tracer position en fonction du temps. Comparer les distances successives et discuter les erreurs de lecture. \textbf{À rendre :} référentiel, tableau, graphique et conclusion. Aucune mesure ne se fait sur une route ouverte.\end{pcEncadre}
\pcSousTitre{Ce que je dois savoir expliquer}\begin{questions}
\item Un mouvement est décrit dans un référentiel annoncé.
\item Trajectoire et évolution de la vitesse sont deux descriptions complémentaires.
\item À temps égaux, des déplacements égaux indiquent une valeur de vitesse constante.
\item La vitesse moyenne vaut distance totale divisée par durée totale.
\item Les unités doivent être converties avant le calcul.
\end{questions}
\begin{pcEncadre}{attention}{Une erreur à éviter}Faire la moyenne de deux vitesses ne donne pas toujours la vitesse moyenne du trajet : il faut utiliser la distance totale et la durée totale.\end{pcEncadre}
\pcRepere{cor-activite-8}{Pistes pour l’investigation}\pcGrille
\clearpage\markright{Je démarre}
\rubrique{Je démarre}
\par\noindent\begin{minipage}{\linewidth}
\begin{exercice}[1]{Le référentiel}
\begin{questions}
\item Un livre posé dans un train est-il immobile par rapport au train ?
\item Peut-il être en mouvement par rapport au sol ?
\end{questions}
\end{exercice}
\end{minipage}
\par\noindent\begin{minipage}{\linewidth}
\begin{exercice}[1]{La trajectoire}
\begin{questions}
\item Un point décrit un cercle autour d’un axe. Comment qualifier la trajectoire ?
\item Un mouvement en ligne droite est-il circulaire ?
\end{questions}
\end{exercice}
\end{minipage}
\par\noindent\begin{minipage}{\linewidth}
\begin{exercice}[1]{Temps égaux}
\pcFigureEx{figures/ex-chronophoto.pdf}{Les prises de vue sont séparées par une même durée.}
\begin{questions}
\item Sur une chronophotographie, les écarts augmentent pour des durées égales. Que peut-on dire de la valeur de la vitesse ?
\item Quelle donnée faut-il vérifier ?
\end{questions}
\end{exercice}
\end{minipage}
\par\noindent\begin{minipage}{\linewidth}
\begin{exercice}[1]{Une vitesse directe}
\begin{questions}
\item Un mobile parcourt $40\unite m$ en $8\unite s$. Calculer sa vitesse moyenne.
\item Donner l’unité.
\end{questions}
\end{exercice}
\end{minipage}
\par\noindent\begin{minipage}{\linewidth}
\begin{exercice}[1]{Convertir}
\begin{questions}
\item Convertir $10\unite{m/s}$ en km/h.
\item Convertir $72\unite{km/h}$ en m/s.
\end{questions}
\end{exercice}
\end{minipage}
\clearpage\markright{Je m’entraîne}
\rubrique{Je m’entraîne}
\par\noindent\begin{minipage}{\linewidth}
\begin{exercice}[2]{Calculer une distance}
\begin{questions}
\item Un chariot roule à $0{,}40\unite{m/s}$ pendant $15\unite s$ à vitesse constante. Quelle distance parcourt-il ?
\item Et pendant une durée double ?
\end{questions}
\end{exercice}
\end{minipage}
\par\noindent\begin{minipage}{\linewidth}
\begin{exercice}[2]{Calculer une durée}
\begin{questions}
\item Une marcheuse parcourt $1{,}8\unite{km}$ à $6\unite{km/h}$ de moyenne. Calculer la durée en heures.
\item Convertir en minutes.
\end{questions}
\end{exercice}
\end{minipage}
\par\noindent\begin{minipage}{\linewidth}
\begin{exercice}[2]{Lire des positions}
\pcFigureEx{figures/ex-position.pdf}{}
\begin{questions}
\item Aux instants 0, 2, 4 et 6 s, les positions sont 0, 3, 6 et 9 m. Décrire le mouvement rectiligne.
\item Calculer la vitesse.
\end{questions}
\end{exercice}
\end{minipage}
\par\noindent\begin{minipage}{\linewidth}
\begin{exercice}[2]{Uniforme sur un cercle}
\begin{questions}
\item Un siège de manège décrit un cercle à valeur de vitesse constante. Sa direction de vitesse reste-t-elle fixe ?
\item Comment qualifier le mouvement ?
\end{questions}
\end{exercice}
\end{minipage}
\par\noindent\begin{minipage}{\linewidth}
\begin{exercice}[2]{Une vidéo étalonnée}
\begin{questions}
\item Une règle de $1\unite m$ occupe $200$ pixels ; un déplacement occupe $80$ pixels. Quelle distance réelle ?
\item Ce déplacement dure $0{,}20\unite s$. Calculer la vitesse moyenne.
\end{questions}
\end{exercice}
\end{minipage}
\clearpage\markright{Je raisonne}
\rubrique{Je raisonne}
\par\noindent\begin{minipage}{\linewidth}
\begin{exercice}[3]{Deux parties d’un trajet}
\begin{questions}
\item Un trajet de $60\unite{km}$ dure $1\unite h$ ; le retour de $60\unite{km}$ dure $2\unite h$. Calculer la vitesse moyenne de l’aller-retour.
\item Pourquoi $45\unite{km/h}$ serait-il faux ?
\end{questions}
\end{exercice}
\end{minipage}
\par\noindent\begin{minipage}{\linewidth}
\begin{exercice}[3]{Un arrêt compte-t-il ?}
\begin{questions}
\item Un bus parcourt $12\unite{km}$ en $20\unite{min}$, puis s’arrête $10\unite{min}$. Calculer sa moyenne sur les 30 minutes.
\item La vitesse pendant l’arrêt vaut-elle cette moyenne ?
\end{questions}
\end{exercice}
\end{minipage}
\par\noindent\begin{minipage}{\linewidth}
\begin{exercice}[3]{Une conclusion trop forte}
\begin{questions}
\item Une voiture parcourt 100 m en 5 s. Peut-on prouver qu’elle était à 20 m/s à chaque instant ?
\item Que permet réellement le calcul ?
\end{questions}
\end{exercice}
\end{minipage}
\par\noindent\begin{minipage}{\linewidth}
\begin{exercice}[3]{Un graphique ambigu}
\pcFigureEx{figures/ex-position-retour.pdf}{Exemple de position décroissant linéairement.}
\begin{questions}
\item Une courbe représente une position qui diminue linéairement au cours du temps. Le mobile ralentit-il forcément ?
\item Quelle distinction faut-il faire ?
\end{questions}
\end{exercice}
\end{minipage}
\par\noindent\begin{minipage}{\linewidth}
\begin{exercice}[3]{Comparer des mesures}
\begin{questions}
\item Trois temps pour $10\unite m$ valent $4{,}8$, $5{,}0$ et $5{,}2\unite s$. Calculer le temps moyen.
\item Donner une vitesse représentative et discuter sa précision.
\end{questions}
\end{exercice}
\end{minipage}
\pcPage{Défi du chapitre}{Décrire et mesurer un mouvement}
\phantomsection\label{defi-ex-8}
\pcDefiTitre{8}{Choisir le meilleur trajet}
\begin{questions}
\item A mesure $1{,}2\unite{km}$ à $4\unite{km/h}$ de moyenne. B mesure $1{,}5\unite{km}$ à $6\unite{km/h}$, avec $3\unite{min}$ d’arrêt supplémentaire. Calculer les durées totales.
\item Quel trajet est le plus rapide ?
\item Quelle hypothèse pourrait changer la conclusion ?
\end{questions}
\pcGrille

\begin{ouverture}[college2]{9}{Interactions, forces et poids}{Partie 2 · Mouvements et interactions}{Comment représenter les actions sur un objet ?}
\objectif{Identifier les interactions de contact et à distance.}
\objectif{Représenter une force à une échelle donnée.}
\objectif{Distinguer masse et poids et exploiter P = mg.}
\prerequis{Masse en kg et vitesse.}
\prerequis{Directions et proportionnalité.}
\end{ouverture}
\pcPage{Le cours}{Interactions, forces et poids}
\phantomsection\label{ess-9}
\essSection{1}{Choisir le système}
Le système est l’objet ou l’ensemble d’objets que l’on étudie. Un diagramme objet-interactions relie ce système à son environnement : support, fil, air, Terre… Une interaction de contact exige un contact matériel ; la gravitation agit à distance. Le diagramme recense les partenaires ; il ne donne pas à lui seul la valeur des forces.\par
\pcFigure{figures/b09-doi.pdf}{Une balle suspendue : interactions avec le fil et la Terre ; l’air est négligé dans ce modèle.}
\essSection{2}{Modéliser par une force}
Une force représente l’action d’un objet sur un autre. On la caractérise par un point d’application, une direction, un sens et une valeur en newtons, N. La flèche part du point choisi ; sa longueur respecte une échelle annoncée. Direction verticale et sens vers le bas sont deux informations différentes. Un dynamomètre mesure une force selon son axe, après contrôle du zéro.\par
\clearpage
\essSection{3}{Poids et masse}
La masse m caractérise la quantité de matière dans les situations étudiées et se mesure en kg avec une balance. Le poids P est une force de pesanteur, en N, mesurée avec un dynamomètre. $P=mg$, avec m en kg et g en N/kg. Près de la surface terrestre, on prend souvent $g=9{,}8$ ou $10\unite{N/kg}$ selon l’énoncé. Le poids est dirigé vers le centre de la Terre dans le modèle local.\par
\pcFigure{figures/b09-poids.pdf}{Même masse, poids différents : les valeurs de g dépendent du lieu.}
\essSection{4}{Effets et équilibre}
Une action peut déformer un objet ou modifier son mouvement, en valeur ou en direction. Dans un modèle d’objet immobile soumis à deux forces qui se compensent, elles ont même direction, sens opposés et même valeur. Elles s’exercent sur le même système. Deux forces d’une interaction entre deux objets s’exercent, elles, sur des systèmes différents : on ne les additionne pas sur un seul objet.\par
\begin{pcEncadre}{method}{Un exemple expliqué}Une trousse de 350 g est immobile sur une table. $m=0{,}350\unite{kg}$, donc $P=0{,}350\times10=3{,}5\unite N$. Dans le modèle à deux forces, le support exerce 3,5 N vers le haut. À 1 cm pour 1 N, chaque flèche mesure 3,5 cm.\end{pcEncadre}
\pcPage{Expérimenter et retenir}{Interactions, forces et poids}
\phantomsection\label{activite-9}
\begin{pcEncadre}{method}{Investigation en classe}Mesurer le poids de plusieurs masses avec un dynamomètre adapté. Tracer P en fonction de m en kg, puis calculer P/m hors zéro. Comparer la pente à la valeur locale de g. Dessiner ensuite les forces sur une masse suspendue immobile, en distinguant les actions sur la masse et sur l’appareil.\end{pcEncadre}
\pcSousTitre{Ce que je dois savoir expliquer}\begin{questions}
\item Nommer l’auteur et le receveur d’une force.
\item Décrire les quatre caractéristiques d’une flèche.
\item Convertir les grammes en kilogrammes avant P = mg.
\end{questions}
\begin{pcEncadre}{attention}{Une erreur à éviter}Un objet posé sur une table a un poids. Son immobilité ne signifie pas absence de forces : dans le modèle, l’action du support compense le poids.\end{pcEncadre}
\pcRepere{cor-activite-9}{Pistes pour l’investigation}\pcGrille
\clearpage\markright{Je démarre}
\rubrique{Je démarre}
\par\noindent\begin{minipage}{\linewidth}
\begin{exercice}[1]{Contact ou distance}
\begin{questions}
\item Classer l’action d’un fil sur une masse suspendue.
\item Classer celle de la Terre sur cette masse.
\end{questions}
\end{exercice}
\end{minipage}
\par\noindent\begin{minipage}{\linewidth}
\begin{exercice}[1]{Deux appareils}
\begin{questions}
\item Quel appareil mesure une masse ?
\item Quel appareil mesure une force ?
\end{questions}
\end{exercice}
\end{minipage}
\par\noindent\begin{minipage}{\linewidth}
\begin{exercice}[1]{Deux unités}
\begin{questions}
\item Donner l’unité SI de la masse.
\item Donner celle du poids.
\end{questions}
\end{exercice}
\end{minipage}
\par\noindent\begin{minipage}{\linewidth}
\begin{exercice}[1]{Direction et sens}
\begin{questions}
\item Le poids est-il « vers le bas » ou « vertical » : quelle information donne chaque expression ?
\item Citer les deux autres caractéristiques de sa représentation.
\end{questions}
\end{exercice}
\end{minipage}
\par\noindent\begin{minipage}{\linewidth}
\begin{exercice}[1]{Un premier calcul}
\begin{questions}
\item Calculer le poids de 2 kg pour $g=10\unite{N/kg}$.
\item Une masse de 200 g vaut-elle 200 kg ?
\end{questions}
\end{exercice}
\end{minipage}
\clearpage\markright{Je m’entraîne}
\rubrique{Je m’entraîne}
\par\noindent\begin{minipage}{\linewidth}
\begin{exercice}[2]{Une petite masse}
\begin{questions}
\item Calculer P pour 150 g et $g=9{,}8\unite{N/kg}$.
\item Donner une valeur arrondie au dixième de N.
\end{questions}
\end{exercice}
\end{minipage}
\par\noindent\begin{minipage}{\linewidth}
\begin{exercice}[2]{Retrouver la masse}
\begin{questions}
\item Un objet pèse 24 N pour $g=10\unite{N/kg}$. Déterminer sa masse.
\item L’exprimer en g.
\end{questions}
\end{exercice}
\end{minipage}
\par\noindent\begin{minipage}{\linewidth}
\begin{exercice}[2]{Changer de lieu}
\begin{questions}
\item Une masse de 60 kg est transportée sur la Lune, avec $g=1{,}6\unite{N/kg}$. Donner masse et poids.
\item Comparer à son poids terrestre avec $g=10\unite{N/kg}$.
\end{questions}
\end{exercice}
\end{minipage}
\par\noindent\begin{minipage}{\linewidth}
\begin{exercice}[2]{Une échelle}
\begin{questions}
\item Une force de 12 N est représentée à 1 cm pour 4 N. Calculer la longueur.
\item Une flèche de 2,5 cm représente quelle valeur à cette échelle ?
\end{questions}
\end{exercice}
\end{minipage}
\par\noindent\begin{minipage}{\linewidth}
\begin{exercice}[2]{Le livre immobile}
\begin{questions}
\item Nommer les deux partenaires principaux d’un livre immobile sur une table horizontale.
\item Décrire leurs forces dans le modèle à deux forces.
\end{questions}
\end{exercice}
\end{minipage}
\clearpage\markright{Je raisonne}
\rubrique{Je raisonne}
\par\noindent\begin{minipage}{\linewidth}
\begin{exercice}[3]{Une pente mesurée}
\pcFigureEx{figures/rev-ex-poids-graphe.pdf}{}
\begin{questions}
\item On mesure 1,0 N pour 0,10 kg et 2,0 N pour 0,20 kg. Estimer g.
\item Que suggère un graphe linéaire passant près de l’origine ?
\end{questions}
\end{exercice}
\end{minipage}
\par\noindent\begin{minipage}{\linewidth}
\begin{exercice}[3]{Une mesure négative}
\begin{questions}
\item Un dynamomètre indique −0,2 N sans charge. Quelle vérification faire ?
\item Pourquoi ce décalage gêne-t-il surtout une petite mesure ?
\end{questions}
\end{exercice}
\end{minipage}
\par\noindent\begin{minipage}{\linewidth}
\begin{exercice}[3]{Des forces confondues}
\begin{questions}
\item Un élève dessine sur le livre son poids et la force qu’il exerce sur la table. Quelle erreur commet-il ?
\item Quelle force doit remplacer la seconde ?
\end{questions}
\end{exercice}
\end{minipage}
\par\noindent\begin{minipage}{\linewidth}
\begin{exercice}[3]{Un virage}
\begin{questions}
\item Une voiture garde une même valeur de vitesse dans un virage. Son mouvement est-il inchangé ?
\item Une force peut-elle modifier seulement la direction de la vitesse ?
\end{questions}
\end{exercice}
\end{minipage}
\par\noindent\begin{minipage}{\linewidth}
\begin{exercice}[3]{Une affirmation à discuter}
\begin{questions}
\item « Une force de 5 N suffit toujours à déplacer n’importe quel objet. » Pourquoi faut-il connaître les autres actions ?
\item Donner un cas où cette force est compensée.
\end{questions}
\end{exercice}
\end{minipage}
\pcPage{Défi du chapitre}{Interactions, forces et poids}
\phantomsection\label{defi-ex-9}
\pcDefiTitre{9}{La caisse suspendue}
\begin{questions}
\item Une caisse de 4,0 kg est suspendue immobile, $g=9{,}8\unite{N/kg}$. Calculer P.
\item Tracer les forces qui s’appliquent à la caisse, à l’échelle 1 cm pour 10 N.
\item Pourquoi ces forces de compensation ne sont-elles pas les deux forces d’une même interaction ?
\end{questions}
\pcGrille

\begin{ouverture}[college2]{10}{Exploiter la gravitation universelle}{Partie 2 · Mouvements et interactions}{Comment comparer l’attraction entre deux objets ?}
\objectif{Utiliser la loi de gravitation lorsqu’elle est fournie.}
\objectif{Convertir distances et masses en unités SI.}
\objectif{Distinguer gravitation, poids et impesanteur.}
\prerequis{Puissances de dix et calculatrice.}
\prerequis{Forces à distance ; P = mg.}
\end{ouverture}
\pcPage{Le cours}{Exploiter la gravitation universelle}
\phantomsection\label{ess-10}
\essSection{1}{Une interaction attractive}
Deux objets possédant une masse s’attirent par gravitation. Les forces sont dirigées selon la droite joignant leurs centres dans le modèle de corps sphériques. Chaque objet exerce une force sur l’autre. Les valeurs sont égales et les sens opposés ; les forces s’appliquent à des objets différents. Ce modèle permet notamment d’étudier l’interaction Terre-Lune.\par
\pcFigure{figures/b10-forces.pdf}{Deux forces sur deux systèmes ; tailles et distance ne sont pas à l’échelle.}
\essSection{2}{Utiliser la relation fournie}
Quand l’énoncé la donne, on exploite la relation \[F=G\dfrac{m_A m_B}{d^2}.\] F est en N, les masses en kg et d en m, distance entre les centres. On emploie par exemple $G=6{,}67\times10^{-11}\unite{N\,m^2/kg^2}$. Écrire les conversions avant le calcul, regrouper le dénominateur et vérifier la touche « puissance » de la calculatrice.\par
\clearpage
\essSection{3}{Comparer sans tout recalculer}
À distance inchangée, doubler une masse double F ; doubler les deux masses multiplie F par quatre. À masses inchangées, doubler d divise F par quatre ; tripler d divise F par neuf. Le carré porte sur toute la distance. Ces comparaisons concernent les conditions annoncées et le même modèle.\par
\pcFigure{figures/b10-distance.pdf}{La force diminue comme l’inverse du carré de la distance, pour des masses fixées.}
\essSection{4}{Du modèle au poids}
Près de la surface d’un astre, son attraction domine et correspond au poids dans le modèle simplifié. Avec un rayon R, on peut comparer $P=mg$ et $F=GMm/R^2$, sans confondre rayon et altitude. Un satellite en orbite reste attiré : l’impesanteur ressentie à bord correspond à une chute commune, et non à une disparition de la gravitation.\par
\begin{pcEncadre}{method}{Un exemple expliqué}Deux masses de 10 kg et 20 kg, distantes de 2 m, sont modélisées par $F=Gm_Am_B/d^2$. $F=6{,}67\times10^{-11}\times10\times20/2^2=3{,}335\times10^{-9}\unite N$. Cette force est très faible à l’échelle d’un dynamomètre scolaire ; elle n’est pas nulle.\end{pcEncadre}
\pcPage{Expérimenter et retenir}{Exploiter la gravitation universelle}
\phantomsection\label{activite-10}
\begin{pcEncadre}{method}{Investigation en classe}Utiliser une simulation ou un tableur, avec la loi fournie. Fixer les masses et calculer F pour d, 2d et 3d ; puis fixer d et changer une seule masse. Tracer des points et expliquer pourquoi F en fonction de d ne forme pas une droite. Aucun montage de gravitation sensible n’est exigé au collège.\end{pcEncadre}
\pcSousTitre{Ce que je dois savoir expliquer}\begin{questions}
\item Identifier d comme distance entre les centres.
\item Calculer avec des kg, m et une loi fournie.
\item Justifier les facteurs 2, 4 ou 9 par la relation.
\end{questions}
\begin{pcEncadre}{attention}{Une erreur à éviter}Une altitude de 400 km n’est pas la distance au centre de la Terre. Dans ce modèle, il faut ajouter le rayon terrestre avant de convertir et de mettre au carré.\end{pcEncadre}
\pcRepere{cor-activite-10}{Pistes pour l’investigation}\pcGrille
\clearpage\markright{Je démarre}
\rubrique{Je démarre}
\par\noindent\begin{minipage}{\linewidth}
\begin{exercice}[1]{Nature de l’action}
\begin{questions}
\item La gravitation est-elle une action de contact ?
\item Est-elle attractive dans le modèle étudié ?
\end{questions}
\end{exercice}
\end{minipage}
\par\noindent\begin{minipage}{\linewidth}
\begin{exercice}[1]{Deux objets}
\begin{questions}
\item La Terre attire la Lune. La Lune attire-t-elle la Terre ?
\item Sur quel objet agit la force exercée par la Lune ?
\end{questions}
\end{exercice}
\end{minipage}
\par\noindent\begin{minipage}{\linewidth}
\begin{exercice}[1]{Unités}
\begin{questions}
\item Dans $F=Gm_A m_B/d^2$, donner les unités de m et d pour obtenir F en N.
\item Convertir 6 400 km en m.
\end{questions}
\end{exercice}
\end{minipage}
\par\noindent\begin{minipage}{\linewidth}
\begin{exercice}[1]{Quelle distance ?}
\begin{questions}
\item Pour deux sphères, quelle distance figure dans la loi donnée ?
\item Pourquoi la distance entre leurs surfaces ne suffit-elle pas ?
\end{questions}
\end{exercice}
\end{minipage}
\par\noindent\begin{minipage}{\linewidth}
\begin{exercice}[1]{Une masse doublée}
\begin{questions}
\item La masse A double, le reste ne change pas. Que devient F ?
\item Si les deux masses doublent ?
\end{questions}
\end{exercice}
\end{minipage}
\clearpage\markright{Je m’entraîne}
\rubrique{Je m’entraîne}
\par\noindent\begin{minipage}{\linewidth}
\begin{exercice}[2]{Une distance doublée}
\begin{questions}
\item Les masses sont fixées et d double. Que devient F ?
\item Et si d est divisée par deux ?
\end{questions}
\end{exercice}
\end{minipage}
\par\noindent\begin{minipage}{\linewidth}
\begin{exercice}[2]{Une force modèle}
\begin{questions}
\item Avec $G=6{,}67\times10^{-11}$ en unités SI, calculer F pour 1 kg, 1 kg et 1 m.
\item Est-ce une force nulle ?
\end{questions}
\end{exercice}
\end{minipage}
\par\noindent\begin{minipage}{\linewidth}
\begin{exercice}[2]{Une comparaison rapide}
\begin{questions}
\item À d, F vaut 36 N. Quelle valeur à 3d avec les mêmes masses ?
\item Quelle valeur à 2d ?
\end{questions}
\end{exercice}
\end{minipage}
\par\noindent\begin{minipage}{\linewidth}
\begin{exercice}[2]{Depuis l’orbite}
\begin{questions}
\item Rayon terrestre donné : 6 400 km ; altitude : 400 km. Déterminer d en km puis m.
\item Quel terme doit-on mettre au carré ?
\end{questions}
\end{exercice}
\end{minipage}
\par\noindent\begin{minipage}{\linewidth}
\begin{exercice}[2]{Un petit calcul}
\begin{questions}
\item Calculer F pour $m_A=100\unite{kg}$, $m_B=50\unite{kg}$, $d=10\unite m$, $G=6{,}67\times10^{-11}$ en SI.
\item Prévoir la valeur si seule $m_B$ devient 100 kg.
\end{questions}
\end{exercice}
\end{minipage}
\clearpage\markright{Je raisonne}
\rubrique{Je raisonne}
\par\noindent\begin{minipage}{\linewidth}
\begin{exercice}[3]{Trois facteurs}
\begin{questions}
\item La masse A double et d double, B reste fixe. Par quel facteur F est-elle multipliée ?
\item Justifier avec la relation.
\end{questions}
\end{exercice}
\end{minipage}
\par\noindent\begin{minipage}{\linewidth}
\begin{exercice}[3]{Un faux équilibre}
\begin{questions}
\item « Les deux forces Terre-Lune s’annulent, donc la Lune n’est pas attirée. » Quelle erreur de système apparaît ?
\item Quelle force faut-il retenir pour étudier seulement la Lune ?
\end{questions}
\end{exercice}
\end{minipage}
\par\noindent\begin{minipage}{\linewidth}
\begin{exercice}[3]{Impesanteur}
\begin{questions}
\item Pourquoi l’impesanteur dans une station orbitale ne signifie-t-elle pas absence d’attraction ?
\item Quelle hypothèse donne l’explication qualitative ?
\end{questions}
\end{exercice}
\end{minipage}
\par\noindent\begin{minipage}{\linewidth}
\begin{exercice}[3]{Retrouver g}
\begin{questions}
\item La loi fournie donne une force de 588 N sur une masse de 60 kg près du sol. Trouver g à partir de P = mg.
\item Quelle masse aurait un poids de 98 N au même lieu ?
\end{questions}
\end{exercice}
\end{minipage}
\par\noindent\begin{minipage}{\linewidth}
\begin{exercice}[3]{La précision du modèle}
\begin{questions}
\item Une calculatrice affiche dix chiffres pour F, mais les distances ont deux chiffres significatifs. Pourquoi éviter une conclusion à dix chiffres ?
\item Citer une hypothèse du modèle gravitationnel utilisé.
\end{questions}
\end{exercice}
\end{minipage}
\pcPage{Défi du chapitre}{Exploiter la gravitation universelle}
\phantomsection\label{defi-ex-10}
\pcDefiTitre{10}{Autour d’une planète}
\begin{questions}
\item Une planète a $M=6{,}0\times10^{24}\unite{kg}$. Un satellite de 500 kg est à $d=7{,}0\times10^6\unite m$. Calculer F avec $G=6{,}67\times10^{-11}$ en SI.
\item Si d double, que devient F ?
\item Comparer les forces planète-satellite et expliquer pourquoi leurs effets ne sont pas identiques.
\end{questions}
\pcGrille

\begin{ouverture}[college3]{11}{Énergie cinétique et bilan mécanique}{Partie 3 · L’énergie, ses transferts et ses conversions}{Pourquoi la vitesse pèse-t-elle autant sur un freinage ?}
\objectif{Calculer une énergie cinétique en unités SI.}
\objectif{Comparer l’effet de la masse et de la vitesse.}
\objectif{Suivre une conversion et distinguer énergie et force.}
\prerequis{Vitesse en m/s ; carré d’un nombre.}
\prerequis{Conservation et formes d’énergie.}
\end{ouverture}
\pcPage{Le cours}{Énergie cinétique et bilan mécanique}
\phantomsection\label{ess-11}
\essSection{1}{L’énergie du mouvement}
Un objet de masse m et de vitesse de valeur v possède, dans le référentiel choisi, une énergie cinétique \[E_c=\dfrac12 mv^2.\] Avec m en kg et v en m/s, $E_c$ est en joules, J. La vitesse se convertit avant le carré : une valeur en km/h doit être divisée par 3,6. Un objet immobile dans le référentiel considéré y a une énergie cinétique nulle.\par
\pcFigure{figures/b11-ec.pdf}{À masse fixée, le graphe d’Ec en fonction de v n’est pas une droite.}
\essSection{2}{Des facteurs à comparer}
À vitesse fixée, doubler la masse double l’énergie cinétique. À masse fixée, doubler la vitesse multiplie cette énergie par quatre ; tripler la vitesse la multiplie par neuf. Ces effets expliquent pourquoi une petite augmentation de vitesse peut avoir un effet énergétique important. Le calcul ne donne pas à lui seul une distance d’arrêt : freinage, adhérence, pente et temps de réaction interviennent aussi.\par
\clearpage
\essSection{3}{Position et conversion}
Une énergie potentielle de pesanteur dépend de la position et du choix d’un niveau de référence. Lors d’une chute, elle peut diminuer tandis que l’énergie cinétique augmente. Si les frottements sont négligés, leur somme, l’énergie mécanique, reste constante. Avec frottements, une partie de l’énergie mécanique est transférée ou convertie en énergie thermique ; l’énergie totale d’un système suffisamment large reste conservée.\par
\pcFigure{figures/b11-chute.pdf}{Modèle de chute : exemple chiffré sans frottement ; le niveau de référence est annoncé.}
\essSection{4}{Freiner et argumenter}
Freiner réduit l’énergie cinétique et augmente notamment l’énergie thermique des freins, pneus et milieu. L’énergie n’est pas détruite. Une force se mesure en N, une énergie en J, une puissance en W : ce sont des grandeurs distinctes. Au Brevet, comparer les carrés des vitesses peut suffire à justifier un rapport d’énergies sans calculer deux grandes valeurs.\par
\begin{pcEncadre}{method}{Un exemple expliqué}Pour un ensemble vélo-cycliste de 80 kg à 18 km/h, $v=18/3{,}6=5\unite{m/s}$ ; $E_c=\tfrac12\times80\times5^2=1\,000\unite J$. À 36 km/h, v double : $E_c=4\,000\unite J$. La comparaison suppose la même masse et le même référentiel.\end{pcEncadre}
\pcPage{Expérimenter et retenir}{Énergie cinétique et bilan mécanique}
\phantomsection\label{activite-11}
\begin{pcEncadre}{method}{Investigation en classe}Au tableur, comparer les énergies d’une même masse pour 2, 4, 6 et 8 m/s. Ajouter une colonne $v^2$ et tracer Ec en fonction de v, puis de $v^2$. Exploiter ensuite une vidéo ou une simulation de chute en distinguant modèle sans frottement et observation réelle. Ne pas provoquer de choc dangereux.\end{pcEncadre}
\pcSousTitre{Ce que je dois savoir expliquer}\begin{questions}
\item Écrire Ec avec m en kg et v en m/s.
\item Justifier l’effet du carré de la vitesse.
\item Faire un bilan comprenant les sorties thermiques.
\end{questions}
\begin{pcEncadre}{attention}{Une erreur à éviter}Doubler la vitesse ne double pas l’énergie cinétique : dans le modèle, cela la multiplie par quatre. Il faut mettre toute la vitesse convertie au carré.\end{pcEncadre}
\pcRepere{cor-activite-11}{Pistes pour l’investigation}\pcGrille
\clearpage\markright{Je démarre}
\rubrique{Je démarre}
\par\noindent\begin{minipage}{\linewidth}
\begin{exercice}[1]{La forme d’énergie}
\begin{questions}
\item Quelle énergie est associée au mouvement d’un objet ?
\item Dans quelle unité SI s’exprime-t-elle ?
\end{questions}
\end{exercice}
\end{minipage}
\par\noindent\begin{minipage}{\linewidth}
\begin{exercice}[1]{Les bonnes unités}
\begin{questions}
\item Quelles unités de m et v employer dans Ec = ½mv² ?
\item Convertir 500 g en kg.
\end{questions}
\end{exercice}
\end{minipage}
\par\noindent\begin{minipage}{\linewidth}
\begin{exercice}[1]{La vitesse}
\begin{questions}
\item Convertir 36 km/h en m/s.
\item Convertir 72 km/h en m/s.
\end{questions}
\end{exercice}
\end{minipage}
\par\noindent\begin{minipage}{\linewidth}
\begin{exercice}[1]{Un calcul direct}
\begin{questions}
\item Calculer Ec pour 2 kg à 3 m/s.
\item Que vaut Ec si cet objet est immobile dans le référentiel ?
\end{questions}
\end{exercice}
\end{minipage}
\par\noindent\begin{minipage}{\linewidth}
\begin{exercice}[1]{Deux facteurs}
\begin{questions}
\item À vitesse identique, la masse double. Que devient Ec ?
\item À masse identique, la vitesse double ?
\end{questions}
\end{exercice}
\end{minipage}
\clearpage\markright{Je m’entraîne}
\rubrique{Je m’entraîne}
\par\noindent\begin{minipage}{\linewidth}
\begin{exercice}[2]{Le ballon}
\begin{questions}
\item Un ballon de 400 g va à 5 m/s. Calculer Ec.
\item Si v devient 10 m/s, calculer la nouvelle valeur.
\end{questions}
\end{exercice}
\end{minipage}
\par\noindent\begin{minipage}{\linewidth}
\begin{exercice}[2]{Une voiture}
\begin{questions}
\item Une voiture de 1 000 kg va à 54 km/h. Déterminer v puis Ec.
\item Donner Ec en kJ.
\end{questions}
\end{exercice}
\end{minipage}
\par\noindent\begin{minipage}{\linewidth}
\begin{exercice}[2]{Ralentir}
\begin{questions}
\item Un ensemble de 60 kg passe de 10 à 5 m/s. Calculer les deux énergies.
\item Quelle énergie cinétique est perdue ?
\end{questions}
\end{exercice}
\end{minipage}
\par\noindent\begin{minipage}{\linewidth}
\begin{exercice}[2]{Une chute idéale}
\begin{questions}
\item Au départ, Ec = 0 J et Ep = 40 J. Sans frottement, que vaut Ec lorsque Ep = 15 J ?
\item Pourquoi faut-il annoncer le modèle ?
\end{questions}
\end{exercice}
\end{minipage}
\par\noindent\begin{minipage}{\linewidth}
\begin{exercice}[2]{Le rapport}
\begin{questions}
\item À masse fixée, comparer Ec à 30 m/s et à 10 m/s.
\item À vitesse fixée, comparer 120 kg et 80 kg.
\end{questions}
\end{exercice}
\end{minipage}
\clearpage\markright{Je raisonne}
\rubrique{Je raisonne}
\par\noindent\begin{minipage}{\linewidth}
\begin{exercice}[3]{Une hausse de 20 \%}
\begin{questions}
\item La vitesse est multipliée par 1,2 à masse fixée. Par quel facteur Ec est-elle multipliée ?
\item Exprimer l’augmentation en pourcentage.
\end{questions}
\end{exercice}
\end{minipage}
\par\noindent\begin{minipage}{\linewidth}
\begin{exercice}[3]{Changer de référentiel}
\begin{questions}
\item Un passager est immobile par rapport à son train. Peut-il avoir une énergie cinétique dans le référentiel du sol ?
\item Quelle précision manque à « cet objet n’a aucune énergie cinétique » ?
\end{questions}
\end{exercice}
\end{minipage}
\par\noindent\begin{minipage}{\linewidth}
\begin{exercice}[3]{La force n’est pas l’énergie}
\begin{questions}
\item Un élève donne une énergie cinétique de 250 N. Corriger l’unité.
\item Peut-on trouver une distance de freinage avec seulement Ec ?
\end{questions}
\end{exercice}
\end{minipage}
\par\noindent\begin{minipage}{\linewidth}
\begin{exercice}[3]{Un bilan réel}
\begin{questions}
\item Une chute commence avec 100 J d’énergie mécanique et finit avec 85 J. Où les 15 J peuvent-ils être passés ?
\item La conservation totale est-elle remise en cause ?
\end{questions}
\end{exercice}
\end{minipage}
\par\noindent\begin{minipage}{\linewidth}
\begin{exercice}[3]{Comparer deux véhicules}
\begin{questions}
\item A : 800 kg à 20 m/s ; B : 1 600 kg à 10 m/s. Lequel a le plus d’énergie cinétique ?
\item Pourquoi comparer seulement les masses donnerait-il une erreur ?
\end{questions}
\end{exercice}
\end{minipage}
\pcPage{Défi du chapitre}{Énergie cinétique et bilan mécanique}
\phantomsection\label{defi-ex-11}
\pcDefiTitre{11}{Une étude de sécurité}
\begin{questions}
\item Une voiture et ses occupants ont une masse de 1 200 kg. Comparer Ec à 50 km/h et 100 km/h.
\item Expliquer le résultat sans recalculer la seconde énergie.
\item Rédiger une conclusion de quatre lignes qui distingue énergie, distance d’arrêt et conditions de freinage.
\end{questions}
\pcGrille

\begin{ouverture}[college3]{12}{L’énergie : transferts et conversions}{Partie 3 · L’énergie, ses transferts et ses conversions}{Où va l’énergie d’une lampe ?}
\objectif{Distinguer source, forme, transfert et conversion.}
\objectif{Construire une chaîne et un bilan énergétiques simples.}
\objectif{Relier rayonnement, climat et choix de ressources.}
\prerequis{Combustions, mouvement et lecture d’un schéma.}
\end{ouverture}
\pcPage{Le cours}{L’énergie : transferts et conversions}
\phantomsection\label{ess-12}
\essSection{1}{Nommer avec précision}
Une source fournit de l’énergie à un système. Une forme décrit comment elle est stockée ou se manifeste : énergie chimique d’une pile, cinétique d’un objet en mouvement, potentielle liée à une position, thermique ou nucléaire. Le Soleil fournit un rayonnement ; une éolienne reçoit de l’énergie du vent. Le mot « électrique » qualifie notamment un transfert électrique dans une chaîne. Un dispositif convertit l’énergie sans la créer.\par
\pcFigure{figures/c12-chaine.pdf}{Pour la lampe : transfert électrique entrant, transferts par rayonnement et thermiques sortants. Les flèches ne représentent pas de la matière.}
\essSection{2}{Transférer et convertir}
Un transfert fait passer de l’énergie d’un système à un autre : électriquement, mécaniquement, thermiquement ou par rayonnement. Une conversion change la forme sous laquelle on décrit l’énergie dans le fonctionnement étudié. Un moteur électrique reçoit de l’énergie électriquement et met une charge en mouvement, tout en échauffant son environnement. Une chaîne nomme les systèmes et le type de chaque transfert.\par
\clearpage
\essSection{3}{Un bilan qui tient compte de tout}
L’énergie se conserve. Pour un système, la variation de l’énergie stockée correspond à l’énergie reçue moins l’énergie cédée. Si un convertisseur fonctionne sans variation de son stockage, l’énergie entrante égale la somme des énergies sortantes. Les « pertes » sont des transferts qui ne répondent pas à l’effet utile recherché ; elles ne constituent pas une disparition. L’énergie s’exprime en joules (J).\par
\pcFigure{figures/c12-terre.pdf}{Bilan qualitatif Terre-atmosphère : rayonnement solaire reçu, part renvoyée et rayonnement thermique émis. Les flèches ne sont pas quantitatives.}
\essSection{4}{Ressources et climat}
Le charbon, le pétrole et le gaz sont des ressources fossiles ; leur combustion émet du dioxyde de carbone. Une ressource renouvelable se renouvelle à une échelle compatible avec l’usage humain, sous conditions d’exploitation. « Renouvelable » ne signifie pas « sans impact ». Certains gaz à effet de serre absorbent et émettent du rayonnement infrarouge terrestre. Si la Terre reçoit plus d’énergie qu’elle n’en cède, son énergie stockée augmente, ce qui peut élever sa température moyenne.\par
\begin{pcEncadre}{method}{Un exemple expliqué}Une lampe reçoit $100\unite J$ électriquement et cède $25\unite J$ par rayonnement utile dans un modèle sans stockage. Les autres transferts totalisent $100-25=75\unite J$. L’énergie totale se conserve, mais toute l’énergie reçue n’est pas transférée sous la forme utile à l’éclairage.\end{pcEncadre}
\pcPage{Expérimenter et retenir}{L’énergie : transferts et conversions}
\phantomsection\label{activite-12}
\begin{pcEncadre}{method}{Investigation en classe}\textbf{Mission.} Étudier une lampe scolaire, un petit moteur et une cellule photovoltaïque, sur documents ou matériel autorisé. Définir chaque système, repérer l’entrée et les sorties puis construire une chaîne. Chercher un indice d’un transfert moins visible, par exemple un échauffement mesuré avec une sonde. Ne pas toucher une pièce chaude. \textbf{À rendre :} chaîne légendée, effet utile choisi, limites et bilan chiffré seulement lorsque toutes les données nécessaires sont fournies.\end{pcEncadre}
\pcSousTitre{Ce que je dois savoir expliquer}\begin{questions}
\item Une source n’est pas la même chose qu’une forme d’énergie.
\item Un transfert relie deux systèmes ; une conversion décrit leur fonctionnement.
\item L’énergie totale se conserve lorsqu’on inclut tous les transferts et stockages.
\item Un transfert non utile ne fait pas disparaître l’énergie.
\item Le rayonnement permet un transfert à travers le vide.
\end{questions}
\begin{pcEncadre}{attention}{Une erreur à éviter}« Le moteur fabrique de l’énergie. » Il convertit de l’énergie reçue ; la production usuelle d’électricité désigne une conversion et un transfert.\end{pcEncadre}
\pcRepere{cor-activite-12}{Pistes pour l’investigation}\pcGrille
\clearpage\markright{Je démarre}
\rubrique{Je démarre}
\par\noindent\begin{minipage}{\linewidth}
\begin{exercice}[1]{Identifier une source}
\begin{questions}
\item Dans une lampe alimentée par une pile, quel système fournit l’énergie à la lampe ?
\item Quelle énergie est stockée dans la pile avant utilisation ?
\end{questions}
\end{exercice}
\end{minipage}
\par\noindent\begin{minipage}{\linewidth}
\begin{exercice}[1]{Nommer une forme}
\begin{questions}
\item Quelle forme associe-t-on au mouvement d’un vélo ?
\item Et à un objet placé en hauteur dans le champ de pesanteur ?
\end{questions}
\end{exercice}
\end{minipage}
\par\noindent\begin{minipage}{\linewidth}
\begin{exercice}[1]{Une unité}
\begin{questions}
\item Quelle est l’unité SI de l’énergie ?
\item Le watt est-il une autre écriture du joule ?
\end{questions}
\end{exercice}
\end{minipage}
\par\noindent\begin{minipage}{\linewidth}
\begin{exercice}[1]{Une sortie moins visible}
\begin{questions}
\item Une lampe éclaire et chauffe son environnement. Nommer deux transferts sortants.
\item L’échauffement signifie-t-il une destruction d’énergie ?
\end{questions}
\end{exercice}
\end{minipage}
\par\noindent\begin{minipage}{\linewidth}
\begin{exercice}[1]{Renouvelable ou fossile}
\begin{questions}
\item Classer vent et charbon.
\item Renouvelable signifie-t-il sans conséquence environnementale ?
\end{questions}
\end{exercice}
\end{minipage}
\clearpage\markright{Je m’entraîne}
\rubrique{Je m’entraîne}
\par\noindent\begin{minipage}{\linewidth}
\begin{exercice}[2]{Un bilan de lampe}
\begin{questions}
\item Une lampe reçoit $200\unite J$ et produit $60\unite J$ de transfert utile sans stockage. Calculer les autres transferts.
\item Vérifier le bilan.
\end{questions}
\end{exercice}
\end{minipage}
\par\noindent\begin{minipage}{\linewidth}
\begin{exercice}[2]{Le moteur}
\begin{questions}
\item Un moteur reçoit $500\unite J$, cède $350\unite J$ mécaniquement et $150\unite J$ thermiquement. Le bilan est-il cohérent sans stockage ?
\item Quel effet utile peut-on choisir ?
\end{questions}
\end{exercice}
\end{minipage}
\par\noindent\begin{minipage}{\linewidth}
\begin{exercice}[2]{Un stockage augmente}
\begin{questions}
\item Un système reçoit $80\unite J$ et en cède $30\unite J$. Quelle variation de stockage ?
\item Peut-on imposer entrée égale sortie ici ?
\end{questions}
\end{exercice}
\end{minipage}
\par\noindent\begin{minipage}{\linewidth}
\begin{exercice}[2]{Une cellule photovoltaïque}
\begin{questions}
\item Décrire l’entrée et une sortie utile d’une cellule photovoltaïque.
\item La lumière transporte-t-elle seulement de l’information ?
\end{questions}
\end{exercice}
\end{minipage}
\par\noindent\begin{minipage}{\linewidth}
\begin{exercice}[2]{Un bilan terrestre}
\begin{questions}
\item Sur une période, la Terre reçoit plus d’énergie qu’elle n’en cède. Que devient son stockage ?
\item Pourquoi un simple dessin ne donne-t-il pas une valeur chiffrée ?
\end{questions}
\end{exercice}
\end{minipage}
\clearpage\markright{Je raisonne}
\rubrique{Je raisonne}
\par\noindent\begin{minipage}{\linewidth}
\begin{exercice}[3]{Comparer deux lampes}
\begin{questions}
\item Pour un même éclairage, A reçoit $100\unite J$ et B $40\unite J$ sur la même durée. Calculer l’économie avec B.
\item Que faut-il vérifier pour que la comparaison soit équitable ?
\end{questions}
\end{exercice}
\end{minipage}
\par\noindent\begin{minipage}{\linewidth}
\begin{exercice}[3]{Une énergie cinétique fournie}
\begin{questions}
\item On utilise $E_c=\frac12mv^2$. Avec $m=2\unite{kg}$ et $v=3\unite{m/s}$, calculer $E_c$.
\item Si la vitesse double, l’énergie double-t-elle ?
\end{questions}
\end{exercice}
\end{minipage}
\par\noindent\begin{minipage}{\linewidth}
\begin{exercice}[3]{Le mot « perte »}
\begin{questions}
\item Pourquoi les pertes d’un moteur peuvent-elles rester importantes malgré la conservation de l’énergie ?
\item Peut-on récupérer toute énergie aussi facilement ?
\end{questions}
\end{exercice}
\end{minipage}
\par\noindent\begin{minipage}{\linewidth}
\begin{exercice}[3]{Un raccourci sur le climat}
\begin{questions}
\item « Les gaz à effet de serre empêchent toute énergie de sortir. » Corriger.
\item Quelle idée de bilan faut-il retenir ?
\end{questions}
\end{exercice}
\end{minipage}
\par\noindent\begin{minipage}{\linewidth}
\begin{exercice}[3]{Une comparaison incomplète}
\begin{questions}
\item Un appareil émet peu de CO2 pendant son usage. Peut-on en déduire qu’il n’a aucun impact ?
\item Quels postes ajouter ?
\end{questions}
\end{exercice}
\end{minipage}
\pcPage{Défi du chapitre}{L’énergie : transferts et conversions}
\phantomsection\label{defi-ex-12}
\pcDefiTitre{12}{Le convertisseur mystérieux}
\begin{questions}
\item Un appareil reçoit $800\unite J$, cède $500\unite J$ mécaniquement et $220\unite J$ thermiquement. Calculer la variation d’énergie stockée.
\item Un élève affirme qu’il a « perdu 80 J ». Corriger.
\item Comment tester l’hypothèse d’un stockage plutôt qu’un transfert oublié ?
\end{questions}
\pcGrille

\begin{ouverture}[college3]{13}{Circuits : intensités et tensions}{Partie 3 · L’énergie, ses transferts et ses conversions}{Quelles lois permettent de contrôler un montage ?}
\objectif{Schématiser série et dérivation.}
\objectif{Brancher les appareils et convertir les mesures.}
\objectif{Exploiter les lois des intensités et des tensions.}
\prerequis{Symboles usuels et circuit fermé.}
\prerequis{A, mA et V.}
\end{ouverture}
\pcPage{Le cours}{Circuits : intensités et tensions}
\phantomsection\label{ess-13}
\essSection{1}{Lire la topologie}
Deux dipôles sont en série s’ils appartiennent à une même branche sans dérivation entre eux. Deux dipôles sont en dérivation s’ils sont reliés aux mêmes deux nœuds. Le schéma indique les connexions ; la position sur la feuille ne suffit pas. Une source alimente le circuit fermé. Un court-circuit peut provoquer une intensité élevée et un échauffement : on coupe l’alimentation avant toute modification.\par
\pcFigure{figures/b13-circuits.pdf}{Topologie de deux circuits : suivre les connexions et repérer les deux nœuds.}
\essSection{2}{Mesurer correctement}
Un ampèremètre mesure une intensité I et se branche en série dans la branche étudiée. Un voltmètre mesure une tension U et se branche en dérivation aux bornes du dipôle. Utiliser les bornes et le calibre adaptés, commencer par un calibre suffisant et respecter le sens de connexion en courant continu. $1\unite A=1\,000\unite{mA}$. Un affichage négatif peut venir d’une inversion des bornes.\par
\clearpage
\essSection{3}{Les lois des intensités}
Dans une unique branche de dipôles en série, l’intensité du courant est la même en tout point en régime permanent. À un nœud, la somme des intensités entrantes égale celle des intensités sortantes ; avec deux branches réceptrices, $I_G=I_1+I_2$. Le courant n’est pas « consommé » par une lampe : c’est l’énergie qui lui est transférée.\par
\pcFigure{figures/b13-mesures.pdf}{Ampèremètre en série et voltmètre en dérivation : connexions distinctes.}
\essSection{4}{Les lois des tensions}
Dans une boucle simple comprenant la source et deux récepteurs en série, avec les valeurs positives et conventions indiquées, $U_G=U_1+U_2$. Deux dipôles reliés aux mêmes deux nœuds ont la même tension : $U_1=U_2=U_G$ dans le montage en dérivation directement alimenté. Ne pas confondre égalité des tensions et égalité des intensités : les courants des branches peuvent différer.\par
\begin{pcEncadre}{method}{Un exemple expliqué}Deux lampes en dérivation sont soumises à 6,0 V. Les intensités sont 120 mA et 180 mA. La source débite $I_G=0{,}120+0{,}180=0{,}300\unite A$. Chaque lampe reçoit 6,0 V ; cette égalité n’oblige pas les deux courants à être égaux.\end{pcEncadre}
\pcPage{Expérimenter et retenir}{Circuits : intensités et tensions}
\phantomsection\label{activite-13}
\begin{pcEncadre}{method}{Investigation en classe}Réaliser sous contrôle un circuit scolaire à deux lampes. Mesurer les intensités en série, puis les trois intensités d’un montage en dérivation. Contrôler ensuite les tensions sur les mêmes nœuds. Comparer aux lois, en tenant compte du calibre, de la lecture et de la stabilité de la source.\end{pcEncadre}
\pcSousTitre{Ce que je dois savoir expliquer}\begin{questions}
\item Choisir le bon montage de mesure.
\item Écrire la loi avant de remplacer les valeurs.
\item Vérifier la cohérence d’une mesure avec les connexions.
\end{questions}
\begin{pcEncadre}{attention}{Une erreur à éviter}Brancher un ampèremètre directement aux bornes d’une source peut créer un court-circuit. Il ne remplace jamais un voltmètre.\end{pcEncadre}
\pcRepere{cor-activite-13}{Pistes pour l’investigation}\pcGrille
\clearpage\markright{Je démarre}
\rubrique{Je démarre}
\par\noindent\begin{minipage}{\linewidth}
\begin{exercice}[1]{Le branchement}
\begin{questions}
\item Comment brancher un ampèremètre ?
\item Comment brancher un voltmètre ?
\end{questions}
\end{exercice}
\end{minipage}
\par\noindent\begin{minipage}{\linewidth}
\begin{exercice}[1]{Conversions}
\begin{questions}
\item Convertir 250 mA en A.
\item Convertir 0,032 A en mA.
\end{questions}
\end{exercice}
\end{minipage}
\par\noindent\begin{minipage}{\linewidth}
\begin{exercice}[1]{En série}
\begin{questions}
\item Dans une seule branche, on mesure 0,18 A avant une lampe. Quelle intensité après ?
\item Le courant est-il consommé ?
\end{questions}
\end{exercice}
\end{minipage}
\par\noindent\begin{minipage}{\linewidth}
\begin{exercice}[1]{Deux nœuds}
\pcFigureEx{figures/ex-circuit-parallele.pdf}{}
\begin{questions}
\item Deux lampes sont reliées aux mêmes deux nœuds. Comment nommer l’association ?
\item Quelle grandeur a la même valeur aux bornes des lampes ?
\end{questions}
\end{exercice}
\end{minipage}
\par\noindent\begin{minipage}{\linewidth}
\begin{exercice}[1]{Une somme}
\pcFigureEx{figures/rev-ex-noeud-035.pdf}{}
\begin{questions}
\item Deux courants de 0,10 A et 0,25 A sortent d’un nœud. Quelle intensité totale entre ?
\item Donner la valeur en mA.
\end{questions}
\end{exercice}
\end{minipage}
\clearpage\markright{Je m’entraîne}
\rubrique{Je m’entraîne}
\par\noindent\begin{minipage}{\linewidth}
\begin{exercice}[2]{Une branche inconnue}
\pcFigureEx{figures/rev-ex-noeud-080.pdf}{}
\begin{questions}
\item $I_G=0{,}80\unite A$, $I_1=250\unite{mA}$. Calculer I2.
\item Quelle erreur ferait-on en calculant $0{,}80-250$ ?
\end{questions}
\end{exercice}
\end{minipage}
\par\noindent\begin{minipage}{\linewidth}
\begin{exercice}[2]{Tensions en série}
\begin{questions}
\item Une source de 9 V alimente deux récepteurs en série. U1 = 3,5 V. Trouver U2.
\item Où brancher le voltmètre pour mesurer U2 ?
\end{questions}
\end{exercice}
\end{minipage}
\par\noindent\begin{minipage}{\linewidth}
\begin{exercice}[2]{Tensions en dérivation}
\begin{questions}
\item Deux résistances sont en dérivation sur une source idéale de 12 V. Donner leurs tensions.
\item Leurs intensités sont-elles forcément égales ?
\end{questions}
\end{exercice}
\end{minipage}
\par\noindent\begin{minipage}{\linewidth}
\begin{exercice}[2]{Une inversion}
\begin{questions}
\item Un voltmètre affiche −4,5 V. Quelle cause de connexion envisager ?
\item L’intensité est-elle pour autant forcément nulle ?
\end{questions}
\end{exercice}
\end{minipage}
\par\noindent\begin{minipage}{\linewidth}
\begin{exercice}[2]{Commander séparément}
\begin{questions}
\item Deux lampes doivent être commandées indépendamment. Quelle association proposer ?
\item Où placer les deux interrupteurs ?
\end{questions}
\end{exercice}
\end{minipage}
\clearpage\markright{Je raisonne}
\rubrique{Je raisonne}
\par\noindent\begin{minipage}{\linewidth}
\begin{exercice}[3]{Une mesure incohérente}
\begin{questions}
\item En série, un élève mesure 0,20 A puis 0,12 A. La loi est-elle invalidée par ces deux nombres seuls ?
\item Citer deux contrôles.
\end{questions}
\end{exercice}
\end{minipage}
\par\noindent\begin{minipage}{\linewidth}
\begin{exercice}[3]{Un nœud à trois sorties}
\pcFigureEx{figures/rev-ex-noeud-090.pdf}{}
\begin{questions}
\item Entrée : 0,90 A ; sorties connues : 0,20 A et 0,35 A. Calculer la troisième.
\item La somme des sorties doit-elle être plus petite que l’entrée en régime permanent ?
\end{questions}
\end{exercice}
\end{minipage}
\par\noindent\begin{minipage}{\linewidth}
\begin{exercice}[3]{Une lampe retirée}
\begin{questions}
\item Dans un montage en dérivation, on ouvre une branche sans toucher l’autre. L’autre lampe peut-elle rester allumée ?
\item Quelle hypothèse sur la source faut-il annoncer pour conserver exactement sa tension ?
\end{questions}
\end{exercice}
\end{minipage}
\clearpage
\par\noindent\begin{minipage}{\linewidth}
\begin{exercice}[3]{Court-circuit}
\pcFigureEx{figures/ex-lampe-court-circuitee.pdf}{Défaut à analyser sur le schéma, sans le réaliser.}
\begin{questions}
\item Un fil relie directement les deux bornes d’une lampe. Pourquoi la situation est-elle différente d’une branche avec une lampe ?
\item Que faire avant de modifier le montage ?
\end{questions}
\end{exercice}
\end{minipage}
\par\noindent\begin{minipage}{\linewidth}
\begin{exercice}[3]{Un cahier des charges}
\begin{questions}
\item Deux lampes doivent recevoir chacune 6 V, l’une sous 0,10 A et l’autre sous 0,20 A. Quelle association et quelle intensité de source prévoir ?
\item Pourquoi une association en série sur 6 V ne répond-elle pas en général à la demande ?
\end{questions}
\end{exercice}
\end{minipage}
\pcPage{Défi du chapitre}{Circuits : intensités et tensions}
\phantomsection\label{defi-ex-13}
\pcDefiTitre{13}{Le montage à vérifier}
\pcFigureEx{figures/rev-ex-montage-9V.pdf}{}
\begin{questions}
\item Deux branches sont reliées à une source de 9 V. La première comprend deux lampes en série sous 4 V et 5 V ; la seconde une résistance. Vérifier les tensions.
\item Les courants de branches valent 150 mA et 60 mA. Trouver IG.
\item Décrire un protocole de mesure de la tension de la résistance et du courant de la première branche.
\end{questions}
\pcGrille

\begin{ouverture}[college3]{14}{Loi d’Ohm et puissance électrique}{Partie 3 · L’énergie, ses transferts et ses conversions}{Comment prévoir le fonctionnement d’un dipôle ?}
\objectif{Exploiter U = RI et une caractéristique.}
\objectif{Calculer P = UI en courant continu.}
\objectif{Contrôler unités, limites et puissance admissible.}
\prerequis{Lois des circuits et montage A/V.}
\prerequis{Proportionnalité.}
\end{ouverture}
\pcPage{Le cours}{Loi d’Ohm et puissance électrique}
\phantomsection\label{ess-14}
\essSection{1}{Un conducteur ohmique}
Pour un conducteur ohmique dans des conditions physiques fixées, notamment de température, la tension U est proportionnelle à l’intensité I : $U=RI$. U en V, I en A et R en ohms, $\Omega$. La caractéristique U en fonction de I est une droite passant par l’origine dans le modèle. La pente U/I vaut R. Tous les dipôles ne sont pas ohmiques.\par
\pcFigure{figures/b14-ohm.pdf}{Résistance de 100 Ω : attention à l’unité de l’axe des intensités.}
\essSection{2}{Isoler l’inconnue}
On peut écrire $I=U/R$ ou $R=U/I$ pour I non nul. Convertir les mA en A et les k$\Omega$ en $\Omega$. À tension fixée, une résistance plus grande laisse passer une intensité plus faible. Une lampe peut chauffer et présenter une caractéristique différente : on ne lui applique pas automatiquement une résistance constante.\par
\clearpage
\essSection{3}{Puissance électrique}
En courant continu, la puissance électrique reçue par un dipôle vaut $P=UI$, avec P en W, U en V et I en A. La puissance est une énergie transférée par unité de durée. Une puissance nominale décrit un fonctionnement prévu dans les conditions indiquées. Pour un conducteur ohmique, combiner U = RI et P = UI permet de vérifier les risques d’échauffement.\par
\pcFigure{figures/b14-montage.pdf}{Mesurer U et I avant de calculer P ; ne pas dépasser les caractéristiques du matériel.}
\essSection{4}{Vérifier un choix}
Une résistance possède une puissance maximale admissible donnée par le fabricant. Un calcul cohérent de R ne garantit pas que l’échauffement soit acceptable. Les montages du collège utilisent une alimentation scolaire adaptée ; les installations domestiques demandent d’autres modèles et des protections. P = UI n’est utilisé ici qu’en courant continu ou dans les cas explicitement simplifiés par l’énoncé.\par
\begin{pcEncadre}{method}{Un exemple expliqué}Une résistance de $200\unite\Omega$ reçoit $6\unite V$. $I=6/200=0{,}030\unite A=30\unite{mA}$. Sa puissance vaut $P=6\times0{,}030=0{,}18\unite W$. Elle reste inférieure à une limite fournie de 0,25 W ; le professeur vérifie aussi les autres contraintes du montage.\end{pcEncadre}
\pcPage{Expérimenter et retenir}{Loi d’Ohm et puissance électrique}
\phantomsection\label{activite-14}
\begin{pcEncadre}{method}{Investigation en classe}Relever U et I pour une résistance avec une alimentation scolaire réglable, sans dépasser les limites données. Faire un nuage de points U(I) avec I en A. Comparer les quotients hors origine et l’alignement, puis calculer la puissance au point le plus élevé. Signaler toute dérive due à l’échauffement.\end{pcEncadre}
\pcSousTitre{Ce que je dois savoir expliquer}\begin{questions}
\item Reconnaître et exploiter une proportionnalité U/I.
\item Écrire et convertir avant de calculer.
\item Comparer la puissance calculée à une limite fournie.
\end{questions}
\begin{pcEncadre}{attention}{Une erreur à éviter}Si l’axe est en mA, sa valeur numérique n’est pas en A : diviser directement 6 V par 30 donne une erreur d’un facteur 1 000 sur la résistance.\end{pcEncadre}
\pcRepere{cor-activite-14}{Pistes pour l’investigation}\pcGrille
\clearpage\markright{Je démarre}
\rubrique{Je démarre}
\par\noindent\begin{minipage}{\linewidth}
\begin{exercice}[1]{Les symboles}
\begin{questions}
\item Dans U = RI, que désignent U et R ?
\item Donner leurs unités.
\end{questions}
\end{exercice}
\end{minipage}
\par\noindent\begin{minipage}{\linewidth}
\begin{exercice}[1]{Convertir}
\begin{questions}
\item Convertir $2{,}2\unite{k\Omega}$ en $\Omega$.
\item Convertir 15 mA en A.
\end{questions}
\end{exercice}
\end{minipage}
\par\noindent\begin{minipage}{\linewidth}
\begin{exercice}[1]{Trouver U}
\begin{questions}
\item Calculer U pour $R=100\unite\Omega$ et $I=0{,}020\unite A$.
\item Donner la relation utilisée.
\end{questions}
\end{exercice}
\end{minipage}
\par\noindent\begin{minipage}{\linewidth}
\begin{exercice}[1]{Trouver I}
\begin{questions}
\item Calculer I sous 6 V pour $R=300\unite\Omega$.
\item Exprimer en mA.
\end{questions}
\end{exercice}
\end{minipage}
\par\noindent\begin{minipage}{\linewidth}
\begin{exercice}[1]{Trouver R}
\begin{questions}
\item Un dipôle ohmique reçoit 3 V sous 10 mA. Calculer R.
\item Peut-on trouver R par U/I à l’origine où U = I = 0 ?
\end{questions}
\end{exercice}
\end{minipage}
\clearpage\markright{Je m’entraîne}
\rubrique{Je m’entraîne}
\par\noindent\begin{minipage}{\linewidth}
\begin{exercice}[2]{Une puissance}
\begin{questions}
\item Une lampe en continu reçoit 12 V sous 0,50 A. Calculer P.
\item Quelle énergie reçoit-elle chaque seconde ?
\end{questions}
\end{exercice}
\end{minipage}
\par\noindent\begin{minipage}{\linewidth}
\begin{exercice}[2]{Retrouver l’intensité}
\begin{questions}
\item Un appareil continu de 24 W fonctionne sous 12 V. Trouver I.
\item Un second appareil de 12 W sous 12 V demande quelle intensité ?
\end{questions}
\end{exercice}
\end{minipage}
\par\noindent\begin{minipage}{\linewidth}
\begin{exercice}[2]{Une caractéristique}
\pcFigureEx{figures/rev-ex-ohm-100.pdf}{}
\begin{questions}
\item On relève (I en A ; U en V) : (0,01 ; 1), (0,02 ; 2), (0,03 ; 3). Trouver R.
\item Quel contrôle graphique effectuer ?
\end{questions}
\end{exercice}
\end{minipage}
\par\noindent\begin{minipage}{\linewidth}
\begin{exercice}[2]{Une limite dépassée}
\begin{questions}
\item Une résistance de $100\unite\Omega$ sous 10 V est limitée à 0,25 W. Calculer I puis P.
\item Le choix respecte-t-il la limite ?
\end{questions}
\end{exercice}
\end{minipage}
\par\noindent\begin{minipage}{\linewidth}
\begin{exercice}[2]{À tension constante}
\begin{questions}
\item R double sous la même tension. Que devient I ?
\item Que devient P = UI ?
\end{questions}
\end{exercice}
\end{minipage}
\clearpage\markright{Je raisonne}
\rubrique{Je raisonne}
\par\noindent\begin{minipage}{\linewidth}
\begin{exercice}[3]{À résistance constante}
\begin{questions}
\item U double et R reste constante. Comment évolue I ?
\item Comment évolue P ?
\end{questions}
\end{exercice}
\end{minipage}
\par\noindent\begin{minipage}{\linewidth}
\begin{exercice}[3]{Une lampe réelle}
\begin{questions}
\item Une lampe donne U/I = 20 Ω puis 35 Ω à deux fonctionnements. Est-elle ohmique à R constante sur cet intervalle ?
\item Quelle cause physique peut contribuer à cette variation ?
\end{questions}
\end{exercice}
\end{minipage}
\par\noindent\begin{minipage}{\linewidth}
\begin{exercice}[3]{Un graphique trompeur}
\pcFigureEx{figures/rev-ex-ohm-inverse.pdf}{Exemple avec la tension en abscisse et l’intensité en ordonnée.}
\begin{questions}
\item Un graphe porte U horizontalement et I verticalement. Sa pente vaut-elle R ?
\item Comment retrouver R sur un point hors origine ?
\end{questions}
\end{exercice}
\end{minipage}
\par\noindent\begin{minipage}{\linewidth}
\begin{exercice}[3]{Deux branches}
\begin{questions}
\item Deux résistances de 100 Ω et 200 Ω sont en dérivation sous 6 V. Trouver leurs courants.
\item Calculer la puissance totale reçue.
\end{questions}
\end{exercice}
\end{minipage}
\par\noindent\begin{minipage}{\linewidth}
\begin{exercice}[3]{Choisir un dipôle}
\begin{questions}
\item Un conducteur doit être parcouru par 20 mA sous 5 V. Trouver R et P.
\item Une limite de 0,125 W convient-elle dans le modèle ? Discuter la marge.
\end{questions}
\end{exercice}
\end{minipage}
\pcPage{Défi du chapitre}{Loi d’Ohm et puissance électrique}
\phantomsection\label{defi-ex-14}
\pcDefiTitre{14}{Une résistance pour le laboratoire}
\begin{questions}
\item On dispose de résistances de 100 Ω ou 220 Ω, chacune limitée à 0,25 W. La tension continue maximale est 6 V. Calculer I et P pour chaque choix.
\item Quel choix respecte la limite donnée ?
\item Décrire les mesures qui permettraient de contrôler le fonctionnement sans confondre les deux appareils.
\end{questions}
\pcGrille

\begin{ouverture}[college3]{15}{Puissance et énergie électrique}{Partie 3 · L’énergie, ses transferts et ses conversions}{Une lampe puissante consomme-t-elle toujours plus ?}
\objectif{Distinguer puissance et énergie.}
\objectif{Calculer avec E = P t et des unités cohérentes.}
\objectif{Comparer des usages sur une durée et un service identifiés.}
\prerequis{Proportionnalité, tension et intensité.}
\end{ouverture}
\pcPage{Le cours}{Puissance et énergie électrique}
\phantomsection\label{ess-15}
\essSection{1}{Puissance et énergie : deux grandeurs}
La puissance $P$ indique une énergie transférée par unité de durée. Son unité est le watt (W) : $1\unite W=1\unite{J/s}$. Une grande puissance n’implique pas une grande énergie totale si la durée est très courte. Il faut distinguer l’énergie disponible dans une réserve, la puissance de fonctionnement et la durée possible. Les relations sont exploitées avec des unités cohérentes et des hypothèses annoncées.\par
\pcFigure{figures/c17-puissance.pdf}{Une même énergie peut être transférée rapidement à forte puissance ou plus lentement à faible puissance.}
\essSection{2}{Relier énergie et durée}
Pour une puissance constante, $E=P\Delta t$. Avec $P$ en W et $\Delta t$ en s, $E$ est en J. Avec $P$ en W et la durée en h, l’énergie est en Wh. $1\unite{Wh}=3\,600\unite J$ et $1\unite{kWh}=1\,000\unite{Wh}$. Ne pas écrire kW/h à la place de kWh : ce ne sont pas les mêmes grandeurs.\par
\clearpage
\essSection{3}{Relier puissance, tension et intensité}
En courant continu, pour un récepteur dans les conditions du modèle, la puissance électrique reçue vaut $P=UI$, avec U en V et I en A. Les indications nominales d’un appareil décrivent un fonctionnement prévu ; elles ne doivent pas être dépassées au hasard. Une résistance peut dissiper de l’énergie thermiquement et chauffer. La valeur maximale admissible du composant compte autant que sa résistance.\par
\pcFigure{figures/c17-energie.pdf}{Le graphe E en fonction du temps est une droite pour chaque puissance constante ; sa pente dépend de P.}
\essSection{4}{Comparer des usages}
Comparer une consommation demande la puissance et la durée d’utilisation, mais aussi un service rendu comparable. Un coût peut être estimé par $\text{énergie en kWh}\times\text{prix par kWh}$, si le prix est fourni ; les tarifs des exercices sont fictifs, pas des prix actuels. La veille, la fabrication et l’alimentation du numérique peuvent ajouter des dépenses hors du simple temps d’usage affiché.\par
\begin{pcEncadre}{method}{Un exemple expliqué}Une lampe de $8\unite W$ fonctionne $3\unite h$. $E=8\times3=24\unite{Wh}=0{,}024\unite{kWh}$. En joules : $24\times3\,600=86\,400\unite J$. Avec le tarif fictif de $0{,}25$ euro par kWh, le coût d’énergie seul est $0{,}024\times0{,}25=0{,}006$ euro ; il ne représente pas le prix total d’un abonnement.\end{pcEncadre}
\pcPage{Expérimenter et retenir}{Puissance et énergie électrique}
\phantomsection\label{activite-15}
\begin{pcEncadre}{method}{Investigation en classe}\textbf{Mission documentaire.} Comparer deux éclairages assurant un service annoncé équivalent. Relever les puissances et fixer une même durée d’usage. Calculer les énergies en Wh puis kWh. Étudier un deuxième scénario où les durées diffèrent. \textbf{À rendre :} hypothèses, tableau de calcul, unités et conclusion. Si un matériel scolaire permet une mesure, le professeur choisit une alimentation protégée et un dispositif compatible ; aucun appareil domestique n’est démonté.\end{pcEncadre}
\pcSousTitre{Ce que je dois savoir expliquer}\begin{questions}
\item La puissance s’exprime en W et l’énergie en J, Wh ou kWh.
\item À puissance constante, E = P t.
\item Les unités du produit déterminent celles de l’énergie.
\item En courant continu, P = U I dans les conditions étudiées.
\item Une comparaison équitable précise durée et service rendu.
\end{questions}
\begin{pcEncadre}{attention}{Une erreur à éviter}« 2 kW pendant une demi-heure font 4 kWh. » Il faut multiplier : $2\times0{,}5=1\unite{kWh}$.\end{pcEncadre}
\pcRepere{cor-activite-15}{Pistes pour l’investigation}\pcGrille
\clearpage\markright{Je démarre}
\rubrique{Je démarre}
\par\noindent\begin{minipage}{\linewidth}
\begin{exercice}[1]{Deux grandeurs}
\begin{questions}
\item Une étiquette indique 10 W. S’agit-il d’une énergie ou d’une puissance ?
\item Quelle unité conviendrait à une énergie ?
\end{questions}
\end{exercice}
\end{minipage}
\par\noindent\begin{minipage}{\linewidth}
\begin{exercice}[1]{Une heure}
\begin{questions}
\item Une lampe de $5\unite W$ fonctionne $1\unite h$. Quelle énergie en Wh ?
\item En joules ?
\end{questions}
\end{exercice}
\end{minipage}
\par\noindent\begin{minipage}{\linewidth}
\begin{exercice}[1]{Convertir}
\begin{questions}
\item Convertir $1{,}5\unite{kWh}$ en Wh.
\item Convertir $30\unite{min}$ en h.
\end{questions}
\end{exercice}
\end{minipage}
\par\noindent\begin{minipage}{\linewidth}
\begin{exercice}[1]{Un produit électrique}
\begin{questions}
\item Un récepteur reçoit $6\unite V$ sous $0{,}20\unite A$. Calculer sa puissance.
\item Donner l’unité.
\end{questions}
\end{exercice}
\end{minipage}
\par\noindent\begin{minipage}{\linewidth}
\begin{exercice}[1]{La durée compte}
\begin{questions}
\item A fonctionne à 20 W pendant 1 h ; B à 10 W pendant 1 h. Lequel reçoit le plus d’énergie ?
\item L’égalité des durées est-elle utile ici ?
\end{questions}
\end{exercice}
\end{minipage}
\clearpage\markright{Je m’entraîne}
\rubrique{Je m’entraîne}
\par\noindent\begin{minipage}{\linewidth}
\begin{exercice}[2]{En secondes}
\begin{questions}
\item Un moteur de $3\unite W$ fonctionne $40\unite s$. Calculer E.
\item Pourquoi la réponse n’est-elle pas en Wh ?
\end{questions}
\end{exercice}
\end{minipage}
\par\noindent\begin{minipage}{\linewidth}
\begin{exercice}[2]{Une demi-heure}
\begin{questions}
\item Un appareil de $60\unite W$ fonctionne $30\unite{min}$. Calculer l’énergie en Wh et kWh.
\item Calculer en J.
\end{questions}
\end{exercice}
\end{minipage}
\par\noindent\begin{minipage}{\linewidth}
\begin{exercice}[2]{Retrouver une durée}
\begin{questions}
\item Une source transfère $180\unite J$ à puissance constante de $6\unite W$. Quelle durée ?
\item Vérifier.
\end{questions}
\end{exercice}
\end{minipage}
\par\noindent\begin{minipage}{\linewidth}
\begin{exercice}[2]{Une intensité nominale}
\begin{questions}
\item Un moteur scolaire de $2{,}4\unite W$ fonctionne sous $6\unite V$ en courant continu. Calculer I dans ce modèle.
\item Exprimer en mA.
\end{questions}
\end{exercice}
\end{minipage}
\par\noindent\begin{minipage}{\linewidth}
\begin{exercice}[2]{Un coût fictif}
\begin{questions}
\item Une énergie de $2{,}5\unite{kWh}$ est facturée au prix fictif de $0{,}20$ euro/kWh. Quel coût d’énergie ?
\item Cela inclut-il forcément l’abonnement ?
\end{questions}
\end{exercice}
\end{minipage}
\clearpage\markright{Je raisonne}
\rubrique{Je raisonne}
\par\noindent\begin{minipage}{\linewidth}
\begin{exercice}[3]{Puissante mais brève}
\begin{questions}
\item A reçoit $100\unite W$ pendant $1\unite{min}$ ; B $10\unite W$ pendant $20\unite{min}$. Comparer les énergies.
\item Pourquoi la puissance seule ne suffit-elle pas ?
\end{questions}
\end{exercice}
\end{minipage}
\par\noindent\begin{minipage}{\linewidth}
\begin{exercice}[3]{Une veille}
\begin{questions}
\item Un dispositif consomme $2\unite W$ en veille pendant $20\unite h$ par jour. Calculer l’énergie quotidienne puis pour 30 jours.
\item Que manque-t-il pour l’énergie totale de l’appareil ?
\end{questions}
\end{exercice}
\end{minipage}
\par\noindent\begin{minipage}{\linewidth}
\begin{exercice}[3]{Une résistance qui chauffe}
\begin{questions}
\item Une résistance de $100\unite\Omega$ est sous $5\unite V$. Trouver I puis P.
\item Une puissance admissible de 0,125 W suffirait-elle ?
\end{questions}
\end{exercice}
\end{minipage}
\par\noindent\begin{minipage}{\linewidth}
\begin{exercice}[3]{Même service}
\begin{questions}
\item Pour un éclairage annoncé équivalent, A vaut 40 W et B 8 W. Calculer l’économie sur 100 h en kWh.
\item Quelle limite ajouter à la conclusion ?
\end{questions}
\end{exercice}
\end{minipage}
\par\noindent\begin{minipage}{\linewidth}
\begin{exercice}[3]{Une unité fausse}
\begin{questions}
\item Une copie écrit « E = 2 kW/h après deux heures ». Identifier le problème.
\item Donner le résultat pour une puissance de 2 kW pendant 2 h.
\end{questions}
\end{exercice}
\end{minipage}
\pcPage{Défi du chapitre}{Puissance et énergie électrique}
\phantomsection\label{defi-ex-15}
\pcDefiTitre{15}{Le projet d’éclairage}
\begin{questions}
\item Dix lampes de $12\unite W$ fonctionnent $4\unite h$ par jour pendant 20 jours. Calculer l’énergie totale.
\item On les remplace par des lampes de 8 W assurant le même service. Calculer l’économie.
\item À un prix fictif de 0,25 euro/kWh, quel gain sur le seul poste énergie ?
\end{questions}
\pcGrille

\begin{ouverture}[college4]{16}{Lumière : propagation et distances}{Partie 4 · Des signaux pour observer et communiquer}{Voir loin, est-ce voir le passé ?}
\objectif{Distinguer source primaire, diffusion et réception.}
\objectif{Exploiter le modèle du rayon et construire une ombre.}
\objectif{Relier distance, durée et vitesse de la lumière.}
\prerequis{Trajets, vitesse moyenne et organisation du Système solaire.}
\end{ouverture}
\pcPage{Le cours}{Lumière : propagation et distances}
\phantomsection\label{ess-16}
\essSection{1}{De la source à l’œil}
Une source primaire émet de la lumière ; un objet diffusant renvoie dans plusieurs directions une partie de la lumière reçue. Pour voir un objet, de la lumière issue de cet objet doit entrer dans l’œil. L’œil ne projette pas un rayon pour « toucher » l’objet. Un écran peut recevoir de la lumière sans la détecter comme un capteur électronique ; on précise le rôle de chaque système.\par
\pcFigure{figures/c18-ombre.pdf}{Pour une source ponctuelle, les rayons limites construisent une zone d’ombre derrière un objet opaque. Schéma géométrique simplifié.}
\essSection{2}{Un modèle de propagation}
Dans le vide et dans un milieu transparent homogène comme l’air calme à l’échelle du dispositif, la lumière se propage en ligne droite. Un rayon est une ligne orientée représentant sa direction et son sens. Il ne dessine pas la largeur réelle d’un faisceau. Un objet opaque peut intercepter certains trajets et créer une ombre. Les sources étendues peuvent aussi produire une pénombre.\par
\clearpage
\essSection{3}{Une vitesse immense mais finie}
Dans le vide, $c=299\,792\,458\unite{m/s}$ exactement ; pour nos calculs, $c\approx300\,000\unite{km/s}$. On utilise $d=c\Delta t$ ou $\Delta t=d/c$ en gardant les unités cohérentes. Dans un matériau transparent, la lumière peut se propager moins vite que dans le vide. Une distance astronomique peut donc correspondre à un délai observable, même si les délais du laboratoire sont minuscules.\par
\pcFigure{figures/c18-trajet.pdf}{Un trajet aller-retour mesure une distance totale égale à deux fois la séparation des objets.}
\essSection{4}{Année-lumière et observation}
Une année-lumière est la distance parcourue par la lumière dans le vide en un an : environ $9{,}46\times10^{12}\unite{km}$. Une étoile à quatre années-lumière est observée avec environ quatre ans de retard sur l’émission. Ce délai ne signifie pas qu’elle est forcément éteinte aujourd’hui. Les observations du Soleil ou d’une source intense exigent des moyens spécialisés ; aucun regard direct n’est proposé dans ce livre.\par
\begin{pcEncadre}{method}{Un exemple expliqué}Pour une distance Terre-Soleil arrondie à $150$ millions de kilomètres, on calcule : \[\Delta t=\dfrac{150\,000\,000}{300\,000}=500\unite s,\] soit $8\unite{min}\,20\unite s$. On observe donc le Soleil avec ce délai approximatif. Le calcul utilise une distance et une vitesse arrondies, sans prédire toutes les variations orbitales.\end{pcEncadre}
\pcPage{Expérimenter et retenir}{Lumière : propagation et distances}
\phantomsection\label{activite-16}
\begin{pcEncadre}{method}{Investigation en classe}\textbf{Question.} Le modèle rectiligne prédit-il l’ombre ? Avec une petite lampe scolaire, un objet opaque et un écran, tracer d’abord les rayons limites sur un schéma. Mesurer les distances relatives puis comparer la prévision à l’ombre observée. Déplacer l’écran et décrire l’effet. Une source étendue donne des bords moins nets : discuter la limite de l’approximation ponctuelle. \textbf{À rendre :} rayons orientés, zones éclairées et zone d’ombre. Ni laser dirigé vers les yeux ni observation du Soleil.\end{pcEncadre}
\pcSousTitre{Ce que je dois savoir expliquer}\begin{questions}
\item Un objet est vu lorsque sa lumière entre dans l’œil.
\item Un rayon est un modèle orienté de propagation.
\item L’ombre résulte de trajets interceptés par un objet opaque.
\item La vitesse de la lumière dans le vide est finie.
\item L’année-lumière est une unité de distance.
\end{questions}
\begin{pcEncadre}{attention}{Une erreur à éviter}Pour un signal réfléchi, la durée mesurée correspond souvent à l’aller et au retour : ne pas attribuer tout $c\Delta t$ à la seule distance source-obstacle.\end{pcEncadre}
\pcRepere{cor-activite-16}{Pistes pour l’investigation}\pcGrille
\clearpage\markright{Je démarre}
\rubrique{Je démarre}
\par\noindent\begin{minipage}{\linewidth}
\begin{exercice}[1]{Source ou diffuseur}
\begin{questions}
\item Une lampe allumée et une feuille éclairée ont-elles le même rôle lumineux ?
\item Lequel est une source primaire ?
\end{questions}
\end{exercice}
\end{minipage}
\par\noindent\begin{minipage}{\linewidth}
\begin{exercice}[1]{Le bon sens}
\pcFigureEx{figures/ex-vision.pdf}{Orienter le segment entre feuille et œil.}
\begin{questions}
\item Dans un dessin de vision d’une feuille, orienter le rayon entre feuille et œil.
\item L’œil émet-il ce rayon pour voir ?
\end{questions}
\end{exercice}
\end{minipage}
\par\noindent\begin{minipage}{\linewidth}
\begin{exercice}[1]{La trajectoire}
\begin{questions}
\item Comment se propage la lumière dans un milieu transparent homogène, selon le modèle ?
\item Un rayon possède-t-il un sens ?
\end{questions}
\end{exercice}
\end{minipage}
\par\noindent\begin{minipage}{\linewidth}
\begin{exercice}[1]{Une grande vitesse}
\begin{questions}
\item Quelle valeur arrondie de c utilise-t-on en km/s ?
\item Est-ce une vitesse infinie ?
\end{questions}
\end{exercice}
\end{minipage}
\par\noindent\begin{minipage}{\linewidth}
\begin{exercice}[1]{La zone d’ombre}
\begin{questions}
\item Qu’est-ce qui crée une ombre derrière un objet opaque ?
\item Un milieu transparent crée-t-il toujours la même ombre ?
\end{questions}
\end{exercice}
\end{minipage}
\clearpage\markright{Je m’entraîne}
\rubrique{Je m’entraîne}
\par\noindent\begin{minipage}{\linewidth}
\begin{exercice}[2]{Un trajet en une seconde}
\begin{questions}
\item Avec $c=300\,000\unite{km/s}$, quelle distance en $2\unite s$ ?
\item Et en $0{,}10\unite s$ ?
\end{questions}
\end{exercice}
\end{minipage}
\par\noindent\begin{minipage}{\linewidth}
\begin{exercice}[2]{Une durée solaire}
\begin{questions}
\item Calculer le délai sur $150$ millions de kilomètres avec la vitesse du cours.
\item Exprimer en minutes et secondes.
\end{questions}
\end{exercice}
\end{minipage}
\par\noindent\begin{minipage}{\linewidth}
\begin{exercice}[2]{Un aller simple lunaire}
\begin{questions}
\item Pour $384\,000\unite{km}$, calculer la durée aller de la lumière.
\item Calculer l’aller-retour idéal.
\end{questions}
\end{exercice}
\end{minipage}
\par\noindent\begin{minipage}{\linewidth}
\begin{exercice}[2]{Une étoile}
\begin{questions}
\item Une étoile est située à 12 années-lumière. Quel retard approximatif a sa lumière ?
\item Cela donne-t-il son état actuel ?
\end{questions}
\end{exercice}
\end{minipage}
\par\noindent\begin{minipage}{\linewidth}
\begin{exercice}[2]{Une unité à convertir}
\begin{questions}
\item On prend $1\unite{al}=9{,}46\times10^{12}\unite{km}$. Quelle distance pour 2 al ?
\item Quelle durée de propagation dans le vide ?
\end{questions}
\end{exercice}
\end{minipage}
\clearpage\markright{Je raisonne}
\rubrique{Je raisonne}
\par\noindent\begin{minipage}{\linewidth}
\begin{exercice}[3]{Une mesure réfléchie}
\pcFigureEx{figures/ex-echo-lumineux.pdf}{}
\begin{questions}
\item Un signal lumineux revient après $0{,}004\unite s$. Avec $c=300\,000\unite{km/s}$, calculer la distance de l’obstacle.
\item Pourquoi diviser par deux ?
\end{questions}
\end{exercice}
\end{minipage}
\par\noindent\begin{minipage}{\linewidth}
\begin{exercice}[3]{Une ombre qui grandit}
\pcFigureEx{figures/ex-ecrans.pdf}{Même source et même objet ; deux positions possibles de l’écran.}
\begin{questions}
\item Avec une petite source et un objet fixes, on éloigne l’écran derrière l’objet. L’ombre géométrique grandit-elle dans le modèle ?
\item Quelle construction justifie la réponse ?
\end{questions}
\end{exercice}
\end{minipage}
\clearpage
\par\noindent\begin{minipage}{\linewidth}
\begin{exercice}[3]{Le rayon invisible}
\begin{questions}
\item Dans l’air très propre, on voit une tache sur l’écran mais pas forcément le trajet du faisceau de côté. Pourquoi ?
\item Que change une légère diffusion ?
\end{questions}
\end{exercice}
\end{minipage}
\par\noindent\begin{minipage}{\linewidth}
\begin{exercice}[3]{Deux échelles}
\begin{questions}
\item Un laboratoire mesure une distance de $3\unite m$. Avec $c\approx3\times10^8\unite{m/s}$, calculer le délai aller.
\item Pourquoi un chronomètre manuel ne convient-il pas ?
\end{questions}
\end{exercice}
\end{minipage}
\par\noindent\begin{minipage}{\linewidth}
\begin{exercice}[3]{Une conclusion médiatique}
\begin{questions}
\item Un titre dit « Cette étoile est loin, donc elle est déjà morte ». Est-ce démontré ?
\item Quelle conclusion est justifiée par la distance seule ?
\end{questions}
\end{exercice}
\end{minipage}
\pcPage{Défi du chapitre}{Lumière : propagation et distances}
\phantomsection\label{defi-ex-16}
\pcDefiTitre{16}{Le miroir spatial}
\begin{questions}
\item Un dispositif reçoit un écho lumineux après $2{,}4\unite s$. Estimer la distance au réflecteur avec $c=300\,000\unite{km/s}$.
\item Un élève annonce 720 000 km. Expliquer son erreur.
\item Si le délai augmente, quelle conclusion est possible sous les mêmes hypothèses ?
\end{questions}
\pcGrille

\begin{ouverture}[college4]{17}{Le son : fréquence, propagation et écho}{Partie 4 · Des signaux pour observer et communiquer}{Pourquoi voit-on l’éclair avant d’entendre le tonnerre ?}
\objectif{Décrire source vibrante et milieu de propagation.}
\objectif{Distinguer fréquence, hauteur et niveau sonore.}
\objectif{Calculer une distance avec un délai de propagation ou un écho.}
\prerequis{Distance, durée, vitesse et lecture de graphiques.}
\end{ouverture}
\pcPage{Le cours}{Le son : fréquence, propagation et écho}
\phantomsection\label{ess-17}
\essSection{1}{Un son a besoin d’un milieu}
Une source sonore vibre et perturbe un milieu matériel. La perturbation se transmet dans un gaz, un liquide ou un solide ; le son ne se propage pas dans le vide. Les particules du milieu oscillent localement : elles ne voyagent pas toutes de la source à l’oreille. Dans l’air vers $20\,^{\circ}\mathrm C$, on utilise ici $v\approx340\unite{m/s}$. La valeur dépend du milieu et de la température.\par
\pcFigure{figures/c19-signaux.pdf}{Deux signaux périodiques de même amplitude schématique : celui qui présente davantage de motifs par seconde a une fréquence plus élevée.}
\essSection{2}{Fréquence et hauteur}
La fréquence $f$ indique le nombre de motifs périodiques par seconde, en hertz (Hz). Une fréquence plus élevée correspond à un son plus aigu pour les sons simples comparés. L’audition humaine jeune couvre approximativement 20 Hz à 20 000 Hz, avec de fortes variations individuelles. En dessous : infrasons ; au-dessus : ultrasons. Ces mots ne désignent pas des niveaux sonores plus ou moins forts.\par
\clearpage
\essSection{3}{Distance et temps de propagation}
Pour un trajet simple, $d=v\Delta t$. Pour un écho, le son parcourt l’aller et le retour : $2d=v\Delta t$. Le délai entre éclair et tonnerre donne une estimation de distance parce que la lumière arrive bien plus vite que le son. On suppose les deux émissions proches dans le temps et on néglige le délai lumineux pour cette estimation.\par
\pcFigure{figures/c19-echo.pdf}{Un émetteur-récepteur chronomètre le retour ; la distance parcourue par le signal est 2d.}
\essSection{4}{Préserver son audition}
Le niveau sonore, souvent exprimé en décibels (dB), décrit une autre propriété que la fréquence. Le risque auditif dépend notamment du niveau, de la durée et de la répétition des expositions. Une gêne n’est pas un instrument de mesure et l’absence de douleur ne prouve pas l’absence de risque. Réduire le niveau, s’éloigner et limiter la durée sont des moyens de prévention. Aucun essai du livre ne cherche un niveau fort.\par
\begin{pcEncadre}{method}{Un exemple expliqué}Un écho revient $0{,}40\unite s$ après l’émission dans l’air. Le trajet total vaut $340\times0{,}40=136\unite m$. La paroi est à $136/2=68\unite m$. Pour un tonnerre entendu 3 secondes après l’éclair, la distance estimée est au contraire $340\times3=1\,020\unite m$ : on ne divise pas par deux, car c’est un trajet simple.\end{pcEncadre}
\pcPage{Expérimenter et retenir}{Le son : fréquence, propagation et écho}
\phantomsection\label{activite-17}
\begin{pcEncadre}{method}{Investigation en classe}\textbf{Question.} Comment distinguer hauteur et niveau ? Le professeur fournit deux enregistrements modérés ou des tracés, l’un avec davantage de motifs sur une même durée, l’autre avec une amplitude différente. Comparer sans augmenter le volume. Compter les motifs sur une durée donnée et calculer une fréquence. \textbf{Étude de documents :} comparer un délai de tonnerre à un délai d’écho ; dessiner le trajet avant de choisir la formule. \textbf{À rendre :} fréquence, description de la hauteur et distinction avec le niveau.\end{pcEncadre}
\pcSousTitre{Ce que je dois savoir expliquer}\begin{questions}
\item Le son nécessite un milieu matériel ; il ne traverse pas le vide.
\item Une source sonore met un milieu en vibration.
\item La fréquence en Hz caractérise la répétition temporelle.
\item La hauteur et le niveau sonore sont deux propriétés différentes.
\item Un écho correspond à un trajet aller-retour.
\end{questions}
\begin{pcEncadre}{attention}{Une erreur à éviter}Ne pas dire « ultrason = son très fort ». Un ultrason est défini par sa fréquence supérieure à la plage audible humaine usuelle.\end{pcEncadre}
\pcRepere{cor-activite-17}{Pistes pour l’investigation}\pcGrille
\clearpage\markright{Je démarre}
\rubrique{Je démarre}
\par\noindent\begin{minipage}{\linewidth}
\begin{exercice}[1]{Le milieu}
\begin{questions}
\item Le son peut-il se propager dans le vide ?
\item Peut-il traverser de l’eau ?
\end{questions}
\end{exercice}
\end{minipage}
\par\noindent\begin{minipage}{\linewidth}
\begin{exercice}[1]{La source}
\begin{questions}
\item Quelle propriété mécanique permet à une source de produire un son ?
\item Les molécules d’air vont-elles toutes jusqu’à l’oreille ?
\end{questions}
\end{exercice}
\end{minipage}
\par\noindent\begin{minipage}{\linewidth}
\begin{exercice}[1]{L’unité}
\begin{questions}
\item Quelle unité pour la fréquence ?
\item Que signifie une fréquence de 100 Hz pour un signal périodique ?
\end{questions}
\end{exercice}
\end{minipage}
\par\noindent\begin{minipage}{\linewidth}
\begin{exercice}[1]{Trois domaines}
\begin{questions}
\item Classer $10\unite{Hz}$, $1\,000\unite{Hz}$ et $40\unite{kHz}$.
\item Cette classification concerne-t-elle le niveau ?
\end{questions}
\end{exercice}
\end{minipage}
\par\noindent\begin{minipage}{\linewidth}
\begin{exercice}[1]{Le plus aigu}
\begin{questions}
\item Comparer deux sons simples de 200 et 800 Hz. Lequel est le plus aigu ?
\item Est-il nécessairement plus fort ?
\end{questions}
\end{exercice}
\end{minipage}
\clearpage\markright{Je m’entraîne}
\rubrique{Je m’entraîne}
\par\noindent\begin{minipage}{\linewidth}
\begin{exercice}[2]{Une distance simple}
\begin{questions}
\item Dans l’air, un son met $2\unite s$ pour atteindre un auditeur. Calculer la distance avec $340\unite{m/s}$.
\item S’agit-il d’un aller-retour ?
\end{questions}
\end{exercice}
\end{minipage}
\par\noindent\begin{minipage}{\linewidth}
\begin{exercice}[2]{Un écho}
\pcFigureEx{figures/ex-echo-sonore.pdf}{}
\begin{questions}
\item Un écho revient après $0{,}50\unite s$. Calculer la distance à la paroi.
\item Quelle longueur a parcourue le son ?
\end{questions}
\end{exercice}
\end{minipage}
\par\noindent\begin{minipage}{\linewidth}
\begin{exercice}[2]{Compter des motifs}
\pcFigureEx{figures/ex-huit-motifs.pdf}{}
\begin{questions}
\item Un tracé comporte $8$ motifs en $0{,}020\unite s$. Calculer f.
\item Comparer à 200 Hz.
\end{questions}
\end{exercice}
\end{minipage}
\clearpage
\par\noindent\begin{minipage}{\linewidth}
\begin{exercice}[2]{Le tonnerre}
\begin{questions}
\item Un délai de $5\unite s$ sépare éclair et tonnerre. Estimer la distance.
\item Quelle durée a-t-on négligée ?
\end{questions}
\end{exercice}
\end{minipage}
\par\noindent\begin{minipage}{\linewidth}
\begin{exercice}[2]{Dans l’eau}
\pcFigureEx{figures/ex-sonar.pdf}{}
\begin{questions}
\item Un sonar reçoit un écho après $0{,}20\unite s$. On fournit $v=1\,500\unite{m/s}$ dans l’eau. Trouver la profondeur.
\item Pourquoi ne pas utiliser 340 m/s ?
\end{questions}
\end{exercice}
\end{minipage}
\clearpage\markright{Je raisonne}
\rubrique{Je raisonne}
\par\noindent\begin{minipage}{\linewidth}
\begin{exercice}[3]{Un signal plus grand}
\pcFigureEx{figures/ex-meme-periode.pdf}{}
\begin{questions}
\item Deux tracés montrent la même période, mais des amplitudes différentes sur la même échelle. Leur fréquence diffère-t-elle ?
\item Que peut-on dire de leur hauteur si ce sont des sons simples ?
\end{questions}
\end{exercice}
\end{minipage}
\par\noindent\begin{minipage}{\linewidth}
\begin{exercice}[3]{Une information incomplète}
\begin{questions}
\item Un document indique un niveau de 80 dB sans durée ni conditions. Peut-on conclure seul sur tout le risque auditif ?
\item Quels renseignements ajouter ?
\end{questions}
\end{exercice}
\end{minipage}
\par\noindent\begin{minipage}{\linewidth}
\begin{exercice}[3]{Deux trajets confondus}
\begin{questions}
\item Un élève calcule $340\times0{,}30=102\unite m$ comme distance d’une paroi pour un écho. Corriger.
\item Le même produit serait-il pertinent pour un trajet simple de 0,30 s ?
\end{questions}
\end{exercice}
\end{minipage}
\par\noindent\begin{minipage}{\linewidth}
\begin{exercice}[3]{Un délai instrumental}
\begin{questions}
\item Un dispositif ajoute $0{,}010\unite s$ au délai d’écho affiché de $0{,}210\unite s$. Calculer la distance dans l’air après correction.
\item Pourquoi corriger avant de diviser par deux ?
\end{questions}
\end{exercice}
\end{minipage}
\par\noindent\begin{minipage}{\linewidth}
\begin{exercice}[3]{Un contrôle de vraisemblance}
\begin{questions}
\item Un capteur dans l’air annonce un obstacle à 3,4 m avec un écho de 2 s. Est-ce cohérent avec 340 m/s ?
\item Quel délai attendrait-on pour 3,4 m ?
\end{questions}
\end{exercice}
\end{minipage}
\pcPage{Défi du chapitre}{Le son : fréquence, propagation et écho}
\phantomsection\label{defi-ex-17}
\pcDefiTitre{17}{Deux capteurs, deux milieux}
\begin{questions}
\item Un capteur A mesure un écho de $0{,}10\unite s$ dans l’air ; B un écho de même durée dans l’eau. Utiliser respectivement 340 et 1 500 m/s. Calculer les distances.
\item Pourquoi un même délai ne donne-t-il pas la même distance ?
\item Quelle donnée faut-il toujours annoncer dans un compte rendu ?
\end{questions}
\pcGrille

\begin{ouverture}[college4]{18}{Signaux et information}{Partie 4 · Des signaux pour observer et communiquer}{Comment un signal devient-il un message ?}
\objectif{Identifier émetteur, support de propagation et récepteur.}
\objectif{Distinguer information et grandeur physique du signal.}
\objectif{Lire un codage et discuter la qualité d’une mesure.}
\prerequis{Sources lumineuses et sonores, graphique et proportionnalité.}
\end{ouverture}
\pcPage{Le cours}{Signaux et information}
\phantomsection\label{ess-18}
\essSection{1}{Une chaîne de communication}
Un émetteur produit un signal ; un récepteur le détecte et permet d’en extraire une information grâce à un code ou un traitement. Un signal peut être lumineux, sonore, électrique ou radio. Les ondes radio appartiennent à la famille des rayonnements électromagnétiques : elles ne sont pas des sons et peuvent se propager dans le vide. Le signal est un phénomène physique ; l’information est le sens ou la donnée qu’on lui associe.\par
\pcFigure{figures/c20-chaine.pdf}{Une communication peut convertir plusieurs fois la nature du signal entre la source du message et son destinataire.}
\essSection{2}{Coder un message}
Un code établit une correspondance connue entre des signaux et des symboles. Deux états distincts permettent de coder 0 et 1 ; une suite de bits peut ensuite représenter plusieurs symboles. Un signal n’a pas de sens unique sans convention : un voyant rouge peut annoncer une alerte ou un état choisi par le concepteur. Le codage doit préciser les durées, l’ordre des symboles et leur séparation.\par
\clearpage
\essSection{3}{Mesurer avec un signal}
Un capteur transforme une grandeur physique en une sortie exploitable. Un étalonnage relie des valeurs connues de la grandeur aux valeurs du signal. Dans une plage donnée, une relation peut être proportionnelle ou affine, mais ce n’est pas automatique. Un signal analogique varie continûment dans son modèle ; un traitement numérique utilise des valeurs discrètes. Cette distinction ne signifie pas que le numérique est toujours exact.\par
\pcFigure{figures/c20-signaux.pdf}{À gauche : variation continue. À droite : états de code 0 ou 1 ; les valeurs et le sens du message dépendent de la convention.}
\essSection{4}{Détecter les limites}
Bruit, obstacle, distance, mauvaise orientation et perturbations peuvent dégrader la réception. Répéter un message ou ajouter une vérification aide parfois à repérer une erreur, sans la supprimer toujours. Pour une mesure, on contrôle étalonnage, plage, résolution et temps de réponse. Une donnée affichée avec beaucoup de chiffres ne garantit pas une grande exactitude. Expliquer la limite du dispositif fait partie d’une conclusion scientifique.\par
\begin{pcEncadre}{method}{Un exemple expliqué}Un capteur scolaire donne $U=0{,}020\,T+0{,}50$, avec U en V et T en degrés Celsius, entre 0 et 50 °C. Pour $U=1{,}10\unite V$, $T=(1{,}10-0{,}50)/0{,}020=30\,^{\circ}\mathrm C$. Ce calcul utilise l’étalonnage fourni ; il ne permet pas d’extrapoler à une température arbitraire.\end{pcEncadre}
\pcPage{Expérimenter et retenir}{Signaux et information}
\phantomsection\label{activite-18}
\begin{pcEncadre}{method}{Investigation en classe}\textbf{Mission.} Inventer un code lumineux à deux états pour quatre messages, puis le transmettre à faible distance avec une petite lampe adaptée. Échanger le code avant la transmission. Tester un message, puis introduire sur document une erreur de symbole. Proposer un moyen de détection et discuter sa limite. \textbf{À rendre :} table du code, chronologie des états, chaîne émetteur-récepteur et bilan de fiabilité. Aucun émetteur radio n’est fabriqué et aucune lumière intense n’est utilisée.\end{pcEncadre}
\pcSousTitre{Ce que je dois savoir expliquer}\begin{questions}
\item Le signal est un phénomène physique portant une information codée.
\item Un récepteur doit disposer d’un code ou d’un traitement.
\item Les ondes radio ne sont pas des ondes sonores.
\item Une mesure par capteur exige un étalonnage et une plage d’utilisation.
\item Les perturbations peuvent modifier le signal et l’information reçue.
\end{questions}
\begin{pcEncadre}{attention}{Une erreur à éviter}Un affichage numérique de 23,456 °C ne prouve pas que la température est connue au millième de degré : la qualité du capteur et de l’étalonnage limite le résultat.\end{pcEncadre}
\pcRepere{cor-activite-18}{Pistes pour l’investigation}\pcGrille
\clearpage\markright{Je démarre}
\rubrique{Je démarre}
\par\noindent\begin{minipage}{\linewidth}
\begin{exercice}[1]{Les rôles}
\begin{questions}
\item Dans une sonnette, identifier un émetteur sonore et un récepteur humain.
\item Quelle information peut être transmise ?
\end{questions}
\end{exercice}
\end{minipage}
\par\noindent\begin{minipage}{\linewidth}
\begin{exercice}[1]{Radio et son}
\begin{questions}
\item Une onde radio est-elle un son très aigu ?
\item Peut-elle traverser le vide ?
\end{questions}
\end{exercice}
\end{minipage}
\par\noindent\begin{minipage}{\linewidth}
\begin{exercice}[1]{Deux états}
\begin{questions}
\item Que signifie un code à deux états ?
\item Le voyant allumé signifie-t-il partout « danger » ?
\end{questions}
\end{exercice}
\end{minipage}
\par\noindent\begin{minipage}{\linewidth}
\begin{exercice}[1]{Le capteur}
\begin{questions}
\item Quel est le rôle d’un capteur ?
\item Pourquoi l’étalonner ?
\end{questions}
\end{exercice}
\end{minipage}
\par\noindent\begin{minipage}{\linewidth}
\begin{exercice}[1]{Une perturbation}
\begin{questions}
\item Nommer une cause de mauvaise réception lumineuse.
\item Et une cause de mauvaise réception sonore ?
\end{questions}
\end{exercice}
\end{minipage}
\clearpage\markright{Je m’entraîne}
\rubrique{Je m’entraîne}
\par\noindent\begin{minipage}{\linewidth}
\begin{exercice}[2]{Lire un code}
\begin{questions}
\item Le code fourni est A = 00, B = 01, C = 10, D = 11. Décoder 01 10 00.
\item Pourquoi les espaces sont-ils utiles ?
\end{questions}
\end{exercice}
\end{minipage}
\par\noindent\begin{minipage}{\linewidth}
\begin{exercice}[2]{Compter des messages}
\begin{questions}
\item Combien de suites différentes peut-on faire avec deux bits ?
\item Donner les suites.
\end{questions}
\end{exercice}
\end{minipage}
\par\noindent\begin{minipage}{\linewidth}
\begin{exercice}[2]{Un étalonnage simple}
\begin{questions}
\item Un capteur donne $U=0{,}10\,x$ avec U en V et x dans son unité. Calculer U pour x = 12.
\item Retrouver x pour U = 0,8 V.
\end{questions}
\end{exercice}
\end{minipage}
\par\noindent\begin{minipage}{\linewidth}
\begin{exercice}[2]{Une sortie affine}
\pcFigureEx{figures/ex-etalonnage.pdf}{}
\begin{questions}
\item Pour $U=0{,}020T+0{,}50$, calculer U à 20 °C.
\item La relation est-elle proportionnelle entre U et T ?
\end{questions}
\end{exercice}
\end{minipage}
\par\noindent\begin{minipage}{\linewidth}
\begin{exercice}[2]{Une valeur reçue}
\begin{questions}
\item Avec l’étalonnage précédent, retrouver T pour $U=1{,}30\unite V$.
\item La valeur est-elle dans la plage 0 à 50 °C ?
\end{questions}
\end{exercice}
\end{minipage}
\clearpage\markright{Je raisonne}
\rubrique{Je raisonne}
\par\noindent\begin{minipage}{\linewidth}
\begin{exercice}[3]{Une erreur de bit}
\begin{questions}
\item Le code de l’exercice 6 est utilisé. B = 01 est reçu 11. Quel message est lu ?
\item Le récepteur peut-il détecter l’erreur avec cette seule table ?
\end{questions}
\end{exercice}
\end{minipage}
\par\noindent\begin{minipage}{\linewidth}
\begin{exercice}[3]{Répéter et comparer}
\begin{questions}
\item Un message est reçu deux fois : 01 10 puis 01 11. Que montre la comparaison ?
\item Permet-elle de savoir quelle version est correcte ?
\end{questions}
\end{exercice}
\end{minipage}
\par\noindent\begin{minipage}{\linewidth}
\begin{exercice}[3]{Hors de la plage}
\pcFigureEx{figures/ex-hors-plage.pdf}{La zone colorée correspond au seul domaine étalonné.}
\begin{questions}
\item Un élève utilise $U=0{,}020T+0{,}50$ pour U = 3 V, alors que l’étalonnage vaut de 0 à 50 °C. Quel calcul obtient-il ?
\item Pourquoi ne doit-il pas annoncer ce résultat comme une mesure fiable ?
\end{questions}
\end{exercice}
\end{minipage}
\par\noindent\begin{minipage}{\linewidth}
\begin{exercice}[3]{Une résolution fournie}
\begin{questions}
\item Un capteur affiche la température par pas de $0{,}5\,^{\circ}\mathrm C$. Peut-il distinguer à l’affichage 20,1 et 20,2 °C ?
\item Cette résolution résume-t-elle toute la qualité de mesure ?
\end{questions}
\end{exercice}
\end{minipage}
\par\noindent\begin{minipage}{\linewidth}
\begin{exercice}[3]{Analogique ou numérique}
\begin{questions}
\item Une tension varie continûment puis un convertisseur donne des nombres arrondis. Décrire les deux représentations.
\item Les arrondis constituent-ils une erreur humaine ?
\end{questions}
\end{exercice}
\end{minipage}
\pcPage{Défi du chapitre}{Signaux et information}
\phantomsection\label{defi-ex-18}
\pcDefiTitre{18}{Le protocole du phare}
\begin{questions}
\item On code quatre messages par deux bits, avec 1 = lampe allumée et 0 = éteinte. Chaque bit dure 1 s. Combien de temps dure un message, hors séparation ?
\item Proposer un signal de début pour éviter de confondre silence et 00.
\item Pourquoi la seule répétition ne garantit-elle pas une correction parfaite ?
\end{questions}
\pcGrille

\pcAnnexe{Corrigés détaillés}{Comprendre les démarches et refaire sans aide}\phantomsection\label{corriges-debut}\pdfbookmark[0]{Corrigés détaillés}{solutions}
\begin{pcEncadre}{method}{Utiliser les corrigés}Cherche avant de lire. Compare le raisonnement et les unités, puis referme le livre et recommence. Les conclusions expérimentales restent limitées aux conditions du protocole.\end{pcEncadre}
\pcSousTitre{Retrouver les quatre thèmes}
\pcRepere{corchap-1}{Matière : mesures, modèles et transformations}
\pcRepere{corchap-8}{Mouvements : vitesses, interactions et poids}
\pcRepere{corchap-11}{Énergie : bilans et lois de l’électricité}
\pcRepere{corchap-16}{Signaux : lumière, son et information}
\pcSousTitre{Faire de l’erreur un outil}
\begin{questions}
\item Repérer la première étape différente de ta démarche : lecture, modèle, conversion ou calcul.
\item Expliquer l’erreur en une phrase et conserver une correction lisible dans le cahier.
\item Refaire sans aide une question analogue, puis revenir au problème initial après quelques jours.
\end{questions}
\begin{pcEncadre}{definition}{Une réponse complète}
La bonne valeur ne suffit pas : une relation justifiée, des unités cohérentes et une conclusion dans le contexte montrent ce que tu as compris. Une réponse expérimentale peut être nuancée si les mesures ne permettent pas de trancher.
\end{pcEncadre}
\corrigesChapitre{1}{Matière : mesurer, séparer, identifier}
\par\noindent\begin{minipage}{\linewidth}
\begin{corrige}{1}{Trois états}
\begin{questions}
\item Le solide, hors déformation imposée.
\item Un gaz se répartit dans tout le volume disponible.
\end{questions}
\end{corrige}
\end{minipage}
\par\noindent\begin{minipage}{\linewidth}
\begin{corrige}{2}{Une eau salée}
\begin{questions}
\item Mélange : eau et espèces dissoutes.
\item Non ; plusieurs espèces peuvent former une seule phase visible.
\end{questions}
\end{corrige}
\end{minipage}
\par\noindent\begin{minipage}{\linewidth}
\begin{corrige}{3}{Les bons mots}
\begin{questions}
\item Vaporisation.
\item Liquéfaction ; on rencontre aussi condensation liquide.
\end{questions}
\end{corrige}
\end{minipage}
\par\noindent\begin{minipage}{\linewidth}
\begin{corrige}{4}{Volumes}
\begin{questions}
\item $350\unite{mL}$.
\item $18\unite{mL}$.
\end{questions}
\end{corrige}
\end{minipage}
\par\noindent\begin{minipage}{\linewidth}
\begin{corrige}{5}{Le quotient}
\begin{questions}
\item $40/50=0{,}80\unite{g/mL}$.
\item $800\unite{kg/m^3}$.
\end{questions}
\end{corrige}
\end{minipage}
\par\noindent\begin{minipage}{\linewidth}
\begin{corrige}{6}{Solide immergé}
\begin{questions}
\item $V=43-28=15\unite{cm^3}$.
\item $\rho=40{,}5/15=2{,}70\unite{g/cm^3}$.
\end{questions}
\par\smallskip\noindent\textbf{Pour comprendre.}\enspace $V=43-28=15\unite{mL}=15\unite{cm^3}$. On calcule ensuite $\rho=40{,}5/15=2{,}70\unite{g/cm^3}$. Utiliser 43 mL compterait aussi l’eau initiale.\par
\pcFigureEx{figures/rev-cor-immersion.pdf}{Le solide déplace 15 mL ; on utilise ce volume avec la seule masse du solide.}
\end{corrige}
\end{minipage}
\par\noindent\begin{minipage}{\linewidth}
\begin{corrige}{7}{Une masse prévue}
\begin{questions}
\item $m=0{,}90\times250=225\unite g$.
\item Elle diminue puisque $m/V$ diminue.
\end{questions}
\end{corrige}
\end{minipage}
\par\noindent\begin{minipage}{\linewidth}
\begin{corrige}{8}{Retrouver un volume}
\begin{questions}
\item $V=178/8{,}9=20\unite{cm^3}$.
\item Le volume extérieur comprend la cavité ; ce n’est pas le seul volume de cuivre.
\end{questions}
\par\smallskip\noindent\textbf{Pour comprendre.}\enspace Partir de $m=\rho V$ et isoler $V=m/\rho=178/8{,}9=20\unite{cm^3}$. Le volume extérieur d’un objet creux ne se substitue pas au volume du matériau.\par
\pcFigureEx{figures/cor-coque.pdf}{Une cavité contribue au volume extérieur, mais pas au volume de cuivre.}
\end{corrige}
\end{minipage}
\par\noindent\begin{minipage}{\linewidth}
\begin{corrige}{9}{Une limite de dissolution}
\begin{questions}
\item Environ 36 g se dissolvent et 4 g restent solides à l’équilibre, dans ce modèle.
\item $36\times2{,}5=90\unite g$.
\end{questions}
\end{corrige}
\end{minipage}
\par\noindent\begin{minipage}{\linewidth}
\begin{corrige}{10}{Interpréter un palier}
\begin{questions}
\item L’énergie transférée sert notamment au changement d’état pendant la coexistence liquide-vapeur.
\item Non : de la vapeur quitte le système pesé ; inclure cette matière dans le bilan.
\end{questions}
\pcFigureEx{figures/rev-cor-ebullition.pdf}{La température reste constante pendant le changement d’état, à la pression annoncée.}
\end{corrige}
\end{minipage}
\par\noindent\begin{minipage}{\linewidth}
\begin{corrige}{11}{Choisir une séparation}
\begin{questions}
\item Filtrer le sable, puis distiller le filtrat avec un dispositif adapté.
\item Non : la potabilité nécessite des critères et des analyses supplémentaires.
\end{questions}
\pcFigureEx{figures/rev-cor-separation.pdf}{Filtration pour retirer le sable ; distillation du filtrat pour recueillir l’eau et laisser le sel non volatil.}
\end{corrige}
\end{minipage}
\par\noindent\begin{minipage}{\linewidth}
\begin{corrige}{12}{Une mesure discutable}
\begin{questions}
\item $54/25=2{,}16\unite{g/mL}$, trop faible.
\item Oui : la bulle peut augmenter le volume déplacé apparent.
\end{questions}
\par\smallskip\noindent\textbf{Pour comprendre.}\enspace Avec le volume apparent, $54/25=2{,}16\unite{g/cm^3}$. Avec le volume réel de 20 cm³, on aurait $54/20=2{,}70\unite{g/cm^3}$. Une mesure biaisée peut donc changer l’identification proposée.\par
\pcFigureEx{figures/cor-bulle.pdf}{Une bulle fait attribuer au solide un volume trop grand et une masse volumique trop petite.}
\end{corrige}
\end{minipage}
\par\noindent\begin{minipage}{\linewidth}
\begin{corrige}{13}{Identifier avec une plage}
\begin{questions}
\item La valeur 2,70 appartient à l’intervalle 2,63 à 2,73 : résultat compatible.
\item Non ; d’autres matériaux ou compositions peuvent convenir. Un second test serait utile.
\end{questions}
\pcFigureEx{figures/rev-cor-intervalle.pdf}{2,70 appartient à l’intervalle 2,63 à 2,73 g/cm³ ; compatibilité ne signifie pas identification certaine.}
\end{corrige}
\end{minipage}
\par\noindent\begin{minipage}{\linewidth}
\begin{corrige}{14}{Glace et océans}
\begin{questions}
\item Cette eau peut rejoindre l’océan depuis un réservoir auparavant sur les continents.
\item La dilatation thermique de l’eau de mer augmente son volume à masse donnée.
\end{questions}
\pcFigureEx{figures/cor-glaces.pdf}{Distinguer apport d’eau par les glaciers continentaux et dilatation de l’eau ; la banquise flottait déjà.}
\end{corrige}
\end{minipage}
\par\noindent\begin{minipage}{\linewidth}
\begin{corrige}{15}{Construire une preuve}
\begin{questions}
\item Mesurer leurs masses pour un volume connu et comparer leurs masses volumiques ; utiliser les protections prescrites.
\item Une différence de température ou de méthode peut produire un écart qui ne vient pas de la composition.
\end{questions}
\end{corrige}
\end{minipage}
\par\noindent\begin{minipage}{\linewidth}
\phantomsection\label{cor-defi-1}\pcSolutionDefi{1}{Le flacon inconnu}\begin{questions}
\item $m=109-45=64\unite g$ ; $\rho=64/80=0{,}80\unite{g/mL}$.
\item L’éthanol est le plus proche, sans preuve de pureté.
\item La température ou une lecture grossière peut expliquer l’écart. Répéter avec verrerie adaptée et comparer plusieurs propriétés.
\end{questions}
\par\smallskip\noindent\textbf{Pour comprendre.}\enspace La masse du liquide vaut $109-45=64\unite g$. Avec 80 mL, $\rho=64/80=0{,}80\unite{g/mL}$. La proximité avec 0,79 g/mL rend l’éthanol compatible ; d’autres liquides ou mélanges peuvent avoir une valeur proche.\par
\pcFigureEx{figures/rev-cor-tare.pdf}{Soustraire la masse du flacon avant de former le quotient masse/volume.}
\end{minipage}
\par\noindent\begin{minipage}{\linewidth}
\phantomsection\label{cor-activite-1}\pcSousTitre{Pistes pour l’investigation}
Une tare correcte soustrait la masse du récipient. Pour un liquide homogène à température stable, m est sensiblement proportionnelle à V ; les écarts doivent être comparés à la lecture. Pour un solide insoluble, mesurer la masse puis le volume déplacé avec immersion complète et sans bulle. Une identification par masse volumique est compatible avec une référence, sans prouver l’unicité ni la pureté.\par\end{minipage}

\corrigesChapitre{2}{Atomes, noyaux et ions}
\par\noindent\begin{minipage}{\linewidth}
\begin{corrige}{1}{Constituants}
\begin{questions}
\item Dans le noyau.
\item L’électron.
\end{questions}
\end{corrige}
\end{minipage}
\par\noindent\begin{minipage}{\linewidth}
\begin{corrige}{2}{Neutralité}
\begin{questions}
\item Huit électrons, car il est neutre.
\item $Z=8$.
\end{questions}
\end{corrige}
\end{minipage}
\par\noindent\begin{minipage}{\linewidth}
\begin{corrige}{3}{Un noyau}
\begin{questions}
\item $Z=11$.
\item $A=11+12=23$.
\end{questions}
\par\smallskip\noindent\textbf{Pour comprendre.}\enspace Le numéro atomique compte les protons : $Z=11$. Le nombre de nucléons est $A=11+12=23$. Un atome neutre possède autant d’électrons que de protons.\par
\pcFigureEx{figures/rev-cor-atome-na.pdf}{Inventaire des constituants ; représentation symbolique sans échelle de taille ni orbites électroniques.}
\end{corrige}
\end{minipage}
\par\noindent\begin{minipage}{\linewidth}
\begin{corrige}{4}{L’eau}
\begin{questions}
\item Deux H et un O.
\item Six H et trois O.
\end{questions}
\pcFigureEx{figures/rev-cor-trois-eaux.pdf}{Le coefficient multiplie la molécule entière : trois molécules contiennent six H et trois O.}
\end{corrige}
\end{minipage}
\par\noindent\begin{minipage}{\linewidth}
\begin{corrige}{5}{Le signe}
\begin{questions}
\item Il a gagné un électron.
\item Il a perdu un électron.
\end{questions}
\pcFigureEx{figures/rev-cor-ions.pdf}{Perdre un électron donne une charge positive ; en gagner un donne une charge négative. Le noyau reste inchangé.}
\end{corrige}
\end{minipage}
\par\noindent\begin{minipage}{\linewidth}
\begin{corrige}{6}{Un ion cuivre}
\begin{questions}
\item $29-2=27$ électrons.
\item 29 protons ; le noyau n’a pas changé.
\end{questions}
\par\smallskip\noindent\textbf{Pour comprendre.}\enspace Le noyau de cuivre conserve ses 29 protons. La charge $2+$ signifie deux électrons de moins que l’atome neutre : $29-2=27$. Ne pas retrancher des protons.\par
\pcFigureEx{figures/rev-cor-atome-cu.pdf}{Inventaire des constituants ; représentation symbolique sans échelle de taille ni orbites électroniques.}
\end{corrige}
\end{minipage}
\par\noindent\begin{minipage}{\linewidth}
\begin{corrige}{7}{Un ion chlorure}
\begin{questions}
\item 18 électrons.
\item Non : une entité monoatomique chargée diffère d’une molécule neutre à deux atomes.
\end{questions}
\pcFigureEx{figures/rev-cor-atome-cl.pdf}{Inventaire des constituants ; représentation symbolique sans échelle de taille ni orbites électroniques.}
\end{corrige}
\end{minipage}
\par\noindent\begin{minipage}{\linewidth}
\begin{corrige}{8}{Inventaire complet}
\begin{questions}
\item 14 neutrons et 13 électrons.
\item Charge $3+$ ; $\ce{Al^{3+}}$.
\end{questions}
\par\smallskip\noindent\textbf{Pour comprendre.}\enspace $N=A-Z=27-13=14$ neutrons. L’atome neutre possède 13 électrons ; l’ion $\ce{Al^3+}$ en a $13-3=10$.\par
\pcFigureEx{figures/rev-cor-atome-al.pdf}{Inventaire des constituants ; représentation symbolique sans échelle de taille ni orbites électroniques.}
\end{corrige}
\end{minipage}
\par\noindent\begin{minipage}{\linewidth}
\begin{corrige}{9}{Indices et coefficient}
\begin{questions}
\item Deux C et huit H.
\item Six N et trois O.
\end{questions}
\pcFigureEx{figures/rev-cor-coefficients.pdf}{Compter chaque type d’atome en multipliant coefficient et indice.}
\end{corrige}
\end{minipage}
\par\noindent\begin{minipage}{\linewidth}
\begin{corrige}{10}{Trois écritures}
\begin{questions}
\item O : symbole de l’oxygène ou atome selon le contexte ; dioxygène : molécule à deux O ; ion oxyde : un O chargé $2-$.
\item Oui, ils ont le même numéro atomique pour chaque noyau d’oxygène.
\end{questions}
\pcFigureEx{figures/rev-cor-oxygene.pdf}{Atome, molécule et ion sont trois entités différentes ; la charge ne se lit pas comme un nombre d’atomes.}
\end{corrige}
\end{minipage}
\par\noindent\begin{minipage}{\linewidth}
\begin{corrige}{11}{Retrouver un élément}
\begin{questions}
\item Deux charges positives non compensées : $2+$.
\item $\ce{Mg^{2+}}$.
\end{questions}
\par\smallskip\noindent\textbf{Pour comprendre.}\enspace La charge en unités élémentaires vaut $12-10=+2$ ; avec $Z=12$, il s’agit de $\ce{Mg^2+}$. Le signe vient du déficit d’électrons.\par
\pcFigureEx{figures/rev-cor-atome-mg.pdf}{Inventaire des constituants ; représentation symbolique sans échelle de taille ni orbites électroniques.}
\end{corrige}
\end{minipage}
\par\noindent\begin{minipage}{\linewidth}
\begin{corrige}{12}{Une fausse justification}
\begin{questions}
\item Un gain ou une perte d’électron ne change pas l’élément. Un ion chlorure provient du chlore ; un ion sodium reste du sodium.
\item Le nombre de protons, donc Z ; cela sort de la transformation chimique usuelle étudiée.
\end{questions}
\end{corrige}
\end{minipage}
\par\noindent\begin{minipage}{\linewidth}
\begin{corrige}{13}{Un changement d’échelle}
\begin{questions}
\item $10^{-10}/10^{-15}=10^5$.
\item $10/10^5=10^{-4}\unite m=0{,}1\unite{mm}$.
\end{questions}
\par\smallskip\noindent\textbf{Pour comprendre.}\enspace $10^{-10}/10^{-15}=10^5$. Si l’atome mesure 10 m dans une maquette, le noyau mesurerait $10/10^5=10^{-4}\unite m=0{,}1\unite{mm}$.\par
\pcFigureEx{figures/rev-cor-echelles-atome.pdf}{Deux zooms séparés : le diamètre du noyau est environ cent mille fois plus petit que celui de l’atome.}
\end{corrige}
\end{minipage}
\par\noindent\begin{minipage}{\linewidth}
\begin{corrige}{14}{Une solution électriquement neutre}
\begin{questions}
\item La somme des charges est nulle.
\item Quatre ions chlorure : $2\times2-4=0$.
\end{questions}
\par\smallskip\noindent\textbf{Pour comprendre.}\enspace Chaque ion chlorure porte la charge $-1$. Deux ions calcium portent ensemble $2\times(+2)=+4$ : il faut donc quatre ions chlorure pour obtenir $+4-4=0$.\par
\pcFigureEx{figures/rev-cor-neutralite.pdf}{Additionner les charges de tous les ions ; conserver les signes et les coefficients.}
\end{corrige}
\end{minipage}
\par\noindent\begin{minipage}{\linewidth}
\begin{corrige}{15}{Deux noyaux du même élément}
\begin{questions}
\item Respectivement 6 et 8 neutrons.
\item Leur même nombre de protons les rattache au carbone ; ce sont deux isotopes, terme fourni ici.
\end{questions}
\pcFigureEx{figures/rev-cor-atome-isotopes.pdf}{Inventaire des constituants ; représentation symbolique sans échelle de taille ni orbites électroniques.}
\end{corrige}
\end{minipage}
\par\noindent\begin{minipage}{\linewidth}
\phantomsection\label{cor-defi-2}\pcSolutionDefi{2}{L’entité inconnue}\begin{questions}
\item $Z=26$ et $A=56$.
\item $\ce{Fe^{3+}}$, car trois électrons manquent par rapport à l’atome neutre.
\item Le nombre de protons reste 26 : l’élément est le fer. La charge vient du déséquilibre entre protons et électrons. L’entité décrite est un ion monoatomique, sans assemblage de plusieurs atomes.
\end{questions}
\par\smallskip\noindent\textbf{Pour comprendre.}\enspace $Z=26$, $A=26+30=56$. La charge vaut $26-23=+3$, donc $\ce{Fe^3+}$ ; les neutrons ne contribuent pas à cette charge.\par
\pcFigureEx{figures/rev-cor-atome-fe3.pdf}{Inventaire des constituants ; représentation symbolique sans échelle de taille ni orbites électroniques.}
\end{minipage}
\par\noindent\begin{minipage}{\linewidth}
\phantomsection\label{cor-activite-2}\pcSousTitre{Pistes pour l’investigation}
Fixer Z et construire un noyau ; l’atome neutre a Z électrons. Retirer deux électrons donne un ion 2+, en ajouter un donne un ion 1−. Les protons ne changent pas. La taille des jetons, leurs distances et leurs positions fixes ne décrivent pas le nuage réel. Le tableau périodique sert à vérifier le symbole et le numéro atomique.\par\end{minipage}

\corrigesChapitre{3}{Transformations chimiques et équations}
\par\noindent\begin{minipage}{\linewidth}
\begin{corrige}{1}{Fondre ou réagir}
\begin{questions}
\item Non : l’eau change d’état.
\item De nouvelles espèces, présentes dans la rouille, se forment en consommant notamment le fer.
\end{questions}
\end{corrige}
\end{minipage}
\par\noindent\begin{minipage}{\linewidth}
\begin{corrige}{2}{Avant et après}
\begin{questions}
\item Carbone et dioxygène.
\item Dioxyde de carbone.
\end{questions}
\end{corrige}
\end{minipage}
\par\noindent\begin{minipage}{\linewidth}
\begin{corrige}{3}{Un test fourni}
\begin{questions}
\item Le dioxyde de carbone.
\item Non ; plusieurs gaz et mécanismes peuvent produire des bulles.
\end{questions}
\end{corrige}
\end{minipage}
\par\noindent\begin{minipage}{\linewidth}
\begin{corrige}{4}{Compter}
\begin{questions}
\item Quatre H de chaque côté.
\item Deux O de chaque côté.
\end{questions}
\pcFigureEx{figures/cor-reaction-eau.pdf}{Les coefficients rétablissent la conservation des atomes ; les formules des espèces restent inchangées.}
\end{corrige}
\end{minipage}
\par\noindent\begin{minipage}{\linewidth}
\begin{corrige}{5}{Masse totale}
\begin{questions}
\item $44\unite g$.
\item Non : ce sont les atomes de chaque élément qui se conservent.
\end{questions}
\end{corrige}
\end{minipage}
\par\noindent\begin{minipage}{\linewidth}
\begin{corrige}{6}{Un gaz s’échappe}
\begin{questions}
\item $2{,}2\unite g$.
\item Non : elle est sortie, par exemple sous forme gazeuse ; il faut l’inclure dans le bilan complet.
\end{questions}
\pcFigureEx{figures/rev-cor-masse-gaz.pdf}{Le système pesé est ouvert : la différence correspond au gaz sorti avec les hypothèses données.}
\end{corrige}
\end{minipage}
\par\noindent\begin{minipage}{\linewidth}
\begin{corrige}{7}{Une masse inconnue}
\begin{questions}
\item $m=22-6=16\unite g$.
\item Un réactif restant devrait figurer dans la masse finale totale.
\end{questions}
\pcFigureEx{figures/rev-cor-masse-carbone.pdf}{Les deux réactifs consommés contribuent à la masse de dioxyde de carbone.}
\end{corrige}
\end{minipage}
\par\noindent\begin{minipage}{\linewidth}
\begin{corrige}{8}{Le bon inventaire}
\begin{questions}
\item Deux C, six H et sept O de chaque côté.
\item Non, ce sont des proportions d’entités dans le modèle de réaction.
\end{questions}
\par\smallskip\noindent\textbf{Pour comprendre.}\enspace À gauche : C = 2, H = 6 et O = $1+3\times2=7$. À droite : C = $2\times1=2$, H = $3\times2=6$ et O = $2\times2+3=7$. Les nombres d’atomes sont égaux pour chaque élément.\par
\pcFigureEx{figures/rev-cor-ethanol.pdf}{Inventaire atomique de chaque côté ; l’oxygène de l’éthanol fait aussi partie du bilan.}
\end{corrige}
\end{minipage}
\par\noindent\begin{minipage}{\linewidth}
\begin{corrige}{9}{Une équation impossible}
\begin{questions}
\item Elle compte deux O à gauche et un à droite.
\item La seconde conserve quatre H et deux O de chaque côté.
\end{questions}
\pcFigureEx{figures/cor-reaction-eau.pdf}{Les coefficients rétablissent la conservation des atomes ; les formules des espèces restent inchangées.}
\end{corrige}
\end{minipage}
\par\noindent\begin{minipage}{\linewidth}
\begin{corrige}{10}{Des réactifs en excès}
\begin{questions}
\item $10+25-5=30\unite g$.
\item Les 35 g finaux comprennent aussi les 5 g non transformés.
\end{questions}
\par\smallskip\noindent\textbf{Pour comprendre.}\enspace La masse initiale totale vaut $10+25=35\unite g$. S’il reste 5 g de réactif, les produits formés représentent $35-5=30\unite g$ ; la masse finale totale reste 35 g.\par
\pcFigureEx{figures/rev-cor-reactif-restant.pdf}{Ne pas confondre masse des produits formés et masse finale totale avec le réactif restant.}
\end{corrige}
\end{minipage}
\par\noindent\begin{minipage}{\linewidth}
\begin{corrige}{11}{Comparer des contrôles}
\begin{questions}
\item Le changement peut venir du réactif, du milieu ou d’une contamination.
\item Non, il faut contrôler la cause et reprendre un protocole interprétable.
\end{questions}
\pcFigureEx{figures/rev-cor-temoin.pdf}{Le changement du témoin empêche d’attribuer avec certitude le résultat au gaz étudié.}
\end{corrige}
\end{minipage}
\par\noindent\begin{minipage}{\linewidth}
\begin{corrige}{12}{Mélange et réaction}
\begin{questions}
\item Non : les molécules de sucre restent présentes, dispersées dans le solvant.
\item Des substances mises en contact peuvent réagir ; il faut examiner les espèces initiales et finales.
\end{questions}
\end{corrige}
\end{minipage}
\par\noindent\begin{minipage}{\linewidth}
\begin{corrige}{13}{Un solide gagne de la masse}
\begin{questions}
\item Le dioxygène incorporé au solide, ici 0,8 g.
\item $2{,}1+0{,}8=2{,}9\unite g$.
\end{questions}
\pcFigureEx{figures/rev-cor-masse-fer.pdf}{L’augmentation de masse vient de matière du milieu incorporée dans le produit.}
\end{corrige}
\end{minipage}
\par\noindent\begin{minipage}{\linewidth}
\begin{corrige}{14}{Un protocole de gaz}
\begin{questions}
\item Identifier le gaz avec une banque de tests, comparer les espèces initiales et finales, et prévoir un témoin.
\item La quantité peut être trop faible ou le réactif inadéquat ; une absence de résultat visible a une portée limitée.
\end{questions}
\end{corrige}
\end{minipage}
\par\noindent\begin{minipage}{\linewidth}
\begin{corrige}{15}{Prise d’initiative}
\begin{questions}
\item Non : le bilan de masse est cohérent mais ne donne pas l’identité.
\item Un test caractéristique, une formule dans une équation fournie ou d’autres propriétés mesurées.
\end{questions}
\end{corrige}
\end{minipage}
\par\noindent\begin{minipage}{\linewidth}
\phantomsection\label{cor-defi-3}\pcSolutionDefi{3}{Le bilan complet}\begin{questions}
\item $m=8+40-8-22=18\unite g$.
\item Avant : 48 g ; après : $8+22+18=48\unite g$.
\item L’enceinte est un contenant supposé inerte. Sa masse intervient dans la pesée totale, mais elle est inchangée et ne représente pas une espèce consommée.
\end{questions}
\par\smallskip\noindent\textbf{Pour comprendre.}\enspace La masse initiale est $8+40=48\unite g$. La masse finale est $22+18+8=48\unite g$. Les 8 g de dioxygène restant doivent figurer dans le bilan ; ils ne font pas partie des produits formés.\par
\pcFigureEx{figures/rev-cor-masse-methane.pdf}{Bilan complet du modèle : réactifs consommés, produits et dioxygène restant.}
\end{minipage}
\par\noindent\begin{minipage}{\linewidth}
\phantomsection\label{cor-activite-3}\pcSousTitre{Pistes pour l’investigation}
Fusion et dissolution sans réaction ne forment pas de nouvelles espèces ; un dégazage peut seulement libérer une espèce déjà présente. Une réaction nécessite une interprétation des espèces consommées et formées. Le test et son témoin renforcent l’argument, sans que les seules bulles suffisent. Identifier le système et les échanges avant tout bilan de masse.\par\end{minipage}

\corrigesChapitre{4}{Solutions, ions et pH}
\par\noindent\begin{minipage}{\linewidth}
\begin{corrige}{1}{Classer}
\begin{questions}
\item pH 2 : acide ; pH 11 : basique.
\item Neutre dans les conditions annoncées.
\end{questions}
\end{corrige}
\end{minipage}
\par\noindent\begin{minipage}{\linewidth}
\begin{corrige}{2}{L’unité}
\begin{questions}
\item Le pH est sans unité.
\item Un volume, dans l’unité de sa graduation.
\end{questions}
\end{corrige}
\end{minipage}
\par\noindent\begin{minipage}{\linewidth}
\begin{corrige}{3}{Ions et acidité}
\begin{questions}
\item $\ce{H+}$.
\item $\ce{OH-}$.
\end{questions}
\end{corrige}
\end{minipage}
\par\noindent\begin{minipage}{\linewidth}
\begin{corrige}{4}{Du sel dissous}
\begin{questions}
\item $\ce{Na+}$ et $\ce{Cl-}$.
\item Ils sont dispersés à l’échelle microscopique et traversent le filtre ordinaire avec la solution.
\end{questions}
\pcFigureEx{figures/rev-cor-ions-filtre.pdf}{Un filtre ordinaire retient les grains, pas les ions dissous de cette solution.}
\end{corrige}
\end{minipage}
\par\noindent\begin{minipage}{\linewidth}
\begin{corrige}{5}{Un test bleu}
\begin{questions}
\item Présence d’ions $\ce{Cu^{2+}}$ dans les conditions du test.
\item Non : employer un prélèvement séparé.
\end{questions}
\pcFigureEx{figures/rev-cor-tests-ions.pdf}{Prélèvements séparés : chaque réactif sert à une recherche précise ; ne pas cumuler les réactifs dans le même tube.}
\end{corrige}
\end{minipage}
\par\noindent\begin{minipage}{\linewidth}
\begin{corrige}{6}{Une dilution}
\begin{questions}
\item Oui, car 4 reste inférieur à 7.
\item Non, elle a diminué.
\end{questions}
\pcFigureEx{figures/rev-cor-ph-acide.pdf}{Le pH augmente de 2 à 4, mais reste inférieur à 7 ; la solution est encore acide.}
\end{corrige}
\end{minipage}
\par\noindent\begin{minipage}{\linewidth}
\begin{corrige}{7}{Une base diluée}
\begin{questions}
\item Il s’atténue.
\item Non dans ce modèle usuel : le pH tend vers 7 sans rendre la solution acide.
\end{questions}
\pcFigureEx{figures/rev-cor-ph-base.pdf}{Le pH baisse de 12 à 10, mais reste supérieur à 7 ; la solution est encore basique.}
\end{corrige}
\end{minipage}
\par\noindent\begin{minipage}{\linewidth}
\begin{corrige}{8}{Trois inconnues}
\begin{questions}
\item A acide, B neutre, C basique vers 25 °C.
\item Non : la seule présence d’ions ne donne pas le caractère acide ou basique.
\end{questions}
\pcFigureEx{figures/rev-cor-ph-classement.pdf}{Classer par comparaison à 7, à température usuelle.}
\end{corrige}
\end{minipage}
\par\noindent\begin{minipage}{\linewidth}
\begin{corrige}{9}{Un résultat vert}
\begin{questions}
\item Ion fer(II).
\item Deux électrons.
\end{questions}
\pcFigureEx{figures/rev-cor-tests-ions.pdf}{Prélèvements séparés : chaque réactif sert à une recherche précise ; ne pas cumuler les réactifs dans le même tube.}
\end{corrige}
\end{minipage}
\par\noindent\begin{minipage}{\linewidth}
\begin{corrige}{10}{Neutraliser}
\begin{questions}
\item $\ce{H+}$ et $\ce{OH-}$.
\item L’un peut être introduit en excès ; les quantités et les espèces présentes déterminent le résultat.
\end{questions}
\pcFigureEx{figures/rev-cor-acide-base.pdf}{Une unité H\textsuperscript{+} réagit avec une unité HO\textsuperscript{-} dans ce modèle ; un éventuel excès reste dans la solution.}
\end{corrige}
\end{minipage}
\par\noindent\begin{minipage}{\linewidth}
\begin{corrige}{11}{Échantillons comparables}
\begin{questions}
\item Non, la précision de la comparaison colorée ne permet pas cette écriture.
\item Employer un pH-mètre réglé, avec rinçage et protocole adaptés, puis répéter.
\end{questions}
\end{corrige}
\end{minipage}
\par\noindent\begin{minipage}{\linewidth}
\begin{corrige}{12}{Conduction}
\begin{questions}
\item Électrons dans le métal ; ions positifs et négatifs dans la solution.
\item Oui, neutralité globale et mobilité des charges sont compatibles.
\end{questions}
\pcFigureEx{figures/rev-cor-conduction.pdf}{Dans la solution, des ions des deux signes se déplacent ; ce ne sont pas des électrons libres qui traversent la solution comme dans le métal.}
\end{corrige}
\end{minipage}
\par\noindent\begin{minipage}{\linewidth}
\begin{corrige}{13}{Sans moyenne}
\begin{questions}
\item Le pH ne se moyenne pas et ne décrit pas à lui seul les quantités réactives disponibles.
\item Il faut connaître les solutions et les quantités, ou mesurer le pH final dans un protocole autorisé.
\end{questions}
\end{corrige}
\end{minipage}
\par\noindent\begin{minipage}{\linewidth}
\begin{corrige}{14}{L’eau et le climat}
\begin{questions}
\item Non, elle reste basique vers 25 °C.
\item Oui : acidification désigne ici une baisse de pH, même sans passage sous 7.
\end{questions}
\pcFigureEx{figures/rev-cor-ph-ocean.pdf}{Le pH diminue de 8,2 à 8,0 : acidification, tout en restant dans le domaine basique.}
\end{corrige}
\end{minipage}
\par\noindent\begin{minipage}{\linewidth}
\begin{corrige}{15}{Concevoir une analyse}
\begin{questions}
\item Sur l’un : nitrate d’argent, rechercher un précipité blanc. Sur l’autre : soude, rechercher un précipité bleu. Prévoir les protections et témoins.
\item Chaque test renseigne sur l’ion associé ; l’un n’identifie pas automatiquement l’autre.
\end{questions}
\par\smallskip\noindent\textbf{Pour comprendre.}\enspace Prélever deux portions de solution dans deux tubes propres. Ajouter quelques gouttes de nitrate d’argent au premier pour rechercher les chlorures ; utiliser la soude dans l’autre pour rechercher notamment les ions cuivre(II). Comparer avec un témoin adapté et conclure dans les limites du test.\par
\pcFigureEx{figures/rev-cor-tests-ions.pdf}{Prélèvements séparés : chaque réactif sert à une recherche précise ; ne pas cumuler les réactifs dans le même tube.}
\end{corrige}
\end{minipage}
\par\noindent\begin{minipage}{\linewidth}
\phantomsection\label{cor-defi-4}\pcSolutionDefi{4}{Le laboratoire enquête}\begin{questions}
\item A est acide ; B basique.
\item Le mélange peut réagir, modifier les espèces et créer un risque. Il faut des prélèvements séparés et un protocole validé.
\item A est acide et contient des ions chlorure selon le test. B est basique. Le test cuivre négatif ne permet pas d’affirmer une absence absolue sans connaître la sensibilité du test.
\end{questions}
\pcFigureEx{figures/rev-cor-ph-solutions.pdf}{A, de pH 3, est acide ; B, de pH 10, est basique. Les tests se font sur des prélèvements séparés.}
\end{minipage}
\par\noindent\begin{minipage}{\linewidth}
\phantomsection\label{cor-activite-4}\pcSousTitre{Pistes pour l’investigation}
Noter les valeurs de pH avec une précision compatible avec l’outil. Le test de précipitation se fait sur un autre prélèvement. Pour chaque test, écrire le réactif, l’observation et l’espèce associée dans la banque. Un témoin contrôle une cause indépendante de l’échantillon. Une conclusion négative reste limitée par la sensibilité et les conditions.\par\end{minipage}

\corrigesChapitre{5}{Métaux : réactions et corrosion}
\par\noindent\begin{minipage}{\linewidth}
\begin{corrige}{1}{Corrosion}
\begin{questions}
\item Chimique.
\item Eau et dioxygène.
\end{questions}
\end{corrige}
\end{minipage}
\par\noindent\begin{minipage}{\linewidth}
\begin{corrige}{2}{Un gaz}
\begin{questions}
\item Dihydrogène, $\ce{H2}$.
\item Ion zinc, $\ce{Zn^{2+}}$.
\end{questions}
\pcFigureEx{figures/rev-cor-zinc.pdf}{Deux H\textsuperscript{+} compensent la charge de Zn²\textsuperscript{+} ; les atomes et la charge totale sont conservés.}
\end{corrige}
\end{minipage}
\par\noindent\begin{minipage}{\linewidth}
\begin{corrige}{3}{Les réactifs}
\begin{questions}
\item Fer et ions H+.
\item Oui, si le contexte fournit une solution chlorhydrique et qu’il ne participe pas à la transformation simplifiée.
\end{questions}
\end{corrige}
\end{minipage}
\par\noindent\begin{minipage}{\linewidth}
\begin{corrige}{4}{Protection}
\begin{questions}
\item Elle constitue une barrière contre le milieu.
\item Elle peut exposer le métal à l’eau et au dioxygène.
\end{questions}
\pcFigureEx{figures/rev-cor-protection.pdf}{Une peinture protège en limitant le contact ; une rayure peut laisser agir le milieu sur le métal exposé.}
\end{corrige}
\end{minipage}
\par\noindent\begin{minipage}{\linewidth}
\begin{corrige}{5}{Le bon test}
\begin{questions}
\item Ion fer(II).
\item Il met en jeu un gaz inflammable et exige matériel, quantités et encadrement adaptés.
\end{questions}
\end{corrige}
\end{minipage}
\par\noindent\begin{minipage}{\linewidth}
\begin{corrige}{6}{Conserver la charge}
\begin{questions}
\item Deux charges positives de chaque côté.
\item Deux atomes H de chaque côté.
\end{questions}
\pcFigureEx{figures/rev-cor-zinc.pdf}{Deux H\textsuperscript{+} compensent la charge de Zn²\textsuperscript{+} ; les atomes et la charge totale sont conservés.}
\end{corrige}
\end{minipage}
\par\noindent\begin{minipage}{\linewidth}
\begin{corrige}{7}{Le pH évolue}
\begin{questions}
\item Elle diminue, mais la solution reste acide.
\item Oui, dans le contexte et le modèle étudiés.
\end{questions}
\pcFigureEx{figures/rev-cor-ph-acide.pdf}{Le pH augmente de 2 à 4, mais reste inférieur à 7 ; la solution est encore acide.}
\end{corrige}
\end{minipage}
\par\noindent\begin{minipage}{\linewidth}
\begin{corrige}{8}{Une pesée ouverte}
\begin{questions}
\item $0{,}06\unite g$.
\item Elles emporteraient aussi de la matière et empêcheraient d’attribuer toute la perte au gaz.
\end{questions}
\end{corrige}
\end{minipage}
\par\noindent\begin{minipage}{\linewidth}
\begin{corrige}{9}{Choisir un couple}
\begin{questions}
\item Dans ces conditions, l’humidité favorise la corrosion ; le contrôle des autres facteurs est nécessaire.
\item De l’eau désaérée maintenue à l’abri du dioxygène, comparée à B.
\end{questions}
\pcFigureEx{figures/rev-cor-corrosion.pdf}{Comparer des conditions contrôlées, avec des clous comparables et une même durée ; ne changer qu’un facteur dans la comparaison retenue.}
\end{corrige}
\end{minipage}
\par\noindent\begin{minipage}{\linewidth}
\begin{corrige}{10}{Ne pas tout changer}
\begin{questions}
\item Non, plusieurs causes changent simultanément.
\item Même métal, dimensions, durée et accès à l’air ; seule la présence de sel varie.
\end{questions}
\par\smallskip\noindent\textbf{Pour comprendre.}\enspace Choisir deux clous comparables, même volume d’eau, même accès à l’air, même température et même durée. Ne faire varier que la salinité. Comparer ensuite le même indicateur de corrosion ; une expérience isolée ne quantifie pas toutes les situations réelles.\par
\pcFigureEx{figures/rev-cor-corrosion.pdf}{Comparer des conditions contrôlées, avec des clous comparables et une même durée ; ne changer qu’un facteur dans la comparaison retenue.}
\end{corrige}
\end{minipage}
\par\noindent\begin{minipage}{\linewidth}
\begin{corrige}{11}{Un bilan de masse}
\begin{questions}
\item $8{,}0\unite g$.
\item La matière supplémentaire vient de l’environnement et doit être comptée dans les réactifs.
\end{questions}
\pcFigureEx{figures/rev-cor-masse-corrosion.pdf}{Une masse de matière issue du milieu s’ajoute au métal pour former les produits.}
\end{corrige}
\end{minipage}
\par\noindent\begin{minipage}{\linewidth}
\begin{corrige}{12}{Une conclusion trop large}
\begin{questions}
\item Non : l’essai ne concerne qu’un métal, un milieu et une durée.
\item Aucun dégagement visible n’a été observé pour ce cuivre dans ces conditions ; d’autres métaux ou acides peuvent se comporter autrement.
\end{questions}
\end{corrige}
\end{minipage}
\par\noindent\begin{minipage}{\linewidth}
\begin{corrige}{13}{Comparer des résultats}
\begin{questions}
\item Non : il faut connaître les conditions et la dispersion de mesure.
\item Répéter sur plusieurs clous comparables, avec témoins, même protocole et estimation des écarts.
\end{questions}
\end{corrige}
\end{minipage}
\par\noindent\begin{minipage}{\linewidth}
\begin{corrige}{14}{Choisir un modèle}
\begin{questions}
\item Non ; elle représente une réaction particulière du fer avec un acide.
\item La corrosion atmosphérique implique notamment eau et dioxygène et produit d’autres composés.
\end{questions}
\end{corrige}
\end{minipage}
\par\noindent\begin{minipage}{\linewidth}
\begin{corrige}{15}{Réparer plutôt que jeter}
\begin{questions}
\item Entretenir ou réparer une protection adaptée ; réemployer ou recycler le métal après tri.
\item Collecte, transport, séparation et remise en forme nécessitent de l’énergie.
\end{questions}
\end{corrige}
\end{minipage}
\par\noindent\begin{minipage}{\linewidth}
\phantomsection\label{cor-defi-5}\pcSolutionDefi{5}{Le portail près de la mer}\begin{questions}
\item Le revêtement isole le métal du milieu ; une rayure rompt cette barrière.
\item Comparer des pièces identiques, même durée et accès à l’air, dans de l’eau non salée ou salée ; répéter les essais.
\item Entretenir un revêtement adapté peut limiter la corrosion. Une expérience courte au laboratoire ne reproduit pas exactement les embruns, les variations de température et les contraintes réelles.
\end{questions}
\pcFigureEx{figures/rev-cor-protection.pdf}{Une peinture protège en limitant le contact ; une rayure peut laisser agir le milieu sur le métal exposé.}
\end{minipage}
\par\noindent\begin{minipage}{\linewidth}
\phantomsection\label{cor-activite-5}\pcSousTitre{Pistes pour l’investigation}
Air sec contre eau au contact de l’air teste le rôle de l’eau ; eau désaérée protégée contre eau au contact de l’air teste celui du dioxygène. L’eau salée doit être comparée à l’eau non salée avec mêmes pièces et mêmes durées. Pour un revêtement, seule la présence de protection varie. Des répétitions et une mesure de corrosion seraient nécessaires pour quantifier une accélération.\par\end{minipage}

\corrigesChapitre{6}{Combustions et conservation de la masse}
\par\noindent\begin{minipage}{\linewidth}
\begin{corrige}{1}{Deux rôles}
\begin{questions}
\item Le charbon joue le rôle de combustible ; le dioxygène, de comburant.
\item Non : l’air est un mélange ; on distingue ses constituants et le dioxygène consommé.
\end{questions}
\end{corrige}
\end{minipage}
\par\noindent\begin{minipage}{\linewidth}
\begin{corrige}{2}{Les produits du méthane}
\begin{questions}
\item Dioxyde de carbone et eau.
\item Oui, une combustion peut transférer de l’énergie notamment thermiquement et par rayonnement.
\end{questions}
\end{corrige}
\end{minipage}
\par\noindent\begin{minipage}{\linewidth}
\begin{corrige}{3}{Une somme de masses}
\begin{questions}
\item $7\unite g$.
\item La mesure doit inclure toute la matière et aucun produit ne doit quitter le système pesé.
\end{questions}
\end{corrige}
\end{minipage}
\par\noindent\begin{minipage}{\linewidth}
\begin{corrige}{4}{Deux gaz à distinguer}
\begin{questions}
\item CO comporte un O pour un C ; CO2 en comporte deux.
\item Ni l’un ni l’autre ne doit être identifié en le sentant ; le CO est notamment inodore et très toxique.
\end{questions}
\end{corrige}
\end{minipage}
\par\noindent\begin{minipage}{\linewidth}
\begin{corrige}{5}{Une condition manquante}
\begin{questions}
\item Combustible, comburant et énergie d’activation.
\item Uniquement dans une activité encadrée par le professeur, avec le dispositif et les protections prévus.
\end{questions}
\end{corrige}
\end{minipage}
\par\noindent\begin{minipage}{\linewidth}
\begin{corrige}{6}{Compter les atomes}
\begin{questions}
\item Quatre H et quatre O de chaque côté.
\item Elles contiennent les quatre H du méthane et deux des quatre O ; les deux autres O sont dans le dioxyde de carbone.
\end{questions}
\pcFigureEx{figures/cor-reaction-methane.pdf}{Un C, quatre H et quatre O se retrouvent de chaque côté ; le nombre de molécules n’est pas imposé constant.}
\end{corrige}
\end{minipage}
\par\noindent\begin{minipage}{\linewidth}
\begin{corrige}{7}{Une masse manquante}
\begin{questions}
\item $22-6=16\unite g$.
\item Tous les produits sont pris en compte et il n’y a ni perte ni autre échange de matière dans le bilan.
\end{questions}
\par\smallskip\noindent\textbf{Pour comprendre.}\enspace Conservation de masse dans le bilan complet : $m_1+m_2=m_{\text{produit}}$. D’où $m_2=22-6=16\unite g$. Le calcul suppose les deux réactifs entièrement consommés et aucun échange non compté.\par
\end{corrige}
\end{minipage}
\par\noindent\begin{minipage}{\linewidth}
\begin{corrige}{8}{Le produit est plus lourd}
\begin{questions}
\item $44-12=32\unite g$.
\item Non : les $32\unite g$ supplémentaires viennent du dioxygène, qui doit être inclus dans les réactifs.
\end{questions}
\par\smallskip\noindent\textbf{Pour comprendre.}\enspace $m_{\text{dioxygène}}=44-12=32\unite g$. Le produit est plus lourd que le carbone seul parce qu’il contient aussi la matière apportée par le dioxygène.\par
\end{corrige}
\end{minipage}
\par\noindent\begin{minipage}{\linewidth}
\begin{corrige}{9}{Un montage ouvert}
\begin{questions}
\item Non : le système limité à la bougie échange de la matière avec l’air.
\item Le dioxygène consommé et les produits gazeux qui s’échappent, en plus de la cire et des éventuels résidus.
\end{questions}
\pcFigureEx{figures/cor-bougie-ouverte.pdf}{Le bilan doit inclure le dioxygène entrant et les produits gazeux sortants.}
\end{corrige}
\end{minipage}
\par\noindent\begin{minipage}{\linewidth}
\begin{corrige}{10}{Un bilan proportionnel}
\begin{questions}
\item On multiplie par 3 : $24\unite g$ de dioxygène et $33\unite g$ de produit.
\item Non : il concerne le dioxygène pur consommé ; l’air comporte aussi d’autres gaz.
\end{questions}
\par\smallskip\noindent\textbf{Pour comprendre.}\enspace Passer de 3 g à 9 g de carbone multiplie les masses du bilan par trois : 24 g de dioxygène et 33 g de produit. Le bilan chimique utilise O2 ; l’air contient aussi d’autres gaz.\par
\end{corrige}
\end{minipage}
\par\noindent\begin{minipage}{\linewidth}
\begin{corrige}{11}{Une affirmation à corriger}
\begin{questions}
\item Elle peut produire du monoxyde de carbone, de la suie et d’autres produits selon les conditions.
\item Le monoxyde de carbone est inodore ; son absence d’odeur n’atteste pas l’absence de danger.
\end{questions}
\end{corrige}
\end{minipage}
\par\noindent\begin{minipage}{\linewidth}
\begin{corrige}{12}{Une mesure presque égale}
\begin{questions}
\item Non : un écart de $0{,}1\unite g$ doit être discuté avec la précision et la variabilité des mesures.
\item La stabilité de la balance, l’étanchéité, les conditions de pesée et la reproductibilité.
\end{questions}
\end{corrige}
\end{minipage}
\par\noindent\begin{minipage}{\linewidth}
\begin{corrige}{13}{Un réactif en excès}
\begin{questions}
\item Les $6\unite g$ de C demandent $16\unite g$ de dioxygène ; il reste $8\unite g$ de dioxygène.
\item $6+16=22\unite g$ de produit ; $22+8=30\unite g$ est la masse finale totale.
\end{questions}
\par\smallskip\noindent\textbf{Pour comprendre.}\enspace Les 6 g de carbone demandent $6\times8/3=16\unite g$ de dioxygène. Sur 24 g, il en reste 8 g ; le produit formé pèse $6+16=22\unite g$. Contrôle : $22+8=6+24=30\unite g$.\par
\pcFigureEx{figures/cor-bilan-carbone-6.pdf}{Séparer matière transformée et réactif restant ; la masse finale totale comprend les deux.}
\end{corrige}
\end{minipage}
\par\noindent\begin{minipage}{\linewidth}
\begin{corrige}{14}{Émissions comparées}
\begin{questions}
\item A : $1\,000\unite g$ ; B : $400\unite g$.
\item Non : cette comparaison porte sur ce seul indicateur et le périmètre de la donnée, sans inclure nécessairement fabrication, ressources et fin de vie.
\end{questions}
\end{corrige}
\end{minipage}
\par\noindent\begin{minipage}{\linewidth}
\begin{corrige}{15}{Conserver quoi ?}
\begin{questions}
\item Non : $2\,\ce{H2}+\ce{O2}\rightarrow2\,\ce{H2O}$ passe de trois molécules à deux dans le bilan.
\item Le nombre d’atomes de chaque élément et la masse totale du système matériel complet.
\end{questions}
\pcFigureEx{figures/cor-reaction-eau.pdf}{Les coefficients rétablissent la conservation des atomes ; les formules des espèces restent inchangées.}
\end{corrige}
\end{minipage}
\par\noindent\begin{minipage}{\linewidth}
\phantomsection\label{cor-defi-6}\pcSolutionDefi{6}{Le bilan sous contrôle}\begin{questions}
\item Les $9\unite g$ de C demanderaient $24\unite g$ de dioxygène : le dioxygène disponible est insuffisant.
\item Les $20\unite g$ correspondent à $20\times3/8=7{,}5\unite g$ de C. Il reste $1{,}5\unite g$ ; on forme $27{,}5\unite g$ de produit dans ce modèle.
\item $27{,}5+1{,}5=29\unite g$, comme $9+20$. Le calcul utilise un bilan de combustion complète idéal ; des conditions réelles pauvres en dioxygène peuvent produire d’autres espèces.
\end{questions}
\pcFigureEx{figures/cor-bilan-carbone-9.pdf}{Le dioxygène disponible limite ici le modèle de transformation complète ; inclure le carbone restant dans la masse finale.}
\end{minipage}
\par\noindent\begin{minipage}{\linewidth}
\phantomsection\label{cor-activite-6}\pcSousTitre{Pistes pour l’investigation}
Dans le dispositif fermé adapté, la masse totale avant et après est compatible avec une conservation à la précision de mesure. Dans un système ouvert, les entrées et sorties gazeuses expliquent une variation de la partie pesée. Le schéma du système doit donc précéder le calcul. On conserve les nombres d’atomes de chaque élément, pas nécessairement le nombre de molécules.\par\end{minipage}

\corrigesChapitre{7}{La matière dans l’Univers}
\par\noindent\begin{minipage}{\linewidth}
\begin{corrige}{1}{Classer des objets}
\begin{questions}
\item Soleil : étoile ; Terre : planète ; Lune : satellite naturel.
\item Oui, comme l’ensemble du Système solaire.
\end{questions}
\end{corrige}
\end{minipage}
\par\noindent\begin{minipage}{\linewidth}
\begin{corrige}{2}{Le bon emboîtement}
\begin{questions}
\item Terre, Système solaire, Voie lactée, Univers observable, en allant vers l’ensemble le plus vaste.
\item Non : elle rassemble notamment de très nombreuses étoiles, du gaz et des poussières.
\end{questions}
\end{corrige}
\end{minipage}
\par\noindent\begin{minipage}{\linewidth}
\begin{corrige}{3}{Une unité trompeuse}
\begin{questions}
\item Une distance.
\item Elle correspond au trajet de la lumière pendant une durée d’un an.
\end{questions}
\end{corrige}
\end{minipage}
\par\noindent\begin{minipage}{\linewidth}
\begin{corrige}{4}{L’ordre des planètes}
\begin{questions}
\item Vénus avant et Mars après.
\item Neptune.
\end{questions}
\end{corrige}
\end{minipage}
\par\noindent\begin{minipage}{\linewidth}
\begin{corrige}{5}{Lire une échelle}
\begin{questions}
\item $50/10=5\unite{cm}$.
\item $150/10=15\unite{cm}$.
\end{questions}
\pcFigureEx{figures/cor-echelle-astronomie.pdf}{Schéma réduit : les longueurs indiquées sont à reporter sur la frise. L’échelle relie chaque centimètre à dix millions de kilomètres ; les diamètres ne sont pas à l’échelle.}
\end{corrige}
\end{minipage}
\par\noindent\begin{minipage}{\linewidth}
\begin{corrige}{6}{Comparer des distances}
\begin{questions}
\item $600/150=4$.
\item $600-150=450$ millions de kilomètres.
\end{questions}
\end{corrige}
\end{minipage}
\par\noindent\begin{minipage}{\linewidth}
\begin{corrige}{7}{Des milliards d’années}
\begin{questions}
\item $4\,600$ millions d’années.
\item C’est une échelle astronomique et géologique, très supérieure à une durée humaine.
\end{questions}
\end{corrige}
\end{minipage}
\par\noindent\begin{minipage}{\linewidth}
\begin{corrige}{8}{Les mêmes éléments}
\begin{questions}
\item Non : le même symbole H désigne le même élément, soumis aux mêmes lois.
\item Les milieux ont des histoires et compositions différentes ; mêmes éléments possibles ne signifie pas mêmes proportions.
\end{questions}
\end{corrige}
\end{minipage}
\par\noindent\begin{minipage}{\linewidth}
\begin{corrige}{9}{Les limites d’une affiche}
\begin{questions}
\item Non : elle peut montrer l’ordre ou les noms.
\item Préciser que les diamètres ne sont pas à l’échelle ; indiquer séparément le choix éventuel d’une échelle de distances.
\end{questions}
\end{corrige}
\end{minipage}
\par\noindent\begin{minipage}{\linewidth}
\begin{corrige}{10}{Passer à une autre échelle}
\begin{questions}
\item $750/25=30\unite{cm}$ du Soleil.
\item $8\times25=200$ millions de kilomètres.
\end{questions}
\pcFigureEx{figures/cor-frise-750.pdf}{On divise la distance réelle par la distance représentée par un centimètre.}
\end{corrige}
\end{minipage}
\par\noindent\begin{minipage}{\linewidth}
\begin{corrige}{11}{Une preuve d’autorité ?}
\begin{questions}
\item Arrondis différents ou distance instantanée comparée à une distance moyenne ; les orbites ne gardent pas toujours la même distance.
\item La définition de la grandeur, la date, l’unité et l’auteur de la source.
\end{questions}
\end{corrige}
\end{minipage}
\par\noindent\begin{minipage}{\linewidth}
\begin{corrige}{12}{Un métal « rare »}
\begin{questions}
\item Non : il peut être très dispersé et difficile à concentrer.
\item La teneur des gisements, l’accessibilité, l’énergie nécessaire, les déchets et les possibilités de recyclage.
\end{questions}
\end{corrige}
\end{minipage}
\par\noindent\begin{minipage}{\linewidth}
\begin{corrige}{13}{Une autre unité}
\begin{questions}
\item $450/150=3\unite{ua}$.
\item $5\times150=750$ millions de kilomètres.
\end{questions}
\end{corrige}
\end{minipage}
\par\noindent\begin{minipage}{\linewidth}
\begin{corrige}{14}{Un alignement dessiné}
\begin{questions}
\item Non : il ne donne pas leurs positions orbitales à une date précise.
\item Il peut représenter l’ordre des distances moyennes et les noms, selon sa légende.
\end{questions}
\end{corrige}
\end{minipage}
\par\noindent\begin{minipage}{\linewidth}
\begin{corrige}{15}{Rédiger une histoire}
\begin{questions}
\item La matière terrestre provient de matière cosmique enrichie au cours de l’évolution stellaire avant la formation du Système solaire.
\item La Terre a hérité de nombreux éléments ; elle n’est pas l’origine unique de ces éléments.
\end{questions}
\end{corrige}
\end{minipage}
\par\noindent\begin{minipage}{\linewidth}
\phantomsection\label{cor-defi-7}\pcSolutionDefi{7}{La frise impossible ?}\begin{questions}
\item $9\unite m=900\unite{cm}$ ; on peut choisir $1\unite{cm}$ pour $5$ millions de kilomètres.
\item $150/5=30\unite{cm}$ du Soleil.
\item Les diamètres sont minuscules devant les distances orbitales. Il faut employer des symboles et annoncer explicitement la différence d’échelle.
\end{questions}
\pcFigureEx{figures/cor-frise-salle.pdf}{Une longueur totale de 900 cm représente 4 500 millions de kilomètres ; la Terre est à 30 cm du Soleil.}
\end{minipage}
\par\noindent\begin{minipage}{\linewidth}
\phantomsection\label{cor-activite-7}\pcSousTitre{Pistes pour l’investigation}
Calculer chaque position en divisant la distance réelle par la valeur représentée par une unité de longueur. Choisir la plus grande distance pour vérifier que toute la frise tient dans la salle. Mentionner séparément les diamètres hors échelle. Une agence spatiale ou un organisme scientifique identifié fournit une source de référence ; les valeurs arrondies ne sont pas des mesures instantanées.\par\end{minipage}

\corrigesChapitre{8}{Décrire et mesurer un mouvement}
\par\noindent\begin{minipage}{\linewidth}
\begin{corrige}{1}{Le référentiel}
\begin{questions}
\item Oui, s’il reste posé à la même place.
\item Oui lorsque le train roule ; la réponse dépend du référentiel.
\end{questions}
\pcFigureEx{figures/cor-referentiels.pdf}{Le mouvement d’un même objet peut être décrit différemment selon le référentiel.}
\end{corrige}
\end{minipage}
\par\noindent\begin{minipage}{\linewidth}
\begin{corrige}{2}{La trajectoire}
\begin{questions}
\item Circulaire.
\item Non : sa trajectoire est rectiligne dans le référentiel considéré.
\end{questions}
\end{corrige}
\end{minipage}
\par\noindent\begin{minipage}{\linewidth}
\begin{corrige}{3}{Temps égaux}
\begin{questions}
\item Elle augmente : le mouvement est accéléré au sens de sa valeur de vitesse.
\item Les intervalles de temps doivent être réellement égaux et les longueurs comparables sur l’image.
\end{questions}
\pcFigureEx{figures/cor-chronophoto.pdf}{Comparer les distances entre positions successives, mesurées pour des durées égales.}
\end{corrige}
\end{minipage}
\par\noindent\begin{minipage}{\linewidth}
\begin{corrige}{4}{Une vitesse directe}
\begin{questions}
\item $v=40/8=5\unite{m/s}$.
\item Mètre par seconde, car les données sont en mètres et secondes.
\end{questions}
\par\smallskip\noindent\textbf{Pour comprendre.}\enspace Choisir $v=d/\Delta t$ avec d en m et la durée en s : $v=40/8=5\unite{m/s}$. Il s’agit de la moyenne sur le trajet, sauf information supplémentaire.\par
\end{corrige}
\end{minipage}
\par\noindent\begin{minipage}{\linewidth}
\begin{corrige}{5}{Convertir}
\begin{questions}
\item $10\times3{,}6=36\unite{km/h}$.
\item $72/3{,}6=20\unite{m/s}$.
\end{questions}
\par\smallskip\noindent\textbf{Pour comprendre.}\enspace $1\unite{m/s}=3{,}6\unite{km/h}$. Multiplier par 3,6 vers km/h et diviser vers m/s : 36 km/h et 20 m/s.\par
\end{corrige}
\end{minipage}
\par\noindent\begin{minipage}{\linewidth}
\begin{corrige}{6}{Calculer une distance}
\begin{questions}
\item $d=0{,}40\times15=6{,}0\unite m$.
\item $12\unite m$ si la vitesse reste la même.
\end{questions}
\end{corrige}
\end{minipage}
\par\noindent\begin{minipage}{\linewidth}
\begin{corrige}{7}{Calculer une durée}
\begin{questions}
\item $\Delta t=1{,}8/6=0{,}30\unite h$.
\item $0{,}30\times60=18\unite{min}$.
\end{questions}
\par\smallskip\noindent\textbf{Pour comprendre.}\enspace Pour une vitesse moyenne fournie, $\Delta t=d/v=1{,}8/6=0{,}30\unite h$. Puis $0{,}30\times60=18\unite{min}$ ; 0,30 h ne signifie pas 30 min.\par
\end{corrige}
\end{minipage}
\par\noindent\begin{minipage}{\linewidth}
\begin{corrige}{8}{Lire des positions}
\begin{questions}
\item Les écarts sont toujours de $3\unite m$ en $2\unite s$ : la valeur de la vitesse est constante.
\item $v=3/2=1{,}5\unite{m/s}$.
\end{questions}
\pcFigureEx{figures/cor-position.pdf}{Des distances de 3 m sont parcourues en 2 s ; la pente traduit une vitesse constante de 1,5 m/s.}
\end{corrige}
\end{minipage}
\par\noindent\begin{minipage}{\linewidth}
\begin{corrige}{9}{Uniforme sur un cercle}
\begin{questions}
\item Non : la direction est tangente au cercle et varie avec la position.
\item C’est un mouvement circulaire uniforme ; le vecteur vitesse n’est pas constant.
\end{questions}
\pcFigureEx{figures/cor-manege.pdf}{Les flèches de vitesse sont tangentes au cercle : leur direction change même si leur valeur reste la même.}
\end{corrige}
\end{minipage}
\par\noindent\begin{minipage}{\linewidth}
\begin{corrige}{10}{Une vidéo étalonnée}
\begin{questions}
\item $80/200=0{,}40\unite m$, si règle et mouvement sont dans le même plan.
\item $0{,}40/0{,}20=2{,}0\unite{m/s}$.
\end{questions}
\par\smallskip\noindent\textbf{Pour comprendre.}\enspace La calibration donne $1/200=0{,}005\unite m$ par pixel. Donc $d=80\times0{,}005=0{,}40\unite m$, puis $v=0{,}40/0{,}20=2{,}0\unite{m/s}$.\par
\pcFigureEx{figures/rev-cor-pixels.pdf}{Convertir le déplacement de l’image en distance réelle avant de diviser par la durée.}
\end{corrige}
\end{minipage}
\par\noindent\begin{minipage}{\linewidth}
\begin{corrige}{11}{Deux parties d’un trajet}
\begin{questions}
\item $d_{\rm total}=120\unite{km}$ et $\Delta t=3\unite h$ ; $v=40\unite{km/h}$.
\item $45$ est la moyenne arithmétique de 60 et 30, mais les durées ne sont pas égales.
\end{questions}
\par\smallskip\noindent\textbf{Pour comprendre.}\enspace Additionner d’abord les distances et les durées : $d=120\unite{km}$, $\Delta t=3\unite h$, donc $v_{\rm moy}=40\unite{km/h}$. La moyenne arithmétique des deux vitesses ne convient pas lorsque les durées diffèrent.\par
\end{corrige}
\end{minipage}
\par\noindent\begin{minipage}{\linewidth}
\begin{corrige}{12}{Un arrêt compte-t-il ?}
\begin{questions}
\item $30\unite{min}=0{,}5\unite h$ ; $v=12/0{,}5=24\unite{km/h}$.
\item Non : pendant l’arrêt, la vitesse par rapport au sol est nulle.
\end{questions}
\par\smallskip\noindent\textbf{Pour comprendre.}\enspace Le trajet complet dure $20+10=30\unite{min}=0{,}50\unite h$. La distance reste 12 km : $v_{\rm moy}=12/0{,}50=24\unite{km/h}$.\par
\pcFigureEx{figures/rev-cor-bus.pdf}{Pour la vitesse moyenne sur le trajet complet, l’arrêt fait partie de la durée totale.}
\end{corrige}
\end{minipage}
\par\noindent\begin{minipage}{\linewidth}
\begin{corrige}{13}{Une conclusion trop forte}
\begin{questions}
\item Non : elle peut avoir accéléré puis ralenti.
\item Il donne une vitesse moyenne de 20 m/s sur ce tronçon, pas toute son évolution temporelle.
\end{questions}
\end{corrige}
\end{minipage}
\par\noindent\begin{minipage}{\linewidth}
\begin{corrige}{14}{Un graphique ambigu}
\begin{questions}
\item Non : une position décroissante peut correspondre à un mouvement uniforme dans le sens opposé à l’axe choisi.
\item Il faut distinguer signe de l’évolution de position et valeur positive de la vitesse.
\end{questions}
\par\smallskip\noindent\textbf{Pour comprendre.}\enspace La pente est négative parce que le mobile se déplace vers des positions plus petites. La droite a une pente de valeur constante : dans cet exemple rectiligne, cela ne traduit pas un ralentissement.\par
\pcFigureEx{figures/cor-position-retour.pdf}{Une position décroissante indique un sens de déplacement ; la valeur de la vitesse peut rester constante.}
\end{corrige}
\end{minipage}
\par\noindent\begin{minipage}{\linewidth}
\begin{corrige}{15}{Comparer des mesures}
\begin{questions}
\item $\bar t=(4{,}8+5{,}0+5{,}2)/3=5{,}0\unite s$.
\item $v\approx10/5{,}0=2{,}0\unite{m/s}$. La dispersion et le déclenchement manuel ne justifient pas de nombreuses décimales.
\end{questions}
\end{corrige}
\end{minipage}
\par\noindent\begin{minipage}{\linewidth}
\phantomsection\label{cor-defi-8}\pcSolutionDefi{8}{Choisir le meilleur trajet}\begin{questions}
\item A : $1{,}2/4=0{,}30\unite h=18\unite{min}$. B : $1{,}5/6=0{,}25\unite h=15\unite{min}$, puis 3 minutes : 18 minutes.
\item Les durées sont égales dans ce modèle.
\item Les vitesses moyennes, la durée d’arrêt ou les distances peuvent varier. La conclusion est conditionnelle aux données.
\end{questions}
\pcFigureEx{figures/rev-cor-trajets.pdf}{Inclure les trois minutes d’arrêt dans la durée du trajet B ; les deux arrivées sont simultanées dans ce modèle.}
\end{minipage}
\par\noindent\begin{minipage}{\linewidth}
\phantomsection\label{cor-activite-8}\pcSousTitre{Pistes pour l’investigation}
Des distances successives sensiblement égales pour des durées égales sont compatibles avec une vitesse de valeur constante. Le graphe de position est alors proche d’une droite sur le déplacement rectiligne étudié. Contrôler l’échelle, la perspective et le même point suivi. Des écarts plus petits que la lecture ne justifient pas une conclusion certaine d’accélération.\par\end{minipage}

\corrigesChapitre{9}{Interactions, forces et poids}
\par\noindent\begin{minipage}{\linewidth}
\begin{corrige}{1}{Contact ou distance}
\begin{questions}
\item Contact.
\item À distance : gravitation.
\end{questions}
\end{corrige}
\end{minipage}
\par\noindent\begin{minipage}{\linewidth}
\begin{corrige}{2}{Deux appareils}
\begin{questions}
\item Une balance.
\item Un dynamomètre.
\end{questions}
\end{corrige}
\end{minipage}
\par\noindent\begin{minipage}{\linewidth}
\begin{corrige}{3}{Deux unités}
\begin{questions}
\item Kilogramme, kg.
\item Newton, N.
\end{questions}
\end{corrige}
\end{minipage}
\par\noindent\begin{minipage}{\linewidth}
\begin{corrige}{4}{Direction et sens}
\begin{questions}
\item Vers le bas : sens ; vertical : direction dans le modèle local.
\item Point d’application et valeur.
\end{questions}
\pcFigureEx{figures/rev-cor-poids.pdf}{Le poids est appliqué au centre de gravité dans ce modèle, de direction verticale et de sens vers le bas.}
\end{corrige}
\end{minipage}
\par\noindent\begin{minipage}{\linewidth}
\begin{corrige}{5}{Un premier calcul}
\begin{questions}
\item $20\unite N$.
\item Non : $200\unite g=0{,}200\unite{kg}$.
\end{questions}
\end{corrige}
\end{minipage}
\par\noindent\begin{minipage}{\linewidth}
\begin{corrige}{6}{Une petite masse}
\begin{questions}
\item $P=0{,}150\times9{,}8=1{,}47\unite N$.
\item Environ $1{,}5\unite N$.
\end{questions}
\par\smallskip\noindent\textbf{Pour comprendre.}\enspace Convertir d’abord : $150\unite g=0{,}150\unite{kg}$. Ensuite $P=mg=0{,}150\times9{,}8=1{,}47\unite N$, soit environ 1,5 N.\par
\end{corrige}
\end{minipage}
\par\noindent\begin{minipage}{\linewidth}
\begin{corrige}{7}{Retrouver la masse}
\begin{questions}
\item $m=P/g=2{,}4\unite{kg}$.
\item $2\,400\unite g$.
\end{questions}
\end{corrige}
\end{minipage}
\par\noindent\begin{minipage}{\linewidth}
\begin{corrige}{8}{Changer de lieu}
\begin{questions}
\item Masse : 60 kg ; poids : 96 N.
\item 600 N sur Terre ; la masse n’a pas changé.
\end{questions}
\par\smallskip\noindent\textbf{Pour comprendre.}\enspace $P_{\rm Terre}=60\times10=600\unite N$ et $P_{\rm Lune}=60\times1{,}6=96\unite N$. Le poids est multiplié par $1{,}6/10=0{,}16$ ; la masse ne change pas.\par
\pcFigureEx{figures/rev-cor-poids-lune.pdf}{La masse reste 60 kg ; les deux poids sont comparés avec la même échelle graphique qualitative.}
\end{corrige}
\end{minipage}
\par\noindent\begin{minipage}{\linewidth}
\begin{corrige}{9}{Une échelle}
\begin{questions}
\item $12/4=3\unite{cm}$.
\item $2{,}5\times4=10\unite N$.
\end{questions}
\par\smallskip\noindent\textbf{Pour comprendre.}\enspace Une longueur se calcule en divisant la force par le nombre de newtons représentés par un centimètre : $12/4=3\unite{cm}$. Inversement, $2{,}5\times4=10\unite N$.\par
\pcFigureExact{figures/rev-cor-force-echelle.pdf}{Flèches de 3 cm et 2,5 cm à l’impression à 100 \% ; 1 cm représente 4 N.}
\end{corrige}
\end{minipage}
\par\noindent\begin{minipage}{\linewidth}
\begin{corrige}{10}{Le livre immobile}
\begin{questions}
\item La Terre et la table.
\item Poids vertical vers le bas, action du support verticale vers le haut, de même valeur.
\end{questions}
\pcFigureEx{figures/cor-livre-forces.pdf}{Ne représenter que les forces appliquées au livre ; l’action du livre sur la table concerne un autre système.}
\end{corrige}
\end{minipage}
\par\noindent\begin{minipage}{\linewidth}
\begin{corrige}{11}{Une pente mesurée}
\begin{questions}
\item $P/m=10\unite{N/kg}$ dans les deux cas.
\item Une proportionnalité entre P et m compatible avec le modèle, à la précision des mesures.
\end{questions}
\pcFigureEx{figures/rev-cor-poids-graphe.pdf}{Les deux mesures sont compatibles avec P = 10 m ; l’origine est une prévision du modèle, pas une mesure ajoutée.}
\end{corrige}
\end{minipage}
\par\noindent\begin{minipage}{\linewidth}
\begin{corrige}{12}{Une mesure négative}
\begin{questions}
\item Contrôler le zéro et l’orientation selon la notice.
\item L’erreur peut représenter une grande part de la petite valeur mesurée.
\end{questions}
\end{corrige}
\end{minipage}
\par\noindent\begin{minipage}{\linewidth}
\begin{corrige}{13}{Des forces confondues}
\begin{questions}
\item La seconde agit sur la table, pas sur le livre.
\item L’action de la table sur le livre.
\end{questions}
\pcFigureEx{figures/cor-livre-forces.pdf}{Ne représenter que les forces appliquées au livre ; l’action du livre sur la table concerne un autre système.}
\end{corrige}
\end{minipage}
\par\noindent\begin{minipage}{\linewidth}
\begin{corrige}{14}{Un virage}
\begin{questions}
\item Non, la direction de la vitesse change.
\item Oui ; une modification de mouvement ne se limite pas à accélérer ou ralentir en valeur.
\end{questions}
\pcFigureEx{figures/cor-manege.pdf}{La direction de la vitesse change dans le mouvement circulaire uniforme ; cela n’impose pas un changement de sa valeur.}
\end{corrige}
\end{minipage}
\par\noindent\begin{minipage}{\linewidth}
\begin{corrige}{15}{Une affirmation à discuter}
\begin{questions}
\item Les effets dépendent de l’ensemble des actions et du système.
\item Une masse immobile suspendue peut subir 5 N vers le bas et 5 N vers le haut.
\end{questions}
\pcFigureEx{figures/rev-cor-forces-5N.pdf}{Ces deux forces opposées agissent sur le même système ; leurs valeurs sont égales.}
\end{corrige}
\end{minipage}
\par\noindent\begin{minipage}{\linewidth}
\phantomsection\label{cor-defi-9}\pcSolutionDefi{9}{La caisse suspendue}\begin{questions}
\item $P=39{,}2\unite N$.
\item Poids vers le bas et action du câble vers le haut, même direction et valeur ; longueur 3,92 cm chacune, environ 3,9 cm au tracé.
\item Elles agissent toutes deux sur la caisse mais proviennent de deux partenaires différents : Terre et câble.
\end{questions}
\par\smallskip\noindent\textbf{Pour comprendre.}\enspace $P=4\times9{,}8=39{,}2\unite N$. À l’équilibre dans le modèle à deux forces, la force du fil a aussi cette valeur. Chaque flèche mesure $39{,}2/10=3{,}92\unite{cm}$. Les deux flèches s’appliquent à la caisse.\par
\pcFigureExact{figures/rev-cor-caisse.pdf}{Deux flèches de 3,92 cm à l’impression à 100 \% ; même système et échelle 1 cm pour 10 N.}
\end{minipage}
\par\noindent\begin{minipage}{\linewidth}
\phantomsection\label{cor-activite-9}\pcSousTitre{Pistes pour l’investigation}
Tracer P en N en fonction de m en kg. Une droite proche de l’origine et des quotients P/m voisins de 9,8 à 10 N/kg sont attendus près du sol terrestre. La masse suspendue subit le poids et l’action du dynamomètre ; si elle est immobile dans le modèle à deux forces, leurs valeurs sont égales et les sens opposés. La force sur l’appareil concerne un autre système.\par\end{minipage}

\corrigesChapitre{10}{Exploiter la gravitation universelle}
\par\noindent\begin{minipage}{\linewidth}
\begin{corrige}{1}{Nature de l’action}
\begin{questions}
\item Non : elle agit à distance.
\item Oui.
\end{questions}
\end{corrige}
\end{minipage}
\par\noindent\begin{minipage}{\linewidth}
\begin{corrige}{2}{Deux objets}
\begin{questions}
\item Oui, avec une force de même valeur.
\item Sur la Terre.
\end{questions}
\pcFigureEx{figures/rev-cor-gravitation.pdf}{Deux flèches opposées, appliquées à deux corps différents ; elles ne se compensent pas sur le seul corps A.}
\end{corrige}
\end{minipage}
\par\noindent\begin{minipage}{\linewidth}
\begin{corrige}{3}{Unités}
\begin{questions}
\item kg et m.
\item $6{,}4\times10^6\unite m$.
\end{questions}
\end{corrige}
\end{minipage}
\par\noindent\begin{minipage}{\linewidth}
\begin{corrige}{4}{Quelle distance ?}
\begin{questions}
\item La distance entre leurs centres dans le modèle.
\item Il faut aussi connaître les rayons pour retrouver cette distance centrale.
\end{questions}
\pcFigureEx{figures/rev-cor-centres.pdf}{La distance d de la loi gravitationnelle est mesurée entre les centres, pas entre les surfaces.}
\end{corrige}
\end{minipage}
\par\noindent\begin{minipage}{\linewidth}
\begin{corrige}{5}{Une masse doublée}
\begin{questions}
\item F double.
\item F est multipliée par quatre.
\end{questions}
\end{corrige}
\end{minipage}
\par\noindent\begin{minipage}{\linewidth}
\begin{corrige}{6}{Une distance doublée}
\begin{questions}
\item F est divisée par quatre.
\item F est multipliée par quatre.
\end{questions}
\pcFigureEx{figures/rev-cor-inverse-carre.pdf}{Le facteur de distance est mis au carré au dénominateur ; à trois fois la distance, la force est divisée par neuf.}
\end{corrige}
\end{minipage}
\par\noindent\begin{minipage}{\linewidth}
\begin{corrige}{7}{Une force modèle}
\begin{questions}
\item $6{,}67\times10^{-11}\unite N$.
\item Non : elle est très faible, même si un appareil scolaire ne la détecte pas.
\end{questions}
\end{corrige}
\end{minipage}
\par\noindent\begin{minipage}{\linewidth}
\begin{corrige}{8}{Une comparaison rapide}
\begin{questions}
\item $36/9=4\unite N$.
\item $36/4=9\unite N$.
\end{questions}
\pcFigureEx{figures/rev-cor-inverse-carre.pdf}{Le facteur de distance est mis au carré au dénominateur ; à trois fois la distance, la force est divisée par neuf.}
\end{corrige}
\end{minipage}
\par\noindent\begin{minipage}{\linewidth}
\begin{corrige}{9}{Depuis l’orbite}
\begin{questions}
\item $d=6\,800\unite{km}=6{,}8\times10^6\unite m$.
\item Toute la distance en mètres, donc $(6{,}8\times10^6)^2$.
\end{questions}
\par\smallskip\noindent\textbf{Pour comprendre.}\enspace $d=R+h=6\,400+400=6\,800\unite{km}=6{,}8\times10^6\unite m$. Employer seulement l’altitude ferait partir la mesure de la surface et non du centre.\par
\pcFigureEx{figures/rev-cor-altitude.pdf}{Additionner rayon et altitude avant de convertir en mètres ; le point de départ est le centre de la Terre.}
\end{corrige}
\end{minipage}
\par\noindent\begin{minipage}{\linewidth}
\begin{corrige}{10}{Un petit calcul}
\begin{questions}
\item $F=3{,}335\times10^{-9}\unite N$.
\item Elle double : $6{,}67\times10^{-9}\unite N$.
\end{questions}
\end{corrige}
\end{minipage}
\par\noindent\begin{minipage}{\linewidth}
\begin{corrige}{11}{Trois facteurs}
\begin{questions}
\item Par $2/4=1/2$.
\item La masse au numérateur apporte 2 ; le carré de la distance au dénominateur apporte 4.
\end{questions}
\par\smallskip\noindent\textbf{Pour comprendre.}\enspace Le facteur de masse est 2, celui du carré de la distance est $2^2=4$ ; le facteur total est donc $2/4=0{,}5$. Traiter séparément les modifications avant de les combiner.\par
\end{corrige}
\end{minipage}
\par\noindent\begin{minipage}{\linewidth}
\begin{corrige}{12}{Un faux équilibre}
\begin{questions}
\item Elles s’exercent sur deux objets distincts et ne se compensent pas sur un seul.
\item L’attraction exercée par la Terre sur la Lune, avec les autres actions éventuelles du modèle.
\end{questions}
\pcFigureEx{figures/rev-cor-gravitation.pdf}{Deux flèches opposées, appliquées à deux corps différents ; elles ne se compensent pas sur le seul corps A.}
\end{corrige}
\end{minipage}
\par\noindent\begin{minipage}{\linewidth}
\begin{corrige}{13}{Impesanteur}
\begin{questions}
\item La station et les occupants restent attirés par la Terre.
\item Ils sont en chute commune autour de la Terre, hors forces de contact de soutien habituelles.
\end{questions}
\pcFigureEx{figures/rev-cor-orbite.pdf}{Le satellite et son contenu sont en chute commune ; l’impesanteur ne signifie pas absence de gravitation.}
\end{corrige}
\end{minipage}
\par\noindent\begin{minipage}{\linewidth}
\begin{corrige}{14}{Retrouver g}
\begin{questions}
\item $g=588/60=9{,}8\unite{N/kg}$.
\item $m=98/9{,}8=10\unite{kg}$.
\end{questions}
\end{corrige}
\end{minipage}
\par\noindent\begin{minipage}{\linewidth}
\begin{corrige}{15}{La précision du modèle}
\begin{questions}
\item La précision numérique affichée ne corrige pas l’approximation des données.
\item Les corps sont assimilés à des sphères ou des masses ponctuelles pour employer la distance entre centres.
\end{questions}
\end{corrige}
\end{minipage}
\clearpage
\par\noindent\begin{minipage}{\linewidth}
\phantomsection\label{cor-defi-10}\pcSolutionDefi{10}{Autour d’une planète}\begin{questions}
\item $F\approx4{,}08\times10^3\unite N$, soit environ $4{,}1\times10^3\unite N$.
\item Environ $1{,}02\times10^3\unite N$, par division par quatre.
\item Les valeurs sont égales et les sens opposés, mais les forces s’appliquent à deux masses très différentes. Égalité de force ne signifie pas évolution de mouvement identique.
\end{questions}
\par\smallskip\noindent\textbf{Pour comprendre.}\enspace $F=6{,}67\times10^{-11}\times(6{,}0\times10^{24})\times500/(7{,}0\times10^6)^2\approx4{,}1\times10^3\unite N$. Doubler la distance divise par quatre : environ $1{,}0\times10^3\unite N$.\par
\pcFigureEx{figures/rev-cor-satellite.pdf}{Distances entre centres et masses en unités SI ; la force diminue lorsque la distance augmente.}
\end{minipage}
\par\noindent\begin{minipage}{\linewidth}
\phantomsection\label{cor-activite-10}\pcSousTitre{Pistes pour l’investigation}
Pour d, 2d et 3d, les forces sont dans les rapports 1, 1/4 et 1/9. Si une seule masse double à distance fixée, F double. Le graphe F(d) est courbe ; un graphe de F en fonction de 1/d² serait linéaire pour des masses fixées. Le tableur vérifie les conséquences du modèle, il ne constitue pas une nouvelle preuve expérimentale de la loi.\par\end{minipage}

\corrigesChapitre{11}{Énergie cinétique et bilan mécanique}
\par\noindent\begin{minipage}{\linewidth}
\begin{corrige}{1}{La forme d’énergie}
\begin{questions}
\item L’énergie cinétique.
\item En joules, J.
\end{questions}
\end{corrige}
\end{minipage}
\par\noindent\begin{minipage}{\linewidth}
\begin{corrige}{2}{Les bonnes unités}
\begin{questions}
\item kg et m/s.
\item $0{,}500\unite{kg}$.
\end{questions}
\end{corrige}
\end{minipage}
\par\noindent\begin{minipage}{\linewidth}
\begin{corrige}{3}{La vitesse}
\begin{questions}
\item $10\unite{m/s}$.
\item $20\unite{m/s}$.
\end{questions}
\end{corrige}
\end{minipage}
\par\noindent\begin{minipage}{\linewidth}
\begin{corrige}{4}{Un calcul direct}
\begin{questions}
\item $\tfrac12\times2\times9=9\unite J$.
\item $0\unite J$.
\end{questions}
\end{corrige}
\end{minipage}
\par\noindent\begin{minipage}{\linewidth}
\begin{corrige}{5}{Deux facteurs}
\begin{questions}
\item Ec double.
\item Ec est multipliée par quatre.
\end{questions}
\pcFigureEx{figures/rev-cor-ec-facteurs.pdf}{Courbe de rapports à masse constante : le facteur appliqué à la vitesse est mis au carré pour l’énergie cinétique.}
\end{corrige}
\end{minipage}
\par\noindent\begin{minipage}{\linewidth}
\begin{corrige}{6}{Le ballon}
\begin{questions}
\item $E_c=0{,}5\times0{,}400\times25=5\unite J$.
\item $20\unite J$.
\end{questions}
\pcFigureEx{figures/rev-cor-ec-facteurs.pdf}{Courbe de rapports à masse constante : le facteur appliqué à la vitesse est mis au carré pour l’énergie cinétique.}
\end{corrige}
\end{minipage}
\par\noindent\begin{minipage}{\linewidth}
\begin{corrige}{7}{Une voiture}
\begin{questions}
\item $v=15\unite{m/s}$ ; $E_c=112\,500\unite J$.
\item $112{,}5\unite{kJ}$.
\end{questions}
\par\smallskip\noindent\textbf{Pour comprendre.}\enspace $54\unite{km/h}=54/3{,}6=15\unite{m/s}$. Alors $E_c=\frac12\times1\,000\times15^2=112\,500\unite J=112{,}5\unite{kJ}$. Mettre 54 au carré sans conversion donnerait une valeur fausse.\par
\end{corrige}
\end{minipage}
\par\noindent\begin{minipage}{\linewidth}
\begin{corrige}{8}{Ralentir}
\begin{questions}
\item $3\,000\unite J$ puis $750\unite J$.
\item $2\,250\unite J$, transférés ou convertis, notamment en thermique.
\end{questions}
\par\smallskip\noindent\textbf{Pour comprendre.}\enspace $E_{c,i}=\frac12\times60\times10^2=3\,000\unite J$, puis $E_{c,f}=\frac12\times60\times5^2=750\unite J$. L’écart est $3\,000-750=2\,250\unite J$.\par
\pcFigureEx{figures/rev-cor-freinage.pdf}{L’écart de 2 250 J ne disparaît pas : il est transféré ou converti, notamment par les frottements.}
\end{corrige}
\end{minipage}
\par\noindent\begin{minipage}{\linewidth}
\begin{corrige}{9}{Une chute idéale}
\begin{questions}
\item $E_c=40-15=25\unite J$.
\item Les frottements réels peuvent transférer une partie de l’énergie mécanique vers le milieu.
\end{questions}
\pcFigureEx{figures/rev-cor-chute.pdf}{Dans le modèle sans frottements, l’énergie mécanique reste 40 J ; les deux formes se redistribuent.}
\end{corrige}
\end{minipage}
\par\noindent\begin{minipage}{\linewidth}
\begin{corrige}{10}{Le rapport}
\begin{questions}
\item Rapport $3^2=9$.
\item Rapport $120/80=1{,}5$.
\end{questions}
\pcFigureEx{figures/rev-cor-ec-facteurs.pdf}{Courbe de rapports à masse constante : le facteur appliqué à la vitesse est mis au carré pour l’énergie cinétique.}
\end{corrige}
\end{minipage}
\par\noindent\begin{minipage}{\linewidth}
\begin{corrige}{11}{Une hausse de 20 \%}
\begin{questions}
\item $1{,}2^2=1{,}44$.
\item 44 \%.
\end{questions}
\par\smallskip\noindent\textbf{Pour comprendre.}\enspace Une hausse de 20 \% donne le facteur de vitesse $1+20/100=1{,}2$. Le facteur d’énergie est $1{,}2^2=1{,}44$ : la hausse d’énergie vaut donc $(1{,}44-1)\times100=44\,\%$.\par
\pcFigureEx{figures/rev-cor-ec-facteurs.pdf}{Courbe de rapports à masse constante : le facteur appliqué à la vitesse est mis au carré pour l’énergie cinétique.}
\end{corrige}
\end{minipage}
\par\noindent\begin{minipage}{\linewidth}
\begin{corrige}{12}{Changer de référentiel}
\begin{questions}
\item Oui, puisqu’il se déplace par rapport au sol.
\item Le référentiel et la vitesse du système dans ce référentiel.
\end{questions}
\end{corrige}
\end{minipage}
\par\noindent\begin{minipage}{\linewidth}
\begin{corrige}{13}{La force n’est pas l’énergie}
\begin{questions}
\item 250 J si le calcul est correct.
\item Non ; il faut un modèle et des données sur le freinage, la route et les autres actions.
\end{questions}
\end{corrige}
\end{minipage}
\par\noindent\begin{minipage}{\linewidth}
\begin{corrige}{14}{Un bilan réel}
\begin{questions}
\item Vers de l’énergie thermique du système et du milieu, notamment par frottement.
\item Non, à condition d’élargir le bilan aux transferts et aux autres formes d’énergie.
\end{questions}
\pcFigureEx{figures/rev-cor-chute-reelle.pdf}{Le bilan mécanique diminue de 15 J ; le bilan total inclut l’énergie transférée au milieu.}
\end{corrige}
\end{minipage}
\par\noindent\begin{minipage}{\linewidth}
\begin{corrige}{15}{Comparer deux véhicules}
\begin{questions}
\item A : 160 000 J ; B : 80 000 J. A a deux fois plus d’énergie.
\item Le carré de la vitesse compense ici plus que la différence de masse.
\end{questions}
\par\smallskip\noindent\textbf{Pour comprendre.}\enspace A : $\frac12\times800\times20^2=160\,000\unite J$. B : $\frac12\times1\,600\times10^2=80\,000\unite J$. Le rapport A/B vaut 2 : la masse seule ne suffit pas à comparer.\par
\pcFigureEx{figures/rev-cor-vehicules.pdf}{Comparer des résultats dans le même référentiel, après avoir pris le carré des vitesses.}
\end{corrige}
\end{minipage}
\par\noindent\begin{minipage}{\linewidth}
\phantomsection\label{cor-defi-11}\pcSolutionDefi{11}{Une étude de sécurité}\begin{questions}
\item À 50 km/h : $v\approx13{,}89\unite{m/s}$ et $E_c\approx1{,}16\times10^5\unite J$. À 100 km/h : environ $4{,}63\times10^5\unite J$.
\item Même masse et vitesse doublée : énergie quadruplée.
\item Une vitesse plus grande augmente fortement l’énergie à dissiper. Le calcul suppose le même ensemble et le même référentiel. Il ne donne pas directement la distance d’arrêt. Celle-ci dépend aussi du temps de réaction, de l’adhérence et du freinage.
\end{questions}
\par\smallskip\noindent\textbf{Pour comprendre.}\enspace Conserver les valeurs non arrondies de $50/3{,}6$ et $100/3{,}6$ jusqu’au calcul final. Le rapport des vitesses est exactement 2, donc celui des énergies exactement 4 dans le modèle, même si les résultats affichés sont arrondis.\par
\pcFigureEx{figures/rev-cor-ec-facteurs.pdf}{Courbe de rapports à masse constante : le facteur appliqué à la vitesse est mis au carré pour l’énergie cinétique.}
\end{minipage}
\par\noindent\begin{minipage}{\linewidth}
\phantomsection\label{cor-activite-11}\pcSousTitre{Pistes pour l’investigation}
À masse fixée, Ec est proportionnelle à v² : le graphe Ec(v) est courbe et Ec(v²) est une droite passant par l’origine dans le modèle. En chute idéale, la baisse d’énergie potentielle est compensée par la hausse d’énergie cinétique. Avec frottements, ajouter les transferts thermiques et préciser les limites du système pour conserver l’énergie totale.\par\end{minipage}

\corrigesChapitre{12}{L’énergie : transferts et conversions}
\par\noindent\begin{minipage}{\linewidth}
\begin{corrige}{1}{Identifier une source}
\begin{questions}
\item La pile, via le circuit.
\item De l’énergie chimique.
\end{questions}
\end{corrige}
\end{minipage}
\par\noindent\begin{minipage}{\linewidth}
\begin{corrige}{2}{Nommer une forme}
\begin{questions}
\item L’énergie cinétique.
\item Une énergie potentielle liée à sa position dans le système Terre-objet.
\end{questions}
\end{corrige}
\end{minipage}
\par\noindent\begin{minipage}{\linewidth}
\begin{corrige}{3}{Une unité}
\begin{questions}
\item Le joule, J.
\item Non : le watt est une unité de puissance, étudiée au chapitre 15.
\end{questions}
\end{corrige}
\end{minipage}
\par\noindent\begin{minipage}{\linewidth}
\begin{corrige}{4}{Une sortie moins visible}
\begin{questions}
\item Un transfert par rayonnement et des transferts thermiques.
\item Non : l’énergie est transférée à l’environnement.
\end{questions}
\end{corrige}
\end{minipage}
\par\noindent\begin{minipage}{\linewidth}
\begin{corrige}{5}{Renouvelable ou fossile}
\begin{questions}
\item Le vent relève d’une ressource renouvelable ; le charbon d’une ressource fossile.
\item Non : fabrication, usages et fin de vie peuvent avoir des impacts.
\end{questions}
\end{corrige}
\end{minipage}
\par\noindent\begin{minipage}{\linewidth}
\begin{corrige}{6}{Un bilan de lampe}
\begin{questions}
\item $200-60=140\unite J$.
\item $60+140=200\unite J$ : le bilan est conservatif.
\end{questions}
\par\smallskip\noindent\textbf{Pour comprendre.}\enspace Écrire le bilan avant de soustraire : $E_{\text{entrée}}=E_{\text{utile}}+E_{\text{autres}}$ sans stockage. Les autres transferts valent $200-60=140\unite J$. Contrôle : $60+140=200\unite J$.\par
\pcFigureEx{figures/cor-energie-lampe.pdf}{Sans stockage, la somme des transferts sortants égale le transfert entrant.}
\end{corrige}
\end{minipage}
\par\noindent\begin{minipage}{\linewidth}
\begin{corrige}{7}{Le moteur}
\begin{questions}
\item Oui : 350 + 150 = 500 J.
\item La mise en mouvement d’une charge, si tel est l’objectif du dispositif.
\end{questions}
\pcFigureEx{figures/cor-energie-moteur.pdf}{Le transfert mécanique peut être utile ; le transfert thermique ne disparaît pas du bilan.}
\end{corrige}
\end{minipage}
\par\noindent\begin{minipage}{\linewidth}
\begin{corrige}{8}{Un stockage augmente}
\begin{questions}
\item Le stockage augmente de $80-30=50\unite J$.
\item Non : cette égalité ne convient que lorsque le stockage ne varie pas.
\end{questions}
\par\smallskip\noindent\textbf{Pour comprendre.}\enspace La conservation s’écrit ici $\Delta E_{\text{stockée}}=E_{\text{entrée}}-E_{\text{sortie}}=80-30=50\unite J$. Ce n’est pas une disparition des 50 J : le système les stocke en plus.\par
\pcFigureEx{figures/cor-energie-stockage.pdf}{Entrée − sorties = variation de stockage : entrée et sortie ne sont pas égales si le stockage change.}
\end{corrige}
\end{minipage}
\par\noindent\begin{minipage}{\linewidth}
\begin{corrige}{9}{Une cellule photovoltaïque}
\begin{questions}
\item Elle reçoit de l’énergie par rayonnement et peut en transférer électriquement à un circuit.
\item Non : un rayonnement transporte aussi de l’énergie.
\end{questions}
\pcFigureEx{figures/cor-photovoltaique.pdf}{La lumière apporte de l’énergie par rayonnement ; la sortie électrique n’est pas la seule sortie réelle.}
\end{corrige}
\end{minipage}
\par\noindent\begin{minipage}{\linewidth}
\begin{corrige}{10}{Un bilan terrestre}
\begin{questions}
\item Son énergie stockée augmente.
\item Il faut connaître les valeurs et le périmètre des transferts ; des flèches qualitatives ne sont pas des mesures.
\end{questions}
\end{corrige}
\end{minipage}
\par\noindent\begin{minipage}{\linewidth}
\begin{corrige}{11}{Comparer deux lampes}
\begin{questions}
\item $100-40=60\unite J$, soit 60 \% de l’énergie de A.
\item L’éclairage utile, la durée et le périmètre de comparaison doivent être comparables.
\end{questions}
\end{corrige}
\end{minipage}
\par\noindent\begin{minipage}{\linewidth}
\begin{corrige}{12}{Une énergie cinétique fournie}
\begin{questions}
\item $E_c=\frac12\times2\times3^2=9\unite J$.
\item Non : avec $v=6\unite{m/s}$, on obtient $36\unite J$, soit quatre fois plus. La formule relie la masse, la vitesse et l’énergie cinétique.
\end{questions}
\end{corrige}
\end{minipage}
\par\noindent\begin{minipage}{\linewidth}
\begin{corrige}{13}{Le mot « perte »}
\begin{questions}
\item Les transferts thermiques vers l’environnement conservent l’énergie, mais ne répondent pas au but de mise en mouvement.
\item Non : les formes et conditions de transfert ne sont pas également utilisables ; récupérer une énergie dispersée peut être difficile.
\end{questions}
\end{corrige}
\end{minipage}
\par\noindent\begin{minipage}{\linewidth}
\begin{corrige}{14}{Un raccourci sur le climat}
\begin{questions}
\item Ils absorbent et émettent une partie du rayonnement infrarouge ; ils ne bloquent pas toute sortie d’énergie.
\item Il faut comparer l’ensemble des transferts entrants et sortants pour le système Terre-atmosphère.
\end{questions}
\end{corrige}
\end{minipage}
\par\noindent\begin{minipage}{\linewidth}
\begin{corrige}{15}{Une comparaison incomplète}
\begin{questions}
\item Non : l’usage n’est qu’une partie du cycle de vie.
\item Extraction des matériaux, fabrication, transport, énergie fournie, durée de vie, réparabilité et traitement final.
\end{questions}
\end{corrige}
\end{minipage}
\clearpage
\par\noindent\begin{minipage}{\linewidth}
\phantomsection\label{cor-defi-12}\pcSolutionDefi{12}{Le convertisseur mystérieux}\begin{questions}
\item $800-500-220=80\unite J$ stockés en plus, si les transferts indiqués sont complets.
\item L’énergie n’a pas disparu ; le système a pu stocker 80 J ou un transfert a été oublié.
\item Examiner un échauffement interne, l’évolution du fonctionnement et les autres sorties ; le seul écart de bilan ne prouve pas un mécanisme.
\end{questions}
\pcFigureEx{figures/cor-convertisseur.pdf}{Le stockage de 80 J est déduit seulement si les transferts indiqués sont complets ; une sortie oubliée reste une autre hypothèse.}
\end{minipage}
\par\noindent\begin{minipage}{\linewidth}
\phantomsection\label{cor-activite-12}\pcSousTitre{Pistes pour l’investigation}
Pour chaque convertisseur, nommer le système, la source et les transferts entrant et sortant. Une lampe reçoit électriquement et cède notamment par rayonnement et thermiquement ; un moteur transfère aussi mécaniquement. Un échauffement indique une sortie ou un stockage à discuter. Sans valeurs complètes, présenter un bilan qualitatif plutôt qu’un partage chiffré inventé.\par\end{minipage}

\corrigesChapitre{13}{Circuits : intensités et tensions}
\par\noindent\begin{minipage}{\linewidth}
\begin{corrige}{1}{Le branchement}
\begin{questions}
\item En série dans la branche.
\item En dérivation aux bornes du dipôle.
\end{questions}
\pcFigureEx{figures/rev-cor-mesures.pdf}{L’ampèremètre est traversé par le courant de la branche ; le voltmètre est relié aux deux bornes du récepteur.}
\end{corrige}
\end{minipage}
\par\noindent\begin{minipage}{\linewidth}
\begin{corrige}{2}{Conversions}
\begin{questions}
\item $0{,}250\unite A$.
\item $32\unite{mA}$.
\end{questions}
\end{corrige}
\end{minipage}
\par\noindent\begin{minipage}{\linewidth}
\begin{corrige}{3}{En série}
\begin{questions}
\item $0{,}18\unite A$ en régime permanent.
\item Non ; la lampe reçoit et convertit de l’énergie.
\end{questions}
\pcFigureEx{figures/rev-cor-serie-018.pdf}{En régime permanent, la lampe ne consomme pas le courant ; l’intensité est la même avant et après.}
\end{corrige}
\end{minipage}
\par\noindent\begin{minipage}{\linewidth}
\begin{corrige}{4}{Deux nœuds}
\begin{questions}
\item En dérivation.
\item La tension, avec la même convention de mesure.
\end{questions}
\pcFigureEx{figures/rev-cor-parallele-12.pdf}{Deux branches reliées aux mêmes deux nœuds ont la même tension avec la même convention de mesure.}
\end{corrige}
\end{minipage}
\par\noindent\begin{minipage}{\linewidth}
\begin{corrige}{5}{Une somme}
\begin{questions}
\item $0{,}35\unite A$.
\item $350\unite{mA}$.
\end{questions}
\pcFigureEx{figures/rev-cor-noeud-035.pdf}{Somme des intensités entrantes égale à la somme des intensités sortantes.}
\end{corrige}
\end{minipage}
\par\noindent\begin{minipage}{\linewidth}
\begin{corrige}{6}{Une branche inconnue}
\begin{questions}
\item $I_2=0{,}80-0{,}250=0{,}550\unite A$.
\item On mélangerait A et mA.
\end{questions}
\par\smallskip\noindent\textbf{Pour comprendre.}\enspace La loi des nœuds donne $I_G=I_1+I_2$. Avec $I_1=0{,}250\unite A$, $I_2=0{,}80-0{,}250=0{,}550\unite A=550\unite{mA}$.\par
\pcFigureEx{figures/rev-cor-noeud-080.pdf}{250 mA = 0,250 A ; le courant de la branche inconnue est 0,550 A.}
\end{corrige}
\end{minipage}
\par\noindent\begin{minipage}{\linewidth}
\begin{corrige}{7}{Tensions en série}
\begin{questions}
\item $U_2=9-3{,}5=5{,}5\unite V$.
\item Aux bornes du second récepteur, en dérivation.
\end{questions}
\par\smallskip\noindent\textbf{Pour comprendre.}\enspace Les tensions des récepteurs en série s’ajoutent : $U_2=U_G-U_1=9-3{,}5=5{,}5\unite V$. Le voltmètre se connecte aux deux bornes de ce second récepteur.\par
\pcFigureEx{figures/rev-cor-serie-9V.pdf}{Le voltmètre mesure la tension du second récepteur en dérivation ; 3,5 + 5,5 = 9 V.}
\end{corrige}
\end{minipage}
\par\noindent\begin{minipage}{\linewidth}
\begin{corrige}{8}{Tensions en dérivation}
\begin{questions}
\item $U_1=U_2=12\unite V$.
\item Non ; elles dépendent des dipôles.
\end{questions}
\pcFigureEx{figures/rev-cor-resistances-12.pdf}{La tension commune ne force pas les intensités à être égales si les résistances diffèrent.}
\end{corrige}
\end{minipage}
\par\noindent\begin{minipage}{\linewidth}
\begin{corrige}{9}{Une inversion}
\begin{questions}
\item Les bornes V et COM peuvent être inversées par rapport à la convention choisie.
\item Non ; le signe de cette mesure n’impose pas un courant nul.
\end{questions}
\pcFigureEx{figures/cor-voltmetre-inverse.pdf}{Inverser les bornes de mesure inverse le signe affiché, sans imposer un courant nul dans le récepteur.}
\end{corrige}
\end{minipage}
\par\noindent\begin{minipage}{\linewidth}
\begin{corrige}{10}{Commander séparément}
\begin{questions}
\item Deux branches en dérivation.
\item Un interrupteur en série avec sa lampe dans chaque branche.
\end{questions}
\pcFigureEx{figures/cor-commandes-independantes.pdf}{Chaque interrupteur se place en série avec sa lampe dans sa propre branche.}
\end{corrige}
\end{minipage}
\par\noindent\begin{minipage}{\linewidth}
\begin{corrige}{11}{Une mesure incohérente}
\begin{questions}
\item Non ; il faut examiner protocole, stabilité et appareils.
\item Vérifier les connexions et calibres, puis mesurer dans les mêmes conditions avec un circuit stable.
\end{questions}
\end{corrige}
\end{minipage}
\par\noindent\begin{minipage}{\linewidth}
\begin{corrige}{12}{Un nœud à trois sorties}
\begin{questions}
\item $0{,}35\unite A$.
\item Non : elles sont égales dans le modèle du nœud.
\end{questions}
\pcFigureEx{figures/rev-cor-noeud-090.pdf}{0,90 − 0,20 − 0,35 = 0,35 A ; contrôler la somme des trois sorties.}
\end{corrige}
\end{minipage}
\par\noindent\begin{minipage}{\linewidth}
\begin{corrige}{13}{Une lampe retirée}
\begin{questions}
\item Oui, son chemin fermé vers la source demeure.
\item Une source idéale ou suffisamment stable dans les conditions du montage.
\end{questions}
\pcFigureEx{figures/cor-branche-ouverte.pdf}{L’autre branche conserve un trajet fermé ; une tension exactement inchangée suppose la source idéale du modèle.}
\end{corrige}
\end{minipage}
\par\noindent\begin{minipage}{\linewidth}
\begin{corrige}{14}{Court-circuit}
\begin{questions}
\item Le fil peut imposer une tension très faible à la lampe et créer une circulation de courant importante dans le circuit.
\item Couper l’alimentation et faire vérifier le montage.
\end{questions}
\pcFigureEx{figures/cor-lampe-court-circuitee.pdf}{Le fil contourne le récepteur ; couper l’alimentation avant toute modification.}
\end{corrige}
\end{minipage}
\par\noindent\begin{minipage}{\linewidth}
\begin{corrige}{15}{Un cahier des charges}
\begin{questions}
\item Dérivation sur 6 V ; source capable de fournir au moins 0,30 A dans ce modèle.
\item En série, les tensions se partagent et le courant est commun : ces conditions ne donnent pas les deux fonctionnements demandés.
\end{questions}
\pcFigureEx{figures/rev-cor-eclairage-6V.pdf}{Les deux lampes reçoivent chacune 6 V et les intensités de branches s’ajoutent au générateur.}
\end{corrige}
\end{minipage}
\par\noindent\begin{minipage}{\linewidth}
\phantomsection\label{cor-defi-13}\pcSolutionDefi{13}{Le montage à vérifier}\begin{questions}
\item La première branche totalise $4+5=9\unite V$ ; la résistance directement en dérivation reçoit aussi 9 V.
\item $I_G=0{,}150+0{,}060=0{,}210\unite A$.
\item Voltmètre aux bornes de la résistance ; ampèremètre inséré en série dans la première branche, alimentation coupée pendant la modification et contrôle des bornes/calibres.
\end{questions}
\par\smallskip\noindent\textbf{Pour comprendre.}\enspace Contrôler séparément chaque branche : $4+5=9\unite V$ pour les lampes et 9 V pour la résistance. Convertir les courants avant la somme : $0{,}150+0{,}060=0{,}210\unite A$.\par
\pcFigureEx{figures/rev-cor-montage-9V.pdf}{Une branche contient deux lampes en série ; A se place dans cette branche, V aux bornes de la résistance.}
\end{minipage}
\par\noindent\begin{minipage}{\linewidth}
\phantomsection\label{cor-activite-13}\pcSousTitre{Pistes pour l’investigation}
En série, les intensités doivent être compatibles entre elles. En dérivation, IG doit être compatible avec I1 + I2 ; les tensions entre les mêmes nœuds sont égales. Les écarts se discutent avec la stabilité de la source et la précision des appareils. Aucune mesure n’exige de brancher un ampèremètre directement sur la source.\par\end{minipage}

\corrigesChapitre{14}{Loi d’Ohm et puissance électrique}
\par\noindent\begin{minipage}{\linewidth}
\begin{corrige}{1}{Les symboles}
\begin{questions}
\item Tension et résistance.
\item Volt, V ; ohm, $\Omega$.
\end{questions}
\end{corrige}
\end{minipage}
\par\noindent\begin{minipage}{\linewidth}
\begin{corrige}{2}{Convertir}
\begin{questions}
\item $2\,200\unite\Omega$.
\item $0{,}015\unite A$.
\end{questions}
\end{corrige}
\end{minipage}
\par\noindent\begin{minipage}{\linewidth}
\begin{corrige}{3}{Trouver U}
\begin{questions}
\item $2{,}0\unite V$.
\item $U=RI$.
\end{questions}
\end{corrige}
\end{minipage}
\par\noindent\begin{minipage}{\linewidth}
\begin{corrige}{4}{Trouver I}
\begin{questions}
\item $I=0{,}020\unite A$.
\item $20\unite{mA}$.
\end{questions}
\end{corrige}
\end{minipage}
\par\noindent\begin{minipage}{\linewidth}
\begin{corrige}{5}{Trouver R}
\begin{questions}
\item $R=3/0{,}010=300\unite\Omega$.
\item Non : le quotient $0/0$ n’est pas défini.
\end{questions}
\end{corrige}
\end{minipage}
\par\noindent\begin{minipage}{\linewidth}
\begin{corrige}{6}{Une puissance}
\begin{questions}
\item $6\unite W$.
\item $6\unite J$ par seconde à puissance constante.
\end{questions}
\end{corrige}
\end{minipage}
\par\noindent\begin{minipage}{\linewidth}
\begin{corrige}{7}{Retrouver l’intensité}
\begin{questions}
\item $I=P/U=2\unite A$.
\item $1\unite A$.
\end{questions}
\end{corrige}
\end{minipage}
\par\noindent\begin{minipage}{\linewidth}
\begin{corrige}{8}{Une caractéristique}
\begin{questions}
\item $U/I=100\unite\Omega$ pour chaque point.
\item Vérifier l’alignement sur une droite proche de l’origine, avec dispersion compatible avec les mesures.
\end{questions}
\par\smallskip\noindent\textbf{Pour comprendre.}\enspace Former $1/0{,}01$, $2/0{,}02$ et $3/0{,}03$ : on obtient chaque fois 100 ohms. Le contrôle graphique porte aussi sur l’origine ; trois quotients cohérents ne valident pas des conditions jamais testées.\par
\pcFigureEx{figures/rev-cor-ohm-100.pdf}{Le quotient U/I vaut 100 ohms pour chacun des trois points ; le modèle est une droite passant par l’origine.}
\end{corrige}
\end{minipage}
\par\noindent\begin{minipage}{\linewidth}
\begin{corrige}{9}{Une limite dépassée}
\begin{questions}
\item $I=0{,}10\unite A$ ; $P=1{,}0\unite W$.
\item Non, la puissance calculée est quatre fois la limite ; ne pas réaliser ce montage.
\end{questions}
\par\smallskip\noindent\textbf{Pour comprendre.}\enspace $I=U/R=10/100=0{,}10\unite A$, puis $P=UI=1{,}0\unite W$. Comparer dans la même unité : $1{,}0/0{,}25=4$, donc la limite est dépassée.\par
\end{corrige}
\end{minipage}
\par\noindent\begin{minipage}{\linewidth}
\begin{corrige}{10}{À tension constante}
\begin{questions}
\item I est divisée par deux.
\item P est également divisée par deux dans ce modèle.
\end{questions}
\end{corrige}
\end{minipage}
\par\noindent\begin{minipage}{\linewidth}
\begin{corrige}{11}{À résistance constante}
\begin{questions}
\item I double.
\item P est multipliée par quatre, car U et I doublent.
\end{questions}
\end{corrige}
\end{minipage}
\par\noindent\begin{minipage}{\linewidth}
\begin{corrige}{12}{Une lampe réelle}
\begin{questions}
\item Non : les quotients diffèrent au-delà des éventuelles erreurs à examiner.
\item L’échauffement du filament peut modifier sa résistance.
\end{questions}
\pcFigureEx{figures/rev-cor-lampe-reelle.pdf}{Exemple de points illustrant les quotients 20 et 35 ohms donnés ; ce ne sont pas de nouvelles mesures de l’énoncé.}
\end{corrige}
\end{minipage}
\par\noindent\begin{minipage}{\linewidth}
\begin{corrige}{13}{Un graphique trompeur}
\begin{questions}
\item Non ; la pente I/U vaut $1/R$ pour le modèle ohmique.
\item Calculer $U/I$ avec I en A.
\end{questions}
\par\smallskip\noindent\textbf{Pour comprendre.}\enspace La pente se forme comme ordonnée/abscisse. Si I est vertical et U horizontal, elle vaut I/U en A/V. La résistance reste U/I en V/A, c’est-à-dire en ohms.\par
\pcFigureEx{figures/rev-cor-ohm-inverse.pdf}{Dans ces axes, la pente est I/U ; on retrouve R en prenant le quotient inverse U/I.}
\end{corrige}
\end{minipage}
\par\noindent\begin{minipage}{\linewidth}
\begin{corrige}{14}{Deux branches}
\begin{questions}
\item $I_1=0{,}060\unite A$ et $I_2=0{,}030\unite A$.
\item $P=6(0{,}060+0{,}030)=0{,}54\unite W$.
\end{questions}
\par\smallskip\noindent\textbf{Pour comprendre.}\enspace $I_1=6/100=0{,}060\unite A$, $I_2=6/200=0{,}030\unite A$. Le générateur fournit $I_G=0{,}090\unite A$, donc $P_G=UI_G=0{,}54\unite W$.\par
\pcFigureEx{figures/rev-cor-resistances-6.pdf}{Même tension de 6 V, courants de 60 mA et 30 mA ; la puissance totale vaut 6 × 0,090 = 0,54 W.}
\end{corrige}
\end{minipage}
\par\noindent\begin{minipage}{\linewidth}
\begin{corrige}{15}{Choisir un dipôle}
\begin{questions}
\item $R=250\unite\Omega$ ; $P=0{,}10\unite W$.
\item Oui, le calcul est inférieur à la limite ; la marge reste faible et les tolérances réelles doivent être vérifiées par le responsable du montage.
\end{questions}
\end{corrige}
\end{minipage}
\par\noindent\begin{minipage}{\linewidth}
\phantomsection\label{cor-defi-14}\pcSolutionDefi{14}{Une résistance pour le laboratoire}\begin{questions}
\item 100 Ω : 60 mA et 0,36 W. 220 Ω : environ 27,3 mA et 0,164 W.
\item 220 Ω respecte la limite ; 100 Ω la dépasse.
\item A en série, V aux bornes de la résistance, calibres et bornes vérifiés ; couper avant modification et rester dans les limites du matériel.
\end{questions}
\par\smallskip\noindent\textbf{Pour comprendre.}\enspace Pour chaque résistance, appliquer $I=U/R$ puis $P=UI$. Sous 6 V : $P_{100}=0{,}36\unite W>0{,}25\unite W$ ; $P_{220}\approx0{,}164\unite W<0{,}25\unite W$. Le contrôle de puissance précède la réalisation.\par
\pcFigureEx{figures/rev-cor-montage-ohm.pdf}{Mesurer U aux bornes de R et I en série ; le modèle retenu est ici 220 ohms sous 6 V.}
\end{minipage}
\par\noindent\begin{minipage}{\linewidth}
\phantomsection\label{cor-activite-14}\pcSousTitre{Pistes pour l’investigation}
Calculer U/I seulement pour I non nul, avec I en A. Un alignement proche de l’origine et des quotients sensiblement constants sont cohérents avec une résistance constante. Calculer P = UI au maximum de fonctionnement et comparer à la puissance admissible. Une dérive thermique peut modifier la résistance ; il faut respecter les limites avant de chercher une belle droite.\par\end{minipage}

\corrigesChapitre{15}{Puissance et énergie électrique}
\par\noindent\begin{minipage}{\linewidth}
\begin{corrige}{1}{Deux grandeurs}
\begin{questions}
\item Une puissance.
\item J, Wh ou kWh, selon le contexte.
\end{questions}
\end{corrige}
\end{minipage}
\par\noindent\begin{minipage}{\linewidth}
\begin{corrige}{2}{Une heure}
\begin{questions}
\item $5\unite{Wh}$.
\item $5\times3\,600=18\,000\unite J$.
\end{questions}
\par\smallskip\noindent\textbf{Pour comprendre.}\enspace À puissance constante, $E=P\Delta t=5\times1=5\unite{Wh}$. Puis $1\unite{Wh}=3\,600\unite J$ : $E=5\times3\,600=18\,000\unite J$.\par
\end{corrige}
\end{minipage}
\par\noindent\begin{minipage}{\linewidth}
\begin{corrige}{3}{Convertir}
\begin{questions}
\item $1\,500\unite{Wh}$.
\item $0{,}5\unite h$.
\end{questions}
\end{corrige}
\end{minipage}
\par\noindent\begin{minipage}{\linewidth}
\begin{corrige}{4}{Un produit électrique}
\begin{questions}
\item $P=UI=6\times0{,}20=1{,}2\unite W$.
\item Le watt.
\end{questions}
\end{corrige}
\end{minipage}
\par\noindent\begin{minipage}{\linewidth}
\begin{corrige}{5}{La durée compte}
\begin{questions}
\item A : 20 Wh contre 10 Wh.
\item Oui : à durée égale, l’énergie est proportionnelle à la puissance.
\end{questions}
\end{corrige}
\end{minipage}
\par\noindent\begin{minipage}{\linewidth}
\begin{corrige}{6}{En secondes}
\begin{questions}
\item $E=3\times40=120\unite J$.
\item La durée est en secondes ; W multiplié par s donne J.
\end{questions}
\end{corrige}
\end{minipage}
\par\noindent\begin{minipage}{\linewidth}
\begin{corrige}{7}{Une demi-heure}
\begin{questions}
\item $E=60\times0{,}5=30\unite{Wh}=0{,}030\unite{kWh}$.
\item $30\times3\,600=108\,000\unite J$.
\end{questions}
\par\smallskip\noindent\textbf{Pour comprendre.}\enspace Choisir la paire d’unités avant le produit : $30\unite{min}=0{,}5\unite h$ donne $60\times0{,}5=30\unite{Wh}$. En secondes, $30\unite{min}=1\,800\unite s$ donne $60\times1\,800=108\,000\unite J$. Les deux écritures sont cohérentes car 1 Wh = 3 600 J.\par
\end{corrige}
\end{minipage}
\par\noindent\begin{minipage}{\linewidth}
\begin{corrige}{8}{Retrouver une durée}
\begin{questions}
\item $\Delta t=E/P=180/6=30\unite s$.
\item $6\times30=180\unite J$.
\end{questions}
\par\smallskip\noindent\textbf{Pour comprendre.}\enspace Isoler la durée : $\Delta t=E/P=180/6=30\unite s$. L’unité J/W donne des secondes ; le contrôle $6\times30=180\unite J$ retrouve l’énergie reçue.\par
\end{corrige}
\end{minipage}
\par\noindent\begin{minipage}{\linewidth}
\begin{corrige}{9}{Une intensité nominale}
\begin{questions}
\item $I=P/U=2{,}4/6=0{,}40\unite A$.
\item $400\unite{mA}$.
\end{questions}
\end{corrige}
\end{minipage}
\par\noindent\begin{minipage}{\linewidth}
\begin{corrige}{10}{Un coût fictif}
\begin{questions}
\item $2{,}5\times0{,}20=0{,}50$ euro.
\item Non : le calcul ne porte que sur le poste défini dans l’énoncé.
\end{questions}
\end{corrige}
\end{minipage}
\par\noindent\begin{minipage}{\linewidth}
\begin{corrige}{11}{Puissante mais brève}
\begin{questions}
\item A : $100\times60=6\,000\unite J$ ; B : $10\times1\,200=12\,000\unite J$.
\item B reçoit deux fois plus d’énergie malgré sa puissance plus faible, car sa durée est beaucoup plus longue.
\end{questions}
\par\smallskip\noindent\textbf{Pour comprendre.}\enspace Convertir les durées avant de comparer : 1 min = 60 s et 20 min = 1 200 s. A reçoit 6 000 J, B 12 000 J. La plus grande puissance ne suffit donc pas à désigner la plus grande énergie.\par
\pcFigureEx{figures/cor-puissance-duree.pdf}{Comparer des énergies exige de tenir compte à la fois de la puissance et de la durée.}
\end{corrige}
\end{minipage}
\par\noindent\begin{minipage}{\linewidth}
\begin{corrige}{12}{Une veille}
\begin{questions}
\item $40\unite{Wh}$ par jour ; $1\,200\unite{Wh}=1{,}2\unite{kWh}$ pour 30 jours.
\item Les autres périodes de fonctionnement et éventuellement les postes inclus dans un périmètre plus large.
\end{questions}
\end{corrige}
\end{minipage}
\par\noindent\begin{minipage}{\linewidth}
\begin{corrige}{13}{Une résistance qui chauffe}
\begin{questions}
\item $I=5/100=0{,}050\unite A$ ; $P=5\times0{,}050=0{,}25\unite W$.
\item Non : $0{,}25>0{,}125\unite W$. Le composant serait sollicité au-delà de sa valeur admissible.
\end{questions}
\par\smallskip\noindent\textbf{Pour comprendre.}\enspace Calculer successivement $I=U/R=5/100=0{,}050\unite A$, puis $P=UI=5\times0{,}050=0{,}25\unite W$. La limite 0,125 W est dépassée d’un facteur deux : ne pas réaliser ce fonctionnement.\par
\end{corrige}
\end{minipage}
\par\noindent\begin{minipage}{\linewidth}
\begin{corrige}{14}{Même service}
\begin{questions}
\item A : $4\unite{kWh}$ ; B : $0{,}8\unite{kWh}$ ; économie : $3{,}2\unite{kWh}$.
\item Vérifier que qualité et quantité de lumière sont comparables ; ce calcul ne comprend pas tout le cycle de vie.
\end{questions}
\end{corrige}
\end{minipage}
\par\noindent\begin{minipage}{\linewidth}
\begin{corrige}{15}{Une unité fausse}
\begin{questions}
\item kW/h n’est pas une unité d’énergie ; il faut multiplier puissance et durée.
\item $E=2\times2=4\unite{kWh}$.
\end{questions}
\end{corrige}
\end{minipage}
\par\noindent\begin{minipage}{\linewidth}
\phantomsection\label{cor-defi-15}\pcSolutionDefi{15}{Le projet d’éclairage}\begin{questions}
\item $E=10\times12\times4\times20=9\,600\unite{Wh}=9{,}6\unite{kWh}$.
\item Nouvelle énergie : $6{,}4\unite{kWh}$ ; économie : $3{,}2\unite{kWh}$.
\item $3{,}2\times0{,}25=0{,}80$ euro. Ce gain ne décide pas seul d’un achat : prix, durée de vie et impacts doivent être considérés.
\end{questions}
\par\smallskip\noindent\textbf{Pour comprendre.}\enspace La durée d’usage vaut $4\times20=80\unite h$. Additionner les dix lampes : $P_A=120\unite W$, $P_B=80\unite W$. L’écart de 40 W sur 80 h donne $3\,200\unite{Wh}=3{,}2\unite{kWh}$ ; multiplier ensuite par le prix d’un kWh.\par
\end{minipage}
\par\noindent\begin{minipage}{\linewidth}
\phantomsection\label{cor-activite-15}\pcSousTitre{Pistes pour l’investigation}
À même durée, multiplier chaque puissance par cette durée ; convertir Wh en kWh en divisant par 1 000. Si les durées diffèrent, les produits doivent être recalculés avant de comparer. Une puissance plus faible ne garantit pas une énergie totale plus faible. Le service rendu et le périmètre de l’étude conditionnent la conclusion.\par\end{minipage}

\corrigesChapitre{16}{Lumière : propagation et distances}
\par\noindent\begin{minipage}{\linewidth}
\begin{corrige}{1}{Source ou diffuseur}
\begin{questions}
\item Non : la lampe émet sa lumière, la feuille diffuse celle qu’elle reçoit.
\item La lampe allumée.
\end{questions}
\end{corrige}
\end{minipage}
\par\noindent\begin{minipage}{\linewidth}
\begin{corrige}{2}{Le bon sens}
\begin{questions}
\item De la feuille vers l’œil.
\item Non : la vision exige la réception de lumière.
\end{questions}
\pcFigureEx{figures/cor-vision.pdf}{La lumière reçue par l’œil vient de la feuille éclairée ; l’œil n’émet pas le rayon nécessaire à cette vision.}
\end{corrige}
\end{minipage}
\par\noindent\begin{minipage}{\linewidth}
\begin{corrige}{3}{La trajectoire}
\begin{questions}
\item En ligne droite.
\item Oui : une flèche indique le sens de propagation.
\end{questions}
\end{corrige}
\end{minipage}
\par\noindent\begin{minipage}{\linewidth}
\begin{corrige}{4}{Une grande vitesse}
\begin{questions}
\item $300\,000\unite{km/s}$.
\item Non : la propagation sur une distance finie nécessite une durée non nulle.
\end{questions}
\end{corrige}
\end{minipage}
\par\noindent\begin{minipage}{\linewidth}
\begin{corrige}{5}{La zone d’ombre}
\begin{questions}
\item L’objet intercepte une partie des trajets lumineux issus de la source.
\item Non : transmission, absorption, diffusion et géométrie du milieu changent la situation.
\end{questions}
\pcFigureEx{figures/cor-ombre.pdf}{L’objet opaque intercepte une partie des rayons issus de la source ponctuelle du modèle.}
\end{corrige}
\end{minipage}
\par\noindent\begin{minipage}{\linewidth}
\begin{corrige}{6}{Un trajet en une seconde}
\begin{questions}
\item $600\,000\unite{km}$.
\item $30\,000\unite{km}$.
\end{questions}
\end{corrige}
\end{minipage}
\par\noindent\begin{minipage}{\linewidth}
\begin{corrige}{7}{Une durée solaire}
\begin{questions}
\item $500\unite s$.
\item $8\unite{min}\,20\unite s$.
\end{questions}
\par\smallskip\noindent\textbf{Pour comprendre.}\enspace $\Delta t=d/c=150\,000\,000/300\,000=500\unite s$. Comme $500=8\times60+20$, cela représente 8 min 20 s. Il s’agit d’un trajet simple Soleil-Terre.\par
\end{corrige}
\end{minipage}
\par\noindent\begin{minipage}{\linewidth}
\begin{corrige}{8}{Un aller simple lunaire}
\begin{questions}
\item $384\,000/300\,000=1{,}28\unite s$.
\item $2{,}56\unite s$ en négligeant les mouvements pendant le trajet.
\end{questions}
\pcFigureEx{figures/cor-lune-trajets.pdf}{L’aller-retour parcourt deux fois la distance Terre-Lune.}
\end{corrige}
\end{minipage}
\par\noindent\begin{minipage}{\linewidth}
\begin{corrige}{9}{Une étoile}
\begin{questions}
\item Environ 12 ans.
\item Non : on observe l’état correspondant à la lumière partie environ 12 ans auparavant.
\end{questions}
\end{corrige}
\end{minipage}
\par\noindent\begin{minipage}{\linewidth}
\begin{corrige}{10}{Une unité à convertir}
\begin{questions}
\item $1{,}892\times10^{13}\unite{km}$.
\item Deux ans, par définition de l’année-lumière.
\end{questions}
\end{corrige}
\end{minipage}
\par\noindent\begin{minipage}{\linewidth}
\begin{corrige}{11}{Une mesure réfléchie}
\begin{questions}
\item $d=c\Delta t/2=600\unite{km}$.
\item La lumière a parcouru l’aller et le retour pendant la durée mesurée.
\end{questions}
\par\smallskip\noindent\textbf{Pour comprendre.}\enspace Distance totale : $c\Delta t=300\,000\times0{,}004=1\,200\unite{km}$. Le réflecteur est à la moitié : $d=600\unite{km}$. Les unités km/s et s donnent des km.\par
\pcFigureEx{figures/cor-echo-lumineux.pdf}{La durée mesurée couvre deux trajets de longueur d : d = cΔt/2.}
\end{corrige}
\end{minipage}
\par\noindent\begin{minipage}{\linewidth}
\begin{corrige}{12}{Une ombre qui grandit}
\begin{questions}
\item Oui dans cette configuration de source ponctuelle divergente.
\item Prolonger les rayons passant par la source et les bords de l’objet jusqu’aux deux positions de l’écran.
\end{questions}
\pcFigureEx{figures/cor-ecrans.pdf}{Prolonger les rayons passant par les bords de l’objet : l’ombre grandit quand l’écran s’éloigne.}
\end{corrige}
\end{minipage}
\par\noindent\begin{minipage}{\linewidth}
\begin{corrige}{13}{Le rayon invisible}
\begin{questions}
\item Peu de lumière est alors renvoyée latéralement vers l’œil ; le rayon dessiné n’est pas une trace matérielle visible.
\item Des particules diffusent une partie de la lumière vers l’œil et peuvent rendre le trajet perceptible, selon le dispositif.
\end{questions}
\pcFigureEx{figures/cor-diffusion.pdf}{Voir le trajet de côté demande que de la lumière soit diffusée vers l’œil.}
\end{corrige}
\end{minipage}
\par\noindent\begin{minipage}{\linewidth}
\begin{corrige}{14}{Deux échelles}
\begin{questions}
\item $\Delta t=3/(3\times10^8)=10^{-8}\unite s$.
\item Ce délai est bien inférieur à la résolution et au temps de réaction d’un chronométrage manuel.
\end{questions}
\end{corrige}
\end{minipage}
\par\noindent\begin{minipage}{\linewidth}
\begin{corrige}{15}{Une conclusion médiatique}
\begin{questions}
\item Non : son évolution dépend notamment de sa nature et de son histoire, pas uniquement du délai lumineux.
\item On reçoit une information ancienne d’un temps approximativement égal à la distance en années-lumière, exprimé en années.
\end{questions}
\end{corrige}
\end{minipage}
\clearpage
\par\noindent\begin{minipage}{\linewidth}
\phantomsection\label{cor-defi-16}\pcSolutionDefi{16}{Le miroir spatial}\begin{questions}
\item $d=300\,000\times2{,}4/2=360\,000\unite{km}$.
\item 720 000 km est la longueur totale aller-retour, pas la séparation.
\item La distance augmente si le trajet et les autres délais du dispositif sont comparables. Un retard instrumental variable doit être exclu.
\end{questions}
\pcFigureEx{figures/cor-miroir-spatial.pdf}{720 000 km est la longueur totale ; la séparation est la moitié.}
\end{minipage}
\par\noindent\begin{minipage}{\linewidth}
\phantomsection\label{cor-activite-16}\pcSousTitre{Pistes pour l’investigation}
Tracer les deux rayons limites de la petite source passant par les bords de l’objet, puis les prolonger jusqu’à l’écran. La zone située entre eux derrière l’objet est dans l’ombre. En éloignant l’écran, l’ombre grandit dans cette configuration. Une source étendue donne une pénombre et limite la netteté du modèle ponctuel.\par\end{minipage}

\corrigesChapitre{17}{Le son : fréquence, propagation et écho}
\par\noindent\begin{minipage}{\linewidth}
\begin{corrige}{1}{Le milieu}
\begin{questions}
\item Non, il nécessite un milieu matériel.
\item Oui, avec une vitesse qui diffère de celle dans l’air.
\end{questions}
\end{corrige}
\end{minipage}
\par\noindent\begin{minipage}{\linewidth}
\begin{corrige}{2}{La source}
\begin{questions}
\item Une vibration transmise au milieu.
\item Non : elles oscillent localement et transmettent la perturbation.
\end{questions}
\end{corrige}
\end{minipage}
\par\noindent\begin{minipage}{\linewidth}
\begin{corrige}{3}{L’unité}
\begin{questions}
\item Le hertz, Hz.
\item Cent motifs ou cycles par seconde.
\end{questions}
\end{corrige}
\end{minipage}
\par\noindent\begin{minipage}{\linewidth}
\begin{corrige}{4}{Trois domaines}
\begin{questions}
\item Infrason, son dans la plage audible humaine usuelle, ultrason.
\item Non : elle concerne la fréquence.
\end{questions}
\end{corrige}
\end{minipage}
\par\noindent\begin{minipage}{\linewidth}
\begin{corrige}{5}{Le plus aigu}
\begin{questions}
\item Celui de 800 Hz.
\item Non : la fréquence ne donne pas à elle seule le niveau sonore.
\end{questions}
\end{corrige}
\end{minipage}
\par\noindent\begin{minipage}{\linewidth}
\begin{corrige}{6}{Une distance simple}
\begin{questions}
\item $d=340\times2=680\unite m$.
\item Non, c’est un trajet simple de la source au récepteur.
\end{questions}
\end{corrige}
\end{minipage}
\par\noindent\begin{minipage}{\linewidth}
\begin{corrige}{7}{Un écho}
\begin{questions}
\item $d=340\times0{,}50/2=85\unite m$.
\item $170\unite m$ au total.
\end{questions}
\par\smallskip\noindent\textbf{Pour comprendre.}\enspace Le produit $340\times0{,}50=170\unite m$ donne le trajet complet. La paroi est à $d=170/2=85\unite m$. Écrire le facteur deux avant d’effectuer le calcul évite de confondre les deux distances.\par
\pcFigureEx{figures/cor-echo-sonore.pdf}{Le son parcourt 170 m au total ; la paroi est à 85 m.}
\end{corrige}
\end{minipage}
\par\noindent\begin{minipage}{\linewidth}
\begin{corrige}{8}{Compter des motifs}
\begin{questions}
\item $f=8/0{,}020=400\unite{Hz}$.
\item Le son simple de 400 Hz est plus aigu, à conditions de comparaison adaptées.
\end{questions}
\par\smallskip\noindent\textbf{Pour comprendre.}\enspace Convertir 20 ms en 0,020 s puis $f=N/\Delta t=8/0{,}020=400\unite{Hz}$. Le nombre de motifs est un comptage, pas une durée. À 400 Hz, le son simple est plus aigu qu’à 200 Hz.\par
\pcFigureEx{figures/cor-huit-motifs.pdf}{On compte huit motifs complets dans 20 ms, soit 0,020 s ; ne pas confondre nombre de crêtes et nombre d’intervalles.}
\end{corrige}
\end{minipage}
\par\noindent\begin{minipage}{\linewidth}
\begin{corrige}{9}{Le tonnerre}
\begin{questions}
\item $d\approx340\times5=1\,700\unite m$.
\item Le temps de propagation de la lumière, très faible devant celui du son sur cette distance.
\end{questions}
\end{corrige}
\end{minipage}
\par\noindent\begin{minipage}{\linewidth}
\begin{corrige}{10}{Dans l’eau}
\begin{questions}
\item $d=1\,500\times0{,}20/2=150\unite m$.
\item Le milieu est l’eau ; la vitesse fournie est adaptée à ce milieu.
\end{questions}
\par\smallskip\noindent\textbf{Pour comprendre.}\enspace Utiliser la vitesse dans l’eau, puis le trajet double : $d=v\Delta t/2=1\,500\times0{,}20/2=150\unite m$. Employer la vitesse dans l’air changerait le modèle physique.\par
\pcFigureEx{figures/cor-sonar.pdf}{La durée est celle d’un aller-retour dans l’eau : la vitesse fournie est 1 500 m/s.}
\end{corrige}
\end{minipage}
\par\noindent\begin{minipage}{\linewidth}
\begin{corrige}{11}{Un signal plus grand}
\begin{questions}
\item Non : la répétition temporelle est identique.
\item La hauteur est identique ; les amplitudes et niveaux peuvent différer selon la mesure et le dispositif.
\end{questions}
\pcFigureEx{figures/cor-meme-periode.pdf}{Même durée d’un motif, donc même fréquence et même hauteur pour ces sons simples ; amplitudes différentes.}
\end{corrige}
\end{minipage}
\par\noindent\begin{minipage}{\linewidth}
\begin{corrige}{12}{Une information incomplète}
\begin{questions}
\item Non : la durée et la répétition comptent aussi, ainsi que les conditions de mesure.
\item Durée d’exposition, répétition, position du récepteur et variations de niveau ; on ne réalise pas de test fort pour vérifier.
\end{questions}
\end{corrige}
\end{minipage}
\par\noindent\begin{minipage}{\linewidth}
\begin{corrige}{13}{Deux trajets confondus}
\begin{questions}
\item Pour l’écho, la distance de la paroi est $102/2=51\unite m$.
\item Oui : sur un trajet simple, $d=v\Delta t=102\unite m$.
\end{questions}
\pcFigureEx{figures/cor-echo-030.pdf}{102 m représente le trajet aller-retour ; la distance à la paroi est 51 m.}
\end{corrige}
\end{minipage}
\par\noindent\begin{minipage}{\linewidth}
\begin{corrige}{14}{Un délai instrumental}
\begin{questions}
\item Délai de propagation : $0{,}200\unite s$ ; $d=340\times0{,}200/2=34\unite m$.
\item Le retard électronique ne correspond pas à un trajet du son et doit être retiré du temps total d’abord.
\end{questions}
\par\smallskip\noindent\textbf{Pour comprendre.}\enspace Durée réelle de propagation : $0{,}210-0{,}010=0{,}200\unite s$. Distance : $d=340\times0{,}200/2=34\unite m$. Le retard électronique n’est pas un temps passé à parcourir l’aller-retour.\par
\pcFigureEx{figures/cor-echo-retard.pdf}{Retirer le retard instrumental d’abord ; seule la durée de propagation corrigée sert à calculer la distance.}
\end{corrige}
\end{minipage}
\par\noindent\begin{minipage}{\linewidth}
\begin{corrige}{15}{Un contrôle de vraisemblance}
\begin{questions}
\item Non : 2 s correspondraient à $340\unite m$ pour un aller-retour.
\item $\Delta t=2d/v=6{,}8/340=0{,}020\unite s$. Vérifier l’unité de temps affichée.
\end{questions}
\par\smallskip\noindent\textbf{Pour comprendre.}\enspace Un écho de 2 s correspondrait à $340\times2/2=340\unite m$, pas 3,4 m. Pour 3,4 m : $\Delta t=2d/v=6{,}8/340=0{,}020\unite s$.\par
\end{corrige}
\end{minipage}
\par\noindent\begin{minipage}{\linewidth}
\phantomsection\label{cor-defi-17}\pcSolutionDefi{17}{Deux capteurs, deux milieux}\begin{questions}
\item A : $17\unite m$ ; B : $75\unite m$.
\item La vitesse de propagation dépend du milieu.
\item Le milieu, la vitesse utilisée et le caractère aller-retour du trajet, ainsi que les éventuels retards instrumentaux.
\end{questions}
\pcFigureEx{figures/cor-deux-milieux.pdf}{Même durée d’aller-retour, vitesses différentes : les deux distances ne sont pas égales.}
\end{minipage}
\par\noindent\begin{minipage}{\linewidth}
\phantomsection\label{cor-activite-17}\pcSousTitre{Pistes pour l’investigation}
Compter les motifs sur une durée connue : $f=N/\Delta t$. À durée égale, plus de motifs signifie une fréquence plus élevée et un son simple plus aigu. Changer l’amplitude ne change pas nécessairement la fréquence. Pour les distances, un écho correspond à deux trajets ; le tonnerre à un trajet simple dont le délai lumineux est négligé.\par\end{minipage}

\corrigesChapitre{18}{Signaux et information}
\par\noindent\begin{minipage}{\linewidth}
\begin{corrige}{1}{Les rôles}
\begin{questions}
\item L’avertisseur émet ; l’oreille de l’utilisateur reçoit le signal.
\item Une demande d’attention ou l’annonce d’une présence, selon la convention.
\end{questions}
\end{corrige}
\end{minipage}
\par\noindent\begin{minipage}{\linewidth}
\begin{corrige}{2}{Radio et son}
\begin{questions}
\item Non : c’est un rayonnement électromagnétique.
\item Oui, contrairement au son qui nécessite un milieu matériel.
\end{questions}
\end{corrige}
\end{minipage}
\par\noindent\begin{minipage}{\linewidth}
\begin{corrige}{3}{Deux états}
\begin{questions}
\item Deux valeurs ou catégories distinguables représentent par exemple 0 et 1.
\item Non : le sens dépend de la convention annoncée.
\end{questions}
\end{corrige}
\end{minipage}
\par\noindent\begin{minipage}{\linewidth}
\begin{corrige}{4}{Le capteur}
\begin{questions}
\item Relier une grandeur physique à un signal exploitable.
\item Pour établir ou vérifier la correspondance entre la grandeur et la sortie.
\end{questions}
\end{corrige}
\end{minipage}
\par\noindent\begin{minipage}{\linewidth}
\begin{corrige}{5}{Une perturbation}
\begin{questions}
\item Un obstacle, une mauvaise orientation ou un éclairage ambiant gênant.
\item Un obstacle, une source de bruit ou une distance trop grande.
\end{questions}
\end{corrige}
\end{minipage}
\par\noindent\begin{minipage}{\linewidth}
\begin{corrige}{6}{Lire un code}
\begin{questions}
\item B C A.
\item Ils indiquent les groupes de deux bits et empêchent une lecture ambiguë ; un protocole peut aussi fixer cette longueur.
\end{questions}
\end{corrige}
\end{minipage}
\par\noindent\begin{minipage}{\linewidth}
\begin{corrige}{7}{Compter des messages}
\begin{questions}
\item Quatre.
\item 00, 01, 10 et 11.
\end{questions}
\end{corrige}
\end{minipage}
\par\noindent\begin{minipage}{\linewidth}
\begin{corrige}{8}{Un étalonnage simple}
\begin{questions}
\item $U=1{,}2\unite V$.
\item $x=0{,}8/0{,}10=8$ dans l’unité de la grandeur donnée.
\end{questions}
\par\smallskip\noindent\textbf{Pour comprendre.}\enspace Pour x = 12, $U=0{,}10\times12=1{,}2\unite V$. Dans l’autre sens, $x=U/0{,}10=0{,}8/0{,}10=8$ dans l’unité de x. Vérifier le domaine d’étalonnage avant d’interpréter.\par
\end{corrige}
\end{minipage}
\par\noindent\begin{minipage}{\linewidth}
\begin{corrige}{9}{Une sortie affine}
\begin{questions}
\item $U=0{,}90\unite V$.
\item Non : à $T=0$, U vaut 0,50 V. C’est une relation affine dans le domaine fourni.
\end{questions}
\par\smallskip\noindent\textbf{Pour comprendre.}\enspace À 20 °C : $U=0{,}020\times20+0{,}50=0{,}90\unite V$. À 0 °C, U vaut déjà 0,50 V ; la droite ne passe donc pas par l’origine : la relation est affine, pas proportionnelle.\par
\pcFigureEx{figures/cor-etalonnage.pdf}{Une droite qui ne passe pas par l’origine n’exprime pas une proportionnalité entre U et T.}
\end{corrige}
\end{minipage}
\par\noindent\begin{minipage}{\linewidth}
\begin{corrige}{10}{Une valeur reçue}
\begin{questions}
\item $T=(1{,}30-0{,}50)/0{,}020=40\,^{\circ}\mathrm C$.
\item Oui, elle appartient à la plage annoncée.
\end{questions}
\par\smallskip\noindent\textbf{Pour comprendre.}\enspace De $U=0{,}020T+0{,}50$, déduire $T=(U-0{,}50)/0{,}020$. Pour 1,30 V, $T=0{,}80/0{,}020=40\,^{\circ}\mathrm C$, dans la plage 0-50 °C.\par
\pcFigureEx{figures/cor-etalonnage-lecture.pdf}{Enlever d’abord la tension de décalage, puis diviser par le coefficient 0,020 V/°C.}
\end{corrige}
\end{minipage}
\par\noindent\begin{minipage}{\linewidth}
\begin{corrige}{11}{Une erreur de bit}
\begin{questions}
\item D.
\item Non : 11 est un code valide. Il faut une information de contrôle supplémentaire pour repérer certains défauts.
\end{questions}
\end{corrige}
\end{minipage}
\par\noindent\begin{minipage}{\linewidth}
\begin{corrige}{12}{Répéter et comparer}
\begin{questions}
\item Au moins une transmission ou lecture est incohérente ; le deuxième groupe diffère.
\item Non à elle seule : la répétition détecte ici un désaccord, mais ne désigne pas la bonne version.
\end{questions}
\end{corrige}
\end{minipage}
\par\noindent\begin{minipage}{\linewidth}
\begin{corrige}{13}{Hors de la plage}
\begin{questions}
\item $T=(3-0{,}50)/0{,}020=125\,^{\circ}\mathrm C$.
\item Ce résultat est hors de la plage validée ; le capteur peut saturer ou changer de comportement. Il faut une autre méthode adaptée.
\end{questions}
\par\smallskip\noindent\textbf{Pour comprendre.}\enspace Le calcul donne $T=(3-0{,}50)/0{,}020=125\,^{\circ}\mathrm C$. Le modèle n’est validé que de 0 à 50 °C : l’extrapolation ne suffit pas à annoncer une température mesurée fiable.\par
\pcFigureEx{figures/cor-hors-plage.pdf}{125 °C est un résultat d’extrapolation, pas une mesure validée par l’étalonnage 0-50 °C.}
\end{corrige}
\end{minipage}
\par\noindent\begin{minipage}{\linewidth}
\begin{corrige}{14}{Une résolution fournie}
\begin{questions}
\item Pas avec cet affichage en pas de 0,5 °C.
\item Non : exactitude, étalonnage, temps de réponse et conditions de contact interviennent aussi.
\end{questions}
\end{corrige}
\end{minipage}
\par\noindent\begin{minipage}{\linewidth}
\begin{corrige}{15}{Analogique ou numérique}
\begin{questions}
\item La tension est modélisée comme analogique ; la sortie chiffrée est discrète et numérique.
\item Non : une résolution limitée fait partie du dispositif. Il faut en tenir compte dans l’interprétation.
\end{questions}
\pcFigureEx{figures/cor-analogique-numerique.pdf}{L’exemple montre des prélèvements et des arrondis prévus par le dispositif, pas des fautes humaines.}
\end{corrige}
\end{minipage}
\par\noindent\begin{minipage}{\linewidth}
\phantomsection\label{cor-defi-18}\pcSolutionDefi{18}{Le protocole du phare}\begin{questions}
\item Deux secondes, puisque chaque message comporte deux bits.
\item Par exemple un allumage plus long de durée convenue, suivi d’une séparation précise, puis les deux bits. Le récepteur doit connaître cette convention.
\item Une même perturbation peut se reproduire, et deux versions différentes ne permettent pas toujours de choisir la bonne. Il faut tester le protocole et discuter ses limites.
\end{questions}
\pcFigureEx{figures/cor-protocole-phare.pdf}{Un début plus long et une séparation convenue distinguent le silence d’un message 00 ; ici, un exemple avec le message 01.}
\end{minipage}
\par\noindent\begin{minipage}{\linewidth}
\phantomsection\label{cor-activite-18}\pcSousTitre{Pistes pour l’investigation}
Une table à deux bits peut coder quatre messages, mais un début et une séparation sont nécessaires pour distinguer 00 d’une absence de transmission. Deux répétitions différentes détectent un désaccord sans indiquer forcément la bonne version. Une erreur transformant un code valide en un autre peut rester invisible ; un protocole doit préciser ses moyens de contrôle et leurs limites.\par\end{minipage}

\global\otCorrigesfalse\global\otSouplefalse\gdef\otAccent{otNavy}\pcAnnexe{Dix brevets blancs}{Volets physique-chimie · Série générale · Préparation 2027}\phantomsection\label{brevets}\pdfbookmark[0]{Dix brevets blancs}{bm-brevets}
\begin{pcEncadre}{method}{Un entraînement au volet de physique-chimie}Chaque sujet est prévu pour 30 minutes et noté sur 10 : 9 points scientifiques et 1 point pour la langue. L’épreuve officielle de sciences dure une heure et comporte deux disciplines retenues parmi trois. Les sujets de ce livre préparent uniquement le volet physique-chimie. Ce sont des créations OptiTAB, sans statut de sujet officiel.\end{pcEncadre}
\pcSousTitre{Conditions de travail}
Prévoir une copie, un brouillon et une calculatrice de collège pour ces entraînements. À l’examen, suivre les consignes du sujet reçu. Chercher sans le cours ; réserver trois minutes à la relecture. Les démarches pertinentes, même incomplètes, peuvent rapporter les points de méthode explicités dans les corrigés.
\pcSousTitre{Le point de langue}
Les quatre dimensions sont appréciées sur la copie entière : orthographe, construction des phrases, vocabulaire scientifique et organisation de la réflexion. Pour l’autoévaluation proposée ici, chacune vaut au plus 0,25 point. Le barème scientifique est séparé. Il ne s’agit pas d’une grille ministérielle imposant cette subdivision.
\pcSousTitre{Parcours conseillé}
Sujets 1 à 3 : commencer après les chapitres correspondants. Sujets 4 à 7 : croiser les thèmes. Sujets 8 à 10 : composer en temps limité sans aide. Tous les documents et nombres d’étude sont originaux ou explicitement annoncés comme références du cours.
\pcRepere{bb-1}{Sujet 01 · Se déplacer en limitant les rejets}
\pcRepere{bb-2}{Sujet 02 · Contrôler une eau au laboratoire}
\pcRepere{bb-3}{Sujet 03 · Éclairer un refuge}
\pcRepere{bb-4}{Sujet 04 · Préparer une mission lunaire}
\pcRepere{bb-5}{Sujet 05 · Entretenir une structure métallique}
\pcRepere{bb-6}{Sujet 06 · Mesurer sous l’eau et préserver l’audition}
\pcRepere{bb-7}{Sujet 07 · Recycler et identifier des matériaux}
\pcRepere{bb-8}{Sujet 08 · Produire et économiser l’électricité}
\pcRepere{bb-9}{Sujet 09 · Observer un système planétaire}
\pcRepere{bb-10}{Sujet 10 · Mener une enquête scientifique}
\pcPage{Brevet blanc 01}{Se déplacer en limitant les rejets}
\phantomsection\label{bb-1}
\pdfbookmark[1]{01 · Se déplacer en limitant les rejets}{bm-bb-1}
\noindent\textbf{Volet PC · 30 min · /10}\hfill Partie 1/2 · 4 points scientifiques
\par\vspace{2mm}Les essais et démarches sont pris en compte. Répondre sur une copie séparée.\par
\begin{pcEncadre}{definition}{Un vélo en mouvement}Un ensemble vélo-cycliste de masse $80\unite{kg}$ parcourt une portion rectiligne de 200 m en 40 s. Le graphique fournit les positions relevées. On donne $E_c=\tfrac12 mv^2$, avec les unités SI. Les valeurs sont des données d’étude.\end{pcEncadre}
\pcFigureEx{figures/bb01-velo.pdf}{}
\par\addvspace{3mm}\noindent\textbf{Question 1 · 1 point}\quad Décrire le mouvement sur cette portion dans le référentiel de la route.\par
\par\addvspace{3mm}\noindent\textbf{Question 2 · 1 point}\quad Calculer la vitesse moyenne en m/s.\par
\par\addvspace{3mm}\noindent\textbf{Question 3 · 2 points}\quad Calculer l’énergie cinétique correspondant à cette vitesse et prévoir sa valeur si la vitesse double.\par
\vfill\pcRepere{cor-bb-1}{Corrigé, barème et remédiation}
\pcPage{Brevet blanc 01}{Se déplacer en limitant les rejets}
\noindent\textbf{Volet PC · 30 min · /10}\hfill Partie 2/2 · 5 points scientifiques
\par\vspace{2mm}Les essais et démarches sont pris en compte. Répondre sur une copie séparée.\par
\begin{pcEncadre}{definition}{Deux solutions énergétiques}Pour une comparaison fictive sur le même trajet : solution A, 20 kWh d’électricité avec une émission de cycle de vie fournie de 50 g de CO2 par kWh ; solution B, combustion émettant 4 000 g de CO2. La réaction modèle du combustible est $\ce{CH4 + 2 O2 -> CO2 + 2 H2O}$. Le tableau ne décrit que l’indicateur CO2 fourni, pas tous les impacts.\end{pcEncadre}
\pcFigureEx{figures/rev-bb1-comparaison.pdf}{Données fictives pour le même trajet et le seul indicateur fourni.}
\par\addvspace{3mm}\noindent\textbf{Question 4 · 1 point}\quad Nommer le comburant de la réaction fournie.\par
\par\addvspace{3mm}\noindent\textbf{Question 5 · 1 point}\quad Compter les atomes d’oxygène de chaque côté de l’équation.\par
\par\addvspace{3mm}\noindent\textbf{Question 6 · 1 point}\quad Calculer la masse de CO2 associée à A en kg.\par
\par\addvspace{3mm}\noindent\textbf{Question 7 · 2 points}\quad En quatre à six lignes, comparer A et B sur l’indicateur donné et préciser une limite environnementale de la conclusion.\par
\vfill\pcRepere{cor-bb-1}{Corrigé, barème et remédiation}
{\small 1 point de langue sur l’ensemble de la copie : voir la grille d’autoévaluation.}
\pcPage{Brevet blanc 02}{Contrôler une eau au laboratoire}
\phantomsection\label{bb-2}
\pdfbookmark[1]{02 · Contrôler une eau au laboratoire}{bm-bb-2}
\noindent\textbf{Volet PC · 30 min · /10}\hfill Partie 1/2 · 4 points scientifiques
\par\vspace{2mm}Les essais et démarches sont pris en compte. Répondre sur une copie séparée.\par
\begin{pcEncadre}{definition}{Masse volumique et mélange}Un récipient vide a une masse de 32,0 g. Avec 50,0 mL d’un liquide, sa masse est 82,0 g. Le liquide est transparent et contient des ions dissous. À température de l’étude, la masse volumique de l’eau pure est voisine de 1,00 g/mL.\end{pcEncadre}
\pcFigureEx{figures/rev-bb2-masse-ex.pdf}{Distinguer récipient rempli et liquide seul.}
\par\addvspace{3mm}\noindent\textbf{Question 1 · 1 point}\quad Calculer la masse de liquide.\par
\par\addvspace{3mm}\noindent\textbf{Question 2 · 1 point}\quad Déterminer sa masse volumique en g/mL.\par
\par\addvspace{3mm}\noindent\textbf{Question 3 · 2 points}\quad Peut-on affirmer que ce liquide est de l’eau pure et potable ? Donner deux arguments distincts.\par
\vfill\pcRepere{cor-bb-2}{Corrigé, barème et remédiation}
\pcPage{Brevet blanc 02}{Contrôler une eau au laboratoire}
\noindent\textbf{Volet PC · 30 min · /10}\hfill Partie 2/2 · 5 points scientifiques
\par\vspace{2mm}Les essais et démarches sont pris en compte. Répondre sur une copie séparée.\par
\begin{pcEncadre}{definition}{Les résultats d’analyse}Le pH vaut 6. Une banque de tests indique : nitrate d’argent + ions chlorure → précipité blanc ; soude + ions cuivre(II) → précipité bleu. Sur deux prélèvements, on obtient un précipité blanc au premier test et aucun bleu au second. Pour Cl, Z = 17.\end{pcEncadre}
\pcFigureEx{figures/rev-bb2-tests-ex.pdf}{Utiliser les observations et la banque de tests du document.}
\par\addvspace{3mm}\noindent\textbf{Question 4 · 1 point}\quad Classer le liquide selon son pH.\par
\par\addvspace{3mm}\noindent\textbf{Question 5 · 1 point}\quad Quel ion est mis en évidence par le précipité blanc ?\par
\par\addvspace{3mm}\noindent\textbf{Question 6 · 1 point}\quad Combien d’électrons possède cet ion monoatomique ?\par
\par\addvspace{3mm}\noindent\textbf{Question 7 · 2 points}\quad En quatre à six lignes, proposer les précautions nécessaires pour comparer les tests et préciser la portée du résultat négatif au cuivre.\par
\vfill\pcRepere{cor-bb-2}{Corrigé, barème et remédiation}
{\small 1 point de langue sur l’ensemble de la copie : voir la grille d’autoévaluation.}
\pcPage{Brevet blanc 03}{Éclairer un refuge}
\phantomsection\label{bb-3}
\pdfbookmark[1]{03 · Éclairer un refuge}{bm-bb-3}
\noindent\textbf{Volet PC · 30 min · /10}\hfill Partie 1/2 · 4 points scientifiques
\par\vspace{2mm}Les essais et démarches sont pris en compte. Répondre sur une copie séparée.\par
\begin{pcEncadre}{definition}{Le circuit d’éclairage}Une source continue idéale de 12 V alimente deux lampes en dérivation. La première est traversée par 0,30 A, la seconde par 0,20 A. Chaque lampe possède son interrupteur en série dans sa branche.\end{pcEncadre}
\pcFigureEx{figures/rev-bb3-ex.pdf}{Deux commandes indépendantes ; interrupteurs fermés pour les mesures.}
\par\addvspace{3mm}\noindent\textbf{Question 1 · 1 point}\quad Donner la tension aux bornes de chaque lampe.\par
\par\addvspace{3mm}\noindent\textbf{Question 2 · 1 point}\quad Calculer l’intensité fournie par la source.\par
\par\addvspace{3mm}\noindent\textbf{Question 3 · 2 points}\quad Décrire où brancher un voltmètre pour la première lampe et un ampèremètre pour la seconde.\par
\vfill\pcRepere{cor-bb-3}{Corrigé, barème et remédiation}
\pcPage{Brevet blanc 03}{Éclairer un refuge}
\noindent\textbf{Volet PC · 30 min · /10}\hfill Partie 2/2 · 5 points scientifiques
\par\vspace{2mm}Les essais et démarches sont pris en compte. Répondre sur une copie séparée.\par
\begin{pcEncadre}{definition}{Une réserve limitée}La puissance totale des lampes est constante et vaut 6,0 W. La batterie peut fournir au plus 36 Wh d’énergie utilisable. On néglige les autres consommations. On donne $E=P\Delta t$ et $1\unite{Wh}=3\,600\unite J$.\end{pcEncadre}
\par\addvspace{3mm}\noindent\textbf{Question 4 · 1 point}\quad Calculer l’énergie utilisée par les deux lampes en 4 h.\par
\par\addvspace{3mm}\noindent\textbf{Question 5 · 1 point}\quad Convertir cette énergie en joules.\par
\par\addvspace{3mm}\noindent\textbf{Question 6 · 1 point}\quad Déterminer la durée maximale théorique à 6,0 W.\par
\par\addvspace{3mm}\noindent\textbf{Question 7 · 2 points}\quad En quatre à six lignes, expliquer pourquoi une durée réelle peut être plus courte et proposer une économie d’énergie justifiée.\par
\vfill\pcRepere{cor-bb-3}{Corrigé, barème et remédiation}
{\small 1 point de langue sur l’ensemble de la copie : voir la grille d’autoévaluation.}
\pcPage{Brevet blanc 04}{Préparer une mission lunaire}
\phantomsection\label{bb-4}
\pdfbookmark[1]{04 · Préparer une mission lunaire}{bm-bb-4}
\noindent\textbf{Volet PC · 30 min · /10}\hfill Partie 1/2 · 4 points scientifiques
\par\vspace{2mm}Les essais et démarches sont pris en compte. Répondre sur une copie séparée.\par
\begin{pcEncadre}{definition}{Peser un équipement}Un équipement de masse 12 kg est transporté de la Terre à la Lune. On donne $g_T=10\unite{N/kg}$, $g_L=1{,}6\unite{N/kg}$ et $P=mg$. L’équipement est posé immobile sur un sol horizontal.\end{pcEncadre}
\pcFigureEx{figures/b09-poids.pdf}{}
\par\addvspace{3mm}\noindent\textbf{Question 1 · 1 point}\quad Donner sa masse sur la Lune.\par
\par\addvspace{3mm}\noindent\textbf{Question 2 · 1 point}\quad Calculer son poids sur la Lune.\par
\par\addvspace{3mm}\noindent\textbf{Question 3 · 2 points}\quad Décrire les deux forces du modèle d’équilibre sur le sol lunaire, avec leur direction, sens et valeur.\par
\vfill\pcRepere{cor-bb-4}{Corrigé, barème et remédiation}
\pcPage{Brevet blanc 04}{Préparer une mission lunaire}
\noindent\textbf{Volet PC · 30 min · /10}\hfill Partie 2/2 · 5 points scientifiques
\par\vspace{2mm}Les essais et démarches sont pris en compte. Répondre sur une copie séparée.\par
\begin{pcEncadre}{definition}{Communiquer avec la Terre}La distance Terre-Lune est prise égale à $3{,}84\times10^8\unite m$. Un signal radio voyage à $3{,}00\times10^8\unite{m/s}$ dans le vide. On néglige tous les retards électroniques. La relation de gravitation fournie est $F=Gm_A m_B/d^2$.\end{pcEncadre}
\pcFigureEx{figures/bb04-radio.pdf}{}
\par\addvspace{3mm}\noindent\textbf{Question 4 · 1 point}\quad Le signal radio est-il sonore ou électromagnétique ?\par
\par\addvspace{3mm}\noindent\textbf{Question 5 · 1 point}\quad Calculer la durée d’un trajet simple.\par
\par\addvspace{3mm}\noindent\textbf{Question 6 · 1 point}\quad Calculer la durée minimale d’un aller-retour.\par
\par\addvspace{3mm}\noindent\textbf{Question 7 · 2 points}\quad En quatre à six lignes, expliquer pourquoi une station en orbite reste soumise à la gravitation et prévoir l’effet d’une distance centrale doublée sur F, à masses fixées.\par
\vfill\pcRepere{cor-bb-4}{Corrigé, barème et remédiation}
{\small 1 point de langue sur l’ensemble de la copie : voir la grille d’autoévaluation.}
\pcPage{Brevet blanc 05}{Entretenir une structure métallique}
\phantomsection\label{bb-5}
\pdfbookmark[1]{05 · Entretenir une structure métallique}{bm-bb-5}
\noindent\textbf{Volet PC · 30 min · /10}\hfill Partie 1/2 · 4 points scientifiques
\par\vspace{2mm}Les essais et démarches sont pris en compte. Répondre sur une copie séparée.\par
\begin{pcEncadre}{definition}{Comparer des protections}Quatre clous identiques sont suivis pendant sept jours. A : air sec, pas de rouille ; B : eau au contact de l’air, rouille ; C : eau salée au contact de l’air, rouille ; D : eau désaérée maintenue sans dioxygène, pas de rouille. Les autres paramètres sont comparables.\end{pcEncadre}
\pcFigureEx{figures/b05-corrosion.pdf}{}
\par\addvspace{3mm}\noindent\textbf{Question 1 · 1 point}\quad La corrosion est-elle une transformation physique ou chimique ?\par
\par\addvspace{3mm}\noindent\textbf{Question 2 · 1 point}\quad Quel facteur est testé en comparant B et D ?\par
\par\addvspace{3mm}\noindent\textbf{Question 3 · 2 points}\quad Ces seuls résultats permettent-ils d’affirmer que le sel double la vitesse de corrosion ? Justifier et proposer une mesure supplémentaire.\par
\vfill\pcRepere{cor-bb-5}{Corrigé, barème et remédiation}
\pcPage{Brevet blanc 05}{Entretenir une structure métallique}
\noindent\textbf{Volet PC · 30 min · /10}\hfill Partie 2/2 · 5 points scientifiques
\par\vspace{2mm}Les essais et démarches sont pris en compte. Répondre sur une copie séparée.\par
\begin{pcEncadre}{definition}{Le fer dans une solution acide}La réaction fournie est $\ce{Fe + 2 H+ -> Fe^{2+} + H2}$. Le fer a Z = 26. Après réaction, un test encadré identifie le gaz ; un précipité vert avec la soude identifie Fe2+. Le pH passe de 2 à 4.\end{pcEncadre}
\par\addvspace{3mm}\noindent\textbf{Question 4 · 1 point}\quad Nommer le gaz produit.\par
\par\addvspace{3mm}\noindent\textbf{Question 5 · 1 point}\quad Déterminer le nombre d’électrons dans Fe2+.\par
\par\addvspace{3mm}\noindent\textbf{Question 6 · 1 point}\quad Le liquide final est-il neutre ? Justifier.\par
\par\addvspace{3mm}\noindent\textbf{Question 7 · 2 points}\quad En quatre à six lignes, relier les observations à la réaction fournie, puis distinguer cette réaction de la corrosion habituelle du fer dans l’air humide.\par
\vfill\pcRepere{cor-bb-5}{Corrigé, barème et remédiation}
{\small 1 point de langue sur l’ensemble de la copie : voir la grille d’autoévaluation.}
\pcPage{Brevet blanc 06}{Mesurer sous l’eau et préserver l’audition}
\phantomsection\label{bb-6}
\pdfbookmark[1]{06 · Mesurer sous l’eau et préserver l’audition}{bm-bb-6}
\noindent\textbf{Volet PC · 30 min · /10}\hfill Partie 1/2 · 4 points scientifiques
\par\vspace{2mm}Les essais et démarches sont pris en compte. Répondre sur une copie séparée.\par
\begin{pcEncadre}{definition}{Le sonar d’un bateau}Un sonar émet un signal ultrasonore de fréquence 40 kHz. Dans l’eau étudiée, sa vitesse est 1 500 m/s. L’écho provenant du fond revient après 0,080 s. Le fond est supposé horizontal et le trajet vertical.\end{pcEncadre}
\pcFigureEx{figures/rev-bb6-sonar-ex.pdf}{Le délai de 0,080 s comprend les deux trajets dans l’eau.}
\par\addvspace{3mm}\noindent\textbf{Question 1 · 1 point}\quad Convertir la fréquence en Hz et la comparer à la limite audible usuelle de 20 000 Hz.\par
\par\addvspace{3mm}\noindent\textbf{Question 2 · 1 point}\quad Quelle distance totale le signal a-t-il parcourue ?\par
\par\addvspace{3mm}\noindent\textbf{Question 3 · 2 points}\quad Calculer la profondeur et justifier le facteur 2.\par
\vfill\pcRepere{cor-bb-6}{Corrigé, barème et remédiation}
\pcPage{Brevet blanc 06}{Mesurer sous l’eau et préserver l’audition}
\noindent\textbf{Volet PC · 30 min · /10}\hfill Partie 2/2 · 5 points scientifiques
\par\vspace{2mm}Les essais et démarches sont pris en compte. Répondre sur une copie séparée.\par
\begin{pcEncadre}{definition}{Alimenter et interpréter}L’appareil fonctionne en courant continu sous 12 V avec une intensité moyenne supposée constante de 0,50 A. Il est utilisé 20 min. On donne $P=UI$, $E=P\Delta t$. Un son peut être intense même si sa fréquence est élevée ou faible.\end{pcEncadre}
\par\addvspace{3mm}\noindent\textbf{Question 4 · 1 point}\quad Calculer la puissance reçue.\par
\par\addvspace{3mm}\noindent\textbf{Question 5 · 1 point}\quad Convertir 20 min en secondes.\par
\par\addvspace{3mm}\noindent\textbf{Question 6 · 1 point}\quad Calculer l’énergie reçue en J.\par
\par\addvspace{3mm}\noindent\textbf{Question 7 · 2 points}\quad En quatre à six lignes, expliquer pourquoi la fréquence seule ne mesure pas le risque auditif et citer deux mesures de prévention adaptées à une écoute sonore.\par
\vfill\pcRepere{cor-bb-6}{Corrigé, barème et remédiation}
{\small 1 point de langue sur l’ensemble de la copie : voir la grille d’autoévaluation.}
\pcPage{Brevet blanc 07}{Recycler et identifier des matériaux}
\phantomsection\label{bb-7}
\pdfbookmark[1]{07 · Recycler et identifier des matériaux}{bm-bb-7}
\noindent\textbf{Volet PC · 30 min · /10}\hfill Partie 1/2 · 4 points scientifiques
\par\vspace{2mm}Les essais et démarches sont pris en compte. Répondre sur une copie séparée.\par
\begin{pcEncadre}{definition}{Un morceau de métal}Un bloc compact de masse 54,0 g déplace 20,0 mL d’eau quand il est complètement immergé sans bulle. Références à température donnée : aluminium 2,70 ; fer 7,8 ; cuivre 8,9 g/cm³. Un noyau d’aluminium décrit ici contient 13 protons et 14 neutrons. Les lectures sont 30,0 mL avant et 50,0 mL après immersion.\end{pcEncadre}
\pcFigureEx{figures/rev-bb7-enonce.pdf}{Les lectures de 30,0 et 50,0 mL correspondent au déplacement de 20,0 mL annoncé.}
\par\addvspace{3mm}\noindent\textbf{Question 1 · 1 point}\quad Calculer la masse volumique du bloc.\par
\par\addvspace{3mm}\noindent\textbf{Question 2 · 1 point}\quad Quelle référence est la plus compatible avec le résultat ?\par
\par\addvspace{3mm}\noindent\textbf{Question 3 · 2 points}\quad Donner Z, A et le nombre d’électrons de l’atome neutre décrit.\par
\vfill\pcRepere{cor-bb-7}{Corrigé, barème et remédiation}
\pcPage{Brevet blanc 07}{Recycler et identifier des matériaux}
\noindent\textbf{Volet PC · 30 min · /10}\hfill Partie 2/2 · 5 points scientifiques
\par\vspace{2mm}Les essais et démarches sont pris en compte. Répondre sur une copie séparée.\par
\begin{pcEncadre}{definition}{Refondre ou transformer}Un atelier pédagogique étudie un modèle de recyclage. Pour un lot, il reçoit 5,0 MJ d’énergie et en transfère 3,5 MJ utilement à la matière. Le reste est transféré vers l’environnement. La fusion est supposée ne former aucune nouvelle espèce. Les données sont fictives.\end{pcEncadre}
\par\addvspace{3mm}\noindent\textbf{Question 4 · 1 point}\quad La fusion décrite est-elle une transformation physique ou chimique ?\par
\par\addvspace{3mm}\noindent\textbf{Question 5 · 1 point}\quad Calculer l’énergie transférée vers l’environnement.\par
\par\addvspace{3mm}\noindent\textbf{Question 6 · 1 point}\quad Une masse de 54,0 g fond sans perte ni réaction : quelle masse liquide obtient-on ?\par
\par\addvspace{3mm}\noindent\textbf{Question 7 · 2 points}\quad En quatre à six lignes, expliquer pourquoi « l’énergie perdue » n’est pas détruite et proposer une limite au slogan « recycler ne coûte aucune ressource ».\par
\vfill\pcRepere{cor-bb-7}{Corrigé, barème et remédiation}
{\small 1 point de langue sur l’ensemble de la copie : voir la grille d’autoévaluation.}
\pcPage{Brevet blanc 08}{Produire et économiser l’électricité}
\phantomsection\label{bb-8}
\pdfbookmark[1]{08 · Produire et économiser l’électricité}{bm-bb-8}
\noindent\textbf{Volet PC · 30 min · /10}\hfill Partie 1/2 · 4 points scientifiques
\par\vspace{2mm}Les essais et démarches sont pris en compte. Répondre sur une copie séparée.\par
\begin{pcEncadre}{definition}{Une journée de production}Le graphique donne la puissance électrique d’un panneau sur une journée modèle. On l’approxime à 0 W avant 10 h, 200 W de 10 h à 14 h, puis 0 W après 14 h. L’énergie est stockée avec un dispositif dont les pertes sont négligées pour ce premier calcul.\end{pcEncadre}
\pcFigureEx{figures/rev-bb8-solaire-ex.pdf}{Repérer les instants de changement de puissance.}
\par\addvspace{3mm}\noindent\textbf{Question 1 · 1 point}\quad Identifier la conversion principale du panneau.\par
\par\addvspace{3mm}\noindent\textbf{Question 2 · 1 point}\quad Lire la durée de production dans le modèle.\par
\par\addvspace{3mm}\noindent\textbf{Question 3 · 2 points}\quad Calculer l’énergie fournie en Wh puis en kWh.\par
\vfill\pcRepere{cor-bb-8}{Corrigé, barème et remédiation}
\pcPage{Brevet blanc 08}{Produire et économiser l’électricité}
\noindent\textbf{Volet PC · 30 min · /10}\hfill Partie 2/2 · 5 points scientifiques
\par\vspace{2mm}Les essais et démarches sont pris en compte. Répondre sur une copie séparée.\par
\begin{pcEncadre}{definition}{Réduire une consommation}Ancienne lampe : 40 W ; nouvelle lampe : 8 W pour le même service d’éclairage étudié. Utilisation : 5 h par jour pendant 100 jours. Prix fictif : 0,25 € par kWh. On donne $E=P\Delta t$.\end{pcEncadre}
\par\addvspace{3mm}\noindent\textbf{Question 4 · 1 point}\quad Calculer la différence de puissance.\par
\par\addvspace{3mm}\noindent\textbf{Question 5 · 1 point}\quad Calculer l’énergie économisée en kWh sur les 100 jours.\par
\par\addvspace{3mm}\noindent\textbf{Question 6 · 1 point}\quad Calculer l’économie monétaire associée.\par
\par\addvspace{3mm}\noindent\textbf{Question 7 · 2 points}\quad En quatre à six lignes, expliquer pourquoi une ressource renouvelable n’assure pas une production constante et proposer un autre critère pour comparer les deux lampes.\par
\vfill\pcRepere{cor-bb-8}{Corrigé, barème et remédiation}
{\small 1 point de langue sur l’ensemble de la copie : voir la grille d’autoévaluation.}
\pcPage{Brevet blanc 09}{Observer un système planétaire}
\phantomsection\label{bb-9}
\pdfbookmark[1]{09 · Observer un système planétaire}{bm-bb-9}
\noindent\textbf{Volet PC · 30 min · /10}\hfill Partie 1/2 · 4 points scientifiques
\par\vspace{2mm}Les essais et démarches sont pris en compte. Répondre sur une copie séparée.\par
\begin{pcEncadre}{definition}{Une lumière ancienne}Une étoile est située à 4 années-lumière. On fournit $1\unite{al}=9{,}46\times10^{12}\unite{km}$ et $c=3{,}00\times10^5\unite{km/s}$. Dans un schéma, une planète reçoit la lumière de l’étoile et en diffuse une partie vers un observateur.\end{pcEncadre}
\pcFigureEx{figures/bb09-etoile.pdf}{}
\par\addvspace{3mm}\noindent\textbf{Question 1 · 1 point}\quad L’année-lumière est-elle une durée ou une distance ?\par
\par\addvspace{3mm}\noindent\textbf{Question 2 · 1 point}\quad Convertir 4 al en km.\par
\par\addvspace{3mm}\noindent\textbf{Question 3 · 2 points}\quad Distinguer étoile et planète comme sources lumineuses et préciser l’ancienneté approximative de la lumière observée de l’étoile.\par
\vfill\pcRepere{cor-bb-9}{Corrigé, barème et remédiation}
\pcPage{Brevet blanc 09}{Observer un système planétaire}
\noindent\textbf{Volet PC · 30 min · /10}\hfill Partie 2/2 · 5 points scientifiques
\par\vspace{2mm}Les essais et démarches sont pris en compte. Répondre sur une copie séparée.\par
\begin{pcEncadre}{definition}{Attraction entre deux corps}Deux corps sont modélisés par $m_A=1{,}0\times10^{20}\unite{kg}$, $m_B=2{,}0\times10^{10}\unite{kg}$ et $d=1{,}0\times10^6\unite m$. On donne $F=Gm_A m_B/d^2$, avec $G=6{,}67\times10^{-11}$ en SI.\end{pcEncadre}
\pcFigureEx{figures/rev-bb9-grav-ex.pdf}{Distance entre centres ; aucune taille n’est à l’échelle.}
\par\addvspace{3mm}\noindent\textbf{Question 4 · 1 point}\quad La force gravitationnelle est-elle attractive ou répulsive ?\par
\par\addvspace{3mm}\noindent\textbf{Question 5 · 1 point}\quad Calculer sa valeur.\par
\par\addvspace{3mm}\noindent\textbf{Question 6 · 1 point}\quad Que devient cette valeur si la distance triple à masses inchangées ?\par
\par\addvspace{3mm}\noindent\textbf{Question 7 · 2 points}\quad En quatre à six lignes, expliquer pourquoi les deux corps subissent des forces de même valeur sans que ces forces se compensent sur un seul corps, puis citer une limite du dessin.\par
\vfill\pcRepere{cor-bb-9}{Corrigé, barème et remédiation}
{\small 1 point de langue sur l’ensemble de la copie : voir la grille d’autoévaluation.}
\pcPage{Brevet blanc 10}{Mener une enquête scientifique}
\phantomsection\label{bb-10}
\pdfbookmark[1]{10 · Mener une enquête scientifique}{bm-bb-10}
\noindent\textbf{Volet PC · 30 min · /10}\hfill Partie 1/2 · 4 points scientifiques
\par\vspace{2mm}Les essais et démarches sont pris en compte. Répondre sur une copie séparée.\par
\begin{pcEncadre}{definition}{Lire une caractéristique}Le graphique U(I) provient d’une résistance scolaire. Un point correspond à I = 0,020 A et U = 4,0 V. On fournit $U=RI$ et $P=UI$ en courant continu. La puissance maximale admissible est de 0,25 W.\end{pcEncadre}
\pcFigureEx{figures/bb10-ohm.pdf}{}
\par\addvspace{3mm}\noindent\textbf{Question 1 · 1 point}\quad Quel aspect du graphe est compatible avec un conducteur ohmique ?\par
\par\addvspace{3mm}\noindent\textbf{Question 2 · 1 point}\quad Calculer la résistance avec le point fourni.\par
\par\addvspace{3mm}\noindent\textbf{Question 3 · 2 points}\quad Calculer la puissance au point donné et comparer à la limite.\par
\vfill\pcRepere{cor-bb-10}{Corrigé, barème et remédiation}
\pcPage{Brevet blanc 10}{Mener une enquête scientifique}
\noindent\textbf{Volet PC · 30 min · /10}\hfill Partie 2/2 · 5 points scientifiques
\par\vspace{2mm}Les essais et démarches sont pris en compte. Répondre sur une copie séparée.\par
\begin{pcEncadre}{definition}{Distinguer les observations}Un professeur présente deux transformations : A, eau liquide qui bout ; B, carbonate réagissant avec un acide et dégageant un gaz. Le gaz B trouble l’eau de chaux. Dans les deux cas, des bulles apparaissent. Une pesée de B en système ouvert passe de 120,0 g à 118,9 g sans projection ni autre échange.\end{pcEncadre}
\pcFigureEx{figures/rev-bb10-transfos-ex.pdf}{Deux transformations différentes peuvent produire des bulles.}
\par\addvspace{3mm}\noindent\textbf{Question 4 · 1 point}\quad Nommer le changement d’état dans A.\par
\par\addvspace{3mm}\noindent\textbf{Question 5 · 1 point}\quad Identifier le gaz mis en évidence dans B.\par
\par\addvspace{3mm}\noindent\textbf{Question 6 · 1 point}\quad Calculer la masse de gaz qui a quitté B avec les hypothèses données.\par
\par\addvspace{3mm}\noindent\textbf{Question 7 · 2 points}\quad En quatre à six lignes, expliquer pourquoi des bulles ne prouvent pas seules une réaction et proposer un contrôle permettant de renforcer l’interprétation de B.\par
\vfill\pcRepere{cor-bb-10}{Corrigé, barème et remédiation}
{\small 1 point de langue sur l’ensemble de la copie : voir la grille d’autoévaluation.}
\global\otCorrigestrue\global\otSoupletrue\gdef\otAccent{otGreenDark}\pcAnnexe{Corrigés des brevets blancs}{Des points de méthode, des réponses et des critères de réussite}\phantomsection\label{cor-brevets}\pdfbookmark[0]{Corrigés des brevets blancs}{bm-cor-brevets}
\begin{pcEncadre}{method}{Se corriger sans masquer ses erreurs}Attribuer les points à partir des critères. Conserver une démarche pertinente même si une erreur de calcul se propage ; éviter de retirer deux fois les mêmes points pour une même erreur. La proposition de correction décrit une répartition pédagogique, pas un barème d’examen.\end{pcEncadre}
\pcSousTitre{Attribuer le point de langue}
Pour chacune des quatre dimensions, attribuer 0 si la compréhension est très gênée ou le critère absent, 0,125 si la maîtrise reste fragile et 0,25 si elle est satisfaisante. Une copie courte correcte n’a pas à devenir artificiellement longue. Prendre en compte les réponses de toute la copie, particulièrement la dernière question rédigée.
\pcSousTitre{Reprendre après le bilan}
Choisir une erreur de connaissance, de modèle ou de conversion. Relire la partie de cours indiquée, refaire un exercice puis reprendre la question du sujet blanc quelques jours plus tard. Une note totale renseigne moins bien qu’une liste précise des points à retravailler.
\pcPage{Corrigé du sujet 01}{Un vélo en mouvement}
\phantomsection\label{cor-bb-1}
\pcRepere{bb-1}{Revenir à l’énoncé}
\par\noindent\begin{minipage}{\linewidth}\pcSousTitre{Question 1 · 1 point}
Rectiligne uniforme : la position croît linéairement avec le temps. \textbf{Barème :} 0,5 point pour rectiligne ; 0,5 pour uniforme et sa lecture graphique.\par
\smallskip\noindent\textbf{Démarche.}\enspace Comparer les déplacements sur des durées égales, dans le référentiel de la route annoncé.\par
\end{minipage}\par
\par\noindent\begin{minipage}{\linewidth}\pcSousTitre{Question 2 · 1 point}
$v=d/\Delta t=200/40=5\unite{m/s}$. \textbf{Barème :} 0,5 pour relation et substitution ; 0,5 pour valeur et unité.\par
\smallskip\noindent\textbf{Démarche.}\enspace Lire les extrémités : $\Delta x=200-0\unite m$ et $\Delta t=40-0\unite s$ ; le quotient donne des m/s.\par
\pcFigureEx{figures/rev-bb1-corrige.pdf}{La pente représente la distance parcourue par unité de temps.}
\end{minipage}\par
\par\noindent\begin{minipage}{\linewidth}\pcSousTitre{Question 3 · 2 points}
$E_c=0{,}5\times80\times5^2=1\,000\unite J$. Si v double à masse fixée, $E_c$ est multipliée par quatre : $4\,000\unite J$. \textbf{Barème :} 1 point pour le premier calcul ; 1 pour facteur et seconde valeur.\par
\smallskip\noindent\textbf{Démarche.}\enspace Employer kg et m/s ; mettre 5 au carré avant le produit. Doubler v multiplie son carré par quatre.\par
\end{minipage}\par
\par\noindent\begin{minipage}{\linewidth}
\begin{pcEncadre}{method}{Ma grille de bilan après les sept questions}Partie 1 : /4 ; partie 2 : /5 ; langue : /1. Total : /10. Après la correction du sujet entier, noter une réussite et une priorité de travail.\end{pcEncadre}
{\small\textbf{Pour reprendre le cours : }\hyperref[chap-3]{chapitre 3 (p.~\pageref*{chap-3})} ; \hyperref[chap-6]{chapitre 6 (p.~\pageref*{chap-6})} ; \hyperref[chap-8]{chapitre 8 (p.~\pageref*{chap-8})} ; \hyperref[chap-11]{chapitre 11 (p.~\pageref*{chap-11})} ; \hyperref[chap-12]{chapitre 12 (p.~\pageref*{chap-12})}.\par}\end{minipage}\par
\pcPage{Corrigé du sujet 01}{Deux solutions énergétiques}
\pcRepere{bb-1}{Revenir à l’énoncé}
\par\noindent\begin{minipage}{\linewidth}\pcSousTitre{Question 4 · 1 point}
Le dioxygène, $\ce{O2}$. \textbf{Barème :} 1 point pour le nom.\par
\smallskip\noindent\textbf{Démarche.}\enspace Le méthane est le combustible ; le dioxygène joue ici le rôle de comburant.\par
\end{minipage}\par
\par\noindent\begin{minipage}{\linewidth}\pcSousTitre{Question 5 · 1 point}
Quatre O de chaque côté : quatre à gauche dans $2\,\ce{O2}$ ; deux dans CO2 et deux dans $2\,\ce{H2O}$ à droite. \textbf{Barème :} 0,5 point par côté.\par
\smallskip\noindent\textbf{Démarche.}\enspace Le coefficient multiplie toute la molécule : $2\,\ce{O2}$ contient $2\times2$ atomes O ; $2\,\ce{H2O}$ en contient $2\times1$.\par
\pcFigureEx{figures/cor-reaction-methane.pdf}{Quatre atomes d’oxygène de chaque côté.}
\end{minipage}\par
\par\noindent\begin{minipage}{\linewidth}\pcSousTitre{Question 6 · 1 point}
$20\times50=1\,000\unite g=1{,}0\unite{kg}$. \textbf{Barème :} 0,5 pour produit ; 0,5 pour conversion et unité.\par
\smallskip\noindent\textbf{Démarche.}\enspace kWh × g/kWh donne g ; diviser ensuite par 1 000 pour obtenir kg.\par
\end{minipage}\par
\par\noindent\begin{minipage}{\linewidth}\pcSousTitre{Question 7 · 2 points}
A est associée à 1,0 kg de CO2 contre 4,0 kg pour B. Dans les conditions et le périmètre des données, A émet quatre fois moins de CO2. Le résultat ne mesure pas tous les impacts, par exemple matériaux, autres polluants ou biodiversité. \textbf{Barème scientifique :} 1 point pour comparaison chiffrée ; 1 pour périmètre et limite pertinente. La langue est appréciée séparément sur l’ensemble du sujet.\par
\smallskip\noindent\textbf{Démarche.}\enspace Commencer par 1,0 kg contre 4,0 kg, puis préciser le périmètre et un impact que l’indicateur ne mesure pas.\par
\end{minipage}\par
\pcPage{Corrigé du sujet 02}{Masse volumique et mélange}
\phantomsection\label{cor-bb-2}
\pcRepere{bb-2}{Revenir à l’énoncé}
\par\noindent\begin{minipage}{\linewidth}\pcSousTitre{Question 1 · 1 point}
$m=82{,}0-32{,}0=50{,}0\unite g$. \textbf{Barème :} 0,5 pour soustraction ; 0,5 pour résultat et unité.\par
\smallskip\noindent\textbf{Démarche.}\enspace La balance a pesé le récipient et le liquide : soustraire la masse du récipient.\par
\end{minipage}\par
\par\noindent\begin{minipage}{\linewidth}\pcSousTitre{Question 2 · 1 point}
$\rho=m/V=50{,}0/50{,}0=1{,}00\unite{g/mL}$. \textbf{Barème :} 0,5 pour quotient ; 0,5 pour valeur et unité.\par
\smallskip\noindent\textbf{Démarche.}\enspace Former $\rho=m/V$ avec la masse du liquide seul ; g et mL donnent g/mL.\par
\pcFigureEx{figures/rev-bb2-masse-cor.pdf}{Soustraire la masse du récipient avant le quotient.}
\end{minipage}\par
\par\noindent\begin{minipage}{\linewidth}\pcSousTitre{Question 3 · 2 points}
Non. Les ions dissous montrent qu’il ne s’agit pas d’un corps pur constitué uniquement d’eau ; une masse volumique voisine de l’eau ne prouve pas la pureté. La potabilité demande des critères chimiques et microbiologiques supplémentaires. \textbf{Barème :} 1 point pour pureté argumentée ; 1 pour potabilité distincte.\par
\smallskip\noindent\textbf{Démarche.}\enspace Les espèces dissoutes concernent la pureté ; les critères sanitaires concernent la potabilité. Donner les deux arguments.\par
\end{minipage}\par
\par\noindent\begin{minipage}{\linewidth}
\begin{pcEncadre}{method}{Ma grille de bilan après les sept questions}Partie 1 : /4 ; partie 2 : /5 ; langue : /1. Total : /10. Après la correction du sujet entier, noter une réussite et une priorité de travail.\end{pcEncadre}
{\small\textbf{Pour reprendre le cours : }\hyperref[chap-1]{chapitre 1 (p.~\pageref*{chap-1})} ; \hyperref[chap-2]{chapitre 2 (p.~\pageref*{chap-2})} ; \hyperref[chap-4]{chapitre 4 (p.~\pageref*{chap-4})}.\par}\end{minipage}\par
\pcPage{Corrigé du sujet 02}{Les résultats d’analyse}
\pcRepere{bb-2}{Revenir à l’énoncé}
\par\noindent\begin{minipage}{\linewidth}\pcSousTitre{Question 4 · 1 point}
Acide, car pH 6 est inférieur à 7 vers 25 °C. \textbf{Barème :} 0,5 pour classement ; 0,5 pour critère.\par
\smallskip\noindent\textbf{Démarche.}\enspace Comparer 6 à 7 ; la proximité avec 7 ne suffit pas à conclure à la neutralité.\par
\end{minipage}\par
\par\noindent\begin{minipage}{\linewidth}\pcSousTitre{Question 5 · 1 point}
L’ion chlorure $\ce{Cl-}$. \textbf{Barème :} 1 point pour nom ou formule non ambiguë.\par
\smallskip\noindent\textbf{Démarche.}\enspace Croiser réactif et couleur du précipité avec la banque de tests ; la couleur initiale du liquide ne remplace pas le test.\par
\end{minipage}\par
\par\noindent\begin{minipage}{\linewidth}\pcSousTitre{Question 6 · 1 point}
18 électrons : Z = 17 protons et un électron de plus pour la charge −1. \textbf{Barème :} 0,5 pour 18 ; 0,5 pour justification.\par
\smallskip\noindent\textbf{Démarche.}\enspace Cl\textsuperscript{-} conserve le même noyau que Cl et gagne un électron : $17+1=18$.\par
\end{minipage}\par
\par\noindent\begin{minipage}{\linewidth}\pcSousTitre{Question 7 · 2 points}
Utiliser des prélèvements séparés, des réactifs adaptés et des témoins. Respecter les protections et éviter de contaminer les flacons. Le test cuivre négatif indique l’absence de précipité détectable dans les conditions ; une quantité inférieure à la sensibilité du test reste possible. \textbf{Barème scientifique :} 1 point pour protocole et témoin ; 1 pour limite du négatif.\par
\smallskip\noindent\textbf{Démarche.}\enspace Prévoir prélèvements séparés, témoin et conditions comparables ; conclure en tenant compte de la limite de détection.\par
\pcFigureEx{figures/rev-bb2-tests-cor.pdf}{Un test négatif ne prouve pas une absence absolue.}
\end{minipage}\par
\pcPage{Corrigé du sujet 03}{Le circuit d’éclairage}
\phantomsection\label{cor-bb-3}
\pcRepere{bb-3}{Revenir à l’énoncé}
\par\noindent\begin{minipage}{\linewidth}\pcSousTitre{Question 1 · 1 point}
12 V pour chaque lampe, car elles sont reliées aux mêmes deux nœuds. \textbf{Barème :} 0,5 pour valeurs ; 0,5 pour loi.\par
\smallskip\noindent\textbf{Démarche.}\enspace Repérer les deux nœuds communs ; chaque lampe est directement reliée à la source par des interrupteurs fermés supposés idéaux.\par
\end{minipage}\par
\par\noindent\begin{minipage}{\linewidth}\pcSousTitre{Question 2 · 1 point}
$I_G=0{,}30+0{,}20=0{,}50\unite A$. \textbf{Barème :} 0,5 pour loi ; 0,5 pour calcul.\par
\smallskip\noindent\textbf{Démarche.}\enspace Les courants sont déjà en ampères : appliquer l’additivité des intensités au nœud.\par
\end{minipage}\par
\par\noindent\begin{minipage}{\linewidth}\pcSousTitre{Question 3 · 2 points}
Voltmètre en dérivation aux bornes de la première ; ampèremètre en série dans la seconde branche, après coupure avant modification. \textbf{Barème :} 1 point par branchement correct.\par
\smallskip\noindent\textbf{Démarche.}\enspace Insérer A dans la branche de L2 hors tension ; relier les deux bornes de V aux deux bornes de L1. A ne se branche pas seul aux bornes de la source.\par
\pcFigureEx{figures/rev-bb3-cor.pdf}{V aux bornes de L1 ; A dans la branche de L2. Chaque interrupteur reste en série avec sa lampe.}
\end{minipage}\par
\par\noindent\begin{minipage}{\linewidth}
\begin{pcEncadre}{method}{Ma grille de bilan après les sept questions}Partie 1 : /4 ; partie 2 : /5 ; langue : /1. Total : /10. Après la correction du sujet entier, noter une réussite et une priorité de travail.\end{pcEncadre}
{\small\textbf{Pour reprendre le cours : }\hyperref[chap-12]{chapitre 12 (p.~\pageref*{chap-12})} ; \hyperref[chap-13]{chapitre 13 (p.~\pageref*{chap-13})} ; \hyperref[chap-14]{chapitre 14 (p.~\pageref*{chap-14})} ; \hyperref[chap-15]{chapitre 15 (p.~\pageref*{chap-15})}.\par}\end{minipage}\par
\pcPage{Corrigé du sujet 03}{Une réserve limitée}
\pcRepere{bb-3}{Revenir à l’énoncé}
\par\noindent\begin{minipage}{\linewidth}\pcSousTitre{Question 4 · 1 point}
$E=6\times4=24\unite{Wh}$. \textbf{Barème :} 0,5 pour relation ; 0,5 pour valeur et unité.\par
\smallskip\noindent\textbf{Démarche.}\enspace Avec W et h, $E=P\Delta t$ donne des Wh.\par
\end{minipage}\par
\par\noindent\begin{minipage}{\linewidth}\pcSousTitre{Question 5 · 1 point}
$24\times3\,600=86\,400\unite J$. \textbf{Barème :} 0,5 pour conversion ; 0,5 pour résultat.\par
\smallskip\noindent\textbf{Démarche.}\enspace Multiplier par $3\,600\unite{J/Wh}$ ; on retrouve le résultat avec $6\times(4\times3\,600)$.\par
\end{minipage}\par
\par\noindent\begin{minipage}{\linewidth}\pcSousTitre{Question 6 · 1 point}
$\Delta t=36/6=6\unite h$. \textbf{Barème :} 0,5 pour division ; 0,5 pour durée et hypothèse.\par
\smallskip\noindent\textbf{Démarche.}\enspace Isoler $\Delta t=E/P$ ; Wh divisés par W donnent h. Contrôler $6\times6=36\unite{Wh}$.\par
\pcFigureEx{figures/rev-bb3-batterie.pdf}{Comparer réserve et consommation dans la même unité.}
\end{minipage}\par
\par\noindent\begin{minipage}{\linewidth}\pcSousTitre{Question 7 · 2 points}
La durée théorique suppose la réserve intégralement utilisable et les seules lampes comme récepteurs. Pertes et autres appareils peuvent réduire l’autonomie. Éteindre une lampe inutile diminue la puissance totale, donc l’énergie reçue pour une même durée. \textbf{Barème scientifique :} 1 point pour une limite liée au modèle ; 1 pour action reliée à E = PΔt.\par
\smallskip\noindent\textbf{Démarche.}\enspace Associer une limite à une hypothèse du modèle, puis une économie à $E=P\Delta t$ ; réduire la durée est aussi pertinent.\par
\end{minipage}\par
\pcPage{Corrigé du sujet 04}{Peser un équipement}
\phantomsection\label{cor-bb-4}
\pcRepere{bb-4}{Revenir à l’énoncé}
\par\noindent\begin{minipage}{\linewidth}\pcSousTitre{Question 1 · 1 point}
12 kg : le transport ne change pas la masse dans ce contexte. \textbf{Barème :} 1 point pour valeur et distinction.\par
\smallskip\noindent\textbf{Démarche.}\enspace La masse caractérise l’équipement ; g et le poids changent avec le lieu.\par
\end{minipage}\par
\par\noindent\begin{minipage}{\linewidth}\pcSousTitre{Question 2 · 1 point}
$P_L=12\times1{,}6=19{,}2\unite N$. \textbf{Barème :} 0,5 pour relation ; 0,5 pour calcul.\par
\smallskip\noindent\textbf{Démarche.}\enspace Choisir $g_L$ ; kg × N/kg donne N. La valeur terrestre ne convient pas ici.\par
\end{minipage}\par
\par\noindent\begin{minipage}{\linewidth}\pcSousTitre{Question 3 · 2 points}
Poids vertical local vers le centre lunaire et action du sol en sens opposé, de valeurs 19,2 N dans le modèle à deux forces. \textbf{Barème :} 1 point pour poids décrit ; 1 pour support et compensation.\par
\smallskip\noindent\textbf{Démarche.}\enspace Nommer les auteurs des actions sur l’équipement : Lune et sol. Dans ce modèle d’équilibre à deux forces, valeurs égales et sens opposés.\par
\pcFigureEx{figures/rev-bb4-forces.pdf}{Deux forces de 19,2 N appliquées à l’équipement ; dessin sans échelle.}
\end{minipage}\par
\par\noindent\begin{minipage}{\linewidth}
\begin{pcEncadre}{method}{Ma grille de bilan après les sept questions}Partie 1 : /4 ; partie 2 : /5 ; langue : /1. Total : /10. Après la correction du sujet entier, noter une réussite et une priorité de travail.\end{pcEncadre}
{\small\textbf{Pour reprendre le cours : }\hyperref[chap-9]{chapitre 9 (p.~\pageref*{chap-9})} ; \hyperref[chap-10]{chapitre 10 (p.~\pageref*{chap-10})} ; \hyperref[chap-16]{chapitre 16 (p.~\pageref*{chap-16})}.\par}\end{minipage}\par
\pcPage{Corrigé du sujet 04}{Communiquer avec la Terre}
\pcRepere{bb-4}{Revenir à l’énoncé}
\par\noindent\begin{minipage}{\linewidth}\pcSousTitre{Question 4 · 1 point}
Électromagnétique ; il peut se propager dans le vide. \textbf{Barème :} 1 point pour nature.\par
\smallskip\noindent\textbf{Démarche.}\enspace Le signal radio est un rayonnement électromagnétique ; un son demande un milieu matériel.\par
\end{minipage}\par
\par\noindent\begin{minipage}{\linewidth}\pcSousTitre{Question 5 · 1 point}
$\Delta t=d/c=1{,}28\unite s$. \textbf{Barème :} 0,5 pour relation ; 0,5 pour résultat.\par
\smallskip\noindent\textbf{Démarche.}\enspace Écrire $\Delta t=d/c$ ; les puissances $10^8$ se simplifient et l’unité obtenue est s.\par
\end{minipage}\par
\par\noindent\begin{minipage}{\linewidth}\pcSousTitre{Question 6 · 1 point}
$2\times1{,}28=2{,}56\unite s$. \textbf{Barème :} 0,5 pour aller-retour ; 0,5 pour durée.\par
\smallskip\noindent\textbf{Démarche.}\enspace Additionner deux trajets de 1,28 s ; ne pas diviser de nouveau une durée calculée pour un aller simple.\par
\pcFigureEx{figures/rev-bb4-radio-cor.pdf}{Chaque trajet mesure 384 000 km ; durée totale : 2,56 s.}
\end{minipage}\par
\par\noindent\begin{minipage}{\linewidth}\pcSousTitre{Question 7 · 2 points}
En orbite, la station et ses occupants restent attirés par l’astre et sont en chute commune. L’impesanteur ne signifie donc pas une force gravitationnelle nulle. Avec les mêmes masses et une distance centrale doublée, le dénominateur est multiplié par quatre ; F est divisée par quatre. \textbf{Barème scientifique :} 1 point pour explication orbitale ; 1 pour facteur justifié.\par
\smallskip\noindent\textbf{Démarche.}\enspace Expliquer la chute commune ; remplacer ensuite d par 2d : $(2d)^2=4d^2$, donc F est divisée par quatre.\par
\end{minipage}\par
\pcPage{Corrigé du sujet 05}{Comparer des protections}
\phantomsection\label{cor-bb-5}
\pcRepere{bb-5}{Revenir à l’énoncé}
\par\noindent\begin{minipage}{\linewidth}\pcSousTitre{Question 1 · 1 point}
Chimique : des espèces de la rouille se forment. \textbf{Barème :} 0,5 pour nature ; 0,5 pour formation d’espèces.\par
\smallskip\noindent\textbf{Démarche.}\enspace La formation de nouvelles espèces distingue cette transformation chimique d’un changement d’état.\par
\end{minipage}\par
\par\noindent\begin{minipage}{\linewidth}\pcSousTitre{Question 2 · 1 point}
Le rôle du dioxygène disponible, l’eau étant présente dans les deux cas. \textbf{Barème :} 1 point pour facteur.\par
\smallskip\noindent\textbf{Démarche.}\enspace B et D contiennent de l’eau ; le facteur comparé est la disponibilité du dioxygène.\par
\end{minipage}\par
\par\noindent\begin{minipage}{\linewidth}\pcSousTitre{Question 3 · 2 points}
Non : « rouille » est qualitatif et ne mesure pas une vitesse. Comparer une quantité de métal corrodé sur des durées identiques, avec répétitions, permettrait d’estimer un effet du sel. \textbf{Barème :} 1 point pour limite ; 1 pour mesure quantitative et conditions.\par
\smallskip\noindent\textbf{Démarche.}\enspace Une observation oui/non ne quantifie pas « deux fois plus vite ». Mesurer une quantité sur une durée avec répétitions et contrôles.\par
\pcFigureEx{figures/rev-cor-corrosion.pdf}{La présence de rouille ne donne pas à elle seule une vitesse.}
\end{minipage}\par
\par\noindent\begin{minipage}{\linewidth}
\begin{pcEncadre}{method}{Ma grille de bilan après les sept questions}Partie 1 : /4 ; partie 2 : /5 ; langue : /1. Total : /10. Après la correction du sujet entier, noter une réussite et une priorité de travail.\end{pcEncadre}
{\small\textbf{Pour reprendre le cours : }\hyperref[chap-2]{chapitre 2 (p.~\pageref*{chap-2})} ; \hyperref[chap-3]{chapitre 3 (p.~\pageref*{chap-3})} ; \hyperref[chap-4]{chapitre 4 (p.~\pageref*{chap-4})} ; \hyperref[chap-5]{chapitre 5 (p.~\pageref*{chap-5})}.\par}\end{minipage}\par
\pcPage{Corrigé du sujet 05}{Le fer dans une solution acide}
\pcRepere{bb-5}{Revenir à l’énoncé}
\par\noindent\begin{minipage}{\linewidth}\pcSousTitre{Question 4 · 1 point}
Dihydrogène, H2. \textbf{Barème :} 1 point.\par
\smallskip\noindent\textbf{Démarche.}\enspace Lire les produits à droite de l’équation : H₂ désigne le dihydrogène.\par
\end{minipage}\par
\par\noindent\begin{minipage}{\linewidth}\pcSousTitre{Question 5 · 1 point}
$26-2=24$ électrons. \textbf{Barème :} 0,5 pour valeur ; 0,5 pour perte de deux électrons.\par
\smallskip\noindent\textbf{Démarche.}\enspace La charge $2+$ signifie deux électrons en moins que l’atome neutre ; Z ne change pas.\par
\pcFigureEx{figures/rev-cor-atome-fe2.pdf}{26 protons et 24 électrons : charge +2.}
\end{minipage}\par
\par\noindent\begin{minipage}{\linewidth}\pcSousTitre{Question 6 · 1 point}
Non : pH 4 est encore acide vers 25 °C. \textbf{Barème :} 0,5 pour réponse ; 0,5 pour critère.\par
\smallskip\noindent\textbf{Démarche.}\enspace Le pH augmente, mais sa valeur finale 4 reste inférieure à 7.\par
\end{minipage}\par
\par\noindent\begin{minipage}{\linewidth}\pcSousTitre{Question 7 · 2 points}
Le gaz H2 et les ions Fe2+ sont les produits de l’équation. La consommation d’ions H+ est cohérente avec une hausse du pH. La corrosion habituelle dans l’air humide met notamment en jeu eau et dioxygène et forme d’autres composés ; l’équation acide-fer ne la décrit pas à elle seule. \textbf{Barème scientifique :} 1 point pour lien observations-modèle ; 1 pour distinction de situations.\par
\smallskip\noindent\textbf{Démarche.}\enspace Relier les observations aux espèces de l’équation, puis distinguer cette réaction de la formation de rouille dans l’air humide.\par
\pcFigureEx{figures/rev-bb5-reaction.pdf}{Les ions H+ sont consommés ; le pH augmente sans atteindre la neutralité.}
\end{minipage}\par
\pcPage{Corrigé du sujet 06}{Le sonar d’un bateau}
\phantomsection\label{cor-bb-6}
\pcRepere{bb-6}{Revenir à l’énoncé}
\par\noindent\begin{minipage}{\linewidth}\pcSousTitre{Question 1 · 1 point}
$40\unite{kHz}=40\,000\unite{Hz}$, au-dessus de la limite usuelle : ultrason. \textbf{Barème :} 0,5 pour conversion ; 0,5 pour comparaison.\par
\smallskip\noindent\textbf{Démarche.}\enspace Kilo signifie mille. Le classement ultrasonore porte sur la fréquence, sans renseigner un niveau en décibels.\par
\end{minipage}\par
\par\noindent\begin{minipage}{\linewidth}\pcSousTitre{Question 2 · 1 point}
$d_{total}=1\,500\times0{,}080=120\unite m$. \textbf{Barème :} 0,5 pour relation ; 0,5 pour calcul.\par
\smallskip\noindent\textbf{Démarche.}\enspace Utiliser la vitesse dans l’eau et la durée totale de 0,080 s pour la longueur totale parcourue.\par
\end{minipage}\par
\par\noindent\begin{minipage}{\linewidth}\pcSousTitre{Question 3 · 2 points}
Le signal parcourt l’aller et le retour, donc la profondeur vaut $120/2=60\unite m$. \textbf{Barème :} 1 point pour trajet double ; 1 pour profondeur et unité.\par
\smallskip\noindent\textbf{Démarche.}\enspace Écrire $2h=v\Delta t$ ; les deux trajets ont la même longueur dans le modèle annoncé.\par
\pcFigureEx{figures/rev-bb6-sonar-cor.pdf}{120 m parcourus au total ; profondeur : 60 m.}
\end{minipage}\par
\par\noindent\begin{minipage}{\linewidth}
\begin{pcEncadre}{method}{Ma grille de bilan après les sept questions}Partie 1 : /4 ; partie 2 : /5 ; langue : /1. Total : /10. Après la correction du sujet entier, noter une réussite et une priorité de travail.\end{pcEncadre}
{\small\textbf{Pour reprendre le cours : }\hyperref[chap-14]{chapitre 14 (p.~\pageref*{chap-14})} ; \hyperref[chap-15]{chapitre 15 (p.~\pageref*{chap-15})} ; \hyperref[chap-17]{chapitre 17 (p.~\pageref*{chap-17})} ; \hyperref[chap-18]{chapitre 18 (p.~\pageref*{chap-18})}.\par}\end{minipage}\par
\pcPage{Corrigé du sujet 06}{Alimenter et interpréter}
\pcRepere{bb-6}{Revenir à l’énoncé}
\par\noindent\begin{minipage}{\linewidth}\pcSousTitre{Question 4 · 1 point}
$P=12\times0{,}50=6{,}0\unite W$. \textbf{Barème :} 0,5 pour relation ; 0,5 pour calcul.\par
\smallskip\noindent\textbf{Démarche.}\enspace En courant continu, $P=UI$ : V × A donne W.\par
\end{minipage}\par
\par\noindent\begin{minipage}{\linewidth}\pcSousTitre{Question 5 · 1 point}
$20\times60=1\,200\unite s$. \textbf{Barème :} 1 point pour conversion.\par
\smallskip\noindent\textbf{Démarche.}\enspace Multiplier par 60 : chaque minute contient 60 secondes.\par
\end{minipage}\par
\par\noindent\begin{minipage}{\linewidth}\pcSousTitre{Question 6 · 1 point}
$E=6\times1\,200=7\,200\unite J$. \textbf{Barème :} 0,5 pour relation ; 0,5 pour résultat.\par
\smallskip\noindent\textbf{Démarche.}\enspace W × s donne J ; ne pas conserver 20 minutes comme si c’étaient 20 secondes.\par
\pcFigureEx{figures/rev-bb6-energie.pdf}{Convertir la durée avant le produit avec la puissance.}
\end{minipage}\par
\par\noindent\begin{minipage}{\linewidth}\pcSousTitre{Question 7 · 2 points}
La fréquence caractérise surtout la hauteur et ne donne pas le niveau sonore ni la durée d’exposition. Le risque dépend notamment de ces deux facteurs. Réduire le niveau, s’éloigner de la source et faire des pauses sont des mesures adaptées ; une protection appropriée peut être nécessaire. \textbf{Barème scientifique :} 1 point pour distinction fréquence/niveau/durée ; 0,5 par mesure cohérente, jusqu’à 1 point.\par
\smallskip\noindent\textbf{Démarche.}\enspace Distinguer hauteur, niveau et durée ; justifier les mesures par la réduction du niveau reçu ou du temps d’exposition.\par
\end{minipage}\par
\pcPage{Corrigé du sujet 07}{Un morceau de métal}
\phantomsection\label{cor-bb-7}
\pcRepere{bb-7}{Revenir à l’énoncé}
\par\noindent\begin{minipage}{\linewidth}\pcSousTitre{Question 1 · 1 point}
$\rho=54{,}0/20{,}0=2{,}70\unite{g/cm^3}$. \textbf{Barème :} 0,5 pour quotient ; 0,5 pour unité et valeur.\par
\smallskip\noindent\textbf{Démarche.}\enspace $V=50-30=20\unite{mL}=20\unite{cm^3}$. Diviser 54 g par ce volume pour obtenir 2,70 g/cm³.\par
\pcFigureEx{figures/rev-bb7-corrige.pdf}{50 − 30 = 20 cm³ ; 54 / 20 = 2,70 g/cm³.}
\end{minipage}\par
\par\noindent\begin{minipage}{\linewidth}\pcSousTitre{Question 2 · 1 point}
L’aluminium ; cette seule propriété ne prouve pas la pureté. \textbf{Barème :} 1 point pour identification compatible.\par
\smallskip\noindent\textbf{Démarche.}\enspace Comparer dans la même unité et à température comparable ; formuler une compatibilité, pas une preuve absolue d’identité.\par
\end{minipage}\par
\par\noindent\begin{minipage}{\linewidth}\pcSousTitre{Question 3 · 2 points}
$Z=13$, $A=27$ et 13 électrons. \textbf{Barème :} 0,5 pour Z ; 0,5 pour A ; 1 pour nombre d’électrons et neutralité.\par
\smallskip\noindent\textbf{Démarche.}\enspace Z compte les protons, A les protons et neutrons. Dans l’atome neutre demandé, le nombre d’électrons égale Z.\par
\pcFigureEx{figures/rev-bb7-atome.pdf}{Inventaire symbolique des constituants de l’atome neutre décrit ; aucune échelle de taille.}
\end{minipage}\par
\par\noindent\begin{minipage}{\linewidth}
\begin{pcEncadre}{method}{Ma grille de bilan après les sept questions}Partie 1 : /4 ; partie 2 : /5 ; langue : /1. Total : /10. Après la correction du sujet entier, noter une réussite et une priorité de travail.\end{pcEncadre}
{\small\textbf{Pour reprendre le cours : }\hyperref[chap-1]{chapitre 1 (p.~\pageref*{chap-1})} ; \hyperref[chap-2]{chapitre 2 (p.~\pageref*{chap-2})} ; \hyperref[chap-3]{chapitre 3 (p.~\pageref*{chap-3})} ; \hyperref[chap-12]{chapitre 12 (p.~\pageref*{chap-12})}.\par}\end{minipage}\par
\pcPage{Corrigé du sujet 07}{Refondre ou transformer}
\pcRepere{bb-7}{Revenir à l’énoncé}
\par\noindent\begin{minipage}{\linewidth}\pcSousTitre{Question 4 · 1 point}
Physique : changement d’état sans nouvelle espèce. \textbf{Barème :} 0,5 pour nature ; 0,5 pour justification.\par
\smallskip\noindent\textbf{Démarche.}\enspace Le document exclut toute nouvelle espèce : la fusion est physique dans ce modèle.\par
\end{minipage}\par
\par\noindent\begin{minipage}{\linewidth}\pcSousTitre{Question 5 · 1 point}
$5{,}0-3{,}5=1{,}5\unite{MJ}$. \textbf{Barème :} 0,5 pour bilan ; 0,5 pour résultat.\par
\smallskip\noindent\textbf{Démarche.}\enspace Sans stockage ni autre sortie, soustraire l’énergie utile de l’entrée. Vérifier $3{,}5+1{,}5=5{,}0\unite{MJ}$.\par
\pcFigureEx{figures/rev-bb7-bilan.pdf}{5,0 MJ = 3,5 MJ utiles + 1,5 MJ vers le milieu.}
\end{minipage}\par
\par\noindent\begin{minipage}{\linewidth}\pcSousTitre{Question 6 · 1 point}
54,0 g, par conservation de la masse dans les conditions données. \textbf{Barème :} 1 point.\par
\smallskip\noindent\textbf{Démarche.}\enspace Un changement d’état sans perte ni réaction conserve la masse.\par
\end{minipage}\par
\par\noindent\begin{minipage}{\linewidth}\pcSousTitre{Question 7 · 2 points}
L’énergie transférée vers l’environnement existe encore, notamment sous forme thermique, mais n’est plus utile au procédé. Le recyclage demande aussi énergie, collecte et tri. Il peut réduire certains prélèvements sans supprimer tous les besoins ni toutes les pertes de matière. \textbf{Barème scientifique :} 1 point pour conservation avec périmètre ; 1 pour besoin de ressources expliqué.\par
\smallskip\noindent\textbf{Démarche.}\enspace Préciser où l’énergie est transférée, puis citer un besoin concret du recyclage. Un transfert thermique ne disparaît pas.\par
\end{minipage}\par
\pcPage{Corrigé du sujet 08}{Une journée de production}
\phantomsection\label{cor-bb-8}
\pcRepere{bb-8}{Revenir à l’énoncé}
\par\noindent\begin{minipage}{\linewidth}\pcSousTitre{Question 1 · 1 point}
Transfert reçu par rayonnement lumineux, conversion en énergie électrique, avec d’autres sorties notamment thermiques dans la réalité. \textbf{Barème :} 1 point pour entrée et sortie principales.\par
\smallskip\noindent\textbf{Démarche.}\enspace Distinguer apport par rayonnement et sortie électrique ; les autres transferts réels restent dans le bilan.\par
\end{minipage}\par
\par\noindent\begin{minipage}{\linewidth}\pcSousTitre{Question 2 · 1 point}
4 h. \textbf{Barème :} 1 point pour lecture de 10 h à 14 h.\par
\smallskip\noindent\textbf{Démarche.}\enspace Lire la largeur du palier : $14-10=4\unite h$ ; 14 h est l’heure de fin, pas la durée.\par
\end{minipage}\par
\par\noindent\begin{minipage}{\linewidth}\pcSousTitre{Question 3 · 2 points}
$E=200\times4=800\unite{Wh}=0{,}800\unite{kWh}$. \textbf{Barème :} 1 point pour calcul Wh ; 1 pour conversion kWh.\par
\smallskip\noindent\textbf{Démarche.}\enspace La puissance constante permet $E=P\Delta t$. Diviser les Wh par 1 000 pour obtenir les kWh.\par
\pcFigureEx{figures/rev-bb8-solaire-cor.pdf}{L’aire du rectangle donne l’énergie : W × h = Wh.}
\end{minipage}\par
\par\noindent\begin{minipage}{\linewidth}
\begin{pcEncadre}{method}{Ma grille de bilan après les sept questions}Partie 1 : /4 ; partie 2 : /5 ; langue : /1. Total : /10. Après la correction du sujet entier, noter une réussite et une priorité de travail.\end{pcEncadre}
{\small\textbf{Pour reprendre le cours : }\hyperref[chap-12]{chapitre 12 (p.~\pageref*{chap-12})} ; \hyperref[chap-13]{chapitre 13 (p.~\pageref*{chap-13})} ; \hyperref[chap-14]{chapitre 14 (p.~\pageref*{chap-14})} ; \hyperref[chap-15]{chapitre 15 (p.~\pageref*{chap-15})}.\par}\end{minipage}\par
\pcPage{Corrigé du sujet 08}{Réduire une consommation}
\pcRepere{bb-8}{Revenir à l’énoncé}
\par\noindent\begin{minipage}{\linewidth}\pcSousTitre{Question 4 · 1 point}
$40-8=32\unite W$. \textbf{Barème :} 1 point.\par
\smallskip\noindent\textbf{Démarche.}\enspace Soustraire les puissances pour le même service : $\Delta P=40-8=32\unite W$.\par
\end{minipage}\par
\par\noindent\begin{minipage}{\linewidth}\pcSousTitre{Question 5 · 1 point}
$E=0{,}032\times5\times100=16\unite{kWh}$. \textbf{Barème :} 0,5 pour unités et durée ; 0,5 pour résultat.\par
\smallskip\noindent\textbf{Démarche.}\enspace La durée vaut $5\times100=500\unite h$ ; convertir 32 W en 0,032 kW avant le produit en kWh.\par
\end{minipage}\par
\par\noindent\begin{minipage}{\linewidth}\pcSousTitre{Question 6 · 1 point}
$16\times0{,}25=4{,}00$ €. \textbf{Barème :} 0,5 pour produit ; 0,5 pour coût.\par
\smallskip\noindent\textbf{Démarche.}\enspace kWh × euros/kWh donne des euros ; ce coût concerne le seul poste fourni.\par
\pcFigureEx{figures/rev-bb8-lampes.pdf}{Chaque lampe fonctionne 500 h ; économie : 16 kWh.}
\end{minipage}\par
\par\noindent\begin{minipage}{\linewidth}\pcSousTitre{Question 7 · 2 points}
La ressource solaire se renouvelle, mais la puissance dépend de l’éclairement, de la météo et de l’heure. Un stockage ou une autre production peut être nécessaire pour certains usages. Durée de vie, qualité de l’éclairage ou impacts de fabrication sont d’autres critères pertinents. \textbf{Barème scientifique :} 1 point pour variabilité expliquée ; 1 pour autre critère cohérent.\par
\smallskip\noindent\textbf{Démarche.}\enspace Relier éclairement et puissance disponible, puis ajouter un critère comme durée de vie ou qualité de lumière.\par
\end{minipage}\par
\pcPage{Corrigé du sujet 09}{Une lumière ancienne}
\phantomsection\label{cor-bb-9}
\pcRepere{bb-9}{Revenir à l’énoncé}
\par\noindent\begin{minipage}{\linewidth}\pcSousTitre{Question 1 · 1 point}
Une distance. \textbf{Barème :} 1 point.\par
\smallskip\noindent\textbf{Démarche.}\enspace L’année-lumière est la distance parcourue par la lumière en un an ; elle reste une unité de distance.\par
\end{minipage}\par
\par\noindent\begin{minipage}{\linewidth}\pcSousTitre{Question 2 · 1 point}
$4\times9{,}46\times10^{12}=3{,}784\times10^{13}\unite{km}$. \textbf{Barème :} 0,5 pour produit ; 0,5 pour résultat et unité.\par
\smallskip\noindent\textbf{Démarche.}\enspace Multiplier la distance d’une année-lumière par quatre ; conserver puissance de dix et unité km.\par
\end{minipage}\par
\par\noindent\begin{minipage}{\linewidth}\pcSousTitre{Question 3 · 2 points}
L’étoile est une source primaire ; la planète est un objet diffusant dans le modèle. La lumière de l’étoile reçue a été émise environ quatre ans auparavant. \textbf{Barème :} 1 point pour distinction ; 1 pour délai.\par
\smallskip\noindent\textbf{Démarche.}\enspace La planète diffuse la lumière reçue. Les quatre ans de retard concernent la distance étoile-observateur fournie.\par
\pcFigureEx{figures/rev-bb9-lumiere.pdf}{La distance en années-lumière indique le retard en années.}
\end{minipage}\par
\par\noindent\begin{minipage}{\linewidth}
\begin{pcEncadre}{method}{Ma grille de bilan après les sept questions}Partie 1 : /4 ; partie 2 : /5 ; langue : /1. Total : /10. Après la correction du sujet entier, noter une réussite et une priorité de travail.\end{pcEncadre}
{\small\textbf{Pour reprendre le cours : }\hyperref[chap-7]{chapitre 7 (p.~\pageref*{chap-7})} ; \hyperref[chap-10]{chapitre 10 (p.~\pageref*{chap-10})} ; \hyperref[chap-16]{chapitre 16 (p.~\pageref*{chap-16})}.\par}\end{minipage}\par
\pcPage{Corrigé du sujet 09}{Attraction entre deux corps}
\pcRepere{bb-9}{Revenir à l’énoncé}
\par\noindent\begin{minipage}{\linewidth}\pcSousTitre{Question 4 · 1 point}
Attractive. \textbf{Barème :} 1 point.\par
\smallskip\noindent\textbf{Démarche.}\enspace La force sur chaque corps est dirigée vers l’autre corps : l’interaction est attractive.\par
\end{minipage}\par
\par\noindent\begin{minipage}{\linewidth}\pcSousTitre{Question 5 · 1 point}
$F=1{,}334\times10^8\unite N$, soit environ $1{,}3\times10^8\unite N$. \textbf{Barème :} 0,5 pour substitution avec d² ; 0,5 pour calcul et unité.\par
\smallskip\noindent\textbf{Démarche.}\enspace Les données sont en SI ; mettre la distance au carré avec sa puissance de dix : $(10^6)^2=10^{12}$.\par
\end{minipage}\par
\par\noindent\begin{minipage}{\linewidth}\pcSousTitre{Question 6 · 1 point}
F est divisée par neuf, soit environ $1{,}48\times10^7\unite N$. \textbf{Barème :} 0,5 pour facteur ; 0,5 pour valeur.\par
\smallskip\noindent\textbf{Démarche.}\enspace Remplacer d par 3d donne F/9 ; utiliser la valeur non arrondie avant de diviser.\par
\end{minipage}\par
\par\noindent\begin{minipage}{\linewidth}\pcSousTitre{Question 7 · 2 points}
Les deux forces appartiennent à la même interaction mais s’exercent sur des corps différents. Elles ne sont donc pas deux forces opposées appliquées au même système. Le dessin n’est pas à l’échelle et le modèle assimile les corps à des masses ponctuelles ou sphériques. \textbf{Barème scientifique :} 1 point pour systèmes distincts ; 1 pour limite annoncée.\par
\smallskip\noindent\textbf{Démarche.}\enspace Identifier le corps qui reçoit chaque force ; une compensation demanderait des forces appliquées au même système.\par
\pcFigureEx{figures/rev-bb9-grav-cor.pdf}{Chaque corps subit 1,334 × 10⁸ N à la distance initiale ; les deux forces s’appliquent à deux systèmes différents.}
\end{minipage}\par
\pcPage{Corrigé du sujet 10}{Lire une caractéristique}
\phantomsection\label{cor-bb-10}
\pcRepere{bb-10}{Revenir à l’énoncé}
\par\noindent\begin{minipage}{\linewidth}\pcSousTitre{Question 1 · 1 point}
Une droite passant par l’origine, à la dispersion des mesures près. \textbf{Barème :} 0,5 pour droite ; 0,5 pour origine.\par
\smallskip\noindent\textbf{Démarche.}\enspace Il faut une droite et le passage par l’origine ; une droite décalée ne suffit pas pour une proportionnalité.\par
\end{minipage}\par
\par\noindent\begin{minipage}{\linewidth}\pcSousTitre{Question 2 · 1 point}
$R=4{,}0/0{,}020=200\unite\Omega$. \textbf{Barème :} 0,5 pour quotient ; 0,5 pour résultat.\par
\smallskip\noindent\textbf{Démarche.}\enspace Choisir un point hors origine, avec I en A ; le quotient U/I donne des ohms.\par
\end{minipage}\par
\par\noindent\begin{minipage}{\linewidth}\pcSousTitre{Question 3 · 2 points}
$P=4{,}0\times0{,}020=0{,}080\unite W$, inférieure à 0,25 W. \textbf{Barème :} 1 point pour calcul ; 1 pour comparaison justifiée.\par
\smallskip\noindent\textbf{Démarche.}\enspace Calculer P au même point de fonctionnement ; comparer 0,080 W à 0,25 W dans la même unité.\par
\pcFigureEx{figures/rev-bb10-ohm-cor.pdf}{R = 200 ohms ; P = 0,080 W, inférieure à 0,25 W.}
\end{minipage}\par
\par\noindent\begin{minipage}{\linewidth}
\begin{pcEncadre}{method}{Ma grille de bilan après les sept questions}Partie 1 : /4 ; partie 2 : /5 ; langue : /1. Total : /10. Après la correction du sujet entier, noter une réussite et une priorité de travail.\end{pcEncadre}
{\small\textbf{Pour reprendre le cours : }\hyperref[chap-3]{chapitre 3 (p.~\pageref*{chap-3})} ; \hyperref[chap-4]{chapitre 4 (p.~\pageref*{chap-4})} ; \hyperref[chap-8]{chapitre 8 (p.~\pageref*{chap-8})} ; \hyperref[chap-14]{chapitre 14 (p.~\pageref*{chap-14})} ; \hyperref[chap-18]{chapitre 18 (p.~\pageref*{chap-18})}.\par}\end{minipage}\par
\pcPage{Corrigé du sujet 10}{Distinguer les observations}
\pcRepere{bb-10}{Revenir à l’énoncé}
\par\noindent\begin{minipage}{\linewidth}\pcSousTitre{Question 4 · 1 point}
Vaporisation, ici ébullition. \textbf{Barème :} 1 point.\par
\smallskip\noindent\textbf{Démarche.}\enspace Nommer le passage liquide vers gaz ; l’ébullition se produit dans le volume du liquide.\par
\end{minipage}\par
\par\noindent\begin{minipage}{\linewidth}\pcSousTitre{Question 5 · 1 point}
Dioxyde de carbone. \textbf{Barème :} 1 point.\par
\smallskip\noindent\textbf{Démarche.}\enspace La conclusion vient du test à l’eau de chaux, pas des seules bulles.\par
\end{minipage}\par
\par\noindent\begin{minipage}{\linewidth}\pcSousTitre{Question 6 · 1 point}
$120{,}0-118{,}9=1{,}1\unite g$. \textbf{Barème :} 0,5 pour bilan ; 0,5 pour masse et unité.\par
\smallskip\noindent\textbf{Démarche.}\enspace La masse sortie vaut masse initiale moins masse finale dans les hypothèses du système ouvert.\par
\end{minipage}\par
\par\noindent\begin{minipage}{\linewidth}\pcSousTitre{Question 7 · 2 points}
Les bulles peuvent provenir d’un changement d’état, d’un dégazage ou d’une réaction. Dans B, le test identifie CO2 ; il faut aussi comparer les espèces initiales et finales. Un témoin sans mise en contact des réactifs et un contrôle de l’eau de chaux renforcent l’attribution. \textbf{Barème scientifique :} 1 point pour mécanismes distincts ; 1 pour contrôle pertinent.\par
\smallskip\noindent\textbf{Démarche.}\enspace Comparer les espèces avant et après ; proposer un témoin sans gaz ajouté et des conditions de test comparables.\par
\pcFigureEx{figures/rev-bb10-transfos-cor.pdf}{Le test au CO2 renforce l’interprétation chimique de B ; les bulles seules ne suffisent pas.}
\end{minipage}\par
\global\otCorrigesfalse\global\otSouplefalse\gdef\otAccent{otNavy}\pcAnnexe{Cinq sujets format annales}{Physique-chimie · 30 minutes · Sujets originaux 11 à 15}\phantomsection\label{annales-originales}\pdfbookmark[0]{Cinq sujets au format des annales}{bm-annales-originales}
\begin{pcEncadre}{method}{Des sujets complets pour composer}Chaque sujet contient trois parties indépendantes, neuf questions et tous les documents nécessaires. Durée conseillée : 30 minutes. Le total est de 9 points scientifiques, complétés par 1 point de langue selon la grille du livre. Les contextes, les mesures, les questions et les figures ont été créés pour ces entraînements. Ce ne sont pas des annales officielles reproduites.\end{pcEncadre}
\pcSousTitre{Un parcours en trois temps}
Prévoir environ 8 minutes par partie, puis 6 minutes pour lire et relire l’ensemble. Une réponse justifiée compte davantage qu’une valeur isolée. Les relations utiles et les hypothèses du modèle figurent dans les documents. Les schémas sans échelle ne doivent pas être mesurés à la règle.
\pcSousTitre{Après la composition}
Pour chaque question, le corrigé distingue la réponse, la démarche et les critères du barème. Garder les étapes justes même si une erreur de calcul se propage. Les points d’une même erreur ne sont pas retirés deux fois. Le point de langue porte sur l’ensemble des réponses, sans imposer un texte inutilement long.
\pcRepere{ann-11}{Sujet 11 · Un monte-charge pour les archives}
\pcRepere{ann-12}{Sujet 12 · Comprendre un détecteur optique de fumée}
\pcRepere{ann-13}{Sujet 13 · Une station météo sur la côte}
\pcRepere{ann-14}{Sujet 14 · Transporter un colis réfrigéré}
\pcRepere{ann-15}{Sujet 15 · Une maquette de dirigeable}
\pcPage{Sujet 11 · Énoncé}{Un monte-charge pour les archives}
\phantomsection\label{ann-11}
\pdfbookmark[1]{11 · Un monte-charge pour les archives}{bm-ann-11}
\textbf{PC · 30 min · /10}\hfill Partie 1/3 · 3 points scientifiques
\begin{pcEncadre}{definition}{Document · Suivre la cabine}Une cabine et son chargement ont une masse totale de 50 kg. La cabine monte verticalement ; son altitude est relevée toutes les secondes. Pendant la portion étudiée, les frottements sont négligés. On modélise la cabine chargée par un point G soumis à son poids et à la force du câble. $g=9{,}8\unite{N/kg}$.\end{pcEncadre}
\pcFigureEx{figures/ann11-positions.pdf}{Document original · données pédagogiques. Schéma sans échelle, sauf axes gradués.}
\par\addvspace{3mm}\noindent\begin{minipage}{\linewidth}\textbf{Question 1 · 0,5 point}\quad Décrire le mouvement de la cabine dans le référentiel du bâtiment en justifiant les deux mots employés.\par\end{minipage}\par
\par\addvspace{3mm}\noindent\begin{minipage}{\linewidth}\textbf{Question 2 · 1 point}\quad Calculer la vitesse moyenne entre 0 s et 3 s.\par\end{minipage}\par
\par\addvspace{3mm}\noindent\begin{minipage}{\linewidth}\textbf{Question 3 · 1,5 points}\quad Calculer le poids de l’ensemble. Pour le mouvement étudié, les forces se compensent. Représenter sans échelle, au point G, les deux forces et indiquer leur valeur.\par\end{minipage}\par
\vfill\pcRepere{cor-ann-11}{Corrigé expliqué, barème et bilan}
\pcPage{Sujet 11 · Énoncé}{Un monte-charge pour les archives}
\textbf{PC · 30 min · /10}\hfill Partie 2/3 · 3 points scientifiques
\begin{pcEncadre}{definition}{Document · Prévoir l’autonomie}Le moteur reçoit une tension continue de 24 V et une intensité de 3,0 A, supposées constantes pendant un cycle complet de 6 min. Une batterie fournit 18 Wh utilisables. Aucun autre appareil ne la sollicite dans le modèle. Relations : $P=UI$, $E=P\Delta t$ ; avec W et h, E est en Wh.\end{pcEncadre}
\par\addvspace{3mm}\noindent\begin{minipage}{\linewidth}\textbf{Question 4 · 1 point}\quad Calculer la puissance électrique reçue par le moteur.\par\end{minipage}\par
\par\addvspace{3mm}\noindent\begin{minipage}{\linewidth}\textbf{Question 5 · 1 point}\quad Calculer l’énergie électrique nécessaire à un cycle en Wh.\par\end{minipage}\par
\par\addvspace{3mm}\noindent\begin{minipage}{\linewidth}\textbf{Question 6 · 1 point}\quad La batterie permet-elle trois cycles complets ? Justifier, puis donner le nombre maximal de cycles dans ce modèle.\par\end{minipage}\par
\vfill\pcRepere{cor-ann-11}{Corrigé expliqué, barème et bilan}
\pcPage{Sujet 11 · Énoncé}{Un monte-charge pour les archives}
\textbf{PC · 30 min · /10}\hfill Partie 3/3 · 3 points scientifiques
\begin{pcEncadre}{definition}{Document · Identifier un liquide de maintenance}Le laboratoire pèse un récipient vide : 28,5 g. Le même récipient contenant 100,0 mL d’un liquide a une masse de 102,5 g. Références à la même température : A, 0,74 g/mL ; B, 0,90 g/mL ; C, 1,15 g/mL. On donne $\rho=m/V$.\end{pcEncadre}
\pcFigureEx{figures/ann11-pesees.pdf}{Document original · données pédagogiques. Schéma sans échelle, sauf axes gradués.}
\par\addvspace{3mm}\noindent\begin{minipage}{\linewidth}\textbf{Question 7 · 0,5 point}\quad Déterminer la masse du liquide seul.\par\end{minipage}\par
\par\addvspace{3mm}\noindent\begin{minipage}{\linewidth}\textbf{Question 8 · 1 point}\quad Calculer la masse volumique du liquide et indiquer la référence compatible.\par\end{minipage}\par
\par\addvspace{3mm}\noindent\begin{minipage}{\linewidth}\textbf{Question 9 · 1,5 points}\quad En trois à cinq lignes, expliquer pourquoi cette compatibilité ne prouve pas l’identité chimique du liquide. Proposer un contrôle supplémentaire pertinent.\par\end{minipage}\par
\vfill\pcRepere{cor-ann-11}{Corrigé expliqué, barème et bilan}
\pcPage{Sujet 12 · Énoncé}{Comprendre un détecteur optique de fumée}
\phantomsection\label{ann-12}
\pdfbookmark[1]{12 · Comprendre un détecteur optique de fumée}{bm-ann-12}
\textbf{PC · 30 min · /10}\hfill Partie 1/3 · 3 points scientifiques
\begin{pcEncadre}{definition}{Document · Un trajet lumineux modifié}Dans ce modèle de détecteur, une DEL émet de la lumière dans une petite chambre. Le capteur est placé hors du faisceau direct. Sans fumée, il ne reçoit presque pas de lumière. Des particules de fumée peuvent diffuser une partie de la lumière vers lui. Le dessin représente la présence de fumée.\end{pcEncadre}
\pcFigureEx{figures/ann12-optique.pdf}{Document original · données pédagogiques. Schéma sans échelle, sauf axes gradués.}
\par\addvspace{3mm}\noindent\begin{minipage}{\linewidth}\textbf{Question 1 · 0,5 point}\quad Identifier la source primaire de lumière et le récepteur.\par\end{minipage}\par
\par\addvspace{3mm}\noindent\begin{minipage}{\linewidth}\textbf{Question 2 · 1 point}\quad Reproduire le schéma simplifié et ajouter les deux flèches nécessaires pour expliquer l’arrivée de lumière au capteur.\par\end{minipage}\par
\par\addvspace{3mm}\noindent\begin{minipage}{\linewidth}\textbf{Question 3 · 1,5 points}\quad En trois à cinq lignes, expliquer pourquoi le capteur est placé hors du faisceau direct. Pourquoi de la poussière peut-elle aussi perturber ce modèle ?\par\end{minipage}\par
\vfill\pcRepere{cor-ann-12}{Corrigé expliqué, barème et bilan}
\pcPage{Sujet 12 · Énoncé}{Comprendre un détecteur optique de fumée}
\textbf{PC · 30 min · /10}\hfill Partie 2/3 · 3 points scientifiques
\begin{pcEncadre}{definition}{Document · Raisonner sur une combustion}Une combustion complète du propane est modélisée par $\ce{C3H8 + 5 O2 -> 3 CO2 + 4 H2O}$. Le document ne modélise pas la formation de fumée ni les produits d’une combustion incomplète.\end{pcEncadre}
\pcFigureEx{figures/ann12-molecules.pdf}{Document original · données pédagogiques. Schéma sans échelle, sauf axes gradués.}
\par\addvspace{3mm}\noindent\begin{minipage}{\linewidth}\textbf{Question 4 · 0,5 point}\quad Nommer le combustible et le comburant de cette réaction.\par\end{minipage}\par
\par\addvspace{3mm}\noindent\begin{minipage}{\linewidth}\textbf{Question 5 · 1 point}\quad Vérifier la conservation des atomes d’oxygène de part et d’autre de l’équation.\par\end{minipage}\par
\par\addvspace{3mm}\noindent\begin{minipage}{\linewidth}\textbf{Question 6 · 1,5 points}\quad Pendant un essai, le réservoir perd 4,4 g de propane. Un élève en déduit que la masse totale des produits formés vaut 4,4 g. Peut-on valider sa conclusion ? Identifier la masse oubliée et écrire le bilan de masse complet, sans calcul numérique.\par\end{minipage}\par
\vfill\pcRepere{cor-ann-12}{Corrigé expliqué, barème et bilan}
\pcPage{Sujet 12 · Énoncé}{Comprendre un détecteur optique de fumée}
\textbf{PC · 30 min · /10}\hfill Partie 3/3 · 3 points scientifiques
\begin{pcEncadre}{definition}{Document · Recevoir le signal sonore}Un essai en plein air utilise un signal bref. Un microphone est situé à 51 m de l’émetteur. On néglige le vent et les réflexions ; la vitesse du son est $340\unite{m/s}$. Deux enregistrements A et B portent sur une même durée de 4 ms.\end{pcEncadre}
\pcFigureEx{figures/ann12-son.pdf}{Document original · données pédagogiques. Schéma sans échelle, sauf axes gradués.}
\par\addvspace{3mm}\noindent\begin{minipage}{\linewidth}\textbf{Question 7 · 1 point}\quad Calculer la durée de propagation directe entre l’émetteur et le microphone.\par\end{minipage}\par
\par\addvspace{3mm}\noindent\begin{minipage}{\linewidth}\textbf{Question 8 · 1 point}\quad Quel enregistrement correspond au son le plus aigu ? Justifier sans confondre fréquence et amplitude.\par\end{minipage}\par
\par\addvspace{3mm}\noindent\begin{minipage}{\linewidth}\textbf{Question 9 · 1 point}\quad Le son pourrait-il se propager de la même manière dans le vide ? Justifier.\par\end{minipage}\par
\vfill\pcRepere{cor-ann-12}{Corrigé expliqué, barème et bilan}
\pcPage{Sujet 13 · Énoncé}{Une station météo sur la côte}
\phantomsection\label{ann-13}
\pdfbookmark[1]{13 · Une station météo sur la côte}{bm-ann-13}
\textbf{PC · 30 min · /10}\hfill Partie 1/3 · 3 points scientifiques
\begin{pcEncadre}{definition}{Document · Connaître l’air prélevé}Un flacon contient 2,0 L d’air sec. Le modèle utilise 78 \% de diazote, 21 \% de dioxygène et 1 \% d’autres gaz en volume. L’azote possède $Z=7$ ; l’isotope étudié possède $A=14$. La formule du diazote est $\ce{N2}$.\end{pcEncadre}
\pcFigureEx{figures/ann13-air.pdf}{Document original · données pédagogiques. Schéma sans échelle, sauf axes gradués.}
\par\addvspace{3mm}\noindent\begin{minipage}{\linewidth}\textbf{Question 1 · 0,5 point}\quad L’air sec est-il un corps pur ou un mélange ? Justifier.\par\end{minipage}\par
\par\addvspace{3mm}\noindent\begin{minipage}{\linewidth}\textbf{Question 2 · 1 point}\quad Calculer le volume de dioxygène du prélèvement, en litre puis en millilitres.\par\end{minipage}\par
\par\addvspace{3mm}\noindent\begin{minipage}{\linewidth}\textbf{Question 3 · 1,5 points}\quad Donner les nombres de protons, de neutrons et d’électrons d’un atome neutre d’azote de cet isotope. Combien d’atomes contient une molécule de diazote ?\par\end{minipage}\par
\vfill\pcRepere{cor-ann-13}{Corrigé expliqué, barème et bilan}
\pcPage{Sujet 13 · Énoncé}{Une station météo sur la côte}
\textbf{PC · 30 min · /10}\hfill Partie 2/3 · 3 points scientifiques
\begin{pcEncadre}{definition}{Document · Tester un composant résistif}Le graphique présente des mesures sur un conducteur ohmique. On utilise $U=RI$ avec U en V, I en A et R en ohms. Le montage de principe du corrigé sera réalisé uniquement avec le matériel scolaire adapté.\end{pcEncadre}
\pcFigureEx{figures/ann13-ohm.pdf}{Document original · données pédagogiques. Schéma sans échelle, sauf axes gradués.}
\par\addvspace{3mm}\noindent\begin{minipage}{\linewidth}\textbf{Question 4 · 0,5 point}\quad Lire la tension correspondant à une intensité de 30 mA.\par\end{minipage}\par
\par\addvspace{3mm}\noindent\begin{minipage}{\linewidth}\textbf{Question 5 · 1 point}\quad Calculer la résistance à partir de ce point.\par\end{minipage}\par
\par\addvspace{3mm}\noindent\begin{minipage}{\linewidth}\textbf{Question 6 · 1,5 points}\quad Schématiser un montage permettant ces mesures avec générateur, résistance, ampèremètre et voltmètre. Préciser leurs branchements.\par\end{minipage}\par
\vfill\pcRepere{cor-ann-13}{Corrigé expliqué, barème et bilan}
\pcPage{Sujet 13 · Énoncé}{Une station météo sur la côte}
\textbf{PC · 30 min · /10}\hfill Partie 3/3 · 3 points scientifiques
\begin{pcEncadre}{definition}{Document · Transmettre un message}La station envoie un signal radio à un relais distant de 120 km. La vitesse retenue est $3{,}0\times10^8\unite{m/s}$. Le message d’étude comporte quatre symboles avec le code A = 10, B = 01, C = 11, D = 00 ; il est composé de huit intervalles de 2 ms. La durée de préparation n’est pas incluse.\end{pcEncadre}
\pcFigureEx{figures/ann13-code.pdf}{Document original · données pédagogiques. Schéma sans échelle, sauf axes gradués.}
\par\addvspace{3mm}\noindent\begin{minipage}{\linewidth}\textbf{Question 7 · 0,5 point}\quad Décoder les quatre groupes représentés sur le document.\par\end{minipage}\par
\par\addvspace{3mm}\noindent\begin{minipage}{\linewidth}\textbf{Question 8 · 1 point}\quad Calculer la durée totale d’émission du message.\par\end{minipage}\par
\par\addvspace{3mm}\noindent\begin{minipage}{\linewidth}\textbf{Question 9 · 1,5 points}\quad Calculer la durée de propagation du début du message jusqu’au relais. Expliquer pourquoi elle diffère de la durée d’émission.\par\end{minipage}\par
\vfill\pcRepere{cor-ann-13}{Corrigé expliqué, barème et bilan}
\pcPage{Sujet 14 · Énoncé}{Transporter un colis réfrigéré}
\phantomsection\label{ann-14}
\pdfbookmark[1]{14 · Transporter un colis réfrigéré}{bm-ann-14}
\textbf{PC · 30 min · /10}\hfill Partie 1/3 · 3 points scientifiques
\begin{pcEncadre}{definition}{Document · Étudier une réserve froide}Un sachet étanche contient 200 g d’eau pure initialement solide. Sous la pression usuelle, le professeur suit sa température pendant un apport d’énergie. Le sachet reste fermé, sans fuite ; on mesure ici la température de l’eau dans ses différents états.\end{pcEncadre}
\pcFigureEx{figures/ann14-fusion.pdf}{Document original · données pédagogiques. Schéma sans échelle, sauf axes gradués.}
\par\addvspace{3mm}\noindent\begin{minipage}{\linewidth}\textbf{Question 1 · 1 point}\quad Nommer le changement d’état au palier de 0 °C et indiquer les états qui coexistent.\par\end{minipage}\par
\par\addvspace{3mm}\noindent\begin{minipage}{\linewidth}\textbf{Question 2 · 1 point}\quad Lire la durée du palier et prévoir la masse d’eau après fusion complète.\par\end{minipage}\par
\par\addvspace{3mm}\noindent\begin{minipage}{\linewidth}\textbf{Question 3 · 1 point}\quad Un élève affirme : « La température ne change plus, donc l’eau ne reçoit plus d’énergie. » Discuter cette affirmation.\par\end{minipage}\par
\vfill\pcRepere{cor-ann-14}{Corrigé expliqué, barème et bilan}
\pcPage{Sujet 14 · Énoncé}{Transporter un colis réfrigéré}
\textbf{PC · 30 min · /10}\hfill Partie 2/3 · 3 points scientifiques
\begin{pcEncadre}{definition}{Document · Analyser un liquide de conditionnement}Un autre prélèvement de liquide présente un pH de 8,0 vers 25 °C. La fiche de test indique : nitrate d’argent + ions chlorure → précipité blanc. Après ajout du réactif à un prélèvement séparé, un précipité blanc apparaît. Pour le chlore, $Z=17$. Les tests sont réalisés au laboratoire sous la conduite du professeur.\end{pcEncadre}
\pcFigureEx{figures/ann14-test.pdf}{Document original · données pédagogiques. Schéma sans échelle, sauf axes gradués.}
\par\addvspace{3mm}\noindent\begin{minipage}{\linewidth}\textbf{Question 4 · 0,5 point}\quad Classer le liquide selon son pH et justifier.\par\end{minipage}\par
\par\addvspace{3mm}\noindent\begin{minipage}{\linewidth}\textbf{Question 5 · 1 point}\quad Identifier l’ion mis en évidence et donner son nombre d’électrons.\par\end{minipage}\par
\par\addvspace{3mm}\noindent\begin{minipage}{\linewidth}\textbf{Question 6 · 1,5 points}\quad En trois à cinq lignes, expliquer ce que permettent les résultats et ce qu’ils ne prouvent pas sur la composition et l’innocuité du liquide.\par\end{minipage}\par
\vfill\pcRepere{cor-ann-14}{Corrigé expliqué, barème et bilan}
\pcPage{Sujet 14 · Énoncé}{Transporter un colis réfrigéré}
\textbf{PC · 30 min · /10}\hfill Partie 3/3 · 3 points scientifiques
\begin{pcEncadre}{definition}{Document · Limiter la vitesse sur un convoyeur}Le colis de masse 2,0 kg se déplace en ligne droite sur un convoyeur. Une phase à vitesse constante est comparée à deux réglages : 0,50 m/s et 1,0 m/s. On donne $E_c=\frac12mv^2$, avec les unités SI.\end{pcEncadre}
\pcFigureEx{figures/ann14-convoyeur.pdf}{Document original · données pédagogiques. Schéma sans échelle, sauf axes gradués.}
\par\addvspace{3mm}\noindent\begin{minipage}{\linewidth}\textbf{Question 7 · 1 point}\quad Quelle distance parcourt le colis en 8,0 s à 0,50 m/s constants ?\par\end{minipage}\par
\par\addvspace{3mm}\noindent\begin{minipage}{\linewidth}\textbf{Question 8 · 1 point}\quad Calculer l’énergie cinétique du colis pour chacun des deux réglages.\par\end{minipage}\par
\par\addvspace{3mm}\noindent\begin{minipage}{\linewidth}\textbf{Question 9 · 1 point}\quad Le doublement de vitesse double-t-il l’énergie cinétique ? Justifier, puis expliquer l’intérêt de limiter la vitesse.\par\end{minipage}\par
\vfill\pcRepere{cor-ann-14}{Corrigé expliqué, barème et bilan}
\pcPage{Sujet 15 · Énoncé}{Une maquette de dirigeable}
\phantomsection\label{ann-15}
\pdfbookmark[1]{15 · Une maquette de dirigeable}{bm-ann-15}
\textbf{PC · 30 min · /10}\hfill Partie 1/3 · 3 points scientifiques
\begin{pcEncadre}{definition}{Document · Décrire le gaz de l’enveloppe}La maquette contient de l’hélium modélisé par des atomes neutres séparés. L’isotope représenté possède $Z=2$ et $A=4$. Les points du modèle ne sont pas dessinés à la taille réelle des atomes.\end{pcEncadre}
\pcFigureEx{figures/ann15-helium.pdf}{Document original · données pédagogiques. Schéma sans échelle, sauf axes gradués.}
\par\addvspace{3mm}\noindent\begin{minipage}{\linewidth}\textbf{Question 1 · 1 point}\quad Donner les nombres de protons, de neutrons et d’électrons de cet atome.\par\end{minipage}\par
\par\addvspace{3mm}\noindent\begin{minipage}{\linewidth}\textbf{Question 2 · 1 point}\quad Le modèle représente-t-il un gaz constitué de molécules de plusieurs atomes ? Justifier.\par\end{minipage}\par
\par\addvspace{3mm}\noindent\begin{minipage}{\linewidth}\textbf{Question 3 · 1 point}\quad Après compression douce sans fuite, à température redevenue identique, le nombre d’atomes diminue-t-il ? Expliquer la modification du modèle.\par\end{minipage}\par
\vfill\pcRepere{cor-ann-15}{Corrigé expliqué, barème et bilan}
\pcPage{Sujet 15 · Énoncé}{Une maquette de dirigeable}
\textbf{PC · 30 min · /10}\hfill Partie 2/3 · 3 points scientifiques
\begin{pcEncadre}{definition}{Document · Équilibrer la maquette}La masse totale est de 0,30 kg. Le modèle retient le poids vertical vers le bas et une force de l’air verticale vers le haut de 3,4 N. Le dirigeable peut être retenu par un fil vertical ancré au sol. On donne $g=10\unite{N/kg}$.\end{pcEncadre}
\pcFigureEx{figures/ann15-forces-enonce.pdf}{Document original · données pédagogiques. Schéma sans échelle, sauf axes gradués.}
\par\addvspace{3mm}\noindent\begin{minipage}{\linewidth}\textbf{Question 4 · 0,5 point}\quad Calculer le poids total de la maquette.\par\end{minipage}\par
\par\addvspace{3mm}\noindent\begin{minipage}{\linewidth}\textbf{Question 5 · 1 point}\quad Sans fil, les deux forces indiquées se compensent-elles ? Justifier.\par\end{minipage}\par
\par\addvspace{3mm}\noindent\begin{minipage}{\linewidth}\textbf{Question 6 · 1,5 points}\quad Le fil permet de maintenir le dirigeable immobile. Déterminer la valeur et le sens de sa force sur la maquette, puis représenter les trois forces sans échelle.\par\end{minipage}\par
\vfill\pcRepere{cor-ann-15}{Corrigé expliqué, barème et bilan}
\pcPage{Sujet 15 · Énoncé}{Une maquette de dirigeable}
\textbf{PC · 30 min · /10}\hfill Partie 3/3 · 3 points scientifiques
\begin{pcEncadre}{definition}{Document · Alimenter les accessoires}Un générateur adapté fournit une tension constante de 6,0 V à deux branches en dérivation : moteur, 0,20 A ; éclairage, 0,10 A. Les courants restent constants dans le modèle. La batterie fournit 1,8 Wh utilisables. On donne $P=UI$ et $E=P\Delta t$.\end{pcEncadre}
\pcFigureEx{figures/ann15-circuit.pdf}{Document original · données pédagogiques. Schéma sans échelle, sauf axes gradués.}
\par\addvspace{3mm}\noindent\begin{minipage}{\linewidth}\textbf{Question 7 · 1 point}\quad Donner la tension aux bornes de chaque branche et calculer l’intensité fournie par le générateur.\par\end{minipage}\par
\par\addvspace{3mm}\noindent\begin{minipage}{\linewidth}\textbf{Question 8 · 1 point}\quad Calculer la puissance totale et la durée maximale d’alimentation dans ce modèle.\par\end{minipage}\par
\par\addvspace{3mm}\noindent\begin{minipage}{\linewidth}\textbf{Question 9 · 1 point}\quad La branche de l’éclairage devient ouverte. Le moteur est-il nécessairement arrêté ? Expliquer avec le schéma.\par\end{minipage}\par
\vfill\pcRepere{cor-ann-15}{Corrigé expliqué, barème et bilan}
\global\otCorrigestrue\global\otSoupletrue\gdef\otAccent{otGreenDark}\pcAnnexe{Corrigés des sujets 11 à 15}{Comprendre le modèle, les calculs et les critères de réussite}\phantomsection\label{cor-annales-originales}\pdfbookmark[0]{Corrigés des sujets 11 à 15}{bm-cor-annales-originales}
\begin{pcEncadre}{method}{Une grille pour progresser}Chaque sujet : partie 1 /3, partie 2 /3, partie 3 /3 et langue /1, soit /10. Relire d’abord la démarche avant d’attribuer les points. Repérer une erreur de modèle, une erreur de lecture et une erreur d’unité si elles existent, puis choisir un exercice du cours pour reprendre la difficulté.\end{pcEncadre}
\pcPage{Corrigé du sujet 11}{Suivre la cabine}
\phantomsection\label{cor-ann-11}
\pcRepere{ann-11}{Revenir à l’énoncé}
\par\noindent\begin{minipage}{\linewidth}\pcSousTitre{Question 1}Le mouvement est rectiligne, car la trajectoire est verticale et droite, et uniforme : l’altitude augmente de 0,60 m pendant chaque seconde.\par\smallskip\textbf{Démarche.}\enspace Séparer la forme de la trajectoire et l’évolution de la valeur de la vitesse. Les relevés sont faits à durées égales.\par\smallskip{\small\textbf{Barème.}\enspace 0,25 pour rectiligne justifié ; 0,25 pour uniforme justifié.\par}\end{minipage}\par
\par\noindent\begin{minipage}{\linewidth}\pcSousTitre{Question 2}$v=d/\Delta t=1{,}80/3{,}0=0{,}60\unite{m/s}$. La distance parcourue est 1,80 m pendant cette montée.\par\smallskip\textbf{Démarche.}\enspace Utiliser la différence entre les altitudes initiale et finale, puis la durée du même intervalle.\par\smallskip{\small\textbf{Barème.}\enspace 0,5 pour quotient et substitution ; 0,5 pour valeur et unité.\par}\end{minipage}\par
\par\noindent\begin{minipage}{\linewidth}\pcSousTitre{Question 3}$P=mg=50\times9{,}8=490\unite N$. La force du câble vaut aussi 490 N : verticale vers le haut, opposée au poids vertical vers le bas. Elles agissent ici sur le même système.\par\smallskip\textbf{Démarche.}\enspace La compensation impose deux forces de même direction, de même valeur et de sens opposés. Elle est compatible avec une vitesse constante non nulle.\par\smallskip{\small\textbf{Barème.}\enspace 0,5 pour le poids ; 0,5 pour directions et sens ; 0,5 pour égalité des valeurs et représentation.\par}\end{minipage}\par
\pcFigureEx{figures/ann11-forces.pdf}{Modèle de correction · les légendes précisent les directions, les connexions ou les résultats.}
\pcPage{Corrigé du sujet 11}{Prévoir l’autonomie}
\pcRepere{ann-11}{Revenir à l’énoncé}
\par\noindent\begin{minipage}{\linewidth}\pcSousTitre{Question 4}$P=24\times3{,}0=72\unite W$.\par\smallskip\textbf{Démarche.}\enspace Multiplier tension en volts et intensité en ampères. Le watt mesure une puissance.\par\smallskip{\small\textbf{Barème.}\enspace 0,5 pour relation et substitution ; 0,5 pour résultat avec unité.\par}\end{minipage}\par
\par\noindent\begin{minipage}{\linewidth}\pcSousTitre{Question 5}$6\unite{min}=6/60\unite h=0{,}10\unite h$ ; $E=72\times0{,}10=7{,}2\unite{Wh}$.\par\smallskip\textbf{Démarche.}\enspace Convertir les minutes en heures avant de multiplier une puissance en watts pour obtenir des Wh.\par\smallskip{\small\textbf{Barème.}\enspace 0,5 pour conversion ; 0,5 pour produit et unité.\par}\end{minipage}\par
\par\noindent\begin{minipage}{\linewidth}\pcSousTitre{Question 6}Trois cycles demandent $3\times7{,}2=21{,}6\unite{Wh}$, au-delà des 18 Wh disponibles. Deux cycles demandent 14,4 Wh : le maximum est donc deux cycles complets.\par\smallskip\textbf{Démarche.}\enspace Comparer des énergies dans la même unité, puis retenir un nombre entier de cycles. Une fraction de cycle ne constitue pas un cycle complet.\par\smallskip{\small\textbf{Barème.}\enspace 0,5 pour comparaison chiffrée ; 0,5 pour deux cycles justifiés.\par}\end{minipage}\par
\pcFigureEx{figures/ann11-autonomie.pdf}{Modèle de correction · les légendes précisent les directions, les connexions ou les résultats.}
\pcPage{Corrigé du sujet 11}{Identifier un liquide de maintenance}
\pcRepere{ann-11}{Revenir à l’énoncé}
\par\noindent\begin{minipage}{\linewidth}\pcSousTitre{Question 7}$m=102{,}5-28{,}5=74{,}0\unite g$.\par\smallskip\textbf{Démarche.}\enspace La masse totale comprend le récipient et son contenu : retrancher la masse du récipient.\par\smallskip{\small\textbf{Barème.}\enspace 0,25 pour soustraction ; 0,25 pour valeur et unité.\par}\end{minipage}\par
\par\noindent\begin{minipage}{\linewidth}\pcSousTitre{Question 8}$\rho=74{,}0/100{,}0=0{,}740\unite{g/mL}$. La valeur est compatible avec A, aux arrondis des références près.\par\smallskip\textbf{Démarche.}\enspace Conserver les unités g et mL pour comparer directement aux références du document.\par\smallskip{\small\textbf{Barème.}\enspace 0,5 pour quotient et unité ; 0,5 pour comparaison à A.\par}\end{minipage}\par
\par\noindent\begin{minipage}{\linewidth}\pcSousTitre{Question 9}Une même masse volumique peut correspondre à plusieurs substances ou à un mélange. La mesure possède aussi une incertitude. Il faut comparer une autre propriété indépendante, par exemple une température de changement d’état dans un protocole adapté, ou réaliser une analyse chimique appropriée au laboratoire.\par\smallskip\textbf{Démarche.}\enspace Distinguer « compatible avec » et « prouvé ». Le contrôle doit porter sur une propriété différente et rester adapté au produit étudié.\par\smallskip{\small\textbf{Barème.}\enspace 0,5 pour possibilité de substances ou mélanges différents ; 0,5 pour limites de mesure ; 0,5 pour contrôle indépendant pertinent.\par}\end{minipage}\par
\par\noindent\begin{minipage}{\linewidth}{\small\textbf{Bilan :} /3 + /3 + /3 + /1 = /10.\par\textbf{Pour reprendre :} \hyperref[chap-1]{chapitre 1 (p.~\pageref*{chap-1})} ; \hyperref[chap-8]{chapitre 8 (p.~\pageref*{chap-8})} ; \hyperref[chap-9]{chapitre 9 (p.~\pageref*{chap-9})} ; \hyperref[chap-14]{chapitre 14 (p.~\pageref*{chap-14})} ; \hyperref[chap-15]{chapitre 15 (p.~\pageref*{chap-15})}.\par}\end{minipage}\par
\pcPage{Corrigé du sujet 12}{Un trajet lumineux modifié}
\phantomsection\label{cor-ann-12}
\pcRepere{ann-12}{Revenir à l’énoncé}
\par\noindent\begin{minipage}{\linewidth}\pcSousTitre{Question 1}La DEL est la source primaire : elle produit la lumière émise. Le capteur est le récepteur de lumière.\par\smallskip\textbf{Démarche.}\enspace Attribuer les rôles en suivant la lumière, de sa production jusqu’à sa détection.\par\smallskip{\small\textbf{Barème.}\enspace 0,25 pour la DEL ; 0,25 pour le capteur.\par}\end{minipage}\par
\par\noindent\begin{minipage}{\linewidth}\pcSousTitre{Question 2}La première flèche va de la DEL aux particules ; la seconde va des particules vers le capteur. Les particules diffusent une partie de la lumière reçue.\par\smallskip\textbf{Démarche.}\enspace Les flèches indiquent le sens de propagation de la lumière. Le capteur n’envoie pas de rayon pour chercher la fumée.\par\smallskip{\small\textbf{Barème.}\enspace 0,5 par trajet fléché correctement orienté.\par}\end{minipage}\par
\par\noindent\begin{minipage}{\linewidth}\pcSousTitre{Question 3}Hors du faisceau direct, le capteur reçoit peu de lumière en l’absence de particules. Une augmentation de lumière reçue peut donc indiquer une diffusion par des particules. La poussière peut également diffuser de la lumière : cette seule mesure ne donne pas l’identité chimique des particules.\par\smallskip\textbf{Démarche.}\enspace Relier le placement du capteur au signal de référence, puis séparer détection d’une présence et identification de sa nature.\par\smallskip{\small\textbf{Barème.}\enspace 0,5 pour faible signal sans particules ; 0,5 pour diffusion ; 0,5 pour limite liée à la poussière.\par}\end{minipage}\par
\pcFigureEx{figures/ann12-optique-cor.pdf}{Modèle de correction · les légendes précisent les directions, les connexions ou les résultats.}
\pcPage{Corrigé du sujet 12}{Raisonner sur une combustion}
\pcRepere{ann-12}{Revenir à l’énoncé}
\par\noindent\begin{minipage}{\linewidth}\pcSousTitre{Question 4}Le propane est le combustible ; le dioxygène est le comburant.\par\smallskip\textbf{Démarche.}\enspace Repérer la substance qui brûle et le réactif qui permet la combustion.\par\smallskip{\small\textbf{Barème.}\enspace 0,25 pour chacun des deux noms.\par}\end{minipage}\par
\par\noindent\begin{minipage}{\linewidth}\pcSousTitre{Question 5}À gauche, 5 molécules de dioxygène apportent 10 atomes O. À droite, 3 molécules de dioxyde de carbone apportent 6 atomes O et 4 molécules d’eau en apportent 4 : total 10.\par\smallskip\textbf{Démarche.}\enspace Multiplier le nombre devant la formule par le nombre d’atomes O dans une molécule, puis additionner les produits.\par\smallskip{\small\textbf{Barème.}\enspace 0,5 pour le décompte à gauche ; 0,5 pour celui à droite.\par}\end{minipage}\par
\par\noindent\begin{minipage}{\linewidth}\pcSousTitre{Question 6}Non : la réaction consomme aussi du dioxygène prélevé dans l’air. Sa masse doit être ajoutée à celle du propane consommé. Pour la transformation complète du document : $m_{\rm propane}+m_{\rm dioxygène}=m_{\rm dioxyde\ de\ carbone}+m_{\rm eau}$. La seule perte du réservoir ne donne pas la masse totale des produits.\par\smallskip\textbf{Démarche.}\enspace Définir la frontière du système pesé : le réservoir ne contient pas tout le dioxygène consommé. La conservation porte sur tous les réactifs et tous les produits, pas sur la seule masse de combustible.\par\smallskip{\small\textbf{Barème.}\enspace 0,5 pour rejet justifié ; 0,5 pour dioxygène oublié ; 0,5 pour bilan complet.\par}\end{minipage}\par
\pcFigureEx{figures/ann12-bilan.pdf}{Modèle de correction · les légendes précisent les directions, les connexions ou les résultats.}
\pcPage{Corrigé du sujet 12}{Recevoir le signal sonore}
\pcRepere{ann-12}{Revenir à l’énoncé}
\par\noindent\begin{minipage}{\linewidth}\pcSousTitre{Question 7}$\Delta t=d/v=51/340=0{,}15\unite s$.\par\smallskip\textbf{Démarche.}\enspace Il s’agit d’un trajet simple : aucune division par deux n’intervient.\par\smallskip{\small\textbf{Barème.}\enspace 0,5 pour quotient ; 0,5 pour durée et unité.\par}\end{minipage}\par
\par\noindent\begin{minipage}{\linewidth}\pcSousTitre{Question 8}B correspond au son le plus aigu : il présente quatre oscillations pendant les 4 ms, contre deux pour A. Les amplitudes dessinées sont égales.\par\smallskip\textbf{Démarche.}\enspace Comparer le nombre d’oscillations sur une même durée. Une plus grande amplitude ne signifie pas un son plus aigu.\par\smallskip{\small\textbf{Barème.}\enspace 0,5 pour B ; 0,5 pour comparaison à durée égale.\par}\end{minipage}\par
\par\noindent\begin{minipage}{\linewidth}\pcSousTitre{Question 9}Non. Le son est une perturbation mécanique qui nécessite un milieu matériel pour se propager. Le vide ne fournit pas ce milieu.\par\smallskip\textbf{Démarche.}\enspace Distinguer la propagation du son de celle de la lumière, qui peut se propager dans le vide.\par\smallskip{\small\textbf{Barème.}\enspace 0,5 pour réponse ; 0,5 pour milieu matériel nécessaire.\par}\end{minipage}\par
\pcFigureEx{figures/ann12-son-cor.pdf}{Modèle de correction · les légendes précisent les directions, les connexions ou les résultats.}
\par\noindent\begin{minipage}{\linewidth}{\small\textbf{Bilan :} /3 + /3 + /3 + /1 = /10.\par\textbf{Pour reprendre :} \hyperref[chap-3]{chapitre 3 (p.~\pageref*{chap-3})} ; \hyperref[chap-6]{chapitre 6 (p.~\pageref*{chap-6})} ; \hyperref[chap-16]{chapitre 16 (p.~\pageref*{chap-16})} ; \hyperref[chap-17]{chapitre 17 (p.~\pageref*{chap-17})}.\par}\end{minipage}\par
\pcPage{Corrigé du sujet 13}{Connaître l’air prélevé}
\phantomsection\label{cor-ann-13}
\pcRepere{ann-13}{Revenir à l’énoncé}
\par\noindent\begin{minipage}{\linewidth}\pcSousTitre{Question 1}C’est un mélange : il contient plusieurs espèces chimiques, dont le diazote et le dioxygène.\par\smallskip\textbf{Démarche.}\enspace La transparence et l’aspect uniforme ne prouvent pas qu’un gaz est un corps pur.\par\smallskip{\small\textbf{Barème.}\enspace 0,25 pour mélange ; 0,25 pour pluralité des espèces.\par}\end{minipage}\par
\par\noindent\begin{minipage}{\linewidth}\pcSousTitre{Question 2}$V=0{,}21\times2{,}0=0{,}42\unite L=420\unite{mL}$.\par\smallskip\textbf{Démarche.}\enspace Appliquer la proportion au volume total, puis convertir les litres en millilitres.\par\smallskip{\small\textbf{Barème.}\enspace 0,5 pour produit ; 0,5 pour valeurs et conversion.\par}\end{minipage}\par
\par\noindent\begin{minipage}{\linewidth}\pcSousTitre{Question 3}L’atome possède 7 protons, 14 − 7 = 7 neutrons et 7 électrons car il est neutre. Une molécule de diazote contient deux atomes d’azote.\par\smallskip\textbf{Démarche.}\enspace Z compte les protons ; A compte protons et neutrons. La neutralité impose autant d’électrons que de protons. L’indice 2 porte sur les atomes de la molécule.\par\smallskip{\small\textbf{Barème.}\enspace 0,25 pour chacun des trois inventaires ; 0,25 pour neutralité justifiée ; 0,5 pour deux atomes.\par}\end{minipage}\par
\pcFigureEx{figures/ann13-atome.pdf}{Modèle de correction · les légendes précisent les directions, les connexions ou les résultats.}
\pcPage{Corrigé du sujet 13}{Tester un composant résistif}
\pcRepere{ann-13}{Revenir à l’énoncé}
\par\noindent\begin{minipage}{\linewidth}\pcSousTitre{Question 4}$U=6{,}0\unite V$ pour $I=30\unite{mA}$.\par\smallskip\textbf{Démarche.}\enspace Partir de 30 mA sur l’axe horizontal, rejoindre la droite puis lire la tension sur l’axe vertical.\par\smallskip{\small\textbf{Barème.}\enspace 0,5 pour lecture et unité.\par}\end{minipage}\par
\par\noindent\begin{minipage}{\linewidth}\pcSousTitre{Question 5}$30\unite{mA}=0{,}030\unite A$ ; $R=6{,}0/0{,}030=200\unite\Omega$.\par\smallskip\textbf{Démarche.}\enspace Convertir les milliampères en ampères avant d’appliquer la loi d’Ohm.\par\smallskip{\small\textbf{Barème.}\enspace 0,5 pour conversion et quotient ; 0,5 pour résultat et unité.\par}\end{minipage}\par
\par\noindent\begin{minipage}{\linewidth}\pcSousTitre{Question 6}L’ampèremètre se place en série dans la branche de la résistance ; le voltmètre se branche en dérivation entre les deux bornes de la résistance. Les bornes A et V sont du côté des potentiels les plus élevés, les bornes COM de l’autre côté, pour des lectures positives.\par\smallskip\textbf{Démarche.}\enspace Mesurer I impose de traverser l’ampèremètre. Mesurer U impose de comparer deux points du circuit : le voltmètre relie ces deux points.\par\smallskip{\small\textbf{Barème.}\enspace 0,5 pour boucle et résistance ; 0,5 pour ampèremètre en série ; 0,5 pour voltmètre en dérivation.\par}\end{minipage}\par
\pcFigureEx{figures/ann13-mesures.pdf}{Modèle de correction · les légendes précisent les directions, les connexions ou les résultats.}
\pcPage{Corrigé du sujet 13}{Transmettre un message}
\pcRepere{ann-13}{Revenir à l’énoncé}
\par\noindent\begin{minipage}{\linewidth}\pcSousTitre{Question 7}Le message est CADB : 11 → C, 10 → A, 00 → D, 01 → B.\par\smallskip\textbf{Démarche.}\enspace Lire chaque groupe depuis le début annoncé ; ne pas décaler le découpage en groupes de deux.\par\smallskip{\small\textbf{Barème.}\enspace 0,5 pour les quatre symboles dans l’ordre.\par}\end{minipage}\par
\par\noindent\begin{minipage}{\linewidth}\pcSousTitre{Question 8}$8\times2=16\unite{ms}=0{,}016\unite s$.\par\smallskip\textbf{Démarche.}\enspace Compter les huit intervalles, et non seulement les intervalles codés 1. Le zéro occupe lui aussi du temps.\par\smallskip{\small\textbf{Barème.}\enspace 0,5 pour huit intervalles ; 0,5 pour produit et unité.\par}\end{minipage}\par
\par\noindent\begin{minipage}{\linewidth}\pcSousTitre{Question 9}$120\unite{km}=120\,000\unite m$ ; $\Delta t=120\,000/(3{,}0\times10^8)=4{,}0\times10^{-4}\unite s=0{,}40\unite{ms}$. Cette durée est le trajet d’un front du signal, tandis que les 16 ms correspondent à l’émission des huit valeurs successives.\par\smallskip\textbf{Démarche.}\enspace Faire le calcul en mètres et secondes. Un message long peut commencer à arriver avant que son émission soit terminée.\par\smallskip{\small\textbf{Barème.}\enspace 0,5 pour conversion ; 0,5 pour quotient, valeur et unité ; 0,5 pour distinction des durées.\par}\end{minipage}\par
\pcFigureEx{figures/ann13-radio.pdf}{Modèle de correction · les légendes précisent les directions, les connexions ou les résultats.}
\par\noindent\begin{minipage}{\linewidth}{\small\textbf{Bilan :} /3 + /3 + /3 + /1 = /10.\par\textbf{Pour reprendre :} \hyperref[chap-2]{chapitre 2 (p.~\pageref*{chap-2})} ; \hyperref[chap-6]{chapitre 6 (p.~\pageref*{chap-6})} ; \hyperref[chap-13]{chapitre 13 (p.~\pageref*{chap-13})} ; \hyperref[chap-14]{chapitre 14 (p.~\pageref*{chap-14})} ; \hyperref[chap-18]{chapitre 18 (p.~\pageref*{chap-18})}.\par}\end{minipage}\par
\pcPage{Corrigé du sujet 14}{Étudier une réserve froide}
\phantomsection\label{cor-ann-14}
\pcRepere{ann-14}{Revenir à l’énoncé}
\par\noindent\begin{minipage}{\linewidth}\pcSousTitre{Question 1}Il s’agit de la fusion. Pendant ce palier, l’eau solide et l’eau liquide coexistent dans le modèle.\par\smallskip\textbf{Démarche.}\enspace Le sens d’évolution compte : la température venait d’une valeur négative pendant un apport d’énergie.\par\smallskip{\small\textbf{Barème.}\enspace 0,5 pour fusion ; 0,5 pour coexistence solide-liquide.\par}\end{minipage}\par
\par\noindent\begin{minipage}{\linewidth}\pcSousTitre{Question 2}Le palier dure 6 − 2 = 4 min. La masse d’eau reste 200 g : la même eau est conservée sans fuite.\par\smallskip\textbf{Démarche.}\enspace Une durée est une différence de dates. La conservation de la masse concerne le système fermé, même quand son état change.\par\smallskip{\small\textbf{Barème.}\enspace 0,5 pour durée ; 0,5 pour masse justifiée.\par}\end{minipage}\par
\par\noindent\begin{minipage}{\linewidth}\pcSousTitre{Question 3}L’affirmation est fausse dans cette expérience. L’apport d’énergie se poursuit et permet le changement d’état ; pendant la fusion d’un corps pur à pression fixée, la température peut rester constante.\par\smallskip\textbf{Démarche.}\enspace Une température constante ne signifie pas l’absence de transfert d’énergie. Relier l’apport au changement d’état observé.\par\smallskip{\small\textbf{Barème.}\enspace 0,5 pour affirmation rejetée ; 0,5 pour énergie du changement d’état.\par}\end{minipage}\par
\pcFigureEx{figures/ann14-fusion-cor.pdf}{Modèle de correction · les légendes précisent les directions, les connexions ou les résultats.}
\pcPage{Corrigé du sujet 14}{Analyser un liquide de conditionnement}
\pcRepere{ann-14}{Revenir à l’énoncé}
\par\noindent\begin{minipage}{\linewidth}\pcSousTitre{Question 4}Il est basique : le pH 8,0 est supérieur à 7 vers 25 °C.\par\smallskip\textbf{Démarche.}\enspace Comparer le pH à 7 à la température précisée. Une valeur voisine de 7 n’est pas forcément exactement neutre.\par\smallskip{\small\textbf{Barème.}\enspace 0,25 pour basique ; 0,25 pour comparaison à 7.\par}\end{minipage}\par
\par\noindent\begin{minipage}{\linewidth}\pcSousTitre{Question 5}Le test met en évidence l’ion chlorure $\ce{Cl-}$. Avec 17 protons et une charge $-1$, il possède 18 électrons.\par\smallskip\textbf{Démarche.}\enspace Un ion négatif possède un excès d’électrons. La charge ne change pas le numéro atomique de l’élément.\par\smallskip{\small\textbf{Barème.}\enspace 0,5 pour identification ; 0,5 pour 18 électrons justifiés.\par}\end{minipage}\par
\par\noindent\begin{minipage}{\linewidth}\pcSousTitre{Question 6}Les résultats indiquent un liquide basique dans ces conditions et la présence d’ions chlorure dans la limite du test. Ils ne donnent pas tous les constituants ni leur quantité. Ils ne suffisent pas à garantir l’innocuité : il faut d’autres analyses et les informations adaptées sur le produit et son usage.\par\smallskip\textbf{Démarche.}\enspace Séparer trois conclusions : acidité/basicité, présence d’une espèce et évaluation d’un risque. Un test positif ne démontre pas la pureté ni l’absence d’autres substances.\par\smallskip{\small\textbf{Barème.}\enspace 0,5 pour les deux résultats exacts ; 0,5 pour composition non exhaustive ; 0,5 pour innocuité non démontrée.\par}\end{minipage}\par
\pcFigureEx{figures/ann14-ion.pdf}{Modèle de correction · les légendes précisent les directions, les connexions ou les résultats.}
\pcPage{Corrigé du sujet 14}{Limiter la vitesse sur un convoyeur}
\pcRepere{ann-14}{Revenir à l’énoncé}
\par\noindent\begin{minipage}{\linewidth}\pcSousTitre{Question 7}$d=v\Delta t=0{,}50\times8{,}0=4{,}0\unite m$.\par\smallskip\textbf{Démarche.}\enspace Multiplier la vitesse par la durée seulement pour la portion où la vitesse garde la valeur annoncée.\par\smallskip{\small\textbf{Barème.}\enspace 0,5 pour relation et produit ; 0,5 pour distance avec unité.\par}\end{minipage}\par
\par\noindent\begin{minipage}{\linewidth}\pcSousTitre{Question 8}À 0,50 m/s : $E_c=0{,}5\times2{,}0\times0{,}50^2=0{,}25\unite J$. À 1,0 m/s : $E_c=1{,}0\unite J$.\par\smallskip\textbf{Démarche.}\enspace Mettre toute la valeur de la vitesse au carré. Une masse en kilogrammes et une vitesse en m/s donnent une énergie en joules.\par\smallskip{\small\textbf{Barème.}\enspace 0,5 pour chaque calcul avec unité.\par}\end{minipage}\par
\par\noindent\begin{minipage}{\linewidth}\pcSousTitre{Question 9}Non : à masse constante, doubler la vitesse multiplie l’énergie cinétique par quatre. À vitesse plus faible, il y a moins d’énergie cinétique à transférer lors d’un arrêt ou d’un choc, toutes choses égales par ailleurs.\par\smallskip\textbf{Démarche.}\enspace Comparer les deux énergies calculées et relier le facteur quatre au carré de la vitesse.\par\smallskip{\small\textbf{Barème.}\enspace 0,5 pour facteur quatre ; 0,5 pour intérêt formulé en termes d’énergie.\par}\end{minipage}\par
\pcFigureEx{figures/ann14-energies.pdf}{Modèle de correction · les légendes précisent les directions, les connexions ou les résultats.}
\par\noindent\begin{minipage}{\linewidth}{\small\textbf{Bilan :} /3 + /3 + /3 + /1 = /10.\par\textbf{Pour reprendre :} \hyperref[chap-1]{chapitre 1 (p.~\pageref*{chap-1})} ; \hyperref[chap-2]{chapitre 2 (p.~\pageref*{chap-2})} ; \hyperref[chap-4]{chapitre 4 (p.~\pageref*{chap-4})} ; \hyperref[chap-8]{chapitre 8 (p.~\pageref*{chap-8})} ; \hyperref[chap-11]{chapitre 11 (p.~\pageref*{chap-11})}.\par}\end{minipage}\par
\pcPage{Corrigé du sujet 15}{Décrire le gaz de l’enveloppe}
\phantomsection\label{cor-ann-15}
\pcRepere{ann-15}{Revenir à l’énoncé}
\par\noindent\begin{minipage}{\linewidth}\pcSousTitre{Question 1}Il possède 2 protons, 4 − 2 = 2 neutrons et 2 électrons puisqu’il est neutre.\par\smallskip\textbf{Démarche.}\enspace Utiliser Z, puis A − Z, puis la neutralité électrique ; A ne compte pas les électrons.\par\smallskip{\small\textbf{Barème.}\enspace 0,25 pour protons ; 0,25 pour neutrons ; 0,5 pour électrons et neutralité.\par}\end{minipage}\par
\par\noindent\begin{minipage}{\linewidth}\pcSousTitre{Question 2}Non : les entités décrites sont des atomes d’hélium isolés. Il s’agit ici d’un gaz monoatomique.\par\smallskip\textbf{Démarche.}\enspace Une entité n’est pas nécessairement une molécule : tenir compte de la nature du gaz indiquée dans le document.\par\smallskip{\small\textbf{Barème.}\enspace 0,5 pour non ; 0,5 pour atomes isolés.\par}\end{minipage}\par
\par\noindent\begin{minipage}{\linewidth}\pcSousTitre{Question 3}Le nombre d’atomes ne diminue pas : aucun gaz n’a quitté le système. On représente les mêmes atomes, de même taille conventionnelle, dans un volume plus petit ; leurs espacements moyens diminuent.\par\smallskip\textbf{Démarche.}\enspace Distinguer la taille des entités, leur nombre et l’espace disponible. Une compression ne détruit pas des atomes.\par\smallskip{\small\textbf{Barème.}\enspace 0,5 pour conservation du nombre ; 0,5 pour réduction des espacements.\par}\end{minipage}\par
\pcFigureEx{figures/ann15-compression.pdf}{Modèle de correction · les légendes précisent les directions, les connexions ou les résultats.}
\pcPage{Corrigé du sujet 15}{Équilibrer la maquette}
\pcRepere{ann-15}{Revenir à l’énoncé}
\par\noindent\begin{minipage}{\linewidth}\pcSousTitre{Question 4}$P=mg=0{,}30\times10=3{,}0\unite N$.\par\smallskip\textbf{Démarche.}\enspace La masse est déjà en kilogrammes ; le produit par g donne un poids en newtons.\par\smallskip{\small\textbf{Barème.}\enspace 0,25 pour produit ; 0,25 pour résultat avec unité.\par}\end{minipage}\par
\par\noindent\begin{minipage}{\linewidth}\pcSousTitre{Question 5}Non : leurs valeurs sont différentes, 3,4 N vers le haut et 3,0 N vers le bas. La résultante du modèle vaut 0,4 N vers le haut.\par\smallskip\textbf{Démarche.}\enspace Deux forces opposées ne se compensent que si leurs valeurs sont égales.\par\smallskip{\small\textbf{Barème.}\enspace 0,5 pour comparaison des valeurs ; 0,5 pour résultante vers le haut.\par}\end{minipage}\par
\par\noindent\begin{minipage}{\linewidth}\pcSousTitre{Question 6}Le fil exerce une force de 0,4 N vers le bas. Avec le poids, la somme des forces vers le bas vaut 3,0 + 0,4 = 3,4 N, égale à la force de l’air vers le haut.\par\smallskip\textbf{Démarche.}\enspace Isoler le dirigeable comme système. Le fil ancré au sol tire la maquette vers le bas ; il ne la pousse pas vers le haut.\par\smallskip{\small\textbf{Barème.}\enspace 0,5 pour sens ; 0,5 pour 0,4 N justifié ; 0,5 pour schéma complet.\par}\end{minipage}\par
\pcFigureEx{figures/ann15-forces-cor.pdf}{Modèle de correction · les légendes précisent les directions, les connexions ou les résultats.}
\pcPage{Corrigé du sujet 15}{Alimenter les accessoires}
\pcRepere{ann-15}{Revenir à l’énoncé}
\par\noindent\begin{minipage}{\linewidth}\pcSousTitre{Question 7}Chaque branche reçoit 6,0 V. L’intensité totale est $I=0{,}20+0{,}10=0{,}30\unite A$.\par\smallskip\textbf{Démarche.}\enspace En dérivation, la tension est commune aux branches et le courant principal est la somme des courants de branche.\par\smallskip{\small\textbf{Barème.}\enspace 0,5 pour tension commune ; 0,5 pour intensité totale.\par}\end{minipage}\par
\par\noindent\begin{minipage}{\linewidth}\pcSousTitre{Question 8}$P=6{,}0\times0{,}30=1{,}8\unite W$ ; $\Delta t=E/P=1{,}8/1{,}8=1{,}0\unite h$.\par\smallskip\textbf{Démarche.}\enspace Une énergie en Wh divisée par une puissance en W donne une durée en heures. La capacité fournie est déjà l’énergie utilisable.\par\smallskip{\small\textbf{Barème.}\enspace 0,5 pour puissance ; 0,5 pour durée avec unité.\par}\end{minipage}\par
\par\noindent\begin{minipage}{\linewidth}\pcSousTitre{Question 9}Non, dans le modèle à générateur adapté et tension maintenue. La branche du moteur conserve une boucle fermée entre les bornes du générateur. Seule la branche d’éclairage est interrompue.\par\smallskip\textbf{Démarche.}\enspace Suivre un trajet complet incluant le générateur pour chaque branche. Une ouverture dans une branche ne coupe pas automatiquement les autres.\par\smallskip{\small\textbf{Barème.}\enspace 0,5 pour réponse ; 0,5 pour boucle du moteur conservée.\par}\end{minipage}\par
\pcFigureEx{figures/ann15-circuit-cor.pdf}{Modèle de correction · les légendes précisent les directions, les connexions ou les résultats.}
\par\noindent\begin{minipage}{\linewidth}{\small\textbf{Bilan :} /3 + /3 + /3 + /1 = /10.\par\textbf{Pour reprendre :} \hyperref[chap-2]{chapitre 2 (p.~\pageref*{chap-2})} ; \hyperref[chap-9]{chapitre 9 (p.~\pageref*{chap-9})} ; \hyperref[chap-13]{chapitre 13 (p.~\pageref*{chap-13})} ; \hyperref[chap-14]{chapitre 14 (p.~\pageref*{chap-14})} ; \hyperref[chap-15]{chapitre 15 (p.~\pageref*{chap-15})}.\par}\end{minipage}\par
\pcPage{Sources des sujets 11 à 15}{Des attendus étudiés, des documents et des problèmes nouveaux}
\phantomsection\label{sources-annales}
\begin{pcEncadre}{method}{Ce qui a été étudié}Les annales officielles servent à analyser la place des documents, les changements d’unité, le raisonnement progressif, les constructions et les réponses argumentées. Le volet PC du sujet 2025 consulté est noté sur 25, celui de 2026 sur 10. Pour ces entraînements 2027, les critères scientifiques et le point de langue suivent le cadre actuel et restent séparés.\end{pcEncadre}
\pcSousTitre{Sources primaires consultées le 24 octobre 2026}
\begin{questions}
\item Ministère de l’Éducation nationale, DNB 2025, sciences, série générale, sujet 25GENSCAN1 : analyse de la progression des questions, des documents et de la tâche de synthèse.
\item Ministère de l’Éducation nationale, DNB 2026, sciences, série générale, sujet 26GENSCAN1 : analyse de l’usage de mesures et de documents indépendants.
\item Éduscol, sujet de référence de physique-chimie pour le DNB 2027 : analyse des attendus de rédaction et des critères de réponse.
\item Bulletin officiel, note du 11 septembre 2026, NOR MENE2623228N, partie I.4 : durée et modalités de sciences applicables à compter de la session 2027.
\end{questions}
\pcSousTitre{Consulter les documents officiels}
{\small\url{https://eduscol.education.gouv.fr/5607/les-epreuves-du-dnb}\par
\url{https://www.education.gouv.fr/bo/2026/Hebdo4/MENE2623228N}\par
\url{https://www.education.gouv.fr/brevet-bac-et-cap-les-sujets-des-examens-2026-504911}\par}
\pcSousTitre{Des choix éditoriaux explicites}
Le monte-charge, le détecteur optique, la station météo, le colis réfrigéré et le dirigeable constituent cinq dossiers originaux. Les séries numériques et les schémas ont été construits pour ces exercices. Aucun sujet n’annonce les thèmes qui seront retenus à l’examen. La concordance générale du programme figure dans les annexes du manuel.
\pcAnnexe{Méthodes et références utiles}{Des fiches à consulter puis à utiliser sans aide}\phantomsection\label{methodes}\pdfbookmark[0]{Méthodes et outils}{bm-methodes}
\begin{pcEncadre}{method}{Calculer en cinq étapes}1. Nommer la grandeur cherchée. 2. Écrire une relation adaptée. 3. Convertir les données. 4. Remplacer les valeurs et calculer. 5. Donner l’unité, un arrondi utile et une conclusion dans le contexte.\end{pcEncadre}
\pcSousTitre{Exemple : une énergie de fonctionnement}
Un appareil de 1,2 kW fonctionne 15 min à puissance constante. En kWh, écrire $\Delta t=15/60=0{,}25\unite h$, puis $E=1{,}2\times0{,}25=0{,}30\unite{kWh}$. En joules, utiliser 1 200 W et 900 s : $E=1\,080\,000\unite J$. Les deux résultats sont cohérents.
\pcSousTitre{Isoler une grandeur}
De $P=mg$ : $m=P/g$ et $g=P/m$, si les dénominateurs sont non nuls. De $E=P\Delta t$ : $P=E/\Delta t$ et $\Delta t=E/P$. Vérifier le sens : à même énergie, une puissance deux fois plus faible demande deux fois plus de temps.
\pcSousTitre{Comparer avant de calculer}
Si $E_c=\tfrac12mv^2$, doubler m double Ec ; doubler v quadruple Ec. Si la gravitation est fournie sous $F=Gm_Am_B/d^2$, doubler d divise F par quatre. Énoncer ce qui reste fixé.
\pcSousTitre{Contrôler un ordre de grandeur}
Une vitesse de vélo de plusieurs milliers de m/s ou une énergie en N signale une erreur probable de conversion ou de grandeur. Le contrôle de plausibilité complète le calcul ; il ne remplace pas les relations.
\pcPage{Les relations de fin de cycle}{Grandeurs, unités et conditions d’emploi}
\phantomsection\label{formules}
\begingroup\renewcommand{\arraystretch}{1.5}\fontsize{10.5}{13}\selectfont
\begin{tabularx}{\linewidth}{@{}p{37mm}p{37mm}Y@{}}\toprule
Relation & Unités cohérentes & Conditions et précautions\\\midrule
$\rho=m/V$ & kg ; m³ ; kg/m³ & Matériau et température précisés\\
$v=d/\Delta t$ & m ; s ; m/s & Moyenne sur le trajet ; uniforme si précisé\\
$P=mg$ & kg ; N/kg ; N & g du lieu ; P est ici le poids\\
$F=Gm_Am_B/d^2$ & kg ; m ; N & Loi fournie ; d entre centres ; G en SI\\
$E_c=\tfrac12mv^2$ & kg ; m/s ; J & Référentiel et système indiqués\\
$U=RI$ & V ; A ; $\Omega$ & Conducteur ohmique, conditions fixées\\
$P=UI$ & V ; A ; W & Courant continu dans le cours\\
$E=P\Delta t$ & W ; s ; J & Puissance constante ou moyenne fournie\\
$E=P\Delta t$ & kW ; h ; kWh & Ne pas mélanger W et kW\\
$d=v\Delta t$ & m ; s ; m/s & Distance totale parcourue par le signal\\
$f=N/\Delta t$ & répétitions ; s ; Hz & Signal périodique ; N motifs complets\\\bottomrule
\end{tabularx}\endgroup
\begin{pcEncadre}{attention}{Le contexte détermine les symboles}
P peut désigner un poids en N ou une puissance en W. Lire la légende et l’unité, puis écrire le nom de la grandeur au début de la démarche.
\end{pcEncadre}
\pcSousTitre{Lois des circuits}
En série dans une branche : même I. À un nœud : somme entrante = somme sortante. Dans une boucle source + deux récepteurs avec conventions adaptées : $U_G=U_1+U_2$. En dérivation entre les mêmes nœuds : même U.
\pcPage{Conversions et puissances de dix}{Une fiche pour éviter les erreurs d’unité}
\phantomsection\label{conversions}
\begin{questions}
\item $1\unite{kg}=1\,000\unite g$ ; $1\unite g=10^{-3}\unite{kg}$.
\item $1\unite L=1\unite{dm^3}$ ; $1\unite{mL}=1\unite{cm^3}$ ; $1\unite{m^3}=1\,000\unite L$.
\item $1\unite{g/mL}=1\,000\unite{kg/m^3}$. Un cube de 1 m de côté a un volume de $10^6\unite{cm^3}$.
\item $1\unite h=3\,600\unite s$ ; $1\unite{min}=60\unite s$. 15 min = 0,25 h, et non 0,15 h.
\item De km/h vers m/s : diviser par 3,6. De m/s vers km/h : multiplier par 3,6.
\item $1\unite A=1\,000\unite{mA}$ ; $1\unite{k\Omega}=1\,000\unite\Omega$ ; $1\unite{kW}=1\,000\unite W$.
\item $1\unite{Wh}=3\,600\unite J$ ; $1\unite{kWh}=3{,}6\times10^6\unite J$. Le kWh est une unité d’énergie.
\item $1\unite{kHz}=1\,000\unite{Hz}$ ; $1\unite{al}\approx9{,}46\times10^{15}\unite m$.
\end{questions}
\begin{pcEncadre}{method}{Le carré concerne toute la valeur}$(2\times10^3)^2=4\times10^6$ et non $2\times10^6$. Sur la calculatrice, saisir les parenthèses autour d’une distance écrite avec une puissance de dix avant de la mettre au carré.\end{pcEncadre}
\pcPage{Lire un graphique}{Exploiter un document et respecter la précision des données}
\phantomsection\label{graphiques}
\begin{questions}
\item Lire le titre, les noms des axes, les unités, l’origine et les graduations. Vérifier si l’axe des intensités est en A ou en mA.
\item Repérer la série et la légende ; distinguer mesure, modèle tracé et valeur de référence.
\item Pour une lecture, projeter le point vers les axes. Donner une valeur arrondie compatible avec la graduation.
\item Une droite ne prouve une proportionnalité que si elle est compatible avec le passage par l’origine et les conditions étudiées.
\item Pour un tableau, citer au moins une valeur pertinente et son unité avant de comparer.
\item Identifier l’auteur d’un document, sa date, son périmètre et les éventuelles données fictives. Une image seule n’est pas une mesure.
\end{questions}
\begin{pcEncadre}{method}{Construire une réponse ouverte}Dire ce que l’on cherche ; choisir les données utiles ; expliquer la relation ; calculer si nécessaire ; comparer ; conclure avec une limite. Si deux options ont des avantages différents, préciser le critère plutôt que proclamer une supériorité absolue.\end{pcEncadre}
\pcSousTitre{Exemple de réponse scientifique courte}
« Dans le modèle fourni, A émet 1 kg de CO2 et B en émet 4 kg. A a donc des émissions quatre fois plus faibles pour ce trajet. Cette comparaison porte seulement sur le périmètre et l’indicateur fournis ; elle ne décrit pas tous les impacts. »\pcPage{Air, solutions et climat}{Révisions transversales du cycle 4}
\phantomsection\label{air}
\pcSousTitre{Composition et pression}
L’air sec contient environ 78 \% de diazote, 21 \% de dioxygène et 1 \% d’autres gaz en volume. La vapeur d’eau est variable. Un gaz possède une masse et se comprime : à quantité de gaz conservée et après retour à la même température, un volume plus petit correspond à une pression plus élevée. La pression se mesure en pascals, Pa ; $1\unite{hPa}=100\unite{Pa}$ et $1\unite{bar}=10^5\unite{Pa}$. La pression atmosphérique usuelle est voisine de $10^5\unite{Pa}$.
\pcSousTitre{Dissolution et traitement des eaux}
Des gaz comme CO2 peuvent se dissoudre dans l’eau. Un gaz qui se dégage pouvait être déjà dissous : le dégazage n’est pas toujours une réaction. La miscibilité concerne deux liquides ; la solubilité caractérise une limite de dissolution dans des conditions données. Décantation, filtration et distillation séparent selon des propriétés différentes. Une eau clarifiée n’est pas automatiquement potable.
\pcSousTitre{Gaz à effet de serre}
CO2, vapeur d’eau, méthane et protoxyde d’azote font partie des gaz à effet de serre. Ils absorbent une partie du rayonnement infrarouge terrestre et émettent aussi du rayonnement. Un déséquilibre durable entre énergie reçue et énergie renvoyée peut modifier l’énergie du système Terre-atmosphère et sa température. L’effet de serre naturel est distinct de son renforcement lié aux activités humaines.
\pcSousTitre{Des conclusions qui respectent le périmètre}
Une absence de CO2 au point d’usage ne prouve pas une absence d’émissions sur tout le cycle de vie. L’acidification est une baisse de pH : un océan peut devenir moins basique tout en gardant un pH supérieur à 7. Comparer les émissions exige une même quantité de service et un périmètre annoncé.
\begin{pcEncadre}{method}{Trois questions flash, corrigées}
Dans 100 L d’air sec : environ 21 L de dioxygène. Une solution pH 8,0 reste basique vers 25 °C. Un écart énergétique reçu − renvoyé de 10 J correspond, dans un bilan sans autre échange, à une hausse d’énergie du système de 10 J.
\end{pcEncadre}
\pcPage{Tests et formules usuelles}{Une banque à exploiter avec le protocole donné}
\phantomsection\label{tests}
\begingroup\renewcommand{\arraystretch}{1.45}\begin{tabularx}{\linewidth}{@{}p{32mm}p{47mm}Y@{}}\toprule
Espèce recherchée & Test fourni & Résultat positif usuel\\\midrule
CO2 & Eau de chaux & Trouble\\
O2 & Bûchette incandescente & Ravivement, test encadré\\
H2 & Flamme, test sécurisé & Petite détonation, test encadré\\
Eau & Sulfate de cuivre anhydre & Blanc vers bleu\\
Cl− & Nitrate d’argent & Précipité blanc\\
Cu2+ & Soude & Précipité bleu\\
Fe2+ & Soude & Précipité vert\\
Fe3+ & Soude & Précipité brun\\\bottomrule
\end{tabularx}\endgroup
\pcSousTitre{Les formules à reconnaître}
Dihydrogène H2 ; dioxygène O2 ; diazote N2 ; eau H2O ; dioxyde de carbone CO2 ; méthane CH4 ; protoxyde d’azote N2O. Ces écritures décrivent la composition ; les noms ne se remplacent pas par « gaz » sans précision.
\begin{pcEncadre}{attention}{Une banque n’est pas une consigne de manipulation libre}
Utiliser les quantités, protections et traitements des déchets indiqués par le professeur. Un test positif s’interprète dans ses conditions ; un résultat négatif dépend aussi de la sensibilité et de l’état du réactif. Faire les tests sur des prélèvements séparés.
\end{pcEncadre}
\pcPage{Repères des éléments}{Z identifie le nombre de protons}
\phantomsection\label{elements}
Sélection de repérage ; ce tableau ne remplace pas une classification périodique complète. A n’est pas donné ici : on ne déduit pas un nombre précis de neutrons du seul Z.\par\vspace{3mm}
\begingroup\fontsize{10.5}{13}\selectfont\renewcommand{\arraystretch}{1.12}\begin{tabularx}{.85\linewidth}{@{}p{20mm}Yp{25mm}@{}}\toprule Z & Élément & Symbole\\\midrule
1 & Hydrogène & H\\
2 & Hélium & He\\
3 & Lithium & Li\\
4 & Béryllium & Be\\
5 & Bore & B\\
6 & Carbone & C\\
7 & Azote & N\\
8 & Oxygène & O\\
9 & Fluor & F\\
10 & Néon & Ne\\
11 & Sodium & Na\\
12 & Magnésium & Mg\\
13 & Aluminium & Al\\
14 & Silicium & Si\\
15 & Phosphore & P\\
16 & Soufre & S\\
17 & Chlore & Cl\\
18 & Argon & Ar\\
19 & Potassium & K\\
20 & Calcium & Ca\\
26 & Fer & Fe\\
29 & Cuivre & Cu\\
30 & Zinc & Zn\\
\bottomrule\end{tabularx}\endgroup
\pcPage{Diagnostic : corrigé}{Un repère pour reprendre le bon chapitre}
\phantomsection\label{cor-diagnostic}
\begin{questions}
\item 250 mL et 0,350 kg. Reprendre la fiche conversions.
\item Non : une solution homogène peut contenir plusieurs espèces ; chapitre 1.
\item Quatre H et deux O ; chapitre 2.
\item 11 électrons, par neutralité ; chapitre 2.
\item Acide, car pH < 7 ; chapitre 4.
\item $v=120/30=4\unite{m/s}$ ; chapitre 8.
\item Poids en N ; masse en kg ; chapitre 9.
\item Ampèremètre en série ; voltmètre en dérivation ; chapitre 13.
\item $0{,}020\unite A$ et $900\unite s$ ; conversions.
\item La distance totale comprend l’aller et le retour ; chapitre 17.
\end{questions}
\begin{pcEncadre}{method}{Décider de la suite}Si une conversion est erronée, refaire deux exemples avec leurs unités. Si la notion est confuse, relire le cours et expliquer le schéma. Si la réponse est juste mais difficile à suivre, retravailler les phrases et l’organisation.\end{pcEncadre}
\pcPage{Lexique scientifique}{Des définitions à replacer dans une situation}
\phantomsection\label{lexique}
\begingroup\fontsize{10.5}{13.5}\selectfont\begin{multicols}{2}
\par\addvspace{2mm}\noindent\begin{minipage}{\linewidth}\textbf{Atome}\enspace Entité neutre comportant un noyau et des électrons.\par\end{minipage}\par
\par\addvspace{2mm}\noindent\begin{minipage}{\linewidth}\textbf{Ion}\enspace Entité chargée ayant perdu ou gagné des électrons.\par\end{minipage}\par
\par\addvspace{2mm}\noindent\begin{minipage}{\linewidth}\textbf{Nucléon}\enspace Proton ou neutron du noyau.\par\end{minipage}\par
\par\addvspace{2mm}\noindent\begin{minipage}{\linewidth}\textbf{Numéro atomique}\enspace Z : nombre de protons ; identifie l’élément.\par\end{minipage}\par
\par\addvspace{2mm}\noindent\begin{minipage}{\linewidth}\textbf{Molécule}\enspace Assemblage neutre d’atomes liés.\par\end{minipage}\par
\par\addvspace{2mm}\noindent\begin{minipage}{\linewidth}\textbf{Espèce chimique}\enspace Ensemble d’entités chimiques identiques.\par\end{minipage}\par
\par\addvspace{2mm}\noindent\begin{minipage}{\linewidth}\textbf{Réactif}\enspace Espèce consommée lors d’une réaction.\par\end{minipage}\par
\par\addvspace{2mm}\noindent\begin{minipage}{\linewidth}\textbf{Produit}\enspace Espèce formée lors d’une réaction.\par\end{minipage}\par
\par\addvspace{2mm}\noindent\begin{minipage}{\linewidth}\textbf{Corrosion}\enspace Transformation d’un métal avec son environnement.\par\end{minipage}\par
\par\addvspace{2mm}\noindent\begin{minipage}{\linewidth}\textbf{Dissolution}\enspace Dispersion d’un soluté dans un solvant.\par\end{minipage}\par
\par\addvspace{2mm}\noindent\begin{minipage}{\linewidth}\textbf{Dilution}\enspace Ajout de solvant sans ajout de soluté.\par\end{minipage}\par
\par\addvspace{2mm}\noindent\begin{minipage}{\linewidth}\textbf{Masse volumique}\enspace Masse divisée par le volume, dans des conditions précisées.\par\end{minipage}\par
\par\addvspace{2mm}\noindent\begin{minipage}{\linewidth}\textbf{Force}\enspace Modèle d’une action : point, direction, sens et valeur.\par\end{minipage}\par
\par\addvspace{2mm}\noindent\begin{minipage}{\linewidth}\textbf{Poids}\enspace Force de pesanteur exercée sur un objet.\par\end{minipage}\par
\par\addvspace{2mm}\noindent\begin{minipage}{\linewidth}\textbf{Référentiel}\enspace Référence choisie pour décrire les positions et leur évolution.\par\end{minipage}\par
\par\addvspace{2mm}\noindent\begin{minipage}{\linewidth}\textbf{Trajectoire}\enspace Ensemble des positions successives d’un point.\par\end{minipage}\par
\par\addvspace{2mm}\noindent\begin{minipage}{\linewidth}\textbf{Énergie cinétique}\enspace Énergie associée au mouvement dans un référentiel.\par\end{minipage}\par
\par\addvspace{2mm}\noindent\begin{minipage}{\linewidth}\textbf{Transfert}\enspace Passage d’énergie entre systèmes.\par\end{minipage}\par
\par\addvspace{2mm}\noindent\begin{minipage}{\linewidth}\textbf{Conversion}\enspace Changement de forme d’énergie.\par\end{minipage}\par
\par\addvspace{2mm}\noindent\begin{minipage}{\linewidth}\textbf{Puissance}\enspace Énergie transférée par unité de durée.\par\end{minipage}\par
\par\addvspace{2mm}\noindent\begin{minipage}{\linewidth}\textbf{Tension}\enspace Grandeur électrique entre deux points, en V.\par\end{minipage}\par
\par\addvspace{2mm}\noindent\begin{minipage}{\linewidth}\textbf{Intensité}\enspace Grandeur caractérisant le courant, en A.\par\end{minipage}\par
\par\addvspace{2mm}\noindent\begin{minipage}{\linewidth}\textbf{Résistance}\enspace Grandeur d’un conducteur ohmique, en Ω.\par\end{minipage}\par
\par\addvspace{2mm}\noindent\begin{minipage}{\linewidth}\textbf{Fréquence}\enspace Nombre de répétitions périodiques par seconde.\par\end{minipage}\par
\par\addvspace{2mm}\noindent\begin{minipage}{\linewidth}\textbf{Signal}\enspace Phénomène détectable pouvant transporter de l’information.\par\end{minipage}\par
\par\addvspace{2mm}\noindent\begin{minipage}{\linewidth}\textbf{Année-lumière}\enspace Distance parcourue par la lumière dans le vide en un an.\par\end{minipage}\par
\par\addvspace{2mm}\noindent\begin{minipage}{\linewidth}\textbf{Témoin}\enspace Comparaison permettant d’examiner une cause possible.\par\end{minipage}\par
\par\addvspace{2mm}\noindent\begin{minipage}{\linewidth}\textbf{Modèle}\enspace Représentation simplifiée dont le domaine d’emploi est précisé.\par\end{minipage}\par
\end{multicols}\endgroup
\pcPage{Programme et concordance}{Repérer les apprentissages du cycle 4}
\phantomsection\label{programme}
\pdfbookmark[0]{Programme et références}{bm-programme}
Les lignes relient les notions du programme en vigueur aux cours, exercices et sujets blancs. Le découpage annuel reste éditorial. La loi de gravitation est utilisée quand elle est fournie ; la formule de l’énergie potentielle n’est pas un prérequis exigé par ce manuel.\par\vspace{2mm}
\par\noindent\begin{minipage}{\linewidth}\textbf{M01} · États, changements d’état et interprétation microscopique. \hyperref[chap-1]{Chapitre 1}, exercices 1, 3, 10, 14. \emph{Cours ; sujet 7.}\par\end{minipage}\par\vspace{2mm}
\par\noindent\begin{minipage}{\linewidth}\textbf{M02} · Corps purs, mélanges, espèces et séparation. \hyperref[chap-1]{Chapitre 1}, exercices 2, 9, 11, 15. \emph{Fiche air ; sujet 2.}\par\end{minipage}\par\vspace{2mm}
\par\noindent\begin{minipage}{\linewidth}\textbf{M03} · Masse volumique, mesure, unités et température. \hyperref[chap-1]{Chapitre 1}, exercices 5, 6, 7, 8, 12, 13. \emph{Sujets 2 et 7.}\par\end{minipage}\par\vspace{2mm}
\par\noindent\begin{minipage}{\linewidth}\textbf{M04} · Solubilité, miscibilité et dissolution de gaz. \hyperref[chap-1]{Chapitre 1}, exercices 9, 11, 15. \emph{Fiche air, solutions et climat.}\par\end{minipage}\par\vspace{2mm}
\par\noindent\begin{minipage}{\linewidth}\textbf{M05} · Atomes, molécules, ions, formules et tableau périodique. \hyperref[chap-2]{Chapitre 2}, exercices 4, 6, 7, 8, 9, 14. \emph{Sujets 2, 5 et 7.}\par\end{minipage}\par\vspace{2mm}
\par\noindent\begin{minipage}{\linewidth}\textbf{M06} · Protons, neutrons, électrons et structure du noyau. \hyperref[chap-2]{Chapitre 2}, exercices 1, 2, 3, 8, 13, 15. \emph{Fiche des éléments ; sujet 7.}\par\end{minipage}\par\vspace{2mm}
\par\noindent\begin{minipage}{\linewidth}\textbf{M07} · Tests caractéristiques et transformations physiques/chimiques. \hyperref[chap-3]{Chapitre 3}, exercices 1, 3, 11, 12, 14. \emph{Sujets 5 et 10.}\par\end{minipage}\par\vspace{2mm}
\par\noindent\begin{minipage}{\linewidth}\textbf{M08} · Équation fournie, redistribution des atomes et conservation. \hyperref[chap-3]{Chapitre 3}, exercices 4, 5, 7, 8, 9, 10, 13. \emph{Sujets 1 et 10.}\par\end{minipage}\par\vspace{2mm}
\par\noindent\begin{minipage}{\linewidth}\textbf{M09} · Acidité, basicité, pH et ions H+ ou OH−. \hyperref[chap-4]{Chapitre 4}, exercices 1, 3, 6, 7, 10, 14. \emph{Sujets 2 et 5.}\par\end{minipage}\par\vspace{2mm}
\par\noindent\begin{minipage}{\linewidth}\textbf{M10} · Réactions acide-base et acide-métal. \hyperref[chap-5]{Chapitre 5}, exercices 2, 3, 6, 7, 12, 14. \emph{Réaction acide-base au chapitre 4.}\par\end{minipage}\par\vspace{2mm}
\par\noindent\begin{minipage}{\linewidth}\textbf{M11} · Corrosion, environnement et protections. \hyperref[chap-5]{Chapitre 5}, exercices 1, 4, 9, 10, 13, 15. \emph{Sujet 5.}\par\end{minipage}\par\vspace{2mm}
\par\noindent\begin{minipage}{\linewidth}\textbf{M12} · Air, combustions, gaz à effet de serre et impacts. \hyperref[chap-6]{Chapitre 6}, exercices 1, 6, 8, 11, 13, 14. \emph{Fiches air/climat ; sujet 1.}\par\end{minipage}\par\vspace{2mm}
\par\noindent\begin{minipage}{\linewidth}\textbf{M13} · Univers, Système solaire, évolution, éléments et ressources. \hyperref[chap-7]{Chapitre 7}, exercices 1, 2, 7, 8, 12, 15. \emph{Année-lumière au chapitre 16 ; sujet 9.}\par\end{minipage}\par\vspace{2mm}
\par\noindent\begin{minipage}{\linewidth}\textbf{V01} · Référentiel, relativité du mouvement et trajectoires. \hyperref[chap-8]{Chapitre 8}, exercices 1, 3, 8, 9, 13, 14. \emph{Sujet 1.}\par\end{minipage}\par\vspace{2mm}
\par\noindent\begin{minipage}{\linewidth}\textbf{V02} · Vitesse : valeur, direction, sens ; distance et durée. \hyperref[chap-8]{Chapitre 8}, exercices 4, 5, 6, 7, 10, 11, 12. \emph{Sujets 1 et 6.}\par\end{minipage}\par\vspace{2mm}
\par\noindent\begin{minipage}{\linewidth}\textbf{V03} · Interactions, forces et quatre caractéristiques. \hyperref[chap-9]{Chapitre 9}, exercices 1, 4, 9, 10, 13, 14. \emph{Sujet 4.}\par\end{minipage}\par\vspace{2mm}
\par\noindent\begin{minipage}{\linewidth}\textbf{V04} · Poids, masse et P = mg. \hyperref[chap-9]{Chapitre 9}, exercices 5, 6, 7, 8, 11, 15. \emph{Sujet 4.}\par\end{minipage}\par\vspace{2mm}
\par\noindent\begin{minipage}{\linewidth}\textbf{V05} · Exploitation de la loi fournie de gravitation. \hyperref[chap-10]{Chapitre 10}, exercices 5, 6, 8, 9, 10, 11, 14. \emph{Sujets 4 et 9.}\par\end{minipage}\par\vspace{2mm}
\par\noindent\begin{minipage}{\linewidth}\textbf{E01} · Énergie cinétique et énergie potentielle de position. \hyperref[chap-11]{Chapitre 11}, exercices 4, 6, 7, 8, 9, 15. \emph{Sujet 1.}\par\end{minipage}\par\vspace{2mm}
\par\noindent\begin{minipage}{\linewidth}\textbf{E02} · Sources, formes, transferts, conversions et conservation. \hyperref[chap-12]{Chapitre 12}, exercices 1, 2, 6, 7, 8, 15. \emph{Sujets 3, 7 et 8.}\par\end{minipage}\par\vspace{2mm}
\par\noindent\begin{minipage}{\linewidth}\textbf{E03} · Rayonnement, absorption et système Terre-atmosphère. \hyperref[chap-12]{Chapitre 12}, exercices 5, 9, 10, 14, 15. \emph{Fiche climat ; sujet 8.}\par\end{minipage}\par\vspace{2mm}
\par\noindent\begin{minipage}{\linewidth}\textbf{E04} · Circuits série/dérivation et cahier des charges. \hyperref[chap-13]{Chapitre 13}, exercices 3, 4, 8, 10, 13, 15. \emph{Sujet 3.}\par\end{minipage}\par\vspace{2mm}
\par\noindent\begin{minipage}{\linewidth}\textbf{E05} · Mesures, unicité et additivité des intensités/tensions. \hyperref[chap-13]{Chapitre 13}, exercices 1, 6, 7, 8, 11, 12. \emph{Sujet 3.}\par\end{minipage}\par\vspace{2mm}
\par\noindent\begin{minipage}{\linewidth}\textbf{E06} · Loi d’Ohm et caractéristique. \hyperref[chap-14]{Chapitre 14}, exercices 3, 4, 5, 8, 9, 13. \emph{Sujet 10.}\par\end{minipage}\par\vspace{2mm}
\par\noindent\begin{minipage}{\linewidth}\textbf{E07} · Puissance électrique P = UI. \hyperref[chap-14]{Chapitre 14}, exercices 6, 7, 9, 14, 15. \emph{Sujets 3, 6 et 10.}\par\end{minipage}\par\vspace{2mm}
\par\noindent\begin{minipage}{\linewidth}\textbf{E08} · Énergie, puissance, durée, consommation et stockage. \hyperref[chap-15]{Chapitre 15}, exercices 2, 3, 4, 6, 7, 11, 15. \emph{Sujets 3, 6 et 8.}\par\end{minipage}\par\vspace{2mm}
\par\noindent\begin{minipage}{\linewidth}\textbf{S01} · Sources, diffusion, rayon lumineux et propagation. \hyperref[chap-16]{Chapitre 16}, exercices 1, 3, 5, 8, 10, 11. \emph{Sujet 9.}\par\end{minipage}\par\vspace{2mm}
\par\noindent\begin{minipage}{\linewidth}\textbf{S02} · Vitesse de lumière, distances et année-lumière. \hyperref[chap-16]{Chapitre 16}, exercices 4, 6, 7, 12, 13, 15. \emph{Sujets 4 et 9.}\par\end{minipage}\par\vspace{2mm}
\par\noindent\begin{minipage}{\linewidth}\textbf{S03} · Propagation sonore, fréquence et plages audibles. \hyperref[chap-17]{Chapitre 17}, exercices 1, 2, 4, 6, 7, 13. \emph{Sujet 6.}\par\end{minipage}\par\vspace{2mm}
\par\noindent\begin{minipage}{\linewidth}\textbf{S04} · Écho, sonar, distance et prévention auditive. \hyperref[chap-17]{Chapitre 17}, exercices 3, 8, 9, 11, 14, 15. \emph{Sujet 6.}\par\end{minipage}\par\vspace{2mm}
\par\noindent\begin{minipage}{\linewidth}\textbf{S05} · Signaux, information, codage et mesures. \hyperref[chap-18]{Chapitre 18}, exercices 1, 4, 6, 7, 10, 12, 15. \emph{Sujets 4 et 6.}\par\end{minipage}\par\vspace{2mm}
\pcPage{Compétences et numérique}{Pratiquer les démarches, les langages et les outils}
\begin{questions}
\item Démarche expérimentale : hypothèse, témoin, répétition et conclusion, chapitres 1, 3, 4 et 5.
\item Concevoir un dispositif : protocole de mesure et montage répondant à une demande, chapitres 9, 13 et 14.
\item Mobiliser le numérique : nuages de points au tableur, chapitres 1, 10, 11 et 14 ; vidéo au chapitre 8 ; signaux enregistrés aux chapitres 17 et 18.
\item Pratiquer les langages : équations, forces, circuits et graphiques ; réponses rédigées de tous les sujets blancs.
\item Adopter une attitude responsable : produits, déchets et ressources ; sécurité électrique ; protection visuelle et auditive.
\item Se situer dans le temps et l’espace : origine des éléments, évolution de l’Univers et ordres de grandeur, chapitres 2, 7 et 16.
\end{questions}
\begin{pcEncadre}{method}{Atelier tableur reproductible}Saisir v = 0, 2, 4, 6, 8 m/s ; fixer m = 2 kg. Calculer Ec par une formule recopiable. Conserver les données d’origine. Produire un nuage de points Ec(v), puis Ec(v²), avec noms des axes et unités. Expliquer la différence des deux graphes.\end{pcEncadre}
\pcPage{Références et choix éditoriaux}{Sources primaires vérifiées le 24 octobre 2026}
\phantomsection\label{references}
\par\noindent\begin{minipage}{\linewidth}\pcSousTitre{Programme en application}
Éduscol, programmes du cycle 4 pour 2026-2027 ; physique-chimie : BO du 30 juillet 2020.\par{\fontsize{8.5}{11}\selectfont\raggedright\url{https://eduscol.education.gouv.fr/4362/enseigner-au-cycle-4}\par}\end{minipage}\par
\par\noindent\begin{minipage}{\linewidth}\pcSousTitre{Texte de référence}
Programme de physique-chimie du cycle 4, pages 96 à 106 ; découpage annuel éditorial.\par{\fontsize{8.5}{11}\selectfont\raggedright\url{https://eduscol.education.fr/document/621/download}\par}\end{minipage}\par
\par\noindent\begin{minipage}{\linewidth}\pcSousTitre{Épreuves du DNB 2027}
Note de service du 11 septembre 2026, NOR MENE2623228N, partie I.4 : sciences.\par{\fontsize{8.5}{11}\selectfont\raggedright\url{https://www.education.gouv.fr/bo/2026/Hebdo4/MENE2623228N}\par}\end{minipage}\par
\par\noindent\begin{minipage}{\linewidth}\pcSousTitre{Présentation des épreuves}
Éduscol : durée, notation et compétences ; page actualisée en octobre 2026.\par{\fontsize{8.5}{11}\selectfont\raggedright\url{https://eduscol.education.gouv.fr/5607/les-epreuves-du-dnb}\par}\end{minipage}\par
\par\noindent\begin{minipage}{\linewidth}\pcSousTitre{Sujet de référence}
Éduscol, physique-chimie, série générale, n° 01 : repérage des attendus et de la rédaction.\par{\fontsize{8.5}{11}\selectfont\raggedright\url{https://eduscol.education.gouv.fr/sites/default/files/document/dnb-sujet-de-reference-physique-chimie-serie-generale-n01-129673.pdf}\par}\end{minipage}\par
\pcPage{Références complémentaires}{Poursuivre la lecture des sources primaires}
\par\noindent\begin{minipage}{\linewidth}\pcSousTitre{Annale officielle 2026}
Ministère, sciences, série générale : lecture critique du format et des démarches attendues. Les énoncés ne sont pas reproduits.\par{\fontsize{8.5}{11}\selectfont\raggedright\url{https://www.education.gouv.fr/sites/default/files/document/diplome-national-du-brevet-2026-sciences-serie-generale-517070.pdf}\par}\end{minipage}\par
\par\noindent\begin{minipage}{\linewidth}\pcSousTitre{Atome et noyau}
CEA : protons, neutrons, électrons et numéro atomique.\par{\fontsize{8.5}{11}\selectfont\raggedright\url{https://www.cea.fr/comprendre/enseignants/pages/ressources-pedagogiques/essentiel-sur/physique-chimie/essentiel-sur-noyaux-des-atomes.aspx}\par}\end{minipage}\par
\par\noindent\begin{minipage}{\linewidth}\pcSousTitre{Unités}
BIPM, Brochure sur le SI : unités, symboles et préfixes.\par{\fontsize{8.5}{11}\selectfont\raggedright\url{https://www.bipm.org/fr/publications/si-brochure}\par}\end{minipage}\par
\par\noindent\begin{minipage}{\linewidth}\pcSousTitre{Chronologie cosmique}
NASA Science : âges approximatifs de l’Univers et du Système solaire.\par{\fontsize{8.5}{11}\selectfont\raggedright\url{https://science.nasa.gov/exoplanets/what-is-the-universe/}\par}\end{minipage}\par
\par\noindent\begin{minipage}{\linewidth}\pcSousTitre{Ressources pédagogiques}
Éduscol : accompagnement du programme de physique-chimie au cycle 4.\par{\fontsize{8.5}{11}\selectfont\raggedright\url{https://eduscol.education.gouv.fr/5739/ressources-d-accompagnement-du-programme-de-physique-chimie-au-cycle-4}\par}\end{minipage}\par
\pcPage{Références complémentaires}{Poursuivre la lecture des sources primaires}
\par\noindent\begin{minipage}{\linewidth}\pcSousTitre{Prévention électrique}
INRS : risques et prévention en électricité.\par{\fontsize{8.5}{11}\selectfont\raggedright\url{https://www.inrs.fr/risques/electriques/risques-electricite.html}\par}\end{minipage}\par
\par\noindent\begin{minipage}{\linewidth}\pcSousTitre{Éléments et symboles}
IUPAC : tableau périodique des éléments.\par{\fontsize{8.5}{11}\selectfont\raggedright\url{https://iupac.org/what-we-do/periodic-table-of-elements/}\par}\end{minipage}\par
\par\noindent\begin{minipage}{\linewidth}\pcSousTitre{Système solaire}
Agence spatiale canadienne : formation des systèmes solaires.\par{\fontsize{8.5}{11}\selectfont\raggedright\url{https://www.asc-csa.gc.ca/fra/blogue/2021/08/06/comment-se-forment-les-systemes-solaires.asp}\par}\end{minipage}\par
\pcSousTitre{Création et réutilisation}
Les sujets blancs, documents d’étude, barèmes et schémas sont des créations éditoriales OptiTAB. Les sujets ministériels ont été consultés pour analyser les attendus, sans reproduction de leurs exercices ni de leurs images. Les activités de révision sont harmonisées avec les autres manuels de la collection. Les nombres d’étude, prix et scénarios sont fictifs sauf indication explicite.

\end{document}
